% concatenated from master_file_8-6-25.tex
\documentclass[10pt, letterpaper]{amsart}

\usepackage{amscd, amssymb}
\usepackage{enumerate}
\usepackage[all]{xypic}
\usepackage{multirow}
\usepackage{hhline}
\usepackage{verbatim}
\usepackage{mathdots}
\usepackage{comment}
\usepackage{tikz-cd}
\usepackage[titletoc]{appendix}
\usepackage{stmaryrd}%this conflicts with extpfeil
\usepackage{graphicx}
\usepackage{mathtools}
%\usepackage{extpfeil}
%\usepackage{amsrefs}
%\usepackage{thmtools}
%\usepackage{float}
%\usepackage{makeidx}

%\usepackage{amsthm}
\usepackage{leftidx}
\usepackage{mathabx }
\usepackage{amsfonts}
\usepackage{mathrsfs}

\usepackage{dynkin-diagrams}
\usepackage{xcolor}


\usepackage[colorlinks=true, linkcolor=blue,
 citecolor=blue, urlcolor=blue]{hyperref}

%\usetikzlibrary{decorations.markings}


%\renewcommand{\thesection}{\Roman{section}}

%%% table of contents %%%
\setcounter{tocdepth}{2}
\let\oldtocsection=\tocsection
\renewcommand{\tocsection}[2]{\hspace{0em}\oldtocsection{#1}{#2}}





\theoremstyle{plain}
\newtheorem{Thm}[equation]{Theorem}
\newtheorem*{Thm*}{Theorem}
\newtheorem{Cor}[equation]{Corollary}
\newtheorem{Prop}[equation]{Proposition}
\newtheorem{Lem}[equation]{Lemma}


\theoremstyle{definition}
\newtheorem*{Rmk}{Remark}
\newtheorem{Def}[equation]{Definition}
\newtheorem*{hypo}{Working
Hypotheses}


\numberwithin{equation}{section}



%\newcommand{\A}{\mathbb{A}}
\newcommand{\C}{\mathbb{C}}
\newcommand{\F}{\mathbb{F}}
%\newcommand{\HH}{\mathbb{H}}
%\newcommand{\K}{\mathbb{K}}
\newcommand{\N}{\mathbb{N}}
\newcommand{\Q}{\mathbb{Q}}
\newcommand{\R}{\mathbb{R}}
%\newcommand{\V}{\mathbb{V}}
%\newcommand{\W}{\mathbb{W}}
%\newcommand{\X}{\mathbb{X}}
%\newcommand{\Y}{\mathbb{Y}}
\newcommand{\Z}{\mathbb{Z}}
%\newcommand{\OOmega}{\mathbf{\Omega}}



\newcommand{\OF}{\mathcal{O}_F}
\newcommand{\OFd}{\mathcal{O}_F^{\diamondsuit}}


%\newcommand{\Gm}{\mathbb{G}_m}
\newcommand{\GL}{\operatorname{GL}}
%\newcommand{\PGL}{\operatorname{PGL}}
\newcommand{\SL}{\operatorname{SL}}
\newcommand{\Sp}{\operatorname{Sp}}
%\newcommand{\GSp}{\operatorname{GSp}}
\newcommand{\SO}{\operatorname{SO}}
\newcommand{\OO}{\operatorname{O}}
%\newcommand{\UU}{\operatorname{U}}
%\newcommand{\SU}{\operatorname{SU}}
\newcommand{\Mp}{\operatorname{Mp}}

\newcommand{\Rad}{\operatorname{rad}}
\newcommand{\pr}{\operatorname{pr}}



\newcommand{\xt}{\tilde{x}}
\newcommand{\yt}{\tilde{y}}
\newcommand{\zt}{\tilde{z}}
\newcommand{\wt}{\tilde{w}}
\newcommand{\htt}{\tilde{h}}


\newcommand{\xa}{x_{\alpha}}
\newcommand{\wa}{w_{\alpha}}
\newcommand{\ha}{h_{\alpha}}
\newcommand{\xb}{x_{\beta}}
\newcommand{\wb}{w_{\beta}}
\newcommand{\hb}{h_{\beta}}

\newcommand{\xta}{\tilde{x}_{\alpha}}
\newcommand{\wta}{\tilde{w}_{\alpha}}
\newcommand{\hta}{\tilde{h}_{\alpha}}
\newcommand{\xtb}{\tilde{x}_{\beta}}
\newcommand{\wtb}{\tilde{w}_{\beta}}
\newcommand{\htb}{\tilde{h}_{\beta}}



\newcommand{\alphac}{{\check{\alpha}}}
\newcommand{\betac}{{\check{\beta}}}
\newcommand{\gammac}{{\check{\gamma}}}

\newcommand{\g}{\mathfrak{G}}
\newcommand{\h}{\mathfrak{h}}
\renewcommand{\sp}{\mathfrak{sp}}


\newcommand{\Spt}{\widetilde{\operatorname{Sp}}}
\newcommand{\Bt}{\widetilde{B}}
\newcommand{\Mt}{\widetilde{M}}
\newcommand{\Nt}{\widetilde{N}}
\newcommand{\Pt}{\widetilde{P}}
\newcommand{\Qt}{\widetilde{Q}}
\newcommand{\Gt}{\widetilde{G}}
\newcommand{\Ht}{\widetilde{H}}
\newcommand{\It}{\widetilde{I}}
\newcommand{\Jt}{\widetilde{J}}
\newcommand{\Kt}{\widetilde{K}}
\newcommand{\Tt}{\widetilde{T}}
\newcommand{\Xt}{\widetilde{X}}
\newcommand{\XXt}{\widetilde{\mathbb{X}}}
\newcommand{\Yt}{\widetilde{Y}}
\newcommand{\Wt}{\widetilde{W}}
\newcommand{\Omegat}{\widetilde{\Omega}}

\newcommand{\Gammat}{\widetilde{\Gamma}}
\newcommand{\Gammaft}{\widetilde{\mathsf{\Gamma}}}
\newcommand{\Gammatp}{\widetilde{\Gamma}'}


\newcommand{\mt}{\tilde{m}}
\newcommand{\gt}{\tilde{g}}
\newcommand{\hti}{\tilde{h}}
\newcommand{\kt}{\tilde{k}}
\newcommand{\pt}{\tilde{p}}
\newcommand{\qt}{\tilde{q}}
\newcommand{\ttt}{\tilde{t}}

\newcommand{\taut}{\tilde{\tau}}
\newcommand{\pit}{\tilde{\pi}}


\newcommand{\nuh}{\hat{\nu}}



\newcommand{\Ind}{\operatorname{Ind}}
\newcommand{\cInd}{\operatorname{c-Ind}}
\newcommand{\ind}{\operatorname{ind}}
\newcommand{\Irr}{\operatorname{Irr}}
\newcommand{\diag}{\operatorname{diag}}
\newcommand{\supp}{\operatorname{supp}}
\newcommand{\ch}{\operatorname{ch}}
\newcommand{\Span}{\operatorname{span}}
\newcommand{\rank}{\operatorname{rank}}
\newcommand{\End}{\operatorname{End}}
\newcommand{\Hom}{\operatorname{Hom}}
\newcommand{\Aut}{\operatorname{Aut}}
\newcommand{\ord}{\operatorname{ord}}
\newcommand{\vol}{\operatorname{vol}}
\newcommand{\sign}{\operatorname{sign}}
\newcommand{\sig}{\operatorname{sig}}
\newcommand{\Herm}{\operatorname{Herm}}
\newcommand{\Quad}{\operatorname{Quad}}
\newcommand{\Sym}{\operatorname{Sym}}
\newcommand{\image}{\operatorname{Im}}
\newcommand{\Int}{\operatorname{Int}}
\newcommand{\Gal}{\operatorname{Gal}}
\newcommand{\Res}{\operatorname{Res}}
\newcommand{\disc}{\operatorname{disc}}
\newcommand{\tr}{\operatorname{tr}}
\newcommand{\Tr}{\operatorname{Tr}}
\newcommand{\rad}{\operatorname{rad}}

\newcommand{\Aff}{\operatorname{Aff}}
\newcommand{\aff}{\mathrm{aff}}
\newcommand{\af}{\mathrm{af}}
\newcommand{\ea}{\mathrm{ea}}



\newcommand{\Sm}{{\mathcal{S}}}

\newcommand{\ee}{\mathbf{e}}
\newcommand{\ff}{\mathbf{f}}
\newcommand{\rr}{\mathbf{r}}
\newcommand{\hh}{\mathbf{h}}
\newcommand{\xx}{\mathbf{x}}
\newcommand{\yy}{\mathbf{y}}
\newcommand{\ww}{\mathbf{w}}


\newcommand{\s}{\mathbf{s}}

\newcommand{\rp}{\mathbf{r}_{\psi}}

\newcommand{\Aa}{A^{(2)}}
\newcommand{\Bb}{B^{(2)}}
\newcommand{\II}{\mathcal{I}}
\newcommand{\JJ}{\mathcal{J}}
\newcommand{\KK}{\mathcal{K}}
\newcommand{\LL}{\mathcal{L}}
\newcommand{\MM}{\mathcal{M}}

\newcommand{\KKt}{\widetilde{\mathcal{K}}}
\newcommand{\IIt}{\widetilde{\mathcal{I}}}
\newcommand{\JJt}{\widetilde{\mathcal{J}}}

\newcommand{\X}{\mathfrak{X}}


\newcommand{\Scal}{\mathcal{S}}
\newcommand{\Ocal}{\mathcal{O}}
\newcommand{\Fcal}{\mathcal{F}}
\newcommand{\Wcal}{\mathcal{W}}
\newcommand{\Hcal}{\mathcal{H}}
\newcommand{\Ncal}{\mathcal{N}}
\newcommand{\Dcal}{\mathcal{D}}


\newcommand{\Bf}{\mathsf{B}}
\newcommand{\Ff}{\mathsf{F}}
\newcommand{\Hf}{\mathsf{H}}
\newcommand{\If}{\mathsf{I}}
\newcommand{\Kf}{\mathsf{K}}
\newcommand{\Mf}{\mathsf{M}}
\newcommand{\Nf}{\mathsf{N}}
\newcommand{\Tf}{\mathsf{T}}
\newcommand{\Uf}{\mathsf{U}}
\newcommand{\Vf}{\mathsf{V}}
\newcommand{\Gammaf}{\mathsf{\Gamma}}

\newcommand{\Hft}{\widetilde{\mathsf{H}}}
\newcommand{\Ift}{\widetilde{\mathsf{I}}}
\newcommand{\Kft}{\widetilde{\mathsf{K}}}
\newcommand{\Mft}{\widetilde{\mathsf{M}}}
\newcommand{\Nft}{\widetilde{\mathsf{N}}}
\newcommand{\Bft}{\widetilde{\mathsf{B}}}
\newcommand{\Pft}{\widetilde{\mathsf{P}}}
\newcommand{\Tft}{\widetilde{\mathsf{T}}}



\newcommand{\la}{\langle}
\newcommand{\ra}{\rangle}
\newcommand{\lla}{\langle\!\langle}
\newcommand{\rra}{\rangle\!\rangle}
\newcommand{\llp}{(\!(}
\newcommand{\rrp}{)\!)}

\newcommand{\ie}{{\it i.e.}\ }
%\newcommand{\apriori}{{\it a priori }}
%\newcommand{\afortiori}{{\it a fortiori }}

\newcommand{\st}{\,:\,}
\newcommand{\qand}{{\quad\text{and}\quad}}
\newcommand{\qqand}{{\qquad\text{and}\qquad}}
\newcommand{\qor}{{\quad\text{or}\quad}}
\newcommand{\by}{{\quad\text{by}\quad}}

\newcommand{\ve}{\varepsilon}


\newcommand{\one}{\mathbf{1}}




\setlength{\oddsidemargin}{0.2in} \setlength{\evensidemargin}{0.2in}
\setlength{\textwidth}{5.8in}

\title[$K$-types for wild double covers]{Minimal depth $K$-types for wild double covers\\
and Shimura correspondences}


\author{Edmund Karasiewicz}

\author{Shuichiro Takeda}
%\address{Mathematics Dept, Univ of Missouri-Columbia, 202
%  Math Sciences Building, Columbia, MO, 65211}
%\email{takedas@missouri.edu}
\subjclass[2020]{11F70, 22E50}
\keywords{Bernstein Components; Hecke Algebra; $p$-adic groups; Metaplectic Group.}


%\pgfkeys{/Dynkin diagram, edge-length=1cm, root-radius=.1cm}



\def\shui#1{\textcolor{blue}{#1}}



\begin{document}


\begin{abstract}
	We construct some Iwahori types, in the sense of Bushnell-Kutzko, for the double cover of an almost simple simply-laced simply-connected Chevalley group $\widetilde{G}$ over any $2$-adic field. These types capture the covering group analog of the Bernstein block of unramified principal series. 
    
    We also prove that the associated Hecke algebra essentially admits an Iwahori-Matsumoto (IM) presentation. The complete presentation is obtained for types $A_{r}$, $D_{2r+1}$, $E_{6}$, $E_{7}$; for the other types, some technical obstacles remain. Those Hecke algebras with the complete IM presentation are isomorphic to Iwahori-Hecke algebras of explicit linear Chevalley groups, giving rise to Shimura correspondences.

    Along the way, we show that the Iwahori type extends to a hyperspecial maximal compact subgroup $\widetilde{K}\subseteq \widetilde{G}$. This extension has minimal depth among the genuine $\widetilde{K}$-representations and allows us to construct a finite Shimura correspondence, generalizing a result of Savin.
\end{abstract}

\maketitle


\setcounter{tocdepth}{1}
\tableofcontents

%%%%%%%%%%%%%%%%%%%%%%%%%%%%%%%%%%%%%%%%%%%%%%%%%%%%%%%%%%%%%%%%%%%%%%%

\section{Introduction}

%%%%%%%%%%%%%%%%%%%%%%%%%%%%%%%%%%%%%%%%%%%%%%%%%%%%%%%%%%%%%%%%%%%%%%%


Our objective in this work is to address some fundamental problems in the representation theory of wild double covers of $2$-adic groups. Our main results include: 
\begin{enumerate}[(1)]
	\item\label{IntroListKStructure} basic structure theory of the hyperspecial maximal compact subgroups of the wild cover;
	\item\label{IntroListITypes} a construction of genuine irreducible Iwahori types that capture the covering group analog of the Bernstein block of unramified principal series;
        \item\label{IntroListIMPres} an Iwahori-Matsumoto presentation for the Hecke algebra associated to the Iwahori types, which implies a Shimura correspondence with an explicit linear group;
        \item\label{IntroListKType} an extension of the Iwahori type to a hyperspecial maximal compact subgroup, which induces a ``finite Shimura correspondence''.
\end{enumerate}
Now we describe the results more precisely. In what follows, for a group $H$, we write $\one=\one_{H}$ for the trivial representation of $H$.

Let $G$ be a reductive group over a non-archimedean local field $F$ of characteristic 0. Let $T\subseteq G$ be a maximal split torus, and $I\subseteq G$ an Iwahori subgroup intersecting $T$ nontrivially. The Iwahori-Hecke algebra $\Hcal(G, I)=C_0^\infty(I\backslash G\slash I)$ plays an important role in the representation theory of $G$. In particular, the Borel-Casselman theorem \cite[Theorem 4.10]{B76} implies that the category of $\Hcal(G, I)$-modules is equivalent to the $(T,\one)$ Bernstein block of unramified principal series representations of $G$, and no other Bernstein block admits Iwahori fixed vectors. In other words, the pair $(I, \one)$ is a type of $G$ in the sense of Bushnell-Kutzko \cite[(4.1) Definition]{BK98} and this type characterizes the $(T,\one)$-block. Furthermore, the Iwahori Hecke algebra $\Hcal(G, I)$ has a very explicit set of generators and relations \cite{IM}, which is known as the Iwahori-Matsumoto presentation.

Assume now that $G$ is a simply-connected Chevalley group over $F$. By the work of Steinberg, Moore, and Matsumoto (see \cite{GGW18} for the history), when $F$ contains $\mu_{n}$, the full group of $n$-th roots of unity, $G$ has an (essentially) unique nontrivial topological central extension $\Gt$, fitting into an exact sequence
\begin{equation}\label{CExtIntro}
    1\longrightarrow\mu_n\longrightarrow\Gt\stackrel{\pr}{\longrightarrow} G\longrightarrow 1.
\end{equation}
For any $H\subseteq G$ we write $\Ht=\pr^{-1}(H)$. A representation of $\Ht$ is called $\epsilon$-genuine if $\mu_{n}$ acts through a faithful character $\epsilon:\mu_{n}\rightarrow \mathbb{C}^{\times}$. We fix this character throughout and use it to identify $\mu_{n}$ with the $n$-th roots of unity in $\mathbb{C}$.

The maximal torus $\Tt\subseteq \Gt$ admits a distinguished genuine representation $\rho$. (See \cite{S04} and \cite[Section 6, Theorems 6.6, 6.8]{GG18}.) We note that such a $\rho$ is not unique, but they are finite in number. When $\rho$ is minimally ramified, the $(\Tt,\rho)$-block of $\Gt$ is a natural analog of the $(T,\one)$-block of $G$. Thus we may ask: is there an analog of the Borel-Casselman theorem for the $(\Tt,\rho)$-block of $\Gt$? (see Section \ref{TtReps} for the precise set of representations we consider.) The pair $(\It, \one)$ does not serve the purpose because $\one_{\It}$ is not genuine, and hence can only detect representations of $\Gt$ that factor through $G$. 

When $\mathrm{GCD}(p,n)=1$, the cover $\Gt$ is called a tame cover. In this case, Savin \cite{S04} established this analog of the Borel-Casselman Theorem when $G$ is simply-laced. The main consequence of the tame assumption for this work is that a hyperspecial maximal compact subgroup $K\subseteq G$ splits the sequence \eqref{CExtIntro}. To highlight the significance of this splitting we make a brief digression. 

Here we are not making any assumption on $\mathrm{GCD}(n,p)$. Suppose the subgroup $H\subset G$ splits into $\Gt$, and let $\s :H\rightarrow \Gt$ be a splitting. Then the map $\mu_{n}\times H\rightarrow \Ht$ defined by $(\zeta,h)\mapsto \zeta\s(h)$ defines a group isomorphism. Thus, there is a transparent correspondence (depending on the splitting) between the representations of $H$ and the $\epsilon$-genuine representations of $\Ht$. Specifically, if $\pi$ is a representation of $H$, then $\epsilon\otimes \pi$ is an $\epsilon$-genuine representation of $\Ht\cong \mu_{n}\times H$.


Now we return to the tame case. Let $\s:K\rightarrow \Gt$ be the unique splitting of $K$. Hence $I\subseteq K$ splits to $\Gt$ by restricting the splitting $\s$. When $\rho$ is unramified with respect to $\mathbf{s}$, a natural candidate for the Iwahori type describing the $(\Tt,\rho)$-block is the pair $(\mu_{n}\times \s(I),\epsilon \otimes \one_{\s(I)})$; the associated Hecke algebra $\Hcal_{\epsilon}(\Gt, \s(I))$ consists of compactly supported smooth functions on $\Gt$ that are $\s(I)$-biinvariant and $\epsilon$-genuine. Using this pair $(\mu_{n}\times \s(I),\epsilon \otimes \one_{\s(I)})$, Savin proves that the $(\Tt,\rho)$-block is equivalent to $\Hcal_{\epsilon}(\Gt, \s(I))$-mod giving an analog of the Borel-Casselman theorem. Savin also shows that $\Hcal_{\epsilon}(\Gt, \s(I))$ admits a Bernstein presentation. This allows Savin to identify $\Hcal_{\epsilon}(\Gt, \s(I))$ with the Iwahori Hecke algebra of an explicit linear group $G^{\prime}$ giving a correspondence between certain representations of $\Gt$ and $G^{\prime}$, called a Shimura correspondence.

When $\mathrm{GCD}(n,p)\neq 1$, we say $\Gt$ is wild. In this case, the Iwahori subgroup $I$ does not split. So, neither does $K$. Thus there is no obvious candidate for a type to describe the $(\Tt,\rho)$-block. A few special cases have been worked out. Loke-Savin \cite{LS10a} resolve the case of the $2$-fold cover of $\SL_{2}$ over $\mathbb{Q}_{2}$. This was generalized by Wood \cite{Wood} to the $2$-fold cover of $\Sp_{2r}$ also over $\mathbb{Q}_{2}$. Takeda-Wood \cite{Takeda_Wood} further extend the work of Wood to include any $2$-adic field. Karasiewicz \cite{Karasiewicz} resolves the case of $2$-fold covers of almost simple simply-laced simply-connected Chevalley groups over $\mathbb{Q}_{2}$. We note that \cite{Wood,LS10a,Karasiewicz} all construct a type on a proper subgroup of the Iwahori. In all of these cases the construction of the type is somewhat ad hoc. The work of Takeda-Wood \cite{Takeda_Wood} constructs a type on the Iwahori. Their construction is conceptually satisfying as it uses the $K$-types of the Weil representation. Unfortunately, the special nature of the Weil representation makes the possibility of generalization unclear. %In particular, we note that a naive generalization of these types to the situation where $F/\mathbb{Q}_{2}$ has ramification results in a Hecke algebra that is too large.

In this paper, our main theorem is an analog of the Borel-Casselman theorem for the nontrivial $2$-fold cover of an almost simple simply-laced simply-connected Chevalley group over an \textit{arbitrary} $2$-adic field $F$. This generalizes the work of Karasiewicz \cite{Karasiewicz} to an arbitrary $2$-adic field. The main novelty of our work is how we address the technical issues arising from the ramification of $F/\mathbb{Q}_{2}$. We will say more about this after stating our results.

To state our main theorem, we must first describe some of our other results on the structure and representation theory of the compact groups $\Tt_{0}:=\Tt\cap \Kt$ and $\Kt$. These results are interesting in their own right. From now on, unless otherwise stated, we assume that $F$ is a $2$-adic field and $G$ is the $F$-points of an almost simple simply-connected simply-laced Chevalley group of rank at least $2$, and $\Gt$ is the unique nontrivial topological $2$-fold cover of $G$. So, $\Gt$ is a wild cover.

We begin with $\Tt_{0}$. Here our main result is an explicit construction of a family of Weyl invariant genuine irreducible $\Tt_{0}$-representations (Theorem \ref{T:PseudoT0}). These $\Tt_{0}$-representations are called pseudo-spherical representations. They are not unique, but are finite in number, and our construction produces all of them. We also prove that a pseudospherical representation $\tau$ is a type for a genuine distinguished $\Tt$-representation $\rho$ (Proposition \ref{TTypes}). Our explicit construction is arithmetic in nature and requires special properties of the $2$-adic quadratic Hilbert symbol.

Now we turn to the structure theory of $\Kt$. For any smooth irreducible genuine $\Kt$-representation $\pi$, the subgroup $\pr(\ker\pi)$ must split into $\Kt$ and the image of this splitting (which is just $\ker\pi$) must be normal in $\Kt$. Our main theorem on the structure of $\Kt$ identifies the maximal principal congruence subgroup of $K$ that splits into $\Kt$ with normal image.

\begin{Thm}[Propositions \ref{Gamma12eSplitting}, \ref{P:SplitNormIm}]
    Let $e$ be the ramification index of $F/\mathbb{Q}_{2}$. Let $\Gamma(k)\subset K$ be the $k$-th principal congruence subgroup of $K$. Then
    \begin{enumerate}
        \item There exists a splitting $\s_{\Gamma(2e)}:\Gamma(2e)\rightarrow \Gt$ such that its image $\s_{\Gamma(2e)}(\Gamma(2e))$ is normal in $\Kt$;
        \item if $\Gamma(2e-1)$ splits into $\Kt$, then the image of this splitting is not normal in $\Kt$, nor $\It$.
    \end{enumerate}
\end{Thm}
This result is also arithmetic in nature and requires special properties of the $2$-adic quadratic Hilbert symbol. We make a few more comments on this theorem.

The notion of depth introduced by Moy-Prasad \cite{MP94,MP96} has been generalized to tame covers by Howard-Weissman \cite{HW09}. The main point is that in a tame topological central extension any pro-$p$ subgroup must split uniquely \cite[Proposition 2.3]{HW09}. The situation is more complicated for wild covers. As the previous theorem demonstrates, a genuine irreducible representation of $\Kt$ cannot have a large kernel; among the principal congruence subgroups, $\Gamma(2e)$ is maximal such that its image under a splitting can be contained in the kernel of a genuine $\Kt$-representation. While we will not define a notion of depth for $\Kt$ or $\Gt$, the genuine irreducible $\Kt$-representations that are trivial on the image of a splitting of $\Gamma(2e)$ can be called the minimal depth representations of $\Kt$. 

With some handle on the structure of $\Kt$ we can turn to representation theory. Ideally we would like to establish some correspondence between representations of $\Kt$ and $K$, as we had in the tame case. Since both $I$ and $K$ do not split into $\Gt$ there is no obvious correspondence between, say, $K$-representations and genuine $\Kt$-representations. However, we are able to construct a correspondence between the depth $0$ representations of $K$ and certain representations of $\Kt$ (Theorem \ref{FiniteShimCor}). This generalizes a result of Savin \cite{S12} for the unramified $F/\mathbb{Q}_{2}$, where the correspondence is called a finite Shimura correspondence. This correspondence is constructed using a very special representation of $\Kt$, which we now describe.

Let $\tau$ be a pseudo-spherical representation of $\Tt_{0}$. In fact, $\tau$ remains irreducible when restricted to $\Tt_{1}=\Tt\cap \Gammat(1)$. Let $\sigma=\sigma_{\tau}$ be the $\Gammat(1)$-representation constructed using the parahoric induction functor of Dat \cite[Section 2.6]{D09}. From general properties of parahoric induction the representation $\sigma$ is an irreducible genuine $\Gammat(1)$-representation that is trivial on $\s_{\Gamma(2e)}(\Gamma(2e))$. A more in-depth study of $\sigma$ reveals that this representation actually extends to $\Kt$.

\begin{Thm}[Theorem \ref{KRep}]\label{P:sigma_extends_to_K_Intro}
    The irreducible genuine $\Gammat(1)$-representation $\sigma$ constructed by parahoric induction from a pseudo-spherical representation of $\Tt_{1}$ extends uniquely to a representation of $\Kt$. In particular, the extension of $\sigma$ is a minimal depth representation of $\Kt$.
\end{Thm}

We also write $\sigma$ for the extension of $\sigma$ to $\Kt$. The minimal depth $\Kt$-representation $\sigma$ actually gives us substantial control over an interesting family of genuine representations of $\Kt$ via a finite Shimura correspondence with the depth $0$ representations of $K$. To state this precisely we introduce some additional notation. 

Let $\kappa$ be the residue field of $F$, and $G_\kappa$ the $\kappa$-points of $G$. For a totally disconnected topological group $H$ let $\mathcal{M}(H)$ be the category of smooth representations of $H$. Let $\mathcal{M}_{(\Gammat(1),\sigma)}(\Kt)$ be the full subcategory of $\mathcal{M}(\Kt)$ consisting of smooth representations of $\Kt$ such that their restriction to $\Gammat(1)$ is $\sigma$-isotypic. Note that any representation of $G_{\kappa}$ inflates to a non-genuine representation of $\Kt$ via the isomorphism $G_{\kappa}\cong K/\Gamma(1)$ induced by the reduction mod $p$ map. Now we can state the finite Shimura correspondence.

\begin{Thm}[Theorem \ref{FiniteShimCor}]
The functor $\mathcal{M}(G_{\kappa})\rightarrow \mathcal{M}_{(\Gammat(1),\sigma)}(\Kt)$ defined by $\pi\mapsto \pi\otimes\sigma$ defines an equivalence of categories. %\S \ref{SS:finite_Shimura}
\end{Thm}
As mentioned earlier, this theorem generalizes Savin \cite[Theorem 4.1]{S12}, which assumes $F$ is unramified over $\Q_2$.

The finite Shimura correspondence provides a construction of a large family of minimal depth representations of $\Kt$ and may be valuable for the construction of minimal depth supercuspidal representations of $\Gt$. It also behaves well with respect to induction and thus can be used to reduce the study of certain Hecke algebras on $\Kt$ to related Hecke algebras on $K$ (Theorem \ref{T:finite_Hecke_algebra_isom}).

Now we can state our main theorem on the analog of the Borel-Casselman Theorem for $\Gt$.

\begin{Thm}[Theorem \ref{IType}]
Let $\tau$ be a pseudo-spherical $\Tt_{0}$-representation and let $\rho$ be a distinguished genuine $\Tt$-representation containing $\tau$. Let $\sigma=\sigma_{\tau}$ be the $\It$-representation parahorically induced from $\tau$.

The pair $(\It, \sigma)$ is a type in the sense of Bushnell-Kutzko \cite[(4.1) Definition]{BK98}. Furthermore, the category of $\Hcal(G, \It; \sigma)$-modules is equivalent to the $(\Tt, \rho)$ Bernstein block of $\Gt$, and no other Bernstein block contains representations with $(\It,\sigma)$-isotypic vectors.
\end{Thm}

As in the linear case, we expect the Hecke algebra $\Hcal(\Gt, \It; \sigma)$ to admit an Iwahori-Matsumoto (IM) presentation \cite{IM}, which essentially describes the algebra as a deformation of the group algebra of an extended affine Weyl group. We mostly prove this, but fall short of the complete IM presentation for certain Cartan types where technical obstructions appear. We briefly indicate the extent of our results here, and refer the reader to the body of the paper for the precise statements of our results. 

First we note that in all Cartan types we can determine a weak IM presentation for $\Hcal(\Gt, \It; \sigma)$. This includes identifying an unnormalized basis for $\Hcal(\Gt, \It; \sigma)$ consisting of invertible elements, identifying a set of generators for the algebra, braid relations and quadratic relations up to nonzero scalar for these generators, and a complete IM presentation for elements supported in $\Kt$ (so in particular the basis elements supported in $\Kt$ can be normalized properly). 

Now we will describe our results for the specific Cartan types. We can obtain the complete IM presentation when $G$ has Cartan type $A_r,D_{2r+1},E_{6},E_{7}$ (Theorem \ref{IMPres}). In type $D_{2r}$ we obtain an IM presentation only up to a $2$-cocycle on the length $0$ elements. In type $E_{8}$ we will not be able to obtain the proper normalization of the one generator corresponding to the simple affine reflection.

We note that when $F/\mathbb{Q}_{2}$ is unramified the techniques of Karasiewicz \cite{Karasiewicz} can be adapted to the present setting to derive the complete IM presentation for all Cartan types. We have opted not to pursue this because we believe that the technical obstructions in types $D_{2r}$ and $E_{8}$ for any $2$-adic field $F$ will naturally be resolved as we refine and expand the techniques used in this paper. For example, the difficulty with $E_{8}$ stems from the fact that its affine Dynkin diagram only has a single special vertex, while all of the other types we consider have at least two. We expect this to be resolved as follows. It seems plausible that one can extend the finite Shimura correspondence to other points in the Bruhat-Tits building of $G$. We expect such a finite Shimura correspondence to encode an isomorphism of Hecke algebras, as we see in this paper for the finite Shimura correspondence at a special vertex (Theorem \ref{T:finite_Hecke_algebra_isom}). This isomorphism of Hecke algebras when applied to a non-special vertex in type $E_{8}$ should determine the proper normalization for the last generator completing the IM presentation in type $E_{8}$.

One of the consequences of the IM presentation is that there is an explicit linear Chevalley group $G^{\prime}$ and Iwahori subgroup $I^{\prime}\subseteq G^{\prime}$ such that $\Hcal(G', I')\cong\Hcal(\Gt, \It; \sigma)$. This Hecke algebra isomorphism gives a Shimura correspondence of representations.

\begin{Thm}[Corollary \ref{C:local_Shimura_Correspondence}]
    For $G$ of Cartan type $A_{r},D_{2r+1},E_{6},E_{7}$, there is an equivalence of categories between the $(I',\one)$-block of $G'$ and the $(\It, \sigma)$-block of $\Gt$.
\end{Thm}
We note that this result will also hold for the other Cartan types, once the the remaining gaps in the IM presentation are filled. 

Another consequence of the IM presentation is that the subalgebra $\Hcal(\Gt,\Kt,\sigma)\subset \Hcal(\Gt,\It,\sigma)$ satisfies a Satake isomorphism (Corollary \ref{C:Satake}). 

%At this point we mention the recent work of Heiermann-Wu \cite{HW25}. They addresses similar questions about Hecke algebras on covering groups, e.g. presentations and relation to Bernstein blocks. In \cite{HW25}, Heiermann-Wu study the structure of Hecke algebras of certain coverings of classical groups via endomorphism algebras of projective generators of Bernstien blocks. They do not make any assumptions on the residual characteristic, but they do assume that the Levi subgroups decompose \cite[Definition 5.1]{HW25}. Our work is disjoint from the work of Heiermann-Wu \cite{HW25}. In a $2$-fold cover, if a split torus decomposes in the sense of Heiermann-We, then it must be abelian; in our setting, all maximal tori are nonabelian.

\quad

Next we highlight several of the ideas and tools that we develop to overcome the technical issues raised by the ramification of $F/\mathbb{Q}_{2}$. The first technical issue is that the structure of the $2$-fold cover $\Gt$ over $F$ is closely connected with the quadratic Hilbert symbol of $F$. When $F$ has odd residual characteristic the quadratic Hilbert symbol is well understood and even obeys a relatively simple formula. 
%(e.g. see \cite[(3.4) Proposition]{Neukirch}.) 
When $F$ has even residual characteristic the quadratic Hilbert symbol is more complicated and its properties seem to be less well-known; the complexity is directly connected to the ramification index. So, to get a handle on the structure of $\Gt$ and its subgroups we must first establish some general properties of the quadratic Hilbert symbol for an arbitrary $2$-adic field. The most important result we prove is an explicit description of the radical of the Hilbert symbol restricted to $\Ocal^{\times}$. We expect that this and the other results we prove about the Hilbert symbol are already known, but we have not been able to locate them in the literature, so we record the proofs in Appendix \ref{S:Appendix_Hilbert_symbol}.

These results on the Hilbert symbol are used repeatedly throughout our paper and are part of the foundation of all of our results. To highlight just two, the construction of the pseudo-spherical representations of $\Tt_{0}$ and the structure theory of $\Kt$ depend crucially on the the results we prove about the Hilbert symbol.

With a pseudo-spherical $\Tt_{0}$-representation $\tau$ that is a type of the $\Tt$-representation $\rho$, we can begin to try to construct a type for for the Bernstein block of $\Gt$ attached to $(\Tt,\rho)$. Here we want to highlight two points. 

The first point is about Bushnell-Kutzko's notion of $G$-covers, which essentially suggests a candidate type for $G$ given a type $\tau$ for $\Tt$. We encountered some difficulty when we applied this theory directly to $\tau$. The theory suggests a candidate type $(\It_{2e},\sigma_{2e})$, where $\It_{2e}\subset \It$ is an open compact subgroup and $\sigma_{2e}$ an irreducible representation of $\It_{2e}$. The trouble is that $\It_{2e}$ is quite small compared to $\It$, even when $e=1$. This makes questions about the support of the associated Hecke algebra difficult to investigate. This seems somewhat similar to the situation in Roche \cite{R98}. In this work Roche investigates types for ramified principal series, so in particular is working in arbitrary depth, using Bushnell-Kutkzo's theory of $G$-covers. Roche's candidate types are on small open compact subgroups. He exerts considerable effort to determine the support of the Hecke algebras and must impose constraints on the residual characteristic, which includes $p=2$ in all Cartan types (though he mentions a way to remove these restrictions in type $A_r$). (See \cite[pg. 370-379]{R98}.)

In our situation we are seeing high-depth phenomenon from the ramification of $F/\mathbb{Q}_{2}$, but since we are exclusively working with $2$-adic fields we must deal with residual characteristic $2$. It is not clear if Roche's techniques adapt to our setting, so we took a different path. To minimize the support difficulties we set out to construct an Iwahori type. This is the approach in Takeda-Wood \cite{Takeda_Wood}, but there the Iwahori type is constructed using the Weil representation. In our situation we construct a candidate  Iwahori type $(\It,\sigma)$ using a pseudo-spherical $\Tt_{0}$-representation $\tau$ and the parahoric induction of Dat \cite{D09}. Now we might try to prove that $\sigma$ is a type using the theory of $G$-covers. Unfortunately the $\It$-representation $\sigma$ does not fit into this theory. However, we are able to adapt the theory to our setting. This seems to suggest that the theory of $G$-covers can be extended by incorporating parahoric induction. (See Lemmas \ref{LemSurj}, \ref{LemInj}.)

The second point is that the statement of our results in terms of Iwahori types of $\Gt$ is misleading, from the standpoint of our arguments. Our arguments crucially use two different models for the Hecke algebra $\Hcal(\Gt, \It; \sigma)$; the other involving a smaller open compact subgroup $\It_{2}\subset \It$ and an irreducible representation $\sigma_{2}$, also constructed from $\tau$ via parahoric induction. In fact, $\sigma\cong \Ind_{\It_{2}}^{\It}\sigma_{2}$ (Proposition \ref{IndSig2}) and so $\Hcal(\Gt, \It; \sigma)\cong \Hcal(\Gt, \It_{2}; \sigma_{2})$, from general theory. Roughly, the advantage of the $\It$-model appears when bounding the support of the Hecke algebra, this allows us to avoid the tedious support calculations in Karasiewicz \cite{Karasiewicz}; the advantage of the $\It_{2}$-model is that it is more compatible with the underlying affine Weyl group structure of $\Gt$ than $\It$ is. We will briefly elaborate on this affine Weyl group comment.

Let $T\subset G$ be a maximal split torus with cocharacter lattice $Y$. Let $\Phi$ be the root system of $(G,T)$. Let $\mathfrak{A}\subset \mathbb{R}\otimes Y$ be an alcove containing $0$ in its closure. The stabilizer of $\mathfrak{A}$ in $G$ is an Iwahori subgroup, say $I$. Now we turn to $\Gt$. The root datum for the Langlands dual group of $\Gt$ suggests that the relevant root system for $\Gt$ is $\frac{1}{2}\Phi$. Thus in place of the alcove $\mathfrak{A}$ we should consider $2\mathfrak{A}$, and the stabilizer of $2\mathfrak{A}$ in $G$ is the group we have been calling $I_{2}$.

\quad


To close this introduction we will summarize the content of each section. In \S \ref{S:Prelim} we set up notation and recall some well known results. We note that \S \ref{SS:HilbSymb} includes the fundamental properties of the Hilbert symbol that we use repeatedly. We expect these results are known, but we include the proofs in Appendix \ref{S:Appendix_Hilbert_symbol}.

In \S \ref{S:2Cover} we recall the presentation of $\Gt$ and prove some basic results on the modified cocharacter lattice $\Yt\subset Y$ and the action of the Weyl group on it.

\S \ref{SecTORep} - \ref{S:IwahoriType} are the core of our work. In \S \ref{SecTORep} we construct the pseudo-spherical representations of $\Tt_{0}$, the maximal compact subgroup of $\Tt$. While $\Tt_{0}$ is a $2$-step nilpotent group with a Stone-von Neumann Theorem, for use in the rest of the paper we need a fairly explicit construction of these representations. To achieve this we describe a factorization of $\Tt_{0}$, heavily using the properties of the Hilbert symbol from \S \ref{SS:HilbSymb}. We note that we construct all of the pseudo-spherical representations of $\Tt_{0}$, show that they are Weyl group invariant (Theorem \ref{T:PseudoT0}), and describe how they branch when restricted to certain subtori of rank one less (Proposition \ref{PseudoSphericalBranch}).

In \S \ref{S:PSKType} we investigate the structure and representation theory of a hyperspecial maximal compact subgroup $\Kt\subset\Gt$ and other open compact subgroups. In \S \ref{SSecSplittings} we examine the structure of $\Kt$. Our main theorem identifies the maximal principal congruence subgroup of $K$ that splits into $\Gt$ with normal image in $\Kt$ (Propositions \ref{Gamma12eSplitting}, \ref{P:SplitNormIm}). In \S \ref{SS:ParaInd} we review the parahoric induction and restriction functors of Dat \cite{D09} and their extension by Crisp-Meir-Onn \cite{CMO19}, and records basic properties of these functors, some from \cite{D09,CMO19} and some new (e.g. Proposition \ref{ResGamma1}). In \S \ref{IwahoriType} we construct a candidate Iwahori type $\sigma$ and several related representations on finite index subgroups of $\It$. The main result in this subsection is Theorem \ref{IRep}, where we show how these representations relate to one another and establish Weyl-invariance, where applicable.

In \S \ref{S:KType} we prove that the $\It$-representation $\sigma$ extends to $\Kt$ (Theorem \ref{KRep}). This is a crucial result. We use it later to construct the finite Shimura correspondence (Theorem \ref{FiniteShimCor}), which itself essentially determines the structure of the subalgebra of $\Hcal(\Gt,\It,\sigma)$ supported on $\Kt$.

In \S \ref{SigProps} we introduce $I_{2}$ and $\sigma_{2}$, mentioned above, and collect various properties of $\sigma$ and $\sigma_{2}$. These properties play a role in our investigation of the Hecke algebras in \S \ref{S:Hecke}.

In \S \ref{SS:finite_Shimura} we prove the finite Shimura correspondence in Theorem \ref{FiniteShimCor}. This result is interesting in its own right, but it also figures into our study of $\Hcal(\Gt,\It,\sigma)$. We return to describe other important properties of this correspondence in \S \ref{FinShimHecke}, after introducing the Hecke algebra.

In \S \ref{SS:WhitInvar} we collect a couple simple corollaries on Whittaker invariants of the representations of $\Kt$ that arise in the finite Shimura correspondence. These results are not used elsewhere in the paper, but they should be important for the study of the Gelfand-Graev representation of $\Gt$. 

In \S \ref{S:Hecke} we begin our study of the Hecke algebra $\Hcal=\Hcal(\Gt,\It,\sigma)$. After reviewing basic properties of Hecke algebras in \S \ref{SS:HeckeReview}, we return in \S \ref{FinShimHecke} to the finite Shimura correspondence. We prove that this correspondence preserves induction in a certain sense and use this to connect the finite Shimura correspondence to Hecke algebras. Consequently, we show that a family of Hecke algebras on $\Kt$ are isomorphic to Hecke algebras studied by Howlett-Lehrer \cite{HL80}.

In \S \ref{HandH2} we prove that the Hecke algebra $\Hcal$ constructed from $\sigma$ and the Hecke algebra $\Hcal_{2}=\Hcal(\Gt,\It_{2},\sigma_{2})$ constructed from $\sigma_{2}$ are in fact isomorphic (Theorem \ref{HIsoH2}). This proves to be quite important since each model offers complementary advantages. 

Over the next six subsections we study the structure of $\Hcal\cong \Hcal_{2}$. In \S \ref{SS:MultOne} we prove in Theorem \ref{HMult1} that every Iwahori double coset supports a space of functions of dimension at most one in $\Hcal$. For this theorem the $\Hcal$-model is quite convenient.

We prove in \S \ref{SuppAfSimpRel} and \ref{SuppLen0}, using the $\Hcal_{2}$-model, that the Hecke algebra supports functions on double cosets represented by certain elements in the extended affine Weyl group. Specifically the simple reflections (Proposition \ref{affSupp}) and the length $0$ elements (Proposition \ref{Len0Supp}).

In \S \ref{SS:QuadRel} we prove the quadratic relations in the Hecke algebra. For the simple reflections represented by elements of $\Kt$ we can prove the exact quadratic relation (Proposition \ref{QuadBraidSimpLinReflect}) using the finite Shimura correspondence, and so in particular, this uses the $\Hcal$-model. For the simple affine reflection we can use the $\Hcal_{2}$-model to prove a weak quadratic relation (Proposition \ref{QuadSimpAff}). Outside of type $E_{8}$ we use an automorphism to upgrade this weak quadratic relation to the expected quadratic relation. We also use the $\Hcal_{2}$-model to prove that the length $0$ elements satisfy a quadratic relation.

Next we turn to the Braid relations in \S \ref{SS:BraidRel}. Using an auxiliary linear Hecke algebra (introduced in \S \ref{SS:LinearHecke}) we can prove that $\Hcal_{2}$ has a Braid relation up to a $2$-cocycle (Proposition \ref{Prop:GenBraid}). The rest of this section is devoted to showing when this $2$-cocycle is trivial. We prove that the $2$-cocycle is trivial in types $A_{r},D_{2r+1},E_{6},E_{7}$, and it is expected to be trivial in general. 

%We will say a bit more about $D_{2r}$ and $E_{8}$ in the next paragraph.

Finally in \S \ref{SS:IMPres} we combine our results from the previous subsections to show that $\Hcal\cong\Hcal_{2}$ has an IM presentation (Theorem \ref{IMPres}). This is complete in types $A_{r},D_{2r+1},E_{6},E_{7}$, but it is expected for all Caratn types. In type $D_{2r}$ the last step is to prove that the $2$-cocycle on the length $0$ elements is trivial; in type $E_{8}$ the last step is to prove that the simple affine reflection satisfies the expected quadratic relation. We end \S \ref{SS:IMPres} with a few comments on $D_{2r}$, and $E_{8}$.

In \S \ref{S:IwahoriType} we prove our main Theorem \ref{TypesThm} that the pair $(\It,\sigma)$ is a type in the sense of Bushnell-Kutzko \cite[(4.1) Definition]{BK98}, and this type captures the Bernstein block associated to $(\Tt,\rho)$. This result has no restrictions on the (irreducible) Cartan type. Given our previous results on the structure of the Hecke algebra $\Hcal$ this argument essentially follows Bushnell-Kutzko \cite{BK98}. However, as mentioned above, the pair $(\It,\sigma)$ does not fit into Bushnell-Kutzko's theory of $G$-covers. So, in this section we also show how the notion of $G$-cover extends to our setting.

This section ends with Corollary \ref{C:local_Shimura_Correspondence}. This states that for the Cartan types for which we know that $\Hcal$ admits a complete IM presentation there is a linear group $G^{\prime}$ with an Iwahori subgroup $I^{\prime}\subset G^{\prime}$ such that $\Hcal\cong \Hcal(G^{\prime},I^{\prime})$. This isomorphism of Hecke algebras implies that the $(\Tt,\rho)$-Bernstein block of $\Gt$ is equivalent to the block of unramified principal series of $G^{\prime}$, giving a Shimura correspondence of representations.

We end with Appendix \ref{S:Appendix_Hilbert_symbol}. It contains proofs of the important properties of the Hilbert symbol that we have been unable to locate in the literature. 



\quad

\begin{center}{\bf Notation and assumptions}\end{center}

Throughout the paper, we let $F$ be a finite extension of $\Q_2$, with ring of integers $\Ocal$ and a chosen uniformizer $\varpi$, and let $\kappa=\Ocal\slash\varpi\Ocal$ be the residue field. Let $e, f\in\Z_{\geq 1}$ be the ramification index and the inertia degree of $F\slash\Q_2$, respectively, so that
\[
[F:\Q_2]=ef,\quad 2\Ocal=\varpi^e\Ocal\qand \kappa=\Ocal\slash\varpi\Ocal\cong \F_{q},
\]
where $q=2^f$.

Let $G$ be any group and $H\subseteq G$ a subgroup. Let $x\in G$. We set
\[
^xH=xHx^{-1}.
\]
For a representation $\pi$ of $H$, we define
\[
^x\pi
\]
to be the representation of ${^xH}$ defined by
\[
^x\pi(h)=\pi(x^{-1}hx),\quad (h\in {^xH}).
\]
Note that $x^{-1}hx\in H$ because $h\in{^xH}=xHx^{-1}$, and hence this is well-defined. More generally, for another group $H'$ and a group isomorphism $\phi:H\to H'$, we define 
\[
^\phi\pi
\]
to be the representation of $H'$ defined by
\begin{equation}\label{E:definition_phi_pi}
^\phi\pi(h')=\pi(\phi^{-1}(h')),\quad(h'\in H').
\end{equation}
We let
\[
\Int(x):H\longrightarrow xHx^{-1},\quad h\mapsto xhx^{-1}.
\]
Then $^x\pi={^{\Int(x)}}\pi$. We denote by $\pi^H$ the space of $H$-invariant vectors in $\pi$.


\begin{comment}
For $x, y\in G$, one has
\[
^{xy}\pi={^x({^y\pi})}.
\]
But it should be mentioned that the left-hand side is defined on $H\cap{^{xy}H}$, whereas on the right-hand side ${^y\pi}$ is defined on $H\cap{^yH}$, and hence {\it a priori} ${^x({^y\pi})}$ is defined on 
\[
(H\cap{^yH})\cap{^x(H\cap{^yH})}=H\cap{^xH}\cap{^yH}\cap{^{xy}H}\subseteq H\cap{^{xy}H}.
\]
Hence the left-hand side of the identity should be interpreted as $^{xy}\pi$ restricted to this possibly smaller group.
\end{comment}


\quad


\begin{center}{\bf Acknowledgments}\end{center}
We would like to thank Gus Lehrer for pointing us to the work of Uri Onn, from which we learned about Dat \cite{D09} and Crisp-Meir-Onn \cite{CMO19}. We also thank Wee Teck Gan for numerous discussions throughout the preparation of this work. The second author was partially supported by the Sumitomo Foundation Fiscal 2024 Grant for Basic Science Research Projects J230603025, and  JSPS KAKENHI grant number 24K06648. Also part of the research was done while both authors were attending the conference ``The 25th Autumn Workshop on Number Theory" at Hokkaido University in fall 2024, and while the second author was visiting the National University of Singapore in Septembers of 2024 and 2025. We would like to thank their hospitality.



\quad\\



%%%%%%%%%%%%%%%%%%%%%%%%%%%%%%%%%%%%%%%%%%%%%%%%%%%%%%%%%%%%%%%%%%%%%%%

\section{Preliminaries}\label{S:Prelim}

%%%%%%%%%%%%%%%%%%%%%%%%%%%%%%%%%%%%%%%%%%%%%%%%%%%%%%%%%%%%%%%%%%%%%%%

In this section, we first establish basic properties of the quadratic Hilbert symbol over our $2$-adic field $F$. After that, we will review the basic theory of Chevalley groups.

\subsection{$2$-adic Hilbert symbols}\label{SS:HilbSymb}
Let
\[
(-,-):F^\times\times F^\times\longrightarrow \{\pm 1\}
\]
be the quadratic Hilbert symbol of the $2$-adic field $F$. (For the general properties of the Hilbert symbol we refer the reader to \cite[Chapter V \S 3]{Neukirch}; for the more specialized results used in this paper see Appendix \ref{S:Appendix_Hilbert_symbol}.) Throughout, we often use
\[
(x, -x)=(x, 1-x)=1\qand (x, x)=(x, -1)
\]
for all $x\in F^\times$ such that the symbols are defined.

To investigate the structure of a double cover of a Chevalley group over a $2$-adic field we need to analyze the Hilbert symbol restricted to $\Ocal^\times\times\Ocal^\times$. Define
\begin{equation}\label{E:integral_radical_def}
R:=\{a\in\Ocal^\times\st \text{$(a, -)$ is trivial on $\Ocal^\times$}\},
\end{equation}
which we call the {\it integral radical} of the Hilbert symbol. 

In the next proposition we collect some properties of the quadratic $2$-adic Hilbert symbol.
\begin{Prop}\label{P:property_Hilber_symbol_not_in_Appendix}$ $
	\begin{enumerate}
		\item $(1+\varpi^{2e}\Ocal)\Ocal^{\times 2}= R$ and $1+\varpi^{2e+1}\Ocal\subseteq \Ocal^{\times 2}$. \label{ORad}
		\item $(1+\varpi^{2e}\Ocal)\slash((1+\varpi^{2e}\Ocal)\cap\Ocal^{\times 2})\cong R\slash\Ocal^{\times 2}\cong\Z\slash 2\Z$.
		\item $[\Ocal^{\times}:R]=2^{ef}$.
	\end{enumerate}
\end{Prop}
\begin{proof}
See Appendix \ref{S:Appendix_Hilbert_symbol}.
\end{proof}

The following corollary will be crucially used in this paper.
\begin{Cor}\label{C:Hilbert_symbol_unique_element_integral_radical}
Let $u\in (1+\varpi^{2e}\Ocal)\smallsetminus\Ocal^{\times 2}$, which exists by the above proposition. Then for all odd $k\in\Z$ and $a\in\Ocal^\times$, we have
\[
(u, a\varpi^k)=-1.
\]
\end{Cor}
\begin{proof}
Since $u\in R$, we know $(a, u)=1$. So $(u, a\varpi^k)=(u, \varpi)$ because $k$ is odd. Since every element in $F^\times$ is written as $a\varpi^n$ for some $a\in \Ocal^\times$ and $n\in\Z$, if $(u, \varpi)=1$, then we would have $(u, x)=1$ for all $x\in F^\times$, which would imply $u\in \Ocal^{\times 2}$.
\end{proof}


The Hilbert symbol restricted to units descends to a non-degenerate bilinear form
\[
\Ocal^\times\slash R\times\Ocal^\times\slash R\longrightarrow\{\pm 1\}.
\]
Since $\Ocal^{\times 2}\subseteq R$, we can consider $\Ocal^\times\slash R$ as an $\F_2$ vector space with a nondegenerate symmetric bilinear form. (We note that this form may not arise from a quadratic form on $\Ocal^\times\slash R$.) For later purposes, we need to decompose this symmetric bilinear form as follows.

\begin{Prop}\label{P:orthogonal_decomposition_of_Hilbert}
Let $S\subseteq \Ocal^\times\slash R$ be the set of all isotropic vectors, namely
\[
S=\{u\in\Ocal^\times\slash R\st (u, u)=1\}.
\]
Then there exists a subspace $D\subseteq \Ocal^\times\slash R$ with $\dim_{\F_2}D\leq 2$ such that
\[
\Ocal^\times\slash R=D\oplus D^\perp\qand D^\perp\subseteq S.
\]
To be more explicit, we have the following.
\begin{enumerate}[(i)]
\item If $-1\in R$, then $D=0$ and $S=\Ocal^\times\slash R$.\label{H1}
\item If $(-1, -1)=-1$, then $D=\la-1\ra$ and $\la-1\ra^{\perp}=S$. \label{H2}
\item If $(-1, -1)=1$ and $-1\notin R$, then $D$ is diagonalizable and 2-dimensional with $-1\in D$ and $D^\perp\subsetneq S$. (Here, we have $D^\perp\subsetneq S$ because $-1\in S$ but $-1\notin D^\perp$.) \label{H3}
\end{enumerate}
\end{Prop}
\begin{proof}
Assume $-1\in R$. Then for all $u\in\Ocal^\times$, we have $(u, u)=(u, -1)=1$. Hence every vector in $\Ocal^\times\slash R$ is isotropic. 

Next assume $(-1, -1)=-1$. Then the 1-dimensional subspace $\la-1\ra$ is non-degenerate. So we have the orthogonal sum decomposition $\Ocal^\times\slash R=\la-1\ra\oplus \la-1\ra^\perp$. (See \cite[Proposition 1.1.1]{Kitaoka} for the existence of orthogonal sum decomposition.) Since $(u, u)=(u, -1)$ for all $u\in\Ocal^\times$, we know that $u\in S$ if and only if $u\in \la-1\ra^\perp$.

Finally, assume $(-1, -1)=1$ and $-1\notin R$. Since $-1\notin R$, there exists $u_0\in\Ocal^\times$ such that $(-1, u_0)=-1$, so $(u_0, u_0)=-1$. One can then check that $(-u_0, -u_0)=-1$. Since $u_0\neq -u_0$, we have the 2-dimensional subspace $D:=\la u_0, -u_0\ra$ generated by $u_0$ and $-u_0$. (Written multiplicatively, $D=\{\pm 1, \pm u_0\}$.) Certainly this space is non-degenerate and contains $-1$. We then have the orthogonal sum decomposition $\Ocal^\times\slash R=D\oplus D^\perp$. Now, let $u\in D^\perp$. Since $-1\in D$, we have $(u, -1)=1$, so $(u, u)=1$, namely $u\in S$. 
\end{proof}

\begin{Rmk}
All the three cases in Proposition \ref{P:orthogonal_decomposition_of_Hilbert} can occur. Case \eqref{H1} occurs whenever $\sqrt{-1}\in F$; Case \eqref{H2} occurs whenever $[F:\Q_{2}]$ is odd because then $(-1, -1)=(-1, N_{F\slash\Q_2}(-1))_{\Q_2}=(-1, -1)_{\Q_2}=-1$, where $(-, -)_{\Q_2}$ is the Hilbert symbol for $\Q_2$; Case \eqref{H3} occurs whenever $\sqrt{-1}\notin F$ and $[F:\Q_{2}]$ is even, because then $(-1, -1)=(-1, N_{F\slash\Q_2}(-1))_{\Q_2}=1$ and $-1=1-2\in 1+\varpi^e\Ocal^\times$ so that $-1\notin R$.
\end{Rmk}

Let $D$ be as in the above proposition, and let $\{u_1,\dots,u_k\}$ be an orthogonal basis for $D$, where actually $k=0, 1, 2$. Note that we have
\begin{equation*}
(u_i, u_j)=\begin{cases}1&\text{if $i\neq j$};\\
-1&\text{if $i=j$}.\end{cases}
\end{equation*}
Also since $D^\perp\subseteq S$, it admits a symplectic basis (\cite[Theorem 1.2.1]{Kitaoka}); namely there exists a basis $\{e_1,\dots, e_{\ell}, f_1,\dots, f_{\ell}\}$ of $D^\perp$ such that
\[
(e_i, f_i)=\begin{cases}1&\text{if $i\neq j$};\\
-1&\text{if $i=j$},\end{cases}
\qand
(e_i, e_j)=(f_i, f_j)=1\; \text{for all $i, j$}.
\]
We fix these basis vectors once and for all, so that we have
\[
\Ocal^\times\slash R=D\oplus D^\perp=\Span_{\F_2}\{u_1,\dots,u_k\}\oplus\Span_{\F_2}\{e_1,\dots,e_\ell, f_1,\dots,f_\ell\}.
\]


\begin{comment}
\subsection{Root Systems}\label{RootSys}

Let $(\Phi,E)$ be an irreducible root system in a Euclidean vector space $(E, \langle-,-\rangle_{E})$ with roots $\Phi$. Let $Q\subseteq E$ be the root lattice, $Q^{\vee}\subseteq E$ the coroot lattice, $P\subseteq E$ the weight lattice, and $P^{\vee}\subseteq E$ the coweight lattice. 

Let $\Delta$ be a set of simple roots. Then let $\Phi^{+}=\Phi^{+}_{\Delta}$ be the set of positive roots determined by $\Delta$ and $\Phi^{-}=\Phi^{-}_{\Delta}=-\Phi^{+}_{\Delta}$ for the set of negative roots. Let $r=|\Delta|$ be the rank $\Phi$, and we index the simple roots $\Delta=\{\alpha_{1},\ldots,\alpha_{r}\}$ using the Bourbaki labeling \cite[Plate I-IX]{BourLGLA46}. Let $\alpha_{0}$ be the lowest root with respect to $\Delta$.

For $\alpha\in\Phi$ let $\alpha^{\vee}:=\frac{2}{\langle\alpha,\alpha\rangle}\alpha$ be the coroot associated to $\alpha$. Then let $\Phi^{\vee}$ be the set of coroots and $\Delta^{\vee}$ the set of simple coroots.

For $\alpha\in\Phi$ and $k\in\mathbb{Z}$ we define the affine hyper plane $H_{\alpha,k}:=\{v\in E|\langle v,\alpha\rangle_{E}=k\}$. Let $w_{\alpha,k}$ be the affine reflection of $E$ defined by $w_{\alpha,k}(v):=v-(\langle v,\alpha\rangle_{E}-k)\alpha^{\vee}$.

Let $\mathfrak{A}=\mathfrak{A}_{\Phi^{+}}:=\{v\in E|0<\langle\alpha,v\rangle<1\text{ for all }\alpha\in \Phi^{+}\}$. 

Let $W_{\af}=W_{\af,\Phi}$ be the affine Weyl group generated by the affine reflections $w_{\alpha,k}$ for all $\alpha\in\Phi$ and $k\in \mathbb{Z}$. Let $W\subseteq W_{\af}$ be the subgroup fixing $0\in E$. Then $W$ is the Weyl group of $(\Phi,E)$ and is generated by the elements $w_{\alpha}:=w_{\alpha,0}$ for all $\alpha\in\Delta$. The reflections $w_{\alpha_{0},-1},w_{\alpha_{1},0}, \ldots, w_{\alpha_{r},0}$ are the reflections through the walls of the chamber $\mathfrak{A}$ and generate $W_{\af}$.

Let $L=L_{\Phi}\subseteq W_{\af}$ be the subgroup of translations. The group homomorphism $L\rightarrow E$ defined by $\lambda\mapsto \lambda(0)$ identifies $L$ with the coroot lattice $\mathrm{span}_{\mathbb{Z}}(\Phi^{\vee})\subseteq E$. The affine Weyl group can be factored as $W_{\af}\cong W\ltimes L$, where $W$ acts on $L$ through its action on $E$.

We need to simultaneously consider the isomorphic root system $(\frac{1}{2}\Phi,E)$. Let $\widetilde{W}_{\af}:=W_{\af,\frac{1}{2}\Phi}$, $\widetilde{L}:=L_{\frac{1}{2}\Phi}$, and $\widetilde{\mathfrak{A}}:=\mathfrak{A}_{\frac{1}{2}\Phi^{+}}$. Note that $\widetilde{W}_{\af}\subseteq W_{\af}$, $\widetilde{L}=2L$, and $\widetilde{\mathfrak{A}}=2\mathfrak{A}$.

The group $\widetilde{W}_{\af}$ is generated by the reflections $w_{\alpha_{0},-2},w_{\alpha_{1},0}, \ldots, w_{\alpha_{r},0}$. For these reflections we simplify notation and write $w_{\alpha_{0}}=w_{\alpha_{0},-2},w_{\alpha_{1}}=w_{\alpha_{1},0}, \ldots, w_{\alpha_{r}}=w_{\alpha_{r},0}$.

Let $\ell$ be the length function of $\Wt_{\af}$ relative to the simple reflections $w_{\alpha_{0}},w_{\alpha_{1}}, \ldots, w_{\alpha_{r}}$.
\end{comment}


\subsection{Chevalley groups}\label{SS:ChevGroup}

\begin{comment}
If $\mathbb{H}$ is a linear algebraic group defined over a ring $R$ and $S$ is an $R$-algebra let $\mathbb{H}(S)$ be the set of $S$-points.

Attached to the root system $(\Phi,E)$ there is a simply-connected split semisimple linear algebraic group $\mathbb{G}$ defined over $\Ocal$ (in fact over $\mathbb{Z}$). Let $\mathbb{T}\subseteq\mathbb{G}$ be a maximal torus and let $\mathbb{B}$ be a Borel subgroup of $\mathbb{G}$ containing $\mathbb{T}$ with unipotent radical $\mathbb{U}$. Let $\mathbb{B}^{-}\supset \mathbb{U}^{-}$ be the opposite Borel and its unipotent radical.
\end{comment}

In this subsection we review the basic structure theory of Chevalley groups and establish our notation. In particular we recall a presentation of such groups.

Let $G$ be the group of $F$-rational points of an almost simple simply-connected Chevalley group. Let $T\subseteq G$ be a maximal split torus and let $\Phi=\Phi(G,T)$ be the irreducible root system associated with the pair $(G,T)$. Let $W=N_G(T)\slash T$ be the Weyl group. Let $X=X^*(T)$ and $Y=X_*(T)$ be the group of rational characters and cocharacters, respectively, so that $(X, \Phi, Y, \Phi^\vee)$ is the root datum for $(G, T)$, where $\Phi^\vee$ is the set of coroots as usual. Note that we have the bijection
\[
\Phi\xrightarrow{\;\sim\;}\Phi^\vee,\quad \alpha\mapsto \alphac.
\]
Let
\[
\la-,-\ra:X\times Y\longrightarrow\Z
\]
be the canonical pairing. It is $W$-invariant and $\la\alpha, \alphac\ra=2$ for all $\alpha\in\Phi$. It should be mentioned that if $\Phi$ is simply-laced we have
\[
\la \alpha, \betac\ra=\la \beta, \alphac\ra
\]
for all $\alpha, \beta\in\Phi$. Since $G$ is simply-connected, we have $Y=\Z\Phi^{\vee}$. So, the identification of $\Phi$ with $\Phi^\vee$ gives rise to a $W$-invariant inner product on $X$, which we denote by the same symbol
\[
\la-,-\ra:X\times X\longrightarrow\Z.
\]
Although the symbol $\la-,-\ra$ has two meanings, the correct one will be clear from context.


%After extension of scalars, $X\otimes\mathbb{R}$ is the Euclidean space of the roots system $\Phi$, as in \S \ref{RootSys}.

Let $B=B^{+}\subseteq G$ be a Borel subgroup containing $T$ with unipotent radical $U=U^{+}$. Let $B^{-}\supset U^{-}$ be the opposite Borel and its unipotent radical. The Borel subgroup $B$ determines a set of simple roots
\[
\Delta=\{\alpha_1,\dots,\alpha_r\}\subseteq \Phi,
\]
which determines positive roots $\Phi^+\subseteq\Phi$. We order the simple roots by Bourbaki's convention. Let
\[
\alpha_0=\text{lowest root with respect to $\Delta$}.
\]

%For each $\alpha\in \Phi^{\pm}$ there is a subgroup $U_{\alpha}\subseteq U^{\pm}$ isomorphic to $F$. Let $x_{\alpha}:F\rightarrow U_{\alpha}$ be this isomorphism. The field $F$ is filtered by subgroups of the form $\varpi^{j}\Ocal$, where $j\in \mathbb{Z}$. Define $U_{\alpha,j}$ to be the image of $x_{\alpha}$ restricted to $\varpi^{j}\Ocal$.

Now we recall a presentation for $G$. The group $G$ is generated by the symbols of the form $x_{\alpha}(t)$, namely
\[
G=\la \xa(t)\st \alpha\in\Phi, t\in F\ra.
\] 
For $t\in F^\times$, define
\begin{align*}
\wa(t)&:=\xa(t)x_{-\alpha}(-t^{-1})\xa(t);\\
\ha(t)&:=w_{\alpha}(t)\wa(-1).
\end{align*}
Then the following is a complete set of relations for the generators $\xa(t)$, when $rk(G)\geq 2$:
\begin{align}
\tag{R1}\xa(t)\xa(u)&=\xa(t+u);\label{R1}\\
\tag{R2}[\xa(t), \xb(u)]&=\prod_{i\alpha+j\beta\in\Phi}x_{i\alpha+j\beta}(c_{ij}t^iu^j),\quad(\alpha+\beta\neq 0);\label{R2}\\
\tag{R3}\ha(t)\ha(u)&=\ha(tu).\label{R3}
\end{align}
where the $c_{ij}\in \{\pm1\}$ are independent of $t$ and $u$ but dependent on $\alpha$ and $\beta$ and the ordering of the roots. Note that, if $i\alpha+j\beta\notin\Phi$, then we interpret $x_{i\alpha+j\beta}(c_{ij}t^iu^j)$ as 1. If $\Phi$ is simply-laced, \eqref{R2} simplifies to 
\begin{equation}
\tag{R2$^{\prime}$}[\xa(t), \xb(u)]=\begin{cases}1&\text{if $\alpha+\beta\notin\Phi$};\\
x_{\alpha+\beta}(ctu)&\text{if $\alpha+\beta\in\Phi$},
\end{cases}
\quad (\alpha+\beta\neq 0),
\end{equation}
where $c=c(\alpha, \beta)\in\{\pm 1\}$. Note also that $\wa(t)^{-1}=\wa(-t)$. 


The following relations follow from \eqref{R1}, \eqref{R2}, \eqref{R3}:
\begin{align*}
\tag{R4}\ha(t)\hb(u)&=\hb(u)\ha(t);\label{R4}\\
\tag{R5}\wa(t)\hb(u)\wa(-t)&=h_{w_\alpha(\beta)}(u);\label{R5}.\\
\tag{R6}\wa(t)\wb(u)\wa(-t)&=w_{w_\alpha(\beta)}(ct^{-\la\alpha, \beta\ra}u)\label{R6},
\end{align*}
where $c=c(\alpha, \beta)\in\{\pm 1\}$.


\subsection{Some notation for subgroups}
Let us establish some notation for subgroups of $G$.

The maximal torus $T$ is generated by the $\ha(t)$'s, namely
\[
T=\la\ha(t)\st t\in F^\times, \alpha\in\Phi\ra.
\]
For $\alpha\in \Phi$ and a subgroup $A\subseteq F^{\times}$, let 
\[
T_{\alpha, A}:=\{h_{\alpha}(a)\st a\in A\},
\]
and
\[
T_A:=\la h_{\alpha}(a)\st a\in A, \alpha\in\Phi\ra,
\]
both of which are subgroups of T. In particular, we often use
\begin{equation}\label{E:def_of_T_R}
T_R:=\la h_{\alpha}(a)\st a\in R, \alpha\in\Phi\ra,
\end{equation}
where $R$ is the integral radical defined in \eqref{E:integral_radical_def}. We also set
\begin{equation}\label{E:def_of_T_0}
T_0=T_{\Ocal^{\times}}=\la h_{\alpha}(a)\st a\in \Ocal^{\times}, \alpha\in\Phi\ra.
\end{equation}
We write
\[
T_{\alpha, k}:=T_{\alpha, 1+\varpi^k\Ocal}=\{h_{\alpha}(a)\st a\in 1+\varpi^k\Ocal\}
\]
and
\begin{equation}\label{E:def_of_T_k}
T_{k}:=T_{1+\varpi^k\Ocal}=\la h_{\alpha}(a)\st a\in 1+\varpi^k\Ocal, \alpha\in\Phi\ra.
\end{equation}
Note that every element $t\in T_{A}$ is uniquely written as $t=\prod_{1\leq i\leq r}h_{\alpha_{i}}(a_{i})$ for some $a_i\in A$ by using only the simple roots. 


Next, for any subgroup $A\subseteq F$ and any $\alpha\in\Phi$, we set 
\[
U_{\alpha}(A)=\la x_{\alpha}(a)\st a\in A\ra
\] 
or more generally for any subset $S\subseteq \Phi$, we set
\[
U_{S}(A)=\la x_{\alpha}(a)\st a\in A, \alpha\in S\ra.
\] 
In particular, we set
\[
U^{+}(A)=U_{\Phi^{+}}(A)\qand U^{-}(A)=U_{\Phi^{-}}(A).
\]
Also we set
\begin{equation}\label{E:notation_U_alpha_j}
U_{\alpha, j}:=U_{\alpha}(\varpi^j\Ocal),\quad U_j^{\pm}:=U_{\Phi^{\pm}}(\varpi^j\Ocal)\qand U^{\pm}=U^{\pm}(F),
\end{equation}
where $j\in\Z$. We often omit the superscript $^+$, so in particular
\begin{equation}\label{E:notation_U_j}
U_j=U_j^+=U_{\Phi^{+}}(\varpi^j\Ocal).
\end{equation}

Set
\[
K:=\la x_{\alpha}(a)\st a\in\Ocal, \alpha\in\Phi\ra,
\]
namely the group generated by the subgroups $U_{\alpha,0}$ for all $\alpha\in\Phi$. This group $K$ is a hyperspecial maximal compact subgroup of $G$. Note that $T_{0}\subseteq K$. Also set
\[
I:=\la U_0, T_0, U_1^-\ra,
\]
namely the group generated by $U_0, T_0, U_1^-$, which is an Iwahori subgroup of $G$.



\subsection{Affine Weyl group of $G$}\label{SS:affine_Weyl_linear_group}
In this subsection, we review the basic theory of the affine Weyl group of the split reductive group $G$ over $F$. (In this subsection, $G$ is not necessarily our Chevalley group but any connected almost simple split reductive group over $F$.)

Let
\[
W_{\ea}:=N_G(T)\slash T_0
\]
be the extended affine Weyl group. The Iwahori-Bruhat decomposition tells that there is a natural bijection
\begin{equation}\label{E:Iwahori-Bruhat_decomp_reductive}
W_{\ea}\xrightarrow{\;\sim\;} I\backslash G\slash I,\quad w\mapsto IwI.
\end{equation}
Note that $W_{\ea}$ acts faithfully on $Y\otimes \R$ as affine transformations via the Tits homomorphism \cite[page 32]{T79}. %This identifies $W_{\ea}$ with a subgroup of the group $\Aff(Y\otimes \R)$ of affine transformations on $Y\otimes \R$.

For each $\alpha\in\Phi$ and $k\in\Z$, let
\[
H_{\alpha, k}:=\{v\in Y\otimes \R\st \la\alpha, v\ra=k\}
\]
be the affine hyperplane, where $\la-,-\ra$ is the ($\R$-valued pairing) on $(X\otimes\R)\times (Y\otimes \R)$ naturally induced from the canonical pairing on $X\times Y$. We then define
\[
w_{\alpha, k}: Y\otimes \R\longrightarrow Y\otimes \R
\]
by
\[
w_{\alpha, k}(v)=v-(\la\alpha, v\ra-k)\alphac,
\]
which is the affine reflection fixing $H_{\alpha, k}$. We let the affine Weyl group
\[
W_{\af}:=\la w_{\alpha, k}\st \alpha\in\Phi, k\in\Z\ra
\]
be the group of generated by the affine reflections $w_{\alpha, k}$. We will write $w_{\alpha}=w_{\alpha,0}$. Let
\[
W=\la w_{\alpha}\ra=\{w\in W_{\af}\st w(0)=0\}.
\]
We call this subgroup the finite Weyl group. It is isomorphic to the Weyl group of $(G,T)$.

With the choice of a set of simple roots $\Delta\subset \Phi$, the affine Weyl group $W_{\af}$ can be realized as a Coxeter group with generators $S_{\af,\Delta}=\{w_{\alpha}|\alpha\in \Delta\}\cup\{w_{\alpha_{0},-1}\}$, where we recall that $\alpha_0$ is the lowest root relative to $\Delta$. Similarly, $W$ is a Coxeter group with generating set $S_{\Delta}=\{w_{\alpha}|\alpha\in \Delta\}$. Note that the generators in $S_{\af,\Delta}$ are the reflections through the walls of the alcove 
\[
\mathfrak{A}:=\{v\in Y\otimes \mathbb{R}\st\alpha(v)>0 \text{ for all }\alpha\in \Delta\text{ and }\alpha_{0}(v)>-1\}.
\]
One can define the length function
\[
\ell=\ell_{\Delta}:W_{\af}\longrightarrow\Z_{\geq 0}
\]
by setting $\ell(w)$ to be the minimum number of generators needed to express $w$. This naturally extends the length function of the finite Weyl group $W$. 

The map $W\ltimes \mathbb{Z}[\Phi^{\vee}]\to W_{\af}$, defined by $(w, y)\mapsto (v\mapsto w(v)+y)$, is a group isomorphism
\[
W_{\af}\cong W\ltimes \Z[\Phi^{\vee}].
\]
Similarly, the group $W_{\ea}$ is generated by $W_{\af}$ and translations by elements of $Y$, and one has
\[
W_{\ea}\cong W\ltimes Y.
\]

Let $\Omega\subset W_{\ea}$ be the subgroup that preserves $\mathfrak{A}$. The action of $\Omega$ on $\mathfrak{A}$ permutes the supporting hyperplanes of $\mathfrak{A}$ and thus induces an automorphism on the affine Dynkin diagram. Thus $\Omega$ permutes the generators of $W_{\af}$ from which it follows that $\Omega$ normalizes $W_{\af}$. It is also known that $W_{\ea}$ is generated by $\Omega$ and $W_{\af}$, and $\Omega\cap W_{\af}=\{1\}$, and thus 
\[
W_{\ea}\cong \Omega\ltimes W_{\af}.
\]
With this decomposition we can extend the length function $\ell$ on $W_{\af}$ to a length function on $W_{\ea}$ by defining $\ell(\omega w)=\ell(w)$, where $w\in W_{\af}$ and $\omega\in\Omega$. Finally, the group $\Omega$ can be shown to be isomorphic to the finite abelian group $Y/\mathbb{Z}[\Phi^{\vee}]$.


%%%%%%%%%%%%%%%%%%%%%%%%%%%%%%%%%%%%%%%%%%%%%%%%%%%%%%%%%%%%%%%%%%%%%%%

\section{Double covers}\label{S:2Cover}

%%%%%%%%%%%%%%%%%%%%%%%%%%%%%%%%%%%%%%%%%%%%%%%%%%%%%%%%%%%%%%%%%%%%%%%
It is known that the Chevalley group $G$ has a unique nontrivial topological $\mu_2$-extension $\Gt$, which fits into an exact sequence
\begin{equation}\label{CExt}
	1\longrightarrow\mu_2\longrightarrow\Gt\stackrel{\pr}{\longrightarrow} G\longrightarrow 1.
\end{equation}
In this section, we will establish various important properties of this $\mu_2$-extension $\Gt$. Although we assume that everything in this section is essentially known to experts, we will give a self-contained exposition and fix our notation. Also from this section on, we assume that $G$ is the almost simple simply-connected Chevalley group associated with a simply-laced root system $\Phi$, unless otherwise stated. Also we exclude the groups of type  $A_1$ and $D_2$. Hence $G$ is a Chevalley group of type $A_r (r\geq 2)$, $D_r (r\geq 3)$, $E_6, E_7$ or $E_8$.


%%%%%%%%%%%%%%%%%%%%%%%%%%%%%%%%%%%%%%%%%%%%%%%%%%%%%%%%%%%%%%%%%%%%%%%

\subsection{Steinberg relations}

%%%%%%%%%%%%%%%%%%%%%%%%%%%%%%%%%%%%%%%%%%%%%%%%%%%%%%%%%%%%%%%%%%%%%%%

In terms of generators and relations, $\Gt$ is given by
\[
\Gt=\la \epsilon\xta(t)\st t\in F, \alpha\in\Phi, \epsilon\in\mu_2\ra
\]
subject to the relations
\begin{align}
\tag{MR1}\xta(t)\xta(u)&=\xta(t+u);\label{MR1}\\
\tag{MR2}[\xta(t), \xtb(u)]&=\prod_{i\alpha+j\beta\in\Phi}\tilde{x}_{i\alpha+j\beta}(c_{ij}t^iu^j),\quad(\alpha+\beta\neq 0);\label{MR2}\\
\tag{MR3}\hta(t)\hta(u)&=(t, u)\hta(tu),\label{MR3}
\end{align}
where $c_{ij}$ is the same as linear case. Again, if $\Phi$ is simply-laced, we have
\begin{equation}\label{E:MR2'}
\tag{MR2$'$}[\xta(t), \xtb(u)]=\begin{cases}1&\text{if $\alpha+\beta\notin\Phi$};\\
\xt_{\alpha+\beta}(ctu)&\text{if $\alpha+\beta\in\Phi$},
\end{cases}
\quad (\alpha+\beta\neq 0),
\end{equation}
where $c\in\{\pm 1\}$. As before, we set
\begin{align*}
\wta(t)&:=\xta(t)\xt_{-\alpha}(-t^{-1})\xta(t);\\
\hta(t)&:=\wt_{\alpha}(t)\wta(-1).
\end{align*}
Let us note that, under the projection $\Gt\xrightarrow{\,\pr\,} G$, the images of the elements $\xta(t), \wta(t)$ and $\hta(t)$ are, respectively, $\xa(t), \wa(t)$ and $\ha(t)$.



Many of the important relations holding for these symbols are summarized in \cite[p.24]{Matsumoto}. In particular, we need to mention the following:
\begin{align}
\tag{MR4}\label{MR4}\hta(t)\htb(u)&=(u, t^{\la\beta,\alpha\ra})\htb(u)\hta(t);\\
\tag{MR5}\label{MR5}\wta(t)\htb(u)\wta(-t)&=\htt_{w_\alpha(\beta)}(c t^{-\la\alpha, \beta\ra}u)\htt_{w_\alpha(\beta)}(c t^{-\la\alpha, \beta\ra})^{-1}\\
\notag&=(u, c t^{-\la\beta,\alpha\ra})\htt_{w_\alpha(\beta)}(u);\\
\tag{MR6}\label{MR6}\wta(t)\wtb(u)\wta(-t)&=\wt_{w_\alpha(\beta)}(ct^{-\la\alpha, \beta\ra}u)
\end{align}
where $c=c(\alpha, \beta)\in\{\pm 1\}$ is as in \eqref{R6}, which is denoted by $\eta$ in \cite[(c), p.24]{Matsumoto}. As a special case of \eqref{MR5}, we have
\begin{gather}\tag{MR5'}\label{MR5'}
\begin{aligned}
\wta(1)\htb(u)\wta(-1)&=\htb(u)\hta(u),\quad (\alpha+\beta\in\Phi);\\
\wta(1)\hta(u)\wta(-1)&=\hta(u^{-1}).
\end{aligned}
\end{gather}

%For any $y\in Y$, we can uniquely write $y=\sum_{\alpha_{j}\in \Delta} c_{j}\alpha_{j}^{\vee}$. For $u\in F^{\times}$ define $\overline{y}(u):=h_{\alpha_{1}}(u)^{c_{1}}\ldots h_{\alpha_{r}}(u)^{c_{r}}$. (Remember that we are using the Bourbaki labeling of the simple roots.)
%
%For each subgroup $H\subseteq G$, we let $\Ht\subseteq\Gt$ be its preimage in $\Gt$. 


We often use the following.

\begin{Prop}\label{PMRootSwap}
Let $\alpha\in \Phi$. Let $t,u\in F^{\times}$ be such that $1+tu\neq 0$. Then
\begin{equation*}
\xt_{\alpha}\,(t)\xt_{-\alpha}(u)=(1+tu,-t)\,\xt_{-\alpha}(\frac{u}{1+tu})\,\htt_{\alpha}(1+tu)\,\xt_{\alpha}(\frac{t}{1+tu}),
\end{equation*}
where we note
\[
(1+tu, -t)=(1+tu, u).
\]
\end{Prop}
\begin{proof}
This is an elementary exercise using the Steinberg relations. Also note that $(1+tu, -t)=(1+tu, 1-(1+tu))(1+tu, -t)=(1+tu, u)$.
\end{proof}


\begin{comment}
\begin{proof}
		\begin{align*}
		x_{\alpha}(t)x_{-\alpha}(u)=&x_{-\alpha}(t^{-1})w_{-\alpha}(-t^{-1})x_{-\alpha}(t^{-1}+u)\\
		=&x_{-\alpha}(t^{-1})w_{-\alpha}(-t^{-1})x_{\alpha}((t^{-1}+u)^{-1})w_{\alpha}(-(t^{-1}+u)^{-1})x_{\alpha}((t^{-1}+u)^{-1})\\
		=&x_{-\alpha}(t^{-1})x_{-\alpha}(-t^{-2}(t^{-1}+u)^{-1})w_{-\alpha}(-t^{-1})w_{\alpha}(-(t^{-1}+u)^{-1})x_{\alpha}((t^{-1}+u)^{-1})\\
		=&x_{-\alpha}(\frac{u}{1+tu})w_{\alpha}(t)w_{\alpha}(-(t^{-1}+u)^{-1})x_{\alpha}((t^{-1}+u)^{-1})\\
		=&(-t,(t^{-1}+u)^{-1})x_{-\alpha}(\frac{u}{1+tu})h_{\alpha}(1+tu)x_{\alpha}(\frac{t}{1+tu})\\
		=&(-t,t^{-1}+u)x_{-\alpha}(\frac{u}{1+tu})h_{\alpha}(1+tu)x_{\alpha}(\frac{t}{1+tu})\\
		=&(-t,1+tu)x_{-\alpha}(\frac{u}{1+tu})h_{\alpha}(1+tu)x_{\alpha}(\frac{t}{1+tu})
	\end{align*}
\end{proof}
\end{comment}

The following corollary is immediate. 

\begin{Cor}\label{PMRootSwapCor}
Let $\alpha\in \Phi$. Let $t\in\mathcal{O}$ and $u\in \varpi^{2e}\mathcal{O}$. Then
\begin{equation*}
		\xt_{\alpha}(t)\,\xt_{-\alpha}(u)=\xt_{-\alpha}(\frac{u}{1+tu})\,\htt_{\alpha}(1+tu)\,\xt_{\alpha}(\frac{t}{1+tu}).
\end{equation*}
\end{Cor}
\begin{proof}
If $t\in\mathcal{O}^{\times}$, then $(1+tu,-t)=1$ because $1+tu\in U_{2e}\subseteq R$, where we recall $R$ is the integral radical.  If $t\in\varpi \mathcal{O}$, then $(1+tu,-t)=1$ because $1+tu\in U_{2e+1}\subseteq\Ocal^{\times 2}$.
\end{proof}

%Corollary \ref{PMRootSwapCor} will allow us to prove the splitting of certain congruence subgroups of $G$.

For any subgroup $H\subseteq G$ let 
\[
\widetilde{H}:=\pr^{-1}(H).
\]
We say that $H$ splits sequence \eqref{CExt} or splits to $\Gt$ if there exists a group homomorphism $\s:H\rightarrow \widetilde{G}$ such that $\pr\circ \s=id_{H}$. In this case, $H\times \mu_{2}\cong \widetilde{H}$ via the map $(h,\pm1)\mapsto \pm\s(h)$. 

From the Steinberg relations \eqref{MR1} and \eqref{MR2} there is a group homomorphism 
\begin{equation}\label{E:section_for_unipotent_radical}
\s^{\pm}:U^{\pm}(F)\longrightarrow \Gt \text{\quad such that\quad} \s^{\pm}(x_{\alpha}(t))=\xt_{\alpha}(t)
\end{equation}
for all $\alpha\in\Phi^{\pm}$ and $t\in F$. This defines a splitting for $U^{\pm}$. 

\begin{comment}
\subsection{Automorphisms of $\Gt$}
Let us discuss some automorphisms of $\Gt$. Let $G_{\mathrm{ad}}$ be the $F$-rational points of the adjoint Chevalley group associated with $\Phi$. Recall that $G_{\mathrm{ad}}$ acts on $G$ via automorphisms. This action of $G_{\mathrm{ad}}$ lifts to an action of $G_{\mathrm{ad}}$ on $\widetilde{G}$, which is trivial on $\mu_{2}$. (See \cite{} for a discussion of this.) We will be specifically concerned with automorphisms induced by the action of $T_{ad}$ on $\widetilde{G}$ and with these we can be very explicit as follows.

Let $P^\vee$ be the coweight lattice of $G_{ad}$, so $P^\vee=\{\omega^\vee\in X_*(T_{ad})\st \la \alpha, \omega^\vee\ra\in \Z\text{ for all $\alpha\in\Phi$}\}$. For each $\omega^{\vee}\in P^\vee$ and $u\in F^{\times}$, we can define an action $\phi_{\omega^{\vee}(u)}$ of $\omega^{\vee}(u)\in T_{ad}$ on $\Gt$ by
\begin{equation*}
	\phi_{\omega^{\vee}(u)}:\xt_{\alpha}(v)\mapsto \xt_{\alpha}(u^{\langle\alpha,\omega^{\vee}\rangle}v),
\end{equation*}
where $\alpha\in \Phi$ and $v\in F$. It can be checked directly that this assignment on the generators of $\Gt$ is compatible with the relations defining $\Gt$ and thus defines a group automorphism $\phi_{\omega^{\vee}(u)}\in\Aut(\Gt)$.

By a direct calculation we can make the action of $\omega^{\vee}(u)$ on $\Tt$ explicit as follows.

\begin{Lem}\label{AdAuts} Let $\omega^{\vee}\in P^\vee$ and $u\in F^{\times}$. Then
	\begin{equation}\label{TadAut}
		\phi_{\omega^{\vee}(u)}(\htt_{\alpha}(v))=(v,u^{\langle\alpha,\omega^{\vee}\rangle})\htt_{\alpha}(v).
	\end{equation}
	Thus, if $\omega^{\vee}\in 2P^{\vee}$, then $\phi_{\omega^{\vee}(u)}$ acts trivially on $\Tt$. Also $\phi_{\omega^{\vee}(u)}$ acts trivially on $\mu_{2}$.
\end{Lem}

\begin{proof}
	Recall that by definition for $\alpha\in\Phi$ and $v\in F^{\times}$ we have
	\begin{equation*}
		\htt_{\alpha}(v)=\wt_{\alpha}(v)\wt_{\alpha}(-1)=\xt_{\alpha}(v)\xt_{-\alpha}(-v^{-1})\xt_{\alpha}(v)\xt_{\alpha}(-1)\xt_{-\alpha}(1)\xt_{\alpha}(-1).
	\end{equation*}
	Thus we have 
	\begin{align*}
		\phi_{\omega^{\vee}(u)}(\htt_{\alpha}(v))
		=&\xt_{\alpha}(u^{\langle\alpha,\omega^{\vee}\rangle}v)\xt_{-\alpha}(-[u^{\langle\alpha,\omega^{\vee}\rangle}v]^{-1})\xt_{\alpha}(u^{\langle\alpha,\omega^{\vee}\rangle}v)\\
		&\quad\cdot\xt_{\alpha}(-u^{\langle\alpha,\omega^{\vee}\rangle})\xt_{-\alpha}(u^{\langle-\alpha,\omega^{\vee}\rangle})\xt_{\alpha}(-u^{\langle\alpha,\omega^{\vee}\rangle})\\
		=&\wt_{\alpha}(u^{\langle\alpha,\omega^{\vee}\rangle}v)\wt_{\alpha}(-u^{\langle\alpha,\omega^{\vee}\rangle})\\
		=&(-u^{\langle\alpha,\omega^{\vee}\rangle}v,u^{\langle\alpha,\omega^{\vee}\rangle})\htt_{\alpha}(v)\\
		=&(v,u^{\langle\alpha,\omega^{\vee}\rangle})\htt_{\alpha}(v).
	\end{align*}
\end{proof}	

\end{comment}

%%%%%%%%%%%%%%%%%%%%%%%%%%%%%%%%%%%%%%%%%%%%%%%%%%%%%%%%%%%%%%%%%%%%%%%

\subsection{The lattice $\Yt$}

%%%%%%%%%%%%%%%%%%%%%%%%%%%%%%%%%%%%%%%%%%%%%%%%%%%%%%%%%%%%%%%%%%%%%%%

In this subsection, we will introduce a certain sublattice $\Yt$ of $Y$ and explicitly describe the quotient $\Yt\slash 2Y$. We will then show how the lattice $\Yt$ is used to describe the center $Z(\Tt)$ of the covering torus $\Tt$. The description of the quotient $\Yt\slash 2Y$ (Proposition \ref{tildYmod2Y}) is essentially \cite[Proposition 4.2]{S04}. For our construction of pseudo-spherical representations, however, it will be convenient for us to reprove this result.

Recall that $Y=X_*(T)$ is the group of cocharacters (\ie cocharacter lattice). Since our group $G$ is simply-connected, we have $Y=\Z[\Phi^\vee]$. Hence each $y\in Y$ is uniquely written as 
\[
y=\sum_{i}c_i\alphac_i.
\] 
For each $t\in F^\times$, we set
\[
y(t)=h_{\alpha_1}(t^{c_1})\cdots h_{\alpha_n}(t^{c_n})\in T.
\]
Then the maximal torus $T$ is generated by $y(t)$'s. Indeed, we have the isomorphism
\[
Y\otimes_{\Z} F^\times\longrightarrow T,\quad y\otimes t\mapsto y(t).
\]

Let us consider the analog for the cover $\Tt$. For each $y=\sum_{i}c_i\alphac_i$ and $t\in F^\times$, set
\[
\yt(t)=\htt_{\alpha_1}(t^{c_1})\cdots \htt_{\alpha_n}(t^{c_n})\in \Tt,
\]
where the order of multiplication is the order of $\Delta^\vee$. However, since $\Tt$ is not commutative due to (MR3) and (MR4), this product is not independent of the order of multiplication. We can, however, choose a sublattice of $Y$ so that the product is independent of the ordering as follows.

Let
\begin{equation}\label{E:definition_of_Yt}
\Yt=\{y\in Y\st \la \alpha, y\ra\in 2\Z\;\text{for all $\alpha \in \Phi$}\}.
\end{equation}
Note that
\[
2Y\subseteq \Yt\subseteq Y.
\]
We will prove that the product of $\yt(t)$ is independent of the ordering of $y\in\Yt$. For this purpose, we explicitly derive the conditions the coefficients $c_i$'s of $y=\sum_ic_i\alphac_i\in \Yt$ have to satisfy modulo $2\Z$, namely the ``parity conditions" to be satisfied by the $c_i$'s, as follows.

\begin{Lem}\label{L:parity_of_adjacent_root}
Let $y=\sum_ic_i\alphac_i\in \Yt$ and fix a simple root $\alpha_k$. Then we have
\[
\sum_{\text{$\alpha_i$ is adjacent to $\alpha_k$}}c_i\in 2\Z.
\]
Note that the number of $c_i$'s in the sum is at most $3$. 

Further, if $\alpha_k$ is adjacent to a root which has no other adjacent root than $\alpha_k$ itself (\ie $\alpha_k$ is adjacent to an ``endpoint" in the Dynkin diagram), then $c_k\in 2\Z$.
\end{Lem}
\begin{proof}
Since 
\[
\la \alpha_k,  y\ra=2c_k-\sum_{\text{$\alpha_i$ is adjacent to $\alpha_k$}}c_i\in 2\Z,
\]
the first part of the lemma follows.

Assume $\alpha_k$ is adjacent to an endpoint $\alpha_\ell$. Then
\[
\la \alpha_\ell, y\ra=2c_\ell-c_k\in 2\Z,
\]
which gives $c_k\in 2\Z$. The lemma follows.
\end{proof}


Now, given $y=\sum_ic_i\alphac_i\in \Yt$, we describe the parity of each $c_i$ by using the Dynkin diagram as follows. If $c_i$ is even, we denote the corresponding node for $\alpha_i$ in the Dynkin diagram by a black dot, and if it is odd, we denote it by a white dot. For example, assume our $\Phi$ is the $E_7$ root system. If $y=\sum_{i=1}^6c_i\alphac_i$ is such that $c_1, c_3, c_4$ and $c_6$ are even and the others are odd, then one can readily see that $y\in\Yt$. We denote this $y$ modulo $2Y$ by 
\[
\pgfkeys{/Dynkin diagram, edge-length=1cm, root-radius=.1cm}
\dynkin[labels={\alpha_1,\alpha_2,\alpha_3,\alpha_4,\alpha_5,\alpha_6,\alpha_7}]  E{*o**o*o}\rlap{\, .}
\]
Furthermore, one can see that this is the only non-zero element in $\Yt\slash 2Y$ for $E_7$ by using Lemma \ref{L:parity_of_adjacent_root} above, and hence $\Yt\slash 2Y=\Z\slash 2\Z$. 


\begin{comment}
Using these two lemmas along with the Dynkin diagrams, one can completely describe the parity conditions each $y=\sum_ic_i\alphac_i\in \Yt$ has to satisfy as follows. For example, consider the type $E_7$ case, so that we have the Dynkin diagram
\[
\pgfkeys{/Dynkin diagram, edge-length=1cm, root-radius=.1cm}
\dynkin[labels={\alpha_1,\alpha_2,\alpha_3,\alpha_4,\alpha_5,\alpha_6,\alpha_7}]  E{ooooooo}\rlap{\, .}
\]
If $y=\sum_{i=1}^7c_i\alphac_i\in \Yt$, then by the first of the above lemma, we always have $c_3, c_4, c_6\in 2\Z$ because $\alpha_3, \alpha_4, \alpha_6$ are adjacent to endpoints. Now, let us adapt the following convention: when the coefficient $c_i$ is even, we denote the corresponding node for $\alpha_i$ in the Dynkin diagram by a black dot. Then in this case we have
\[
\pgfkeys{/Dynkin diagram, edge-length=1cm, root-radius=.1cm}
\dynkin[labels={\alpha_1,\alpha_2,\alpha_3,\alpha_4,\alpha_5,\alpha_6,\alpha_7}]  E{oo**o*o}\rlap{\, .}
\]
Next, by using the second one of the above two lemmas with $\alpha_k=\alpha_3$, we obtain $c_1+c_4\in 2\Z$. But we already know $c_4\in 2\Z$, which gives $c_1\in 2\Z$. Hence the above Dynkin diagram is further written as
\[
\pgfkeys{/Dynkin diagram, edge-length=1cm, root-radius=.1cm}
\dynkin[labels={\alpha_1,\alpha_2,\alpha_3,\alpha_4,\alpha_5,\alpha_6,\alpha_7}]  E{*o**o*o}\rlap{\, .}
\]
Further, by using the second one of the above two lemmas with $\alpha_k=\alpha_4$ and $\alpha_6$, we obtain
\[
c_2+c_5\in 2\Z\qand c_5+c_7\in 2\Z.
\]
These are the parity conditions all the coefficients $c_i$'s have to satisfy. But this implies that if one of the three $c_2, c_5, c_7$ is even, then the other two will be automatically even. Hence, for $y=\sum_{i=1}^7c_i\alphac_i\in\Yt\slash 2Y$, either all the coefficients $c_i$ are even, or the three $c_2, c_5, c_7$ are odd and the others are even as in the above diagram.  In particular, we have $\Yt\slash 2Y=\Z\slash 2\Z$.

\end{comment}

With this convention, we have
\begin{Prop}\label{tildYmod2Y}
Depending on the type of the Dynkin diagram, we have
\[
\Yt\slash 2Y=\begin{cases}\Z\slash 2\Z\times \Z\slash 2\Z&\text{Type $D_{2r}$};\\
\Z\slash 2\Z&\text{Type $A_{2r}, D_{2r+1}$ and $E_7$};\\
1&\text{otherwise}.
\end{cases}
\]
The results are summarized in Table \ref{T:table}, where the nontrivial elements in $\Yt\slash 2Y$ are described by Dynkin diagrams with white dots indicating odd and black dots even.
\end{Prop}
\begin{proof}
This can be verified case by case using Lemma \ref{L:parity_of_adjacent_root}. Since the verification is elementary, we omit the details.
\end{proof}

{
%\pgfkeys{/Dynkin diagram, edge-length=.7cm, root-radius=.07cm}

\begin{table}[h]
\caption{Lattice $\Yt$}\label{T:table}
\centering
\renewcommand{\arraystretch}{1.5}
\begin{tabular}{cccl}
\hline
Type & $\Yt\slash 2Y$ & $\left|Y\slash \Yt\right|$ &non-trivial elements in $\Yt\slash 2Y$ \\\hline\hline
$A_{2r+1}$ & $\Z\slash 2\Z$ & $2^{2r}$ & \dynkin A{o*o*.*o}\\
$A_{2r}$ & 1 & $2^{2r}$ &\\\hline
$D_{2r+1}$  & $\Z\slash 2\Z$ & $2^{2r}$ & \dynkin D{****.**oo}\\
$D_{2r}$ & $\Z\slash 2\Z\times \Z\slash 2\Z$  & $2^{2r-2}$ & \dynkin D{o*o*.o**o}\;\dynkin D{o*o*.o*o*}\;\dynkin D{****.**oo}\\\hline
 $E_6$ & 1 & $2^6$ & \\\hline
 $E_7$ & $\Z\slash 2\Z$ & $2^6$ &\dynkin E{*o**o*o} \\\hline
 $E_8$ & 1 & $2^8$ &\\\hline
\end{tabular}
\end{table}

}





As a small corollary, let us mention
\begin{Cor}\label{C:parity_of_adjacent_root3}
Let $y=\sum_ic_i\alphac_i\in \Yt$. If $\alpha_i$ and $\alpha_j$ are adjacent to each other, then $c_ic_j\in 2\Z$.
\end{Cor}
\begin{proof}
This can be read off from the Dynkin diagrams in Table \ref{T:table}.


%By the above lemma, the lemma holds if one of $\alpha_i$ and $\alpha_j$ is either an endpoint or adjacent to an endpoint. Hence we assume neither of them is an endpoint or adjacent to an endpoint. By looking at all the simply-laced Dynkin diagrams, one can see that $\alpha_i$ and $\alpha_j$ are located in the Dynkin diagram as follows
%\[
%\dynkin[labels={\alpha_k,\alpha_i,\alpha_j,\alpha_{\ell}}] A{*.oo.*}
%\]
%where the black dots ($\alpha_k$ and $\alpha_{\ell}$) are adjacent to an endpoint and hence $c_k, c_{\ell}\in 2\Z$.
%
%Now, let $\alpha_{k+1}$ be the root adjacent to $\alpha_k$ from the right, and $\alpha_{k+2}$ be the root adjacent to $\alpha_{k+1}$ from the right, so that we have
%\[
%\dynkin[labels={\alpha_k,\alpha_{k+1},\alpha_{k+2},\alpha_\ell}] A{*oo.*}.
%\]
%Then
%\[
%\la y, \alphac_{k+1}\ra=-c_k+2c_{k+1}-c_{k+2}\in 2\Z,
%\]
%which implies $c_{k+2}\in 2\Z$. Arguing inductively, we have $c_{k+2m}\in 2\Z$ for all $m\in\Z$ with $k+2m\leq \ell$. Namely, we have
%\[
%\dynkin[labels={\alpha_k,,,,,\alpha_\ell}] A{*o*o*o.*}
%\]
%where the black dots indicate that the corresponding coefficients in $y=\sum_ic_i\alphac_i$ is always even; namely the black dots and white dots alternate. Hence at least one of $c_i$ and $c_j$ are even. The lemma follows.
\end{proof}


This corollary implies
\begin{Prop}\label{ZTbiject}
Let $y=\sum_ic_i\alphac_i\in\Yt$ and $t_i\in F^\times$. The product
\[
\prod_{i}\htt_{\alpha_i}(t_i^{c_i})
\]
is independent of the ordering.
\end{Prop}
\begin{proof}
We have only to show that if $\alpha_i$ and $\alpha_j$ are adjacent then
\[
\htt_{\alpha_i}(t_i^{c_i})\htt_{\alpha_j}(t_j^{c_j})=\htt_{\alpha_j}(t_j^{c_j})\htt_{\alpha_i}(t_i^{c_i}).
\]
But
\begin{align*}
\htt_{\alpha_i}(t_i^{c_i})\htt_{\alpha_j}(t_j^{c_j})&=(t_i^{c_i}, t_j^{-c_j})\,\htt_{\alpha_j}(t_j^{c_j})\htt_{\alpha_i}(t_i^{c_i})\quad\text{by (MR4)}\\
&=(t_i, t_j)^{c_ic_j}\htt_{\alpha_j}(t_j^{c_j})\htt_{\alpha_i}(t_i^{c_i}),
\end{align*}
and $(t_i, t_j)^{c_ic_j}=1$ by the above corollary.
\end{proof}


In particular, for each $y=\sum_ic_i\alphac_i\in\Yt$ and $t\in F^\times$, we can define
\[
\yt(t):=\prod_{i}\htt_{\alpha_i}(t^{c_i}),
\]
which is independent of the ordering by Proposition \ref{ZTbiject} with $t=t_i$ for all $i$.

\begin{comment}
We then have the following.
\begin{Prop}\label{P:Center_Z(Tt)}
Let $Z(\Tt)$ be the center $\Tt$. The map
\[
\mu_2\times\Yt\otimes_{\Z}F^\times \longrightarrow Z(\Tt),\quad (\epsilon, y\otimes t)\mapsto \epsilon\, \yt(t),
\]
where $\epsilon\in\mu_2$ and $y\otimes t\in \Yt\otimes_{\Z}F^\times$, is a bijection.
\end{Prop}
\begin{proof}
First, we will show that $\epsilon\, \yt(t)$ is indeed in $Z(\Tt)$. For this, it suffices to show that $\yt(t)\htb(s)=\htb(s)\yt(t)$ for all $\htb(s)\in \Tt$, where $\beta\in\Phi$ and $s\in F^\times$. But by (MR4), for each $\alpha_i$,
\[
\htt_{\alpha_i}(t^{c_i})\htb(s)=(t^{c_i}, s^{\la \alphac_i, \beta\ra})\,\htb(s)\htt_{\alpha_i}(t^{c_i})=(t, s^{\la \beta, c_i\alphac_i\ra})\,\htb(s)\htt_{\alpha_i}(t^{c_i}).
\]
Hence one can readily see that
\[
\yt(t)\htb(s)=(t, s^{\la \beta, y\ra})\,\htb(s)\yt(t).
\]
Since $\la\beta, y\ra\in2\Z$, we have $\yt(t)\htb(s)=\htb(s)\yt(t)$.

To show that the map is surjective, note that each element in $\Tt$ is (not necessarily uniquely) written as
\[
\epsilon\prod_i\htt(t^{d_i})
\]
for some $\sum_id_i\alphac_i\in Y$. If this element is in the center $Z(\Tt)$, it has to commute with all the $\htb(s)$. But by the same computation as above, we have
\[
\prod_i\htt(t^{d_i})\cdot\htb(s)=(t, s^{\la\beta, \sum_id_i\alphac_i\ra})\,\htb(s) \cdot \prod_i\htt(t^{d_i}). 
\]
Hence we must have either $\sum_id_i\alphac_i\in \Yt$ or $t\in F^{\times 2}$. If $t\in F^{\times 2}$, so $t=u^2$ for some $u\in F^\times$, then
\[
(\sum_id_i\alphac_i)\otimes u^2=(\sum_i2d_i\alphac_i)\otimes u\in \Yt\otimes_{\Z}F^\times.
\]
Hence the given map is surjective.

To show it is injective, compose the map with $Z(\Tt)\slash\mu_2$. This is a homomorphism whose kernel is $\mu_2$. This shows that the given map is injective.
\end{proof}

\begin{Rmk}
The bijection in Proposition \ref{P:Center_Z(Tt)} may not be a group isomorphism. Indeed, if the group is of type $A_{2m}$ with $m$ odd, then this bijection fails to be a homomorphism. 
\end{Rmk}

\end{comment}


%%%%%%%%%%%%%%%%%%%%%%%%%%%%%%%%%%%%%%%%%%%%%%%%%%%%%%%%%%%%%%%%%%%%%%%

\subsection{The Weyl group action on $\Yt$}

%%%%%%%%%%%%%%%%%%%%%%%%%%%%%%%%%%%%%%%%%%%%%%%%%%%%%%%%%%%%%%%%%%%%%%%

The Weyl group $W$ acts on $\Yt$ because the inner product $\la-,-\ra$ is $W$-invariant. Also $2Y$ is $W$-invariant. Hence the action of $W$ on $\Yt$ descends to $\Yt\slash 2Y$. But this action of $W$ on $\Yt\slash 2Y$ is trivial; namely we have
\begin{Lem}\label{L:Yt_is_Weyl_invariant_mod_2Y}
If $y\in \Yt$, then $w y=y$ in $\Yt\slash 2Y$ for all $w\in W$.
\end{Lem}
\begin{proof}
Let $\alpha\in\Delta$ be fixed. We have to show $w_\alpha y=y\pmod{2Y}$. Note that
\begin{equation*}
    w_{\alpha}y = y-\langle\alpha,y\rangle\alphac.
\end{equation*}
Now the result follows because $y\in \Yt$.

\begin{comment}
\[
w_\alpha\alphac_i=\begin{cases} -\alphac_i&\text{if $\alpha_i=\alpha$};\\
\alphac_i+\alphac&\text{if $\alpha$ and $\alpha_i$ are adjacent, \ie $\la\alpha_i,\alpha\ra=-1$};\\
\alphac_i&\text{otherwise, \ie $\la\alpha_i,\alpha\ra=0$}.
\end{cases}
\]
But viewed modulo $2Y$, we have $w_\alpha\alphac=-\alphac=\alphac$. Hence $w_\alpha$ fixes all the simple roots possibly except the adjacent ones.

Hence, modulo $2Y$, we have
\begin{align*}
w_\alpha y&= \sum_{\text{$\alpha_k$ is adjacent to $\alpha$}}w_\alpha c_k\alphac_k+\sum_{\text{$\alpha_i$ is not adjacent to $\alpha$}}c_i\alphac_i\\
&=\sum_{\text{$\alpha_k$ is adjacent to $\alpha$}}(c_k\alphac_k+c_k\alphac)+\sum_{\text{$\alpha_i$ is not adjacent to $\alpha$}}c_i\alphac_i\\
&=\sum_{\text{$\alpha_k$ is adjacent to $\alpha$}}c_k\alphac+\sum_{\text{$\alpha_k$ is adjacent to $\alpha$}}c_k\alphac_k+\sum_{\text{$\alpha_i$ is not adjacent to $\alpha$}}c_i\alphac_i\\
&=\left(\sum_{\text{$\alpha_k$ is adjacent to $\alpha$}}c_k\right)\alphac+\sum_{i}c_i\alphac_i\\
&=\sum_{i}c_i\alphac_i,
\end{align*}
where we used Lemma \ref{L:parity_of_adjacent_root} for the last step. Hence $w_\alpha y=y$.

%Assume there is only one root $\alpha_k$ adjacent to $\alpha$. Then 
%\begin{align*}
%w_\alpha y&=w_\alpha c_k\alphac_k+\sum_{i\neq k}c_i\alphac_i\\
%&=c_k\alphac_k+c_k\alphac+\sum_{i\neq j}c_i\alphac_i.
%\end{align*}
%Since $y\in\Yt$, we have
%\[
%\la y,\alphac\ra=-c_k+2c_\alpha\in 2\Z,
%\]
%where $c_\alpha$ is the coefficient in $y$ of this fixed $\alpha$. Hence $c_k\in 2\Z$, which implies $w_\alpha y=y\pmod{2Y}$.
%
%Next assume there are exactly two roots $\alpha_k$ and $\alpha_\ell$ adjacent to $\alpha$. Then
%\begin{align*}
%w_\alpha y&=w_\alpha c_k\alphac_k+w_\alpha c_\ell\alphac_\ell+\sum_{i\neq k, \ell}c_i\alphac_i\\
%&=c_k\alphac_k+c_\ell\alphac_\ell+(c_k+c_\ell)\alphac+\sum_{i\neq k, \ell}c_i\alphac_i.
%\end{align*}
%Since $y\in\Yt$, we have
%\[
%\la y,\alphac\ra=2c_\alpha-c_k-c_\ell\in 2\Z,
%\]
%where $c_\alpha$ is the coefficient of this fixed $\alpha$. Hence $c_k+c_\ell\in 2\Z$, which implies $w_\alpha y=y\pmod{2Y}$.
%
%Finally, assume there are three roots $\alpha_k, \alpha_\ell, \alpha_m$ adjacent to $\alpha$. Then
%\begin{align*}
%w_\alpha y&=w_\alpha c_k\alphac_k+w_\alpha c_\ell\alphac_\ell+w_\alpha c_m\alphac_m+\sum_{i\neq k, \ell, m}c_i\alphac_i\\
%&=c_k\alphac_k+c_\ell\alphac_\ell+c_m\alphac_m+(c_k+c_\ell+c_m)\alphac+\sum_{i\neq k, \ell, m}c_i\alphac_i.
%\end{align*}
%Since $y\in\Yt$, we have
%\[
%\la y,\alphac\ra=2c_\alpha-c_k-c_\ell-c_m\in 2\Z,
%\]
%where $c_\alpha$ is the coefficient of this fixed $\alpha$. Hence $c_k+c_\ell+c_m\in 2\Z$, which implies $w_\alpha y=y\pmod{2Y}$.
%
%
%The lemma is proven.
\end{comment}
\end{proof}


%%%%%%%%%%%%%%%%%%%%%%%%%%%%%%%%%%%%%%%%%%%%%%%%%%%%%%%%%%%%%%%%%%%%%%%

\subsection{The group $\Wcal$}

%%%%%%%%%%%%%%%%%%%%%%%%%%%%%%%%%%%%%%%%%%%%%%%%%%%%%%%%%%%%%%%%%%%%%%%
For each $\alpha\in\Phi$, let
\[
\wta:=\wta(1).
\]
We then define
\begin{equation}\label{E:definition_of_Wcal}
\Wcal:=\la \wta\st\alpha\in\Phi\ra,
\end{equation}
namely the group generated by the $\wta$'s.

Note that
\[
\wta(1)^{-1}=\wta(-1).
\]
Since $\wta(-1)^{2}=\hta(-1)$ and $\hta(-1)^{-1}=(-1,-1)\hta(-1)$ we have
\begin{equation}\label{E:square_of_w}
\wta(1)^2=(-1,-1)\cdot\hta(-1).
\end{equation}
Hence each $\wta$ has order 4 or 8, depending on $(-1,-1)=1$ or $-1$ because $\hta(-1)^2=(-1,-1)$. The following should be mentioned : if $\alpha+\beta$ is a root, then
\begin{gather}\label{E:Savin_relation}
\begin{aligned}
\wta(1)\wtb(1)\wta(1)&=\wtb(1)\wta(1)\wtb(1)\quad\text{(braid relation)};\\
\wta(1)\htb(t)\wta(-1)&=\htb(t)\hta(t),
\end{aligned}
\end{gather}
where the second one is actually \eqref{MR5'} listed earlier. (See \cite[Th\'eor\`eme 6]{Matsumoto} and \cite[Lemma 3.2]{S04}.) Set
\[
\Tt(\Z):=\la\hta(-1)\st\alpha\in\Delta\ra.
\]
Then we have the exact sequence
\begin{equation}\label{E:exact_sequence_of_Weyl_group}
1\longrightarrow \Tt(\Z)\longrightarrow\Wcal\longrightarrow W\longrightarrow 1,
\end{equation}
where $W$ is the Weyl group of $G$ and the surjection sends each $\wta(1)$ to the corresponding Weyl group element.



\subsection{Affine Weyl group of $\Gt$}\label{SS:affine_Weyl_covering_group}

In this section, we modify the theory of affine Weyl groups of the linear group $G$ described in \S \ref{SS:affine_Weyl_linear_group} to that of the covering group $\Gt$.

\begin{comment}
Recall for the linear case that we have
\[
W_{\ea}\cong W\ltimes Y\qand W_{\af}\cong W\ltimes \Z[\Phi^\vee],
\]
so that 
\[
1\longrightarrow W_{\af}\longrightarrow W_{\ea}\longrightarrow Y\slash \Z[\Phi^\vee]\longrightarrow 1,
\]
which gives $W_{\ea}=W_{\af}\rtimes Y\slash\Z[\Phi^\vee]$, where $Y\slash\Z[\Phi^\vee]$ is the set of length 0 elements. In particular, if $G$ is simply-connected, we have $W_{\af}=W_{\ea}$. Also as in \eqref{E:W_af_generators}, the group $W_{\af}$ is generated by the simple reflections $w_{\alpha_0, -1}, w_{\alpha_1}, \dots, w_{\alpha_n}$.
\end{comment}

First, note that the root datum $(X, \Phi, Y, \Phi^\vee)$ is isomorphic to the root datum
\[
(\frac{1}{2}X, \frac{1}{2}\Phi, 2Y, 2\Phi^\vee).
\]
We also have a modified root datum of $\Gt$ \cite{Mc12,W14}
\[
(\Xt, \frac{1}{2}\Phi, \Yt, 2\Phi^\vee),
\]
where $\Yt$ is as before and $\Xt=\Hom(\Yt,\mathbb{Z})$. This modified datum is that of the linear group
\[
G':=G\slash Z_2,
\]
where $Z_2$ is the 2-torsion of the center of $G$. Note that
\[
Z_2\cong \Yt\slash \Z[2\Phi^{\vee}].
\]

Let $\Wt_{\af}$ be the affine Weyl group associated with the root system $\frac{1}{2}\Phi$, which acts on $Y\otimes \mathbb{R}$. Fix a subset of simple roots $\Delta\subseteq \Phi$. Then $\frac{1}{2}\Delta $ is a subset of simple roots for $\frac{1}{2}\Phi$. Let
\begin{equation}\label{E:w_aff_definition}
    w_{\af}:=w_{\alpha_0, -2}=w_{\frac{1}{2}\alpha_{0},-1}.
\end{equation}
One can then see that
\[
\Wt_{\af}=\la w_{\af}, w_{\alpha_1}, \dots, w_{\alpha_n}\ra.
\]

Also let $\Wt_{\ea}$ be the extended affine Weyl group associated with $G'$. We then have
\[
\Wt_{\ea}\cong W\ltimes \Yt\qand \Wt_{\af}\cong W\ltimes 2Y,
\]
so that 
\[
1\longrightarrow \Wt_{\af}\longrightarrow \Wt_{\ea}\longrightarrow \Yt\slash 2Y\longrightarrow 1.
\]
(Note that in \cite{Karasiewicz}, our $\Wt_{\ea}$ is denoted by $\Wt_{\aff}$, and our $\Wt_{\af}$ by $\widehat{W}_{\aff}$.) 

Let 
\[
\tilde{\ell}=\tilde{\ell}_{\frac{1}{2}\Delta}:\Wt_{\ea}\cong W\ltimes \Yt\longrightarrow\Z_{\geq 0},
\]
be the length function, which can be also computed by the formula of \cite[Proposition 1.23]{IM}; namely
\[
\tilde{\ell}(w, y)=\sum_{\alpha\in\frac{1}{2}\Phi^+\cap w\frac{1}{2}\Phi^+}|\la\alpha, y\ra|+\sum_{\alpha\in\frac{1}{2}\Phi^-\cap w\frac{1}{2}\Phi^-}|\la\alpha, y\ra+1|,
\]
where $(w, y)\in W\ltimes\Yt$. Let $\Omegat$ be the set of length 0 elements. Then 
\[
\Omegat\cong \Yt\slash 2Y;
\]
indeed, we have
\[
\tilde{\ell}:\Omegat\ltimes \Wt_{\af}\longrightarrow\Z_{\geq 0},\quad \tilde{\ell}(w, \varepsilon)=\tilde{\ell}(w),
\]
for $(w, \varepsilon)\in \Omegat\ltimes \Wt_{\af}$. Note that $\Omegat$ is the subgroup that preserves the alcove $2\mathfrak{A}$, where $\mathfrak{A}$ is defined relative to $\Delta\subseteq\Phi$, as in \S \ref{SS:affine_Weyl_linear_group}.

Since $G$ is simply-connected and 
\[
\It\backslash\Gt\slash \It = I\backslash G\slash I,
\]
the Iwahori-Bruhat decomposition \eqref{E:Iwahori-Bruhat_decomp_reductive} implies
\[
W_{\ea}=W_{\af}\cong \It\backslash\Gt\slash \It.
\]
As we will see, however, the Iwahori Hecke algebra we will be considering is not supported on all of $W_{\af}$, but only on $\Wt_{\ea}$. This is precisely the rationale for introducing the groups $\Wt_{\af}$ and $\Wt_{\ea}$. 

%%%%%%%%%%%%%%%%%%%%%%%%%%%%%%%%%%%%%%%%%%%%%%%%%%%%%%%%%%%%%%%%%%%%%%%

\section{Pseudo-spherical representations of $\Tt_{0}$}\label{SecTORep}

%%%%%%%%%%%%%%%%%%%%%%%%%%%%%%%%%%%%%%%%%%%%%%%%%%%%%%%%%%%%%%%%%%%%%%%

Our goal in this section is to construct the pseudo-spherical representations of $\Tt_{0}$ (Definition \ref{Def:PSTRep}) and show that they are $W$-invariant and types of $\Tt$ in the sense of Bushnell-Kutzko \cite[(4.1) Definition]{BK98}.


%%%%%%%%%%%%%%%%%%%%%%%%%%%%%%%%%%%%%%%%%%%%%%%%%%%%%%%%%%%%%%%%%%%%%%%

\subsection{Structure of $\Tt_{0}$}

%%%%%%%%%%%%%%%%%%%%%%%%%%%%%%%%%%%%%%%%%%%%%%%%%%%%%%%%%%%%%%%%%%%%%%%


Let us recall from \eqref{E:def_of_T_0} and \eqref{E:def_of_T_R} that we have defined
\[
T_0=T_{\Ocal}=\{h_{\alpha}(a)\st a\in \Ocal, \alpha\in\Phi\}
\qand
T_R=\{h_{\alpha}(a)\st a\in R, \alpha\in\Phi\},
\]
where $R$ is the integral radical as in \eqref{E:integral_radical_def}. We let $\Tt_0$ be the preimage of $T_0$ in $\Tt$, which is the maximal compact subgroup of $\Tt$.  

Let 
\[
\s_{T_{R}}:T_{R}\longrightarrow \Tt,\quad \prod_{i}h_{\alpha_{i}}(u_{i})\mapsto \prod_{i}\htt_{\alpha_{i}}(u_{i}),
\]
where the product on the right-hand side is independent of the ordering because $u_i\in R$ for all $i$. By comparing the relations \eqref{R3} and \eqref{R4} with \eqref{MR3} and \eqref{MR4}, one can readily see that this map is a homomorphism (\ie it splits the sequence \eqref{CExt}), and moreover its image $\s_{T_R}(T_R)$ is normal in $\Tt_0$, so that we can view $T_R$ as a subgroup of $\Tt_0$.

\begin{comment}
\begin{Lem}\label{TorSplit}
\quad
\begin{enumerate}
	\item The map $\s_{T_{R}}$ splits sequence \eqref{CExt}, \ie it is a homomorphism. 
	\item The image of $\s_{T_{R}}$ is normal in $\Tt_{0}$.
	%\item Among split subgroups of $T_{0}$, the subgroup $T_{R}$ is maximal. (True for SL3, but might not be true in general. Try SL4.)
	%\item The image of a splitting of $T_{2e-1}$ (if it exists) cannot be normal in $\Tt_{0}$, where we recall $T_{2e-1}$ is defined in \eqref{E:def_of_T_k}.\label{T2e-1NotNorm}
\end{enumerate}
\end{Lem}
\begin{proof}
The lemma follows by comparing the relations \eqref{R3} and \eqref{R4} with \eqref{MR3} and \eqref{MR4}.
\end{proof}
\end{comment}


In the rest of this subsection, we will explicitly describe the structure of the finite quotient $\Tt_0\slash T_R$. As it turns out, the quotient $\Tt_0\slash T_R$ admits a factorization in which each factor has a simple structure, and moreover each factor is preserved by the action of the Weyl group. To show the existence of the factorization, we will use the structure of the quotient $\Ocal^\times\slash R$ viewed as a symmetric bilinear form over $\F_2$ given by the Hilbert symbol, which is mainly developed in Proposition \ref{P:orthogonal_decomposition_of_Hilbert}. We will construct a pseudo-spherical representation of $\Tt_{0}$ by constructing a Weyl-invariant representation for each of the factors in this factorization.

\begin{comment}
\begin{Lem}$ $
	\begin{enumerate}
		\item If $H_{1}\subseteq H_{2}\subseteq G$ are subgroups such that $H_{1}$ does not split, then $H_{2}$ does not split. 
		\item If the subgroup $H\subseteq T$ splits, then $\widetilde{H}$ is abelian.
		\item If the subgroup $H\subseteq T_{0}$ splits and the image of the splitting is normal in $\Tt_{0}$, then $\widetilde{H}\subseteq Z(\Tt_{0})$.
	\end{enumerate}
\end{Lem}
\end{comment} 


Now, recall from Proposition \ref{P:orthogonal_decomposition_of_Hilbert} that we have the orthogonal sum decomposition
\[
\Ocal^\times\slash R=D\oplus D^\perp,
\]
where $\dim_{\F_2}D\leq 2$ and $D^\perp$ consists of isotropic vectors. Also recall we have chosen an orthogonal set $\{u_1,\dots, u_k\}$ of $D$ and a symplectic basis $\{e_1,\dots, e_\ell,f_1,\dots,f_\ell\}$ of $D^\perp$, so that we have
\[
\Ocal^\times\slash R=D\oplus D^\perp=\Span_{\F_2}\{u_1,\dots,u_k\}\oplus\Span_{\F_2}\{e_1,\dots,e_\ell, f_1,\dots,f_\ell\}.
\]

For each $u_i\in \{u_1,\dots, u_k\}$ let us define
\[
\Mt_i:=\la \hta(u_i)\st\alpha\in\Phi\ra\subseteq \Tt_{0}\slash T_{R},
\]
namely the group generated by the $\hta(u_i)$'s viewed in $\Tt_{0}\slash T_{R}$. Note that by \eqref{MR3} we have
\[
\hta(u_i)^2=(u_i, u_i)\hta(u_i^2)=(u_i,u_i)=-1,
\]
because $\hta(u_i^2)=1$ in $\Tt_{0}\slash T_{R}$. Hence $\mu_2\subseteq\Mt_i$. 

For each $\{e_i, f_i\}\subseteq \{e_1,\dots, e_\ell,f_1,\dots,f_\ell\}$, let us define
\[
N_i:=\la \hta(e_i),\hta(f_i)\st\alpha\in\Phi\ra\subseteq \Tt_{0}\slash T_{R},
\]
namely the group generated by the $\hta(e_i)$'s and $\hta(f_i)$'s viewed in $\Tt_{0}\slash T_{R}$. Note that
\[
\hta(e_i)\htb(f_i)=(f_i, e_i^{\la\beta, \alpha\ra})\, \htb(f_i)\hta(e_i)
\]
by \eqref{MR4}. Hence if we choose $\alpha$ and $\beta$ so that $\la\alpha, \beta\ra=-1$ (such $\alpha$ and $\beta$ exist because $G$ is not of type $A_1$ or $D_2$, and $\Phi$ is simply-laced), we have $\hta(e_i)\htb(f_i)=-\htb(f_i)\hta(e_i)$. Hence $\mu_2\subseteq\Nt_i$.

Now, let us define
\[
\Mt:=\Mt_1\cdots\Mt_k\qand
\widetilde{N}:=\widetilde{N}_{1}\cdots \widetilde{N}_{\ell}.
\]
By the orthogonality of the basis vectors $u_i$'s, $e_i$'s and $f_i$'s and relations \eqref{MR3} and \eqref{MR4}, we know that the $\Mt_i$'s and $\Nt_i$'s mutually commute. Hence $\Mt$ is the group generated by all the elements of the form $\hta(u_i)$ viewed in $\Tt_{0}\slash T_{R}$ and $\Nt$ is the one generated by all the elements of the form $\hta(e_i)$ and $\hta(f_i)$ viewed in $\Tt_{0}\slash T_{R}$. We then have
\begin{equation}\label{E:exact_sequence_Mt}
1\longrightarrow(\underbrace{\mu_2\times\cdots\times\mu_2}_{\text{$k+\ell$-times}})^{\Delta}\longrightarrow\Mt_1\times\cdots\times\Mt_k\times \Nt_1\times\cdots\times\Nt_{\ell}\longrightarrow\Mt\cdot \Nt\longrightarrow 1,
\end{equation}
where
\[
(\mu_2\times\cdots\times\mu_2)^{\Delta}=\{(\epsilon_1,\dots,\epsilon_{k+\ell})\st \epsilon_1\cdots\epsilon_{k+\ell}=1\}\subseteq \mu_2\times\cdots\times\mu_2.
\]
Further, the group $\Mt\cdot\Nt$ is the group generated by all the elements $\hta(u_i)$, $\hta(e_i)$ and $\hta(f_i)$, which is the image of $\Tt_{0}$ in $\Tt_{0}\slash T_{R}$. Hence we have
\begin{equation}\label{E:exact_sequence_Tt}
1\longrightarrow T_{R}\longrightarrow\Tt_{0}\longrightarrow \Mt\cdot \widetilde{N}\longrightarrow 1.
\end{equation}



%%%%%%%%%%%%%%%%%%%%%%%%%%%%%%%%%%%%%%%%%%%%%%%%%%%%%%%%%%%%%%%%%%%%%%%

%\subsection{Weyl group action on $\Tt_{0}$}

%%%%%%%%%%%%%%%%%%%%%%%%%%%%%%%%%%%%%%%%%%%%%%%%%%%%%%%%%%%%%%%%%%%%%%%
 
The groups $T_{R}$, $\Mt_i$ and $\Nt_i$ are all invariant under the action of $\Wcal$ as follows.
\begin{Lem}\label{L:action_of_W_on_T(O)}
Let $\wt\in\Wcal$.  Then we have
\begin{enumerate}[(a)]
\item $\wt T_{R}\wt^{-1}=T_{R}$, so that the conjugation action of $\Wcal$ on $\Tt_{0}$ descends to $\Mt\cdot\Nt=\Tt_{0}\slash T_{R}$;
\item $\wt\Mt_i\wt^{-1}=\Mt_i$ and $\wt\Nt_i\wt^{-1}=\Nt_i$.
\end{enumerate}
\end{Lem}
\begin{proof}
We may assume that $\wt$ is of the form $\wta:=\wt_{\alpha}(1)$ for some $\alpha\in\Phi$. For each generator $\htb(t)$ with $t\in\Ocal^\times$, we have
\begin{align*}
\wta\htb(t)\wta^{-1}&=\wta(1)\htb(t)\wta(-1)\\
&=(t, c)\cdot\htt_{w_{\alpha}(\beta)}(t),
\end{align*}
where $c=c(\alpha, \beta)\in\{\pm 1\}$. Hence, if $t\in R$, then $(t, c)=1$, so $\wt T_{R}\wt^{-1}=T_{R}$. This proves (a). If $\htb(t)$ is a generator of $\Mt_i$ or $\Nt_i$, then $\wta\htb(t)\wta^{-1}=(t, c)\cdot\htt_{\gamma}(t)$ is also in $\Mt_i$ or $\Nt_i$ because $\htt_{\gamma}(t)$ is also a generator of $\Mt_i$ or $\Nt_i$. This proves (b).
\end{proof}

The lemma implies that for each $\wt\in\Wcal$ we have the isomorphisms $\Int(\wt):\Mt_i\to\Mt_i$ and $\Int(\wt):\Nt_i\to\Nt_i$ given by conjugation, which gives rise to the following commutative diagram
\begin{equation}\label{D:commutativity_of_Weyl_action}
\begin{tikzcd}
\Mt_1\times\cdots\times\Mt_k\times\Nt_1\cdots\times\Nt_\ell \ar[d, "\Int(\wt)"']\ar[r]&\Mt\cdot\Nt\ar[d, "\Int(\wt)"]\\
\Mt_1\times\cdots\times\Mt_k\times\Nt_1\cdots\times\Nt_\ell\ar[r]&\Mt\cdot\Nt
\rlap{\, ,}
\end{tikzcd}
\end{equation}
where the horizontal arrows are the surjections of \eqref{E:exact_sequence_Mt} and the vertical arrows are the conjugation by $\wt$.





%%%%%%%%%%%%%%%%%%%%%%%%%%%%%%%%%%%%%%%%%%%%%%%%%%%%%%%%%%%%%%%%%%%%%%%

\subsection{Definition}

%%%%%%%%%%%%%%%%%%%%%%%%%%%%%%%%%%%%%%%%%%%%%%%%%%%%%%%%%%%%%%%%%%%%%%%

Let us now define what a pseudo-spherical representation is:
\begin{Def}\label{Def:PSTRep}
	A {\it pseudo-spherical representation} of $\Tt_{0}$ is an irreducible genuine representation of $\Tt_{0}$ which is trivial on $T_{R}$. 
\end{Def}

\begin{Rmk}This terminology is inspired by \cite{ABPTV,LS10a,LS10b}.\end{Rmk}

\begin{comment}
\begin{Rmk}
Let us explain why we call such representations ``pseudo-spherical". For the linear torus $T(F)$, an irreducible representation (\ie a character) of $T(F)$ is spherical (or unramified) if it is trivial on the compact torus $T(\Ocal)$. Hence one might as well call the trivial representation of $T(\Ocal)$ a spherical representation of $T(\Ocal)$. We would like to have its analog for the covering $\Tt_{0}$. Certainly, we are interested in {\it genuine} representations for the covering case. So the covering analog has to be a genuine representation which is as close to the trivial representation as possible. Now, the group $T_{R}$ is the largest subgroup of $T(\Ocal)$ that splits \textcolor{red}{(Not true, need to fix. Maybe just need to add $W$-invar?} to the covering group. Hence the pseudo-spherical representations are the ones that are as close to the trivial representation as possible. 
\end{Rmk}
\end{comment}

%For a subgroup $H\subseteq \Ocal^{\times}$ let $\Tt_{H}\subseteq \Tt_{0}$ be the subgroup generated by the elements of the form $\hta(u)$, where $u\in H$.

Note that the pseudo-spherical representations of $\Tt_{0}$ remain irreducible when restricted to the first congruence subgroup.

\begin{Lem}\label{RedtoU1}
	The natural inclusion $\Tt_{1}\subseteq \Tt_{0}$ descends to a group isomorphism 
	\begin{equation*}
		\Tt_{1}/T_{R}\cap \Tt_{1}\xrightarrow{\;\sim\;}\Tt_{0}/T_{R}.
	\end{equation*}
	Consequently, if $\tau$ is a pseudo-spherical representation of $\Tt_{0}$, then $\tau$ remains irreducible when restricted to $\Tt_{1}$.
\end{Lem}

\begin{proof}
	Recall that $\Ocal^{\times}=\Ocal^{\times 2}(1+\varpi\Ocal)$ by Lemma \ref{L:square_root_exits_in_redidue_field}. Thus the natural inclusion $1+\varpi\Ocal\subseteq\Ocal^{\times}$ descends to an isomorphism $(1+\varpi\Ocal)/R\cap (1+\varpi\Ocal)\xrightarrow{\sim}\Ocal^{\times}/R$. The result follows.
\end{proof}



%%%%%%%%%%%%%%%%%%%%%%%%%%%%%%%%%%%%%%%%%%%%%%%%%%%%%%%%%%%%%%%%%%%%%%%

\subsection{A general construction}

%%%%%%%%%%%%%%%%%%%%%%%%%%%%%%%%%%%%%%%%%%%%%%%%%%%%%%%%%%%%%%%%%%%%%%%

With the short exact sequence \eqref{E:exact_sequence_Tt} in mind, to construct a pseudo-spherical representation of $\Tt_{0}$ we have to construct a genuine irreducible representation for each $\Mt_i$ and $\Nt_i$. For this purpose, let us start with the following general lemma.
\begin{Lem}\label{L:rep_of_2-step_unipotent_group}
Let $H\subseteq G$ be a compact abelian subgroup, and $\Ht$ its preimage in $\Gt$, so that we have
\[
1\longrightarrow\mu_2\longrightarrow\Ht\longrightarrow H\longrightarrow 1.
\]
Let $\Irr_{gen}(\Ht)$ be the set of (equivalence classes of) finite dimensional irreducible genuine representations of $\Ht$ and $\Irr_{gen}(Z(\Ht))$ the set of genuine characters on the center $Z(\Ht)$. Then there is a bijection
\[
\Irr_{gen}(Z(\Ht))\xrightarrow{\;\sim\;}\Irr_{gen}(\Ht),\quad \chi\mapsto \pi(\chi),
\]
where $\pi(\chi)$ is the unique constituent of the isotypic representation $\Ind_{Z(\Ht)}^{\Ht}\chi$. 

Moreover, we have
\[
\Ind_{Z(\Ht)}^{\Ht}\chi=n\pi(\chi),
\]
where
\[
n=\dim\pi(\chi)=\left|\Ht\slash Z(\Ht)\right|^{\frac{1}{2}}.
\]
\end{Lem}
\begin{proof}
This is essentially \cite[Proposition 2.2, p.706]{ABPTV} except that the group $H$ there is an abelian Lie group with $H\slash H^\circ$ finite. Their proof works for our setting. But, to show $\Ind_{Z(\Ht)}^{\Ht}\chi$ is completely reducible, one has to use that $\Ht$ is compact. The details are left to the reader.
\end{proof}

Also if $\Wcal$ acts on $\Ht$ through automorphisms, then $\Wcal$ also acts on $Z(\Ht)$. Then for each $\wt\in\Wcal$ we have
\[
^{\wt}\left(\Ind_{Z(\Ht)}^{\Ht}\chi\right)=\Ind_{Z(\Ht)}^{\Ht}({^{\wt}\chi}).
\]
Hence the bijection of the lemma restricts to the bijection
\begin{equation}\label{E:bijection_Weyl_invariant_rep}
(\Irr_{gen}Z(\Ht))^{\Wcal}\xrightarrow{\;\sim\;}(\Irr_{gen}\Ht)^{\Wcal},
\end{equation}
where $(\Irr_{gen}Z(\Ht))^{\Wcal}$ is the set of all $\Wcal$-invariant genuine characters on $Z(\Ht)$ and $(\Irr_{gen}\Ht)^{\Wcal}$ is the set of (equivalence classes of) $\Wcal$-invariant genuine irreducible representations of $\Ht$.



%%%%%%%%%%%%%%%%%%%%%%%%%%%%%%%%%%%%%%%%%%%%%%%%%%%%%%%%%%%%%%%%%%%%%%%

\subsection{Representations of $\Mt_i$}

%%%%%%%%%%%%%%%%%%%%%%%%%%%%%%%%%%%%%%%%%%%%%%%%%%%%%%%%%%%%%%%%%%%%%%%

Let us construct (finite dimensional) irreducible genuine representations of $\Mt_i$ and show that they are all $\Wcal$-invariant. For this purpose, we describe the center $Z(\Mt_i)$ concretely as follows.
\begin{Lem}
We have a bijection
\[
\mu_2\times \Yt\slash 2Y\longrightarrow Z(\Mt_i)
\]
given by
\[
(\epsilon,\, \sum_{\alpha\in\Delta}c_\alpha\alphac)\mapsto \epsilon\prod_{\alpha}\hta(u_i^{c_\alpha}),
\]
where the product is independent of the ordering. (Note that this does not have to be a group isomorphism.)
\end{Lem}
\begin{proof}
In this proof, we omit the subscript $i$ for $\Mt_i$ and $u_i$, so we write $\Mt=\Mt_i$ and $u=u_i$. Also we write $\sum_{\alpha\in\Delta}c_\alpha\alphac=\sum_{i}c_i\alphac_i$, so that $\prod_{\alpha}\hta(u_i^{c_\alpha})$ is $\prod_{i}\htt_{\alpha_i}(u^{c_i})$.

Note that since $\htt_{\alpha_i}(u^2)=1$, the map is independent of the representative $\sum_{i}c_i\alphac_i$.

Let us show that $\epsilon\prod_{i}\hta(u^{c_i})$ is indeed in $Z(\Mt)$. Let $\htt_\beta(u)\in\Mt$, where $\beta\in\Phi$. Then for each $\alphac_i$ we have
\[
\htt_{\alpha_i}(u^{c_i})\htt_\beta(u)=(u, u^{\la \beta, \alphac_i\ra})\htt_\beta(u)\htt_{\alpha_i}(u^{c_i})=\htt_\beta(u)\htt_{\alpha_i}(u^{c_i})
\]
by \eqref{MR4}, where we used $\la \beta, \alphac_i\ra\in 2\Z$. Hence
\[
\left(\prod_{i}\htt_{\alpha_i}(u^{c_i})\right)\htt_\beta(u)=\htt_\beta(u)\left(\prod_{i}\htt_{\alpha_i}(u^{c_i})\right),
\]
so $\epsilon\prod_{i}\hta(u^{c_i})$ is indeed in $Z(\Mt)$.

To show that the map is surjective, note that each element in $\Mt$ is (not necessarily uniquely) written as
\[
\epsilon\prod_i\htt(u^{d_i})
\]
for some $\sum_id_i\alphac_i\in Y$. If this element is in the center $Z(\Mt)$, it has to commute with all the $\htt_\beta(u)$. But by the same computation as above, we have
\[
\prod_i\htt(u^{d_i})\cdot\htt_\beta(u)=(u, u^{\la\beta,\; \sum_id_i\alphac_i\ra})\,\htt_\beta(u) \cdot \prod_i\htt(u^{d_i}). 
\]
Hence we must have $\sum_id_i\alphac_i\in \Yt$ because $(u, u)=-1$. Hence the given map is surjective.

To show it is injective, compose the map with $Z(\Mt)\slash\mu_2$. This is a homomorphism whose kernel is $\mu_2$. This shows that the given map is injective.
\end{proof}

Once we can describe the center $Z(\Mt_i)$, we can show
\begin{Lem}
The group $\Wcal$ acts trivially on $Z(\Mt_i)$.
\end{Lem}
\begin{proof}
In this proof, we omit the subscript $i$ for $\Mt_i$ and $u_i$, so we write $\Mt=\Mt_i$ and $u=u_i$. Also we write $\sum_{\alpha\in\Delta}c_\alpha\alphac=\sum_{i}c_i\alphac_i$. 

To show the lemma, it suffices to show that for each simple root $\alpha_j$ we have
\[
\wt_{\alpha_j}(1) \prod_{i}\htt_{\alpha_i}(u^{c_i}) \wt_{\alpha_j}(-1)=\prod_{i}\htt_{\alpha_i}(u^{c_i}).
\]
To work this out, let us recall from \eqref{MR5} and \eqref{MR5'} that we have
\[
\wt_{\alpha}(1)\htt_{\beta}(t)\wt_{\alpha}(-1)=\begin{cases}\htt_{\alpha}(t^{-1})&\text{if $\beta=\alpha$};\\
\htt_{\beta}(t)\htt_{\alpha}(t)&\text{if $\la\alpha,\beta\ra=-1$};\\
\htt_{\beta}(t)&\text{if $\la\alpha, \beta\ra=0$},
\end{cases}
\]
for all $t\in F^\times$. Also noting $\htt_\alpha(u^2)=1\pmod{T_{R}}$, we always have
\[
\wt_{\alpha}(1)\htt_{\alpha}(u^c)\wt_{\alpha}(-1)=\htt_{\alpha}(u^c)
\]
for any $c\in\Z$ and any $\alpha\in\Phi$. Also note that if $c=d\pmod{2}$ then $\htt_{\alpha}(u^c)=\htt_{\alpha}(u^d)\mod{T_R}$.

Assume that $\alpha_j$ is adjacent to exactly one simple root $\alpha_k$, namely $\alpha_j$ is an endpoint of the Dynkin diagram. We then know from Lemma \ref{L:parity_of_adjacent_root} that $c_k\in 2\Z$. We can then write
\begin{align*}
&\wt_{\alpha_j}(1) \prod_{i}\htt_{\alpha_i}(u^{c_i}) \wt_{\alpha_j}(-1)\\
&=\left(\wt_{\alpha_j}(1)\htt_{\alpha_k}(u^{c_k})\wt_{\alpha_j}(-1)\right)\left(\wt_{\alpha_j}(1)\htt_{\alpha_j}(u^{c_j})\wt_{\alpha_j}(-1)\right) \prod_{i\neq j, k}\htt_{\alpha_i}(u^{c_i})\\
&=\htt_{\alpha_k}(u^{c_k})\htt_{\alpha_j}(u^{c_k})\htt_{\alpha_j}(u^{c_j})\prod_{i\neq j, k}\htt_{\alpha_i}(u^{c_i})\\
&=\htt_{\alpha_k}(u^{c_k})\htt_{\alpha_j}(u^{c_j})\prod_{i\neq j, k}\htt_{\alpha_i}(u^{c_i})\\
&=\prod_{i}\htt_{\alpha_i}(u^{c_i}),
\end{align*}
where we used $\htt_{\alpha_j}(u^{c_k})=1$.

Next, assume that $\alpha_j$ is adjacent to exactly two roots $\alpha_k$ and $\alpha_\ell$. By the analogous computation, we have
\begin{align*}
&\wt_{\alpha_j}(1) \prod_{i}\htt_{\alpha_i}(u^{c_\alpha}) \wt_{\alpha_j}(-1)\\
&=\htt_{\alpha_k}(u^{c_k})\htt_{\alpha_j}(u^{c_k})\htt_{\alpha_\ell}(u^{c_\ell})\htt_{\alpha_j}(u^{c_\ell})\htt_{\alpha_j}(u^{c_j})\prod_{i\neq j, k, \ell}\htt_{\alpha_i}(u^{c_i}).
\end{align*}
Here, by moving round $\htt$'s, we can further write
\begin{align*}
&\htt_{\alpha_k}(u^{c_k})\htt_{\alpha_j}(u^{c_k})\htt_{\alpha_\ell}(u^{c_\ell})\htt_{\alpha_j}(u^{c_\ell})\htt_{\alpha_j}(u^{c_j})\\
&=(u^{c_k},u^{c_{\ell}})\,\htt_{\alpha_k}(u^{c_k})\htt_{\alpha_\ell}(u^{c_\ell})\htt_{\alpha_j}(u^{c_k})\htt_{\alpha_j}(u^{c_\ell})\htt_{\alpha_j}(u^{c_j})\\
&=(u^{c_k},u^{c_{\ell}})\,(u^{c_k},u^{c_{\ell}})\,\htt_{\alpha_k}(u^{c_k})\htt_{\alpha_\ell}(u^{c_\ell})\htt_{\alpha_j}(u^{c_k+c_\ell})\htt_{\alpha_j}(u^{c_j})\\
&=(u^{c_k+c_\ell}, u^{c_j})\,\htt_{\alpha_k}(u^{c_k})\htt_{\alpha_\ell}(u^{c_\ell})\htt_{\alpha_j}(u^{c_k+c_\ell+c_j})\\\
&=(u,u)^{c_kc_j+c_\ell c_j}\,\htt_{\alpha_k}(u^{c_k})\htt_{\alpha_\ell}(u^{c_\ell})\htt_{\alpha_j}(u^{c_k+c_\ell+c_j})\\
&=\htt_{\alpha_k}(u^{c_k})\htt_{\alpha_\ell}(u^{c_\ell})\htt_{\alpha_j}(u^{c_j}),
\end{align*}
where for the last step we used that $c_kc_j, c_\ell c_j\in 2\Z$ by Corollary \ref{C:parity_of_adjacent_root3}, and $c_k+c_\ell\in 2\Z$ by Lemma \ref{L:parity_of_adjacent_root}.

Finally, if $\alpha_k$ has three adjacent roots, then the calculation is essentially the same and left to the reader.
\end{proof}

We then have
\begin{Prop}\label{P:W-invariant_rep_of_Mt}
Every irreducible genuine representation of $\Mt_i$ is $\Wcal$-invariant and 
\[
\left|\Irr_{gen}Z(\Mt_i)\right|=\left|\Yt\slash 2Y\right|.
\]
Further, for each $\pi\in \Irr_{gen}(\Mt_i)$, we have
\[
\dim\pi=\left|\Mt_i\slash Z(\Mt_i)\right|^{\frac{1}{2}}=\left|Y\slash\Yt\right|^{\frac{1}{2}},
\]
where the cardinality $\left|Y\slash\Yt\right|$ is listed in Table \ref{T:table}.
\end{Prop}
\begin{proof}
This follows from Lemma \ref{L:rep_of_2-step_unipotent_group} and \eqref{E:bijection_Weyl_invariant_rep}.
\end{proof}



\begin{Rmk}
When $F=\R$, the above proposition is proven in \cite{ABPTV}, where the corresponding group is denoted by $\Mt$. One can actually show that our $\Mt_i$ is isomorphic to their $\Mt$. In \cite{Karasiewicz}, the corresponding group is denoted by $\Tt^{\diamond}$, and a $\Wcal$-invariant representation is constructed by constructing an isomorphism between $\Mt$ and $\Tt^{\diamond}$. One can readily verify that our $\Mt_i$ is isomorphic to $\Tt^{\diamond}$ of \cite{Karasiewicz} via the map $\Mt_i\to\Tt^{\diamond}$, $\hta(u_i)\mapsto\hta(-1)$.
\end{Rmk}


%%%%%%%%%%%%%%%%%%%%%%%%%%%%%%%%%%%%%%%%%%%%%%%%%%%%%%%%%%%%%%%%%%%%%%%

\subsection{Representations of $\Nt_i$}

%%%%%%%%%%%%%%%%%%%%%%%%%%%%%%%%%%%%%%%%%%%%%%%%%%%%%%%%%%%%%%%%%%%%%%%

We can construct pseudo-spherical representations of $\Nt_i$ and show that they are all $\Wcal$-invariant essentially in the same way.

Let us first mention
\begin{Lem}
The map
\[
\mu_2\times \Yt\slash 2Y\times\Yt\slash 2Y\longrightarrow Z(\Nt_i)
\]
given by
\[
(\epsilon, y_1, y_2)\mapsto \epsilon\, \yt_1(e_i)\,\yt_2(f_i)
\]
is an isomorphism. Further, the group $\Wcal$ acts trivially on $Z(\Nt_i)$.
\end{Lem}
\begin{proof}
Let us omit the subscript $i$ from the notation for $\Nt_i$, so $\{e, f\}$ is a symplectic basis of a 2-dimensional subspace in $\Ocal^\times \slash R$ such that $(e, f)=-1$ and $(e, e)=(f, f)=1$, and $\Nt$ is the group generated by $\hta(e)$ and $\htb(f)$ for all $\alpha\in \Phi$. Now, since $(e, e)=(f, f)=1$, we know from \eqref{MR4} that $\hta(e)\htb(e)=\htb(e)\hta(e)$ and $\hta(f)\htb(f)=\htb(f)\hta(f)$ for all $\alpha, \beta\in\Phi$. Hence each element in $\Nt$ is uniquely written as
\[
\epsilon\, \yt_1(e)\,\yt_2(f)=\epsilon\cdot \prod_{i}\htt_{\alpha_i}(e^{c_i})\cdot \prod_{i}\htt_{\alpha_i}(f^{d_i}),
\]
where $y_1=\sum_{i}c_i\alphac_i$ and $y_2=\sum_{i}d_i\alphac_i$.

Assume $\epsilon\, \yt_1(e)\,\yt_2(f)\in Z(\Nt)$. Then it commutes with all the elements of the form $\hta(f)$. But by (MR4) we know that
\[
\epsilon\, \yt_1(e)\,\yt_2(f)\,\hta(f)=(e, f^{\la \alpha, y_1\ra})\hta(f)\cdot \epsilon\, \yt_1(e)\,\yt_2(f).
\]
Hence we must have $\la \alpha, y_1\ra\in2\Z$ for all $\alpha\in\Phi$, which implies $y_1\in\Yt$. Similarly, we have $y_2\in\Yt$. If $y_1, y_2\in\Yt$ then we already know that both $\yt_1(e)$ and $\yt_2(f)$ are in $Z(\Tt)$ and hence in $Z(\Nt)$. Hence we have the surjective map
\[
\mu_2\times \Yt\times\Yt\longrightarrow Z(\Nt),\quad (\epsilon, y_1, y_2)\mapsto \epsilon\, \yt_1(e)\,\yt_2(f).
\]
To show this is a homomorphism, notice that for each $(\epsilon, y_1, y_2)$ and $(\epsilon', z_1, z_2)$ we know
\[
\epsilon\, \yt_1(e)\,\yt_2(f)\cdot \epsilon'\, \zt_1(e)\,\zt_2(f)=\epsilon\epsilon'\, \yt_1(e)\zt_1(e)\,\yt_2(f)\zt_2(f),
\]
because $\yt_2(f)\in Z(\Nt)$. We then have $\yt_1(e)\zt_1(e)=\widetilde{y_1+z_1}(e)$ because $(e,e)=1$ and so all the $\hta(e)$'s commute with each other, and similarly we have $\yt_2(f)\zt_2(f)=\widetilde{y_2+z_2}(f)$.

We need to show that the kernel of this map is $\{1\}\times 2Y\times 2Y$. Let $(1, y_1, y_2)$ be in the kernel. Suppose $y_1=\sum_{i}c_i\alphac_i$. Then $\prod_{i}\htt_{\alpha_i}(e^{c_i})=1$, which implies $\htt_{\alpha_i}(e^{c_i})=1$ for each $i$. Then we must have $e^{c_i}\in R$, which implies $c_i\in 2\Z$ because $e\notin R$. Hence $y_1\in 2Y$. Similarly, we have $y_2\in 2Y$. Hence the kernel is $\{1\}\times 2Y\times 2Y$. 

It remains to show that $\Wcal$ acts trivially on $Z(\Nt)$. This can be proven in the same way as $\Mt_i$; the proof is indeed simpler because we always have $(e, e)=(f, f)=1$, and hence all the Hilbert symbols in the calculation for the $\Mt_i$ case will automatically become $1$. The detail is left to the reader.
\end{proof}

We then have
\begin{Prop}\label{P:W-invariant_rep_of_Nt}
Every irreducible genuine representation of $\Nt_i$ is $\Wcal$-invariant and 
\[
\left|\Irr_{gen}Z(\Nt_i)\right|=\left|\Yt\slash 2Y\right|^2.
\]
Further, for each $\pi\in \Irr_{gen}(\Nt_i)$, we have
\[
\dim\pi=\left|\Nt_i\slash Z(\Nt_i)\right|=\left|Y\slash\Yt\right|,
\]
where the cardinality $\left|Y\slash\Yt\right|$ is listed in Table \ref{T:table}.
\end{Prop}
\begin{proof}
This follows from the above lemma, Lemma \ref{L:rep_of_2-step_unipotent_group}, and \eqref{E:bijection_Weyl_invariant_rep}.
\end{proof}

%%%%%%%%%%%%%%%%%%%%%%%%%%%%%%%%%%%%%%%%%%%%%%%%%%%%%%%%%%%%%%%%%%%%%%%

\subsection{Construction of pseudo-spherical representations}\label{SSPSpherT}

%%%%%%%%%%%%%%%%%%%%%%%%%%%%%%%%%%%%%%%%%%%%%%%%%%%%%%%%%%%%%%%%%%%%%%%

By putting together genuine representations of $\Mt_i$'s and $\Nt_i$'s, one can construct a genuine representation of $\Mt\cdot\Nt$ as follows. Let $\delta_i$ be a genuine representation of $\Mt_i$ and $\tau_i$ one of $\Nt_i$. Then consider the tensor product representation
\[
\delta_1\otimes\cdots\delta_k\otimes\tau_1\otimes\cdots\otimes \tau_{\ell}
\]
of the direct product
\[
\Mt_1\times\cdots\times\Mt_k\times \Nt_1\times\cdots\times\Nt_{\ell}.
\]
Since all the $\delta_i$'s and $\tau_i$'s are genuine, this tensor product representation is trivial on the kernel of the surjection 
\[
\Mt_1\times\cdots\times\Mt_k\times \Nt_1\times\cdots\times\Nt_{\ell}\longrightarrow \Mt\cdot\Nt
\]
of \eqref{E:exact_sequence_Mt}, descending to a representation $\tau$ of $\Mt\cdot\Nt$. Conversely, every genuine irreducible representation of $\Mt\cdot\Nt$ arises in this way.


Assume $\delta_i$ and $\tau_i$ are all irreducible, so that they are all $\Wcal$-invariant. Then $\tau$ is also $\Wcal$-invariant by \eqref{D:commutativity_of_Weyl_action}. By pulling back $\tau$ to a representation of $\Tt_{0}$ via the surjection
\[
\Tt_{0}\longrightarrow \Mt\cdot\Nt=\Tt_{0}\slash T_{R},
\]
we obtain an irreducible genuine representation of $\Tt_{0}$ trivial on $T_{R}$ (\ie a pseudo-spherical representation of $\Tt_{0}$). Since $T_{R}$ is invariant under $\Wcal$, the pseudo-spherical representation is $\Wcal$-invariant. We denote this representation of $\Tt_{0}$ by $\tau$. 

To summarize, we have proven
\begin{Thm}\label{T:PseudoT0}
There exists a pseudo-spherical representation of $\Tt_{0}$. Moreover, every pseudo-spherical representation $\tau$ is $\Wcal$-invariant, namely ${^{\wt}\tau}\cong\tau$ for all $\wt\in\Wcal$.
\end{Thm}


%%%%%%%%%%%%%%%%%%%%%%%%%%%%%%%%%%%%%%%%%%%%%%%%%%%%%%%%%%%%%%%%%%%%%%%

\subsection{Dimension and number of pseudo-spherical representations}

%%%%%%%%%%%%%%%%%%%%%%%%%%%%%%%%%%%%%%%%%%%%%%%%%%%%%%%%%%%%%%%%%%%%%%%

The dimension of each pseudo-spherical representation $\tau$ is computed as follows.
\begin{Prop}
Let $\tau$ be a pseudo-spherical representation. Then
\[
\dim\tau=\left|Y\slash \Yt\right|^{\frac{ef}{2}},
\]
where the cardinality $\left|Y\slash \Yt\right|$ is as in Table \ref{T:table}.
\end{Prop}
\begin{proof}
From Propositions \ref{P:W-invariant_rep_of_Mt} and \ref{P:W-invariant_rep_of_Nt}, we already know
\[
\dim\delta_i=\left|Y\slash \Yt\right|^{\frac{1}{2}}\qand \dim\tau_i=\left|Y\slash \Yt\right|
\]
for each $\delta_i\in\Irr_{gen}(\Mt_i)$ and $\tau_i\in\Irr_{gen}(\Nt_i)$. Hence to compute $\dim\tau$, we have to compute the number of $\Mt_i$'s and $\Nt_i$'s.

We know from Proposition \ref{P:property_Hilber_symbol_not_in_Appendix} that
\[
[\Ocal^{\times}: R]=2^{ef}.
\]
So, if we let $m$ be the number of $\Mt_i$'s and $n$ the number of $\Nt_i$'s, then $ef=m+2n$. Hence 
\[
\dim\tau=(\dim\delta_i)^{m}\cdot(\dim\tau_i)^n=\left|Y\slash \Yt\right|^{\frac{m+2n}{2}}=\left|Y\slash \Yt\right|^{\frac{ef}{2}}.
\]
\end{proof}

Similarly, we have
\begin{Prop}
The number of the pseudo-spherical representations of $\Tt_{0}$ is $\left|\Yt\slash 2Y\right|^{ef}$.
\end{Prop}
\begin{proof}
We know that 
\[
\left|\Irr_{gen}(\Mt_i)\right|=\left|\Yt\slash 2Y\right|\qand \left|\Irr_{gen}(\Nt_i)\right|=\left|\Yt\slash 2Y\right|^2.
\]
One can then prove the proposition in the same way as the previous proposition.
\end{proof}

The dimension and the number of the pseudo-spherical representations of $\Tt_{0}$ is summarized in Table \ref{T:table2}, where $A_{1}$ and $D_2$ are excluded.

\begin{table}[h]
\caption{dimension and number of pseudo-spherical representations}\label{T:table2}
\centering
\renewcommand{\arraystretch}{1.5}
\begin{tabular}{ccccc}
\hline
Type & $\left| \Yt\slash 2Y\right|$ & $\left|Y\slash \Yt\right|$ &dimension &number \\\hline\hline
$A_{2r+1}$ & 2 & $2^{2r}$& $2^{efr}$ & $2^{ef}$\\
$A_{2r}$ &  1 & $2^{2r}$ & $2^{efr}$ & 1\\\hline
$D_{2r+1}$  & 2 & $2^{2r}$ & $2^{efr}$ & $2^{ef}$\\
$D_{2r}$ & 4  & $2^{2r-2}$ & $2^{ef(r-1)}$ & $4^{ef}$\\\hline
 $E_6$ & 1 & $2^6$ & $2^{3ef}$ & 1\\\hline
 $E_7$ & 2 & $2^6$ & $2^{3ef}$ & $2^{ef}$ \\\hline
 $E_8$ & 1 & $2^8$ & $2^{4ef}$ & 1\\\hline
\end{tabular}
\end{table}


\subsection{Branching}

For later purposes, we need the following.

\begin{Prop}\label{PseudoSphericalBranch}
	Suppose that $\Phi$ is a root system of type $A_r$ or of type $D_r$. Let $\Phi^{\prime}\subseteq\Phi$ be a subroot system of rank $r-1$ and of the same type as $\Phi$. Let $\Tt_{0}^{\prime}\subseteq \Tt_{0}$ be the normal subgroup generated by the elements $\htt_{\alpha}(u)$, where $u\in \Ocal$ and $\alpha\in\Phi^{\prime}$. 
    \begin{enumerate}[(1)]
    \item For a pseudo-spherical representation $\tau$ of $\Tt_{0}$, the restriction $\tau|_{\Tt_0^{\prime}}$ to $\Tt_{0}^{\prime}$ is a multiplicity-free direct sum of pseudo-spherical representations of $\Tt_{0}^{\prime}$.
    \item If $\tau'$ is a pseudo-spherical representation of $\Tt_0'$, then there exists a pseudo-spherical representation $\tau$ of $\Tt_0$ such that $\tau'\subseteq\tau|_{\Tt_0}$.
    \end{enumerate}
\end{Prop}

\begin{proof}
	The restriction $\tau|_{\Tt_0^{\prime}}$ is a direct sum of pseudo-spherical representations of $\Tt_{0}^{\prime}$ {\it by definition} because $T_R'\subseteq T_R$. Now we will break the proof of (1) up into cases.
	
	First suppose that $\Phi$ is of type $A_{2r+1}$ or $D_{2r}$. From the dimension column of Table \ref{T:table2}, we see that this restriction to $\Tt_{0}^{\prime}$ must remain irreducible. Thus the result follows in this case.
	
	Second, suppose that $\Phi$ is of type $A_{2r}$ or $D_{2r+1}$. Let $\Delta\subseteq\Phi$ be a set of simple roots. Then $\Delta^{\prime}:=\Delta\cap\Phi^{\prime}$ is a set of simple roots in $\Phi^{\prime}$. Let $\alpha\in \Delta\smallsetminus\Delta^{\prime}$. Let $\Tt_{\alpha,0}\subseteq \Tt_{0}$ be the subgroup generated by the elements of the form $\htt_{\alpha}(u)$, where $u\in \Ocal^{\times}$. Note that $\Tt_{\alpha,0}$ is a normal subgroup of $\Tt_{0}$. By definition the restriction of $\tau$ to $\Tt_{\alpha,0}$ factors through the finite abelian group $\Tt_{\alpha,0}/(\Tt_{\alpha,0}\cap T_R)$, which has size $2[\Ocal^{\times}:R]=2^{ef+1}$.
	
	Let $\tau^{\prime}$ be a pseudo-spherical representation of $\Tt_{0}^{\prime}$. By a direct computation with the central character it follows that $\tau^{\prime}\cong \,^{\htt_{\alpha}(u)}\tau^{\prime}$ if and only if $u\in R$. From this it follows that the restriction $\tau|_{\Tt_{0}^{\prime}}$ must contain at least $[\Ocal:R]=2^{ef}$ distinct pseudo-spherical $\Tt_{0}^{\prime}$-modules. By comparing the dimensions in Table \ref{T:table2} it follows that the restriction of $\tau$ to $\Tt_{0}^{\prime}$ is a multiplicity-free direct sum of pseudo-spherical representations of $\Tt_{0}^{\prime}$.
    Item (1) is proven. 
    
    To prove (2), let $\tau'$ be a pseudo-spherical representation of $\Tt'_0$. Since $T_R'=\Tt_0'\cap T_R$ acts trivially and $T_R$ commutes with $\Tt_0'$ pointwise, we can extend $\tau'$ to a representation of $T_R\Tt_0'$ by letting $T_R$ act trivially. Consider $\Ind_{T_R\Tt_0'}^{\Tt_0}\tau'$. Since $T_R$ is in the center of $\Tt_0$, one can see that $T_R$ is in the kernel of this induced representation. Hence any irreducible constituent of this induced representation is pseudo-spherical and contains $\tau'$.
\end{proof}

%%%%%%%%%%%%%%%%%%%%%%%%%%%%%%%%%%%%%%%%%%%%%%%%%%%%%%%%%%%%%%%%%%%%%%%

\subsection{Square roots of $\htt_{\alpha}(-1)$}\label{SSSqRoot}

%%%%%%%%%%%%%%%%%%%%%%%%%%%%%%%%%%%%%%%%%%%%%%%%%%%%%%%%%%%%%%%%%%%%%%%

In this subsection we record a few identities in the group algebra $\mathbb{C}[\Tt_{0}]$ and their relation to each pseudo-spherical representation $\tau$ of $\Tt_{0}$. In particular, we will compute a square root of $\tau(\htt_{\alpha}(-1))$. Later these will be important to our investigation of representations of $\It$. Also in this subsection, we set
\begin{equation}\label{E:U_k_def_subset_of_Ocal}
U_k:=1+\varpi^k\Ocal
\end{equation}
for $k\in\Z_{\geq 0}$ with the convention that $U_0=\Ocal$.

Let $\alpha\in\Phi$, $j\in \mathbb{Z}_{\geq 0}$ and $\chi\in \Hom(U_{j}/(U_{j}\cap R),\mathbb{C}^{\times})$. Define
\begin{equation*}
	X_{\alpha,\chi,j}:=\frac{1}{\left|U_j/U_{2e}\right|}\sum_{v\in U_{j}/U_{2e}}\chi(v)\htt_{\alpha}(v).
\end{equation*}
If $\chi$ is such that $\chi(-)=(-, u)$ for some $u\in F^{\times}$, then we write 
\begin{equation}\label{E:X_alpha_u_j_def}
X_{\alpha,u,j}:=X_{\alpha,\chi,j}.
\end{equation}

\begin{Lem}\label{Xpmalpha}
	\begin{equation*}
		X_{-\alpha,u,j}=X_{\alpha,-u,j}
	\end{equation*}
\end{Lem}

\begin{proof}
	This follows because $\htt_{-\alpha}(v)=(v,-1)\htt_{\alpha}(v^{-1})$.
\end{proof}

The next lemma shows that the elements $X_{\alpha,u,j}$ define nontrivial endomorphisms of $\tau$.

\begin{Lem} Let $\tau$ be a pseudo-spherical representation of $\Tt_{0}$. Then
	\begin{equation*}
		\Tr(\tau (X_{\alpha,u,j})) = |(U_{j}\cap R)/U_{j}|\dim\tau\neq 0.
	\end{equation*}
\end{Lem}

\begin{proof}
	For $t\in \Tt_{0}$, we have $\Tr(\tau(t))\neq 0$ if and only if $t\in Z(\Tt_{0})$. (This follows because if $t\in\Tt_0\smallsetminus Z(\Tt_0)$, then there exists $s\in\Tt_0$ such that $sts^{-1}=-t$, so $\Tr(\tau(sts^{-1}))=-\Tr(\tau(t))$.) The result now follows because $\htt_{\alpha}(v)\in Z(\Tt_{0})$ if and only if $v\in R$. (Here we use that $\Phi$ is simply-laced and not of type $A_{1}$ or $D_2$.)
\end{proof}

\begin{Lem} If $j=0$ or $1$, then
	\begin{equation*}
		X_{\alpha,\chi,j}^{2}=\chi(-1)|U_{j}/U_{2e}|^{-1}\htt_{\alpha}(-1)\sum_{v\in U_{j}\cap R/U_{2e}}\htt_{\alpha}(v).
	\end{equation*}
\end{Lem}

\begin{proof} This follows by a direct calculation:
	\begin{align*}
		\left|U_j/U_{2e}\right|^{2}X_{\alpha,\chi,j}^{2}=&\sum_{v,w\in U_{j}/U_{2e}}\chi(vw)\htt_{\alpha}(v)\htt_{\alpha}(w)\\
		=&\sum_{v,w}\chi(vw)(v,w)\htt_{\alpha}(vw)\\
		=&\sum_{v,w}\chi(v)(vw^{-1},w)\htt_{\alpha}(v)\\
		=&\sum_{v}\chi(v)\htt_{\alpha}(v)\sum_{w}(-v,w),
	\end{align*}
where we used $(vw^{-1}, w)=(vw,w)=(v, w)(w, w)=(v, w)(-1, w)=(-v, w)$.

	Because $j=0,1$ we have
	\begin{align*}
		\sum_{v}\chi(v)\htt_{\alpha}(v)\sum_{w}(-v,w)=&\chi(-1)|U_{j}/U_{2e}|\sum_{-v\in U_{j}\cap R/U_{2e}}\htt_{\alpha}(v)\\
		=&\chi(-1)|U_{j}/U_{2e}|\sum_{v\in U_{j}\cap R/U_{2e}}\htt_{\alpha}(-v)\\
		=&\chi(-1)|U_{j}/U_{2e}|\htt_{\alpha}(-1)\sum_{v\in U_{j}\cap R/U_{2e}}\htt_{\alpha}(v).
	\end{align*}
    The lemma follows.
\end{proof}

\begin{Cor}\label{XInvert}
	Let $\alpha\in\Phi$, $\chi\in \Hom(U_{j}/U_{j}\cap R,\mathbb{C}^{\times})$, and $j=0,1$. If $\tau$ is a pseudo-spherical representation, then
	\begin{equation*}
		\tau(X_{\alpha,\chi,j}^{2}) = \chi(-1)|U_{j}\cap R/U_{j}|\tau(\htt_{\alpha}(-1)).
	\end{equation*} 
In particular, $\tau(X_{\alpha,\chi,j})$ is invertible.
\end{Cor}

\begin{Rmk}
For $\alpha\in \Delta$, the element $\tau(X_{\alpha,1,1})$ is proportional to the extension of the pseudo-spherical $\Tt$-module $\tau$ to the element $w_{\alpha}(-1)\in\Wcal$.
\end{Rmk}

%%%%%%%%%%%%%%%%%%%%%%%%%%%%%%%%%%%%%%%%%%%%%%%%%%%%%%%%%%%%%%%%%%%%%%%

\subsection{Representation theory of $\Tt$}\label{TtReps}

%%%%%%%%%%%%%%%%%%%%%%%%%%%%%%%%%%%%%%%%%%%%%%%%%%%%%%%%%%%%%%%%%%%%%%%

Let $\tau$ be a pseudo-spherical representation of $\Tt_0$. In this subsection, we will show that the pair $(\Tt_0, \tau)$ is a type of $\Tt$ in the sense of Bushnell-Kutzko \cite[(4.1) Definition]{BK98}. For this purpose, we will construct a family of irreducible representations of $\Tt$ which contain the pseudo-spherical representation $\tau$.

First, we need to extend $\tau$ to an irreducible representation of $Z(\Tt)\Tt_0$ by adding a central character. Namely, we need a character on $Z(\Tt)$ which agrees with $\tau$ on $Z(\Tt)\cap\Tt_0$. This can be done by using the following.
\begin{Lem}
We have the short exact sequence
\[
1\longrightarrow\mu_2 \xrightarrow{\;\Delta\;}\mu_2\Yt(\varpi)\times \Tt_0\longrightarrow Z(\Tt)\Tt_0\longrightarrow 1,
\]
where the injection embeds $\mu_2$ diagonally and the surjection is given by multiplication.
\end{Lem}
\begin{proof}
One can readily see that $\Yt(\varpi)\subseteq Z(\Tt)$, so we have the map $\mu_2\Yt(\varpi)\times \Tt_0\to Z(\Tt)\Tt_0$. To show that this map is surjective, it suffices to show that each $z\in Z(\Tt)$ is written as $z=\epsilon y(\varpi)t_0$ for some $\epsilon \in\mu_2$, $y\in \Yt$ and $t_0\in\Tt_0$. But we may assume $z$ is of the form $\prod_{i}\htt_{\alpha_i}(t_i)$, where the order of the product is fixed. For each $i$ we can write $t_i=\varpi^{k_i}a_i$ for some $k_i\in\Z$ and $a_i\in\Ocal^\times$, so
\[
\htt_{\alpha_i}(t_i)=(\varpi^{k_i}, a_i)\cdot\htt_{\alpha_i}(\varpi^{k_i})\cdot\htt_{\alpha_i}(a_i).
\]
Let us set $y=\sum_ik_i\alphac_i$. Then by taking the product over $i$, we have
\[
\prod_i\htt_{\alpha_i}(t_i)=\epsilon\cdot\prod_i\htt_{\alpha_i}(\varpi^{k_i})\cdot\prod_i\htt_{\alpha_i}(a_i)=\epsilon\cdot \yt(\varpi)\cdot\prod_i\htt_{\alpha_i}(a_i)
\]
for some $\epsilon\in\mu_2$. Thus it suffices to show $y\in\Yt$.

Note that by Corollary \ref{C:Hilbert_symbol_unique_element_integral_radical} there exists $r\in R$ such that $(r, a_i)=1$ for all $i$ and $(r, \varpi)=-1$. Then for each $\beta\in\Phi$ we have
\[
\htt_{\alpha_i}(\varpi^{k_i})\cdot\htt_\beta(r)=(r, \varpi^{k_i\la \alpha_i, \beta\ra})\cdot \htt_\beta(r)\cdot\htt_{\alpha_i}(\varpi^{k_i}) 
\qand
\htt_{\alpha_i}(a_i)\cdot \htt_\beta(r)=\htt_\beta(r)\cdot\htt_{\alpha_i}(a_i).
\]
Hence for each $i$, we have
\[
\htt_{\alpha_i}(t_i)\cdot\htt_{\beta}(r)=(r, \varpi^{k_i\la \alpha_i, \beta\ra})\cdot \htt_{\beta}(r)\cdot \htt_{\alpha_i}(t_i).
\]
Thus we have
\[
\left(\prod_i\htt_{\alpha_i}(t_i)\right)\cdot\htt_{\beta}(r)=(r, \varpi^{\sum_ik_i\la \alpha_i, \beta\ra})\cdot\htt_{\beta}(r)\cdot\left(\prod_i \htt_{\alpha_i}(t_i)\right).
\]
Note that
\[
(r, \varpi^{\sum_ik_i\la \alpha_i,\; \beta\ra})=(r, \varpi)^{\sum_ik_i\la \alpha_i,\; \beta\ra}=(-1)^{\sum_ik_i\la \alpha_i,\; \beta\ra}=(-1)^{\la \beta,\; y\ra}.
\]
Hence in order for $\prod_{i}\htt_{\alpha_i}(t_i)$ to be in $Z(\Tt)$, we must have $\la \beta,\; y\ra\in2\Z$ for all $\beta\in\Phi$. Therefore we have $y\in\Yt$ by the definition \eqref{E:definition_of_Yt} of $\Yt$. Hence the map is surjective.

To show that the kernel is the diagonal $\mu_2$, let $(\epsilon \yt(\varpi), t)\in (\mu_2\Yt(\varpi))\times \Tt_0$ be such that $\epsilon\yt(\varpi)t=1$. Then $t=\epsilon\yt(\varpi)^{-1}$. But $t\in \Tt_0$ is uniquely written as
\[
t=\epsilon'\prod_i\htt_{\alpha_i}(t_i)\quad \text{for some $t_i\in\Ocal^\times$}.
\]
Hence we must have $y=0$, which gives $t=\epsilon$. Hence the kernel is the diagonal $\mu_2$.
\end{proof}


Let 
\[
\eta:\mu_2\Yt(\varpi)\longrightarrow \C^\times
\]
be a genuine $W$-invariant character trivial on elements of the form $h_{\alpha}(\varpi^{2})$, where $\alpha\in \Phi$. (For example, one can restrict the character of $Z(\Tt)$ constructed in \cite[page 118]{S04} or \cite[\S 7]{GG18}.) Any two such characters differ by a quadratic character of $\Yt/2Y$. Consider the representation $\eta\otimes\tau$ of $\mu_2\Yt(\varpi)\times\Tt_0$. Since both $\eta$ and $\tau$ are genuine, the diagonal $\mu_2$ acts trivially, which allows
\[
\tau_\eta:=\eta\otimes\tau
\]
to descend to a representation of $Z(\Tt)\Tt_0$, which we also denote by the same symbol $\tau_\eta$.

Now, for each unramified character $\chi:T\to\C^\times$, we define
\[
i_{\eta,\tau}(\chi):=\chi\otimes\Ind_{Z(\Tt)\Tt_0}^{\Tt}\tau_{\eta}.
\]
Note the following. First, for any $\alpha\in\Phi$ and $t\in F^{\times}$ the elements $h_{\alpha}(t^{2})\in Z(\Tt)$ act trivially on $i_{\eta,\tau}(\chi)$. Thus $i_{\eta,\tau}(\chi)$ is a distinguished genuine $\Tt$-representation (\cite[Section 6]{GG18}). Furthermore, by choosing $\eta$ and $\tau$ we can construct $|\Yt/2Y|^{1+ef}$ distinguished genuine representations of $\Tt$. There are $|\Yt/2Y|^{2+ef}$ in total (\cite[Theorem 6.6]{GG18}). Second, the two representations $i_{\eta,\tau}(\chi)$ and $i_{\eta^{\prime},\tau^{\prime}}(\chi^{\prime})$ can be related by an unramified twist if and only if $\tau\cong\tau^{\prime}$ and $\eta^{-1}\otimes \eta^{\prime}$ defines a quadratic character of $\Yt/2Y$. Thus the Bernstein block of $i_{\eta,\tau}(\chi)$ depends only on $\tau$.

In what follows, we will show that this induced representation is irreducible, and every irreducible representation of $\Tt$ which contains the pseudo-spherical representation $\tau$ is of this form, which will prove that the pair $(\Tt_0, \tau)$ is a type of $\Tt$.

Let us start with
\begin{Lem}\label{L:tau_twist}
Let $y\in Y$. Then $^{\yt(\varpi)}\tau=\tau$ if and only if $y\in\Yt$.
\end{Lem}
\begin{proof}
We know that there exists $u\in 1+\varpi^{2e}\Ocal$ such that $(u, \varpi)=-1$ by Corollary \ref{C:Hilbert_symbol_unique_element_integral_radical}. But since $1+\varpi^{2e}\Ocal\subseteq R$, we have $\tau(\htt_{\alpha}(u))=1$ for each root $\alpha\in\Phi$. Now for each $y\in Y$
\[
\yt(\varpi)^{-1}\htt_{\alpha}(u)\yt(\varpi)=(u, \varpi)^{\la \alpha, y\ra}\htt_{\alpha}(u)=(-1)^{\la \alpha, y\ra}\htt_{\alpha}(u)
\]
by \eqref{MR4}.
Hence one can see that $^{\yt(\varpi)}\tau=\tau$ if and only if $\la\alpha, y\ra\in 2\Z$, namely $y\in \Yt$.
\end{proof}

We can then prove

\begin{Prop}\label{IndTau}
The $\Tt$-representation $i_{\eta,\tau}(\chi)$ is irreducible and contains $\tau$ with multiplicity one when restricted to $\Tt_0$. Moreover, $\tau$ is the unique pseudo-spherical $\Tt_{0}$-constituent of $i_{\eta,\tau}(\chi)$. 
\end{Prop}	
\begin{proof}
The irreducibility, multiplicity one property, and the uniqueness property follow from Mackey theory together with Lemma \ref{L:tau_twist} and the fact that every element in $\Tt\slash Z(\Tt)\Tt_0$ is represented by an element of the form $\yt(\varpi)$ for some $y\in Y$. The details are left to the reader.
\end{proof}

Finally the following proposition shows that $(\Tt_{0},\tau)$ is a type of $\Tt$.

\begin{Prop}\label{TTypes}
Let $\pi$ be a smooth irreducible $\Tt$-module. The $\tau$-isotypic subspace of $\pi|_{\Tt_{0}}$ is nonzero if and only if $\pi\cong i_{\eta,\tau}(\chi)$ for some unramified character $\chi:T\rightarrow \mathbb{C}^{\times}$.
\end{Prop}
\begin{proof}
	The if direction is clear by Proposition \ref{IndTau}.
    Suppose that the $\tau$-isotypic subspace of $\pi|_{\Tt_{0}}$ is nonzero. Thus by Frobenius reciprocity we have a surjective $\Tt$-module homomorphism $\Ind_{\Tt_{0}}^{\Tt}(\tau)\twoheadrightarrow \pi$. Let $\chi^{\prime}$ be the central character of $\pi$ restricted to $\mu_{2}\Yt(\varpi)$. Then the surjective map must factor through the $\chi^{\prime}$-coinvariants 
    \[
    \Ind_{\Tt_{0}}^{\Tt}(\tau)_{\chi^{\prime}}\cong \Ind_{Z(\Tt)\Tt_{0}}^{\Tt}((\chi^{\prime}\cdot \eta^{-1})\otimes \tau_{\eta}).
    \]
    Now let $\chi:T\rightarrow \mathbb{C}^{\times}$ be an unramified character extending the non-genuine character $\chi^{\prime}\cdot \eta^{-1}$. By comparison of central characters and the Stone-Von Neumann Theorem for $\Tt$ we see that 
	\begin{equation*}
		\Ind_{Z(\Tt)\Tt_{0}}^{\Tt}((\chi^{\prime}\cdot \eta^{-1})\otimes \tau_{\eta})\cong \chi\otimes \Ind_{Z(\Tt)\Tt_{0}}^{\Tt}\tau_{\eta}.
	\end{equation*}
	Now the result follows because this induced representation is irreducible and maps onto $\pi$.
\end{proof}



\begin{comment}
%%%%%%%%%%%%%%%%%%%%%%%%%%%%%%%%%%%%%%%%%%%%%%%%%%%%%%%%%%%%%%%%%%%%%%%

\subsection{Extending pseudo-spherical representations}\textcolor{red}{This subsection might not be needed. -ST}

%%%%%%%%%%%%%%%%%%%%%%%%%%%%%%%%%%%%%%%%%%%%%%%%%%%%%%%%%%%%%%%%%%%%%%%

In this subsection, we will show that the $\Wcal$-invariance of each pseudo-spherical representation $\tau$ allows us to extend it to a representation of the group generated by $\Tt_{0}$ and $\Wcal$.

First, recall from \eqref{E:square_of_w} that for each $\alpha\in\Phi$, we have
\[
\wta(1)^2=(-1,-1)\cdot\hta(-1)\in\Tt_{0}.
\]
Hence for each fixed $\alpha\in\Phi$, the group
\[
\Tt_{0}_{\wta}:=\la \Tt_{0},\, \wta\ra,
\]
generated by $\Tt_{0}$ and $\wta$ has $\Tt_{0}$ as a subgroup of index $2$, where $\wt_\alpha=\wt_{\alpha}(1)$ as before. Then by the  $\Wcal$-invariance of $\tau$ and basic Mackey theory, 
for each $\alpha\in\Phi$, the representation $\tau$ extends to a representation of $\Tt_{0}_{\wta}$ exactly in two different ways. Let us denote the two extensions by $\tau_\alpha^+$ and $\tau_\alpha^-$. Then
\[
\tau_\alpha^+(\wta)=-\tau_\alpha^-(\wta).
\]

Let us mention the following commutativity relation.
\begin{Lem}\label{L:Ta_and_Tb_commute}
Let $\alpha, \beta\in\Phi$ be such that $\la\alpha, \beta\ra=0$. Then $\tau_\alpha^{\epsilon}(\wta)\tau_\beta^{\delta}(\wtb)=\tau_\beta^{\delta}(\wtb)\tau_\alpha^{\epsilon}(\wta)$ for all $\epsilon, \delta\in\{\pm 1\}$.
\end{Lem}
\begin{proof}
If $\la\alpha, \beta\ra=0$, then $\wta\wtb=\wtb\wta$. Hence the group generated by $\Tt_{0}_{\wta}$ and $\wtb$ has $\Tt_{0}_{\wta}$ as an index 2 subgroup. Thus $\tau^{\epsilon}_{\alpha}$ extends to a representation of this group, and there are exactly two such extensions. One of them is an extension of $\tau_\beta^{\delta}$. The lemma follows because $\wta\wtb=\wtb\wta$.
\end{proof}


Now, let
\[
\Tt_{0}_{\Wcal}:=\la\Tt_{0},\, \Wcal\ra,
\]
namely the group generated by $\Tt_{0}$ and $\Wcal$. Since $\Wcal$ acts on $\Tt_{0}$ by conjugation, we have the natural surjection
\begin{equation}\label{E:group_generated_by_T_and W}
\Tt_{0}\rtimes\Wcal\longrightarrow \Tt_{0}_{\Wcal},\quad (\ttt, \wt)\mapsto \ttt\wt,
\end{equation}
for $\ttt\in\Tt_{0}$ and $\wt\in\Wcal$. In what follows, we will show that the representation $\tau$ extends even to $\Tt_{0}_{\Wcal}$. 

Let us first mention the following elementary lemma.
\begin{Lem}\label{L:extends_to_projective_rep}
The representation $\tau$ extends to a projective representation of $\Tt_{0}_{\Wcal}$. Further, when restricted to $\Tt_{0}_{\wta}$, this projective representation coincides with $\tau_{\alpha}^{\pm}$ as projective representations.
\end{Lem}
\begin{proof}
This follows from the $\Wcal$-invariance of $\tau$. 
\end{proof}


Next, we need to mention the following.
\begin{Lem}\label{L:minimum_relations_of_W}
The group $\Wcal$ is isomorphic to the group generated by the elements $\hta(-1)$ for all $\alpha\in\Delta$ and the ``symbols" $\tilde{\mathbf{w}}_\alpha$ indexed by $\alpha\in\Delta$ subject to the following three relations:
\begin{enumerate}[(i)]
\item $\tilde{\mathbf{w}}_\alpha^2=(-1,-1)\, \hta(-1)$;
\item if $\la\alpha, \beta\ra=-1$ then $\tilde{\mathbf{w}}_\alpha\tilde{\mathbf{w}}_\beta\tilde{\mathbf{w}}_\alpha=\tilde{\mathbf{w}}_\beta\tilde{\mathbf{w}}_\alpha\tilde{\mathbf{w}}_\beta$ (braid relation);
\item if $\la\alpha, \beta\ra=0$ then $\tilde{\mathbf{w}}_\alpha\tilde{\mathbf{w}}_\beta=\tilde{\mathbf{w}}_\beta\tilde{\mathbf{w}}_\alpha$.
\end{enumerate}
\end{Lem}
\begin{proof}
This is precisely the minimum set of relations the generators $\wta$ of $\Wcal$ satisfy. See \cite[Lemma 3.2]{S04} for the braid relation.
\end{proof}

We can then prove
\begin{Lem}
The representation $\tau|_{\Tt_{0}\cap\Wcal}$ extends to a representation of $\Wcal$.
\end{Lem}
\begin{proof}
For each $\alpha\in\Delta$, let 
\[
T_\alpha:=\tau^{\epsilon_\alpha}_\alpha(\wta),
\] 
where for now we choose the sign $\epsilon_\alpha\in\{\pm\}$ arbitrarily. It suffices to show that the signs $\pm$ can be chosen for each $\alpha$ consistently with the three relations of the above lemma. Since $\tau_{\alpha}^{\epsilon_\alpha}|_{\Tt_{0}}=\tau$, the first relation $T_\alpha^2=(-1,-1)\,\tau(\hta(-1))$ is satisfied for any choice of the signs. Also by Lemma \ref{L:Ta_and_Tb_commute}, the third relation is satisfied, again, for any choice of the signs, namely $T_\alpha T_\beta=T_\beta T_\alpha$ if $\la\alpha,\beta\ra=0$. So we need to show that we can choose the signs so that the second relation (braid relation) is satisfied, namely $T_\alpha T_\beta T_\alpha=T_\beta T_\alpha T_\beta$ if $\la\alpha, \beta\ra=-1$.

So, assume $\la\alpha, \beta\ra=-1$. First note that
\[
\wtb\hta(-1)=\hta(-1)\htb(-1)\wtb
\]
by \eqref{E:Savin_relation}. Hence by applying $\tau_\beta^{\epsilon_\alpha}$ to both sides, we have
\begin{equation}\label{E:switch_T_alpha_and_h_beta}
T_\beta\tau(\hta(-1))=\tau(\hta(-1)\htb(-1)) T_\beta
\end{equation}
because $\tau|_{\Tt_{0}}=\tau_\alpha^{\epsilon_\alpha}|_{\Tt_{0}}$. Also note that
\[
\wta\hta(-1)=\hta(-1)\wta
\]
by \cite[(R') in p.113]{S04}. Hence by applying $\tau_{\alpha}^{\epsilon_\alpha}$ to both sides, we have
\begin{equation}\label{E:switch_T_alpha_and_h_alpha}
T_\alpha\tau(\hta(-1))=\tau(\hta(-1)) T_\alpha.
\end{equation}


% Also note that
%\[
%\hta(-1)\htb(-1)=(-1,-1)\htb(-1)\hta(-1).
%\]

Now, we already know that the braid relation $\wta\wtb\wta=\wtb\wta\wtb$ implies that $T_\alpha T_\beta T_\alpha=cT_\beta T_\alpha T_\beta$ for some $c\in\C^\times$ by Lemma \ref{L:extends_to_projective_rep}. We will show $c=\pm 1$. For this, it suffices to show $(T_\alpha T_\beta T_\alpha)^2=(T_\beta T_\alpha T_\beta)^2$. Note that
\begin{align*}
(T_\alpha T_\beta T_\alpha)^2&=T_\alpha T_\beta T_\alpha^2 T_\beta T_\alpha\\
&=(-1,-1)\,T_\alpha T_\beta\, \tau(\hta(-1))\, T_\beta T_\alpha\\
&=(-1,-1)\,T_\alpha\, \tau(\hta(-1)\htb(-1))\, T_\beta^2 T_\alpha\quad\text{by \eqref{E:switch_T_alpha_and_h_beta}}\\
&=T_\alpha\, \tau(\hta(-1)\htb(-1))\, \tau(\htb(-1))\, T_\alpha\\
&=(-1,-1)\, T_\alpha\, \tau(\hta(-1))\, T_\alpha\\
&=(-1,-1)\, \tau(\hta(-1))\, T_\alpha^2\quad\text{by \eqref{E:switch_T_alpha_and_h_alpha}}\\
&=\tau(\hta(-1))\, \tau(\hta(-1))\\
&=(-1, -1).
\end{align*}
By switching the roles of $\alpha$ and $\beta$, we also obtain $(T_\beta T_\alpha T_\beta)^2=(-1,-1)$. Hence 
\[
T_\alpha T_\beta T_\alpha=\pm T_\beta T_\alpha T_\beta.
\]


Fix $\alpha_0\in\Delta$ and choose the sign for $T_{\alpha_0}$ arbitrarily. Then, for each $\beta\in\Delta$, there exists a unique choice of the sign for $T_\beta$ such that the braid relation $T_\alpha T_\beta T_\alpha= T_\beta T_\alpha T_\beta$ is satisfied for all $\alpha, \beta\in\Delta$ with $\la\alpha, \beta\ra=-1$. Indeed, first choose the signs for $T_\beta$ arbitrarily for each $\beta\in\Delta\smallsetminus{\alpha_0}$. If $\la \alpha_0, \beta\ra=-1$ and $T_{\alpha_0} T_\beta T_{\alpha_0}=- T_\beta T_{\alpha_0} T_\beta$, then re-choose our $T_\beta$ to be $-T_\beta$. Namely, once a sign for $T_{\alpha_0}$ is fixed, it determines the sign for each of the adjacent roots $\beta$. Then the sign for each of the roots adjacent to each such $\beta$ is also determined. Repeating the process, the signs for all $T_\alpha$ are determined.
\end{proof}

This lemma immediately implies 
\begin{Prop}
The representation $\tau$ of $\Tt_{0}$ extend to a representation of $\Tt_{0}_{\Wcal}$.
\end{Prop}
\begin{proof}
Let $\tau_{\Wcal}$ be a representation of $\Wcal$ as constructed in the above lemma. Since each element in $\Tt_{0}_{\Wcal}$ is written as $\ttt\wt$ for $\ttt\in\Tt_{0}$ and $\wt\in\Wcal$ by \eqref{E:group_generated_by_T_and W}, we can define a representation on $\Tt_{0}_{\Wcal}$ simply by setting $\tau(\ttt)\tau_{\Wcal}(\wt)$. This is indeed well defined because $\tau$ and $\tau_{\Wcal}$ agree on the overlap $\Tt_{0}\cap\Wcal$.
\end{proof}
\end{comment}







%%%%%%%%%%%%%%%%%%%%%%%%%%%%%%%%%%%%%%%%%%%%%%%%%%%%%%%%%%%%%%%%%%%%%%%

\section{Pseudo-spherical representations of $\Kt$}\label{S:PSKType}

%%%%%%%%%%%%%%%%%%%%%%%%%%%%%%%%%%%%%%%%%%%%%%%%%%%%%%%%%%%%%%%%%%%%%%%

Now that we have constructed a pseudo-spherical representation $\tau$ of $\Tt_0$, the next task is to construct an irreducible representation $\sigma$ of the preimage $\It$ of the Iwahori subgroup $I\subseteq G$ that is to be used to define our Iwahori Hecke algebra. We achieve this by applying what we call {\it parahoric induction}, which is developed by Crisp, Meir and Onn in \cite{CMO19}\footnote{In \cite{CMO19}, this induction is called a variant of the Harish-Chandra induction instead of parahoric induction.}, based on the parahoric induction of Dat \cite{D09}. We call this representation $\sigma$ of $\It$ a pseudo-spherical representation of $\It$. This representation is essentially ``Weyl invariant", which allows us to extend it to a representation of $\Kt$, the preimage of the hyperspecial maximal compact subgroup $K\subseteq G$. We also call this extension a pseudo-spherical representation of $\Kt$.

Also at the end, we establish a finite group Shimura correspondence, generalizing the work of Savin \cite{S12} which proves the case when $F$ is unramified over $\Q_2$. More precisely, tensoring by the pseudo-spherical representation defines an equivalence of categories between the category of representations of $G(\Ocal/\varpi\Ocal)$ and a certain subcategory of $\widetilde{{\Kf}}$-modules. 


%%%%%%%%%%%%%%%%%%%%%%%%%%%%%%%%%%%%%%%%%%%%%%%%%%%%%%%%%%%%%%%%%%%%%%%

\subsection{Structure of $\Kt$}\label{SSecSplittings}

%%%%%%%%%%%%%%%%%%%%%%%%%%%%%%%%%%%%%%%%%%%%%%%%%%%%%%%%%%%%%%%%%%%%%%%
In this subsection we prove that certain subgroups of $K$ split to $\Gt$. In particular, we identify the maximal principal congruence subgroup of $K$ which splits with a normal image in $\Kt$.

To begin we set up some notation for congruence subgroups of $G$. For $k\in\mathbb{Z}_{\geq1}$, define
\begin{align*}
\Gamma(k)&:=\la T_{k},\, \xa(t)\st t\in\varpi^k\Ocal,\, \alpha\in\Phi\ra;\\
\Gamma_1(k)&:=\la\Gamma(k),\, \xa(t)\st t\in\Ocal,\, \alpha\in\Phi^+\ra;\\
\Gamma_0(k)&:=\la \Gamma_1(k),\, T_{0}\ra,
\end{align*}
where we recall
\[
T_k=T_{1+\varpi^k\Ocal}=\la\htt_{\alpha}(t)\st t\in 1+\varpi^k\Ocal, \alpha\in\Phi\ra.
\]
Note that 
\[
\Gamma(k)\subseteq\Gamma_1(k)\subseteq\Gamma_0(k).
\] 
Alternatively, if $\mathbb{G}$ is an $\Ocal$-model for $G$ such that $K=\mathbb{G}(\Ocal)$, then one can consider the surjective mod $\varpi^k$ map (\cite[Proposition 1.6]{A69})
\[
\mathbb{G}(\Ocal)\longrightarrow \mathbb{G}(\Ocal\slash\varpi^k\Ocal),
\]
and the groups $\Gamma(k), \Gamma_1(k)$ and $\Gamma_0(k)$ are the preimages of the trivial group, the unipotent radical $U(\Ocal\slash\varpi^k\Ocal)$ and the Borel subgroup $B(\Ocal\slash\varpi^k\Ocal)$, respectively. 


We are mostly concerned with the groups
\[
\Gamma(2e),\quad \Gamma(1), \quad \Gamma_1(2e),\quad   \Gamma_0(1),\quad\Gamma_0(2),
\]
where we often write
\begin{equation}\label{E:I_and_I_2_def}
    I:=\Gamma_0(1)\qand I_2:=\Gamma_0(2).
\end{equation}
The former is, of course, the Iwahori subgroup. Note the inclusions
\[
\Gamma_1(2e)\subseteq I_2\subseteq I\subseteq K
\qand \Gamma(2e)\subseteq\Gamma(1)\subseteq I\subseteq K.
\]

Let $\Gamma$ be any of the groups $\Gamma_{0}(k),\Gamma_{1}(k),\Gamma(k)$. The group $\Gamma$ admits an Iwahori factorization, meaning the multiplication map 
\begin{equation*}
	\mu_{\Gamma}: (U^{-}_{0}\cap \Gamma)\times (T_{0}\cap \Gamma)\times (U_{0}\cap \Gamma)\rightarrow \Gamma
\end{equation*}
is a homeomorphism. Using this factorization we can define a set-theoretic section of the exact sequence \eqref{CExt} for these congruence subgroups.

First, recall the the sections $\s^{\pm}$ as in \eqref{E:section_for_unipotent_radical}. Second, recall that every element of $T$ can be uniquely written in the form $\prod_{\alpha}h_{\alpha}(u_{\alpha})$, where $u_{\alpha}\in F^{\times}$. If we fix a linear order on $\Delta$ we can define a (set-theoretic) section $\s_{T}:T\rightarrow \Tt$ by $\prod_{\alpha}h_{\alpha}(u_{\alpha})\mapsto \prod_{\alpha}\htt_{\alpha}(u_{\alpha})$.
%Unless otherwise stated we use the Bourbaki labeling on $\Delta$. %If $H\subseteq T$ is a subset we will write $\s_{H}$ for $\s_{T}$ restricted to $H$.

Now we can define a section of $\Gamma$ by
\[
\s_{\Gamma}:=(\s^{-}\times \s_{T}\times \s^{+})\circ \mu_{\Gamma}^{-1}:\Gamma \rightarrow \Gt.
\]
The map $\s_{\Gamma}$ is not a group homomorphism in general, but for $\Gamma=\Gamma_{1}(2e)$ it is a homomorphism.
\begin{Prop}\label{Gamma12eSplitting}
	Let $\Gamma=\Gamma_{1}(2e)$. The section 
    \[
    \s_{\Gamma}:=(\s^{-}\times \s_{T}\times \s^{+})\circ \mu_{\Gamma}^{-1}:\Gamma \longrightarrow \Gt
    \]
    is a splitting of the exact sequence \eqref{CExt}.
\end{Prop}
\begin{proof}
    We already know that each of $\s^{+}, \s_{T}|_{T_{2e}}, \s^{-}$ is a homomorphism. Hence, it suffices to show that the image is closed under multiplication. Note that $\Gamma$ is generated by the subgroups $T_{2e}$, $U_{\alpha_{0},2e}$ and $U_{\alpha,0}$ with $\alpha\in \Delta$, where we recall the notation $U_{\alpha_{0},2e}$ and $U_{\alpha,0}$ from \eqref{E:notation_U_alpha_j}. It is clear that the image is closed under multiplication on the left by $T_{2e}$ and $U_{\alpha_{0},2e}$, and under multiplication on the right by $T_{2e}$ and $U_{0}$. So it remains to show that for any $\xt_{\alpha}(a)$, where $\alpha\in\Delta$ and $a\in\Ocal$, and any $y\in U^{-}_{2e}$, we have
	\begin{equation*}
		\xt_{\alpha}(a)\s_{\Gamma}(y)\in \s_{\Gamma}(\Gamma).
	\end{equation*}
	
	Let $V=U_{\Phi^{-}\smallsetminus\{-\alpha\}}(4\Ocal)\subseteq U^{-}_{2e}$. Since $\alpha$ is a simple root, the multiplication map defines a set bijection $V\times U_{-\alpha,2e}\cong U^{-}_{2e}$, and $\xt_{\alpha}(a)$ normalizes $V$. 
	
	Now let $v\in V$ and $b\in \Ocal$ be such that $y=v\xt_{-\alpha}(4b)$. Because $\xt_{\alpha}(a)$ normalizes $\s^{-}(V)$ by \eqref{MR2}, we have
	\begin{align*}
		\xt_{\alpha}(a)\s_{\Gamma}(y) =& [\xt_{\alpha}(a)\s^{-}(v)\xt_{\alpha}(-a)] \xt_{\alpha}(a)\xt_{-\alpha}(4b)\\
		=& [\xt_{\alpha}(a)\s^{-}(v)\xt_{\alpha}(-a)] \xt_{-\alpha}(\frac{4b}{1+4ab})\htt_{\alpha}(1+4ab)\xt_{\alpha}(\frac{a}{1+4ab})\\
		=& \s_{\Gamma}\left(x_{\alpha}(a)vx_{\alpha}(-a) x_{-\alpha}(\frac{4b}{1+4ab})\right)\htt_{\alpha}(1+4ab)\xt_{\alpha}(\frac{a}{1+4ab}),
	\end{align*}
    where we used Lemma \ref{PMRootSwapCor} for the second equality. The result follows.
\end{proof}

\begin{Rmk}
Proposition \ref{Gamma12eSplitting} generalizes the splitting results in type $A_2$ of \cite[Theorem 3.3]{Karasiewicz}, which are local counterparts to the splitting studied in \cite{K22}.
\end{Rmk}


Now, let us denote the splitting $\s_{\Gamma_1(2e)}$ simply by $\s$, so we have a homomorphism 
\[
\s:\Gamma_1(2e)\longrightarrow\Kt.
\]

The following proposition follows directly from the Steinberg relations, Corollary \ref{PMRootSwapCor}, and our results on the Hilbert symbol.

\begin{Prop}\label{P:SplitNormIm}$ $
    \begin{enumerate}
        \item The subgroup $\s(\Gamma_1(2e))\subseteq \Gammat_{0}(2e)$ is normal.
        \item\label{Gam2eSplitNorm} The subgroup $\s(\Gamma(2e))\subseteq \Kt$ is normal.
        \item If $\Gamma(2e-1)$ splits, then its image cannot be normal in $\Kt$.
    \end{enumerate}
\end{Prop}

\begin{comment}
\begin{Prop}\label{PropSplitPCong}$ $
	\begin{enumerate}
		\item The group $\s_{\Gamma(2e)}(\Gamma(2e))\subseteq \Kt$ is normal subgroup of $\Kt$.\label{Gam2eSplitNorm}
		\item The group $\s_{\Gamma_{1}(2e)}(\Gamma_{1}(2e))\subseteq \widetilde{\Gamma}_{0}(2e)$ is normal subgroup of $\widetilde{\Gamma}_{0}(2e)$.\label{Gam02eSplitNorm}
		\item The image of any splitting of $\Gamma(2e-1)$ (if it exists) cannot be normal in $\Kt$.\label{Gam2e-1NotNorm}
	\end{enumerate}
\end{Prop}

\begin{proof}
	First we prove Item \eqref{Gam2eSplitNorm}. By Proposition \ref{Gamma12eSplitting} the set $\s_{\Gamma(2e)}(\Gamma(2e))$ is a subgroup. We can directly check that the subgroup is normal using the Steinberg relations and Corollary \ref{PMRootSwapCor}.
	
	The proof of Item \eqref{Gam02eSplitNorm} is similar.
	
	Item \eqref{Gam2e-1NotNorm} follows from Lemma \ref{TorSplit}, Item \eqref{T2e-1NotNorm}. 
\end{proof}
\end{comment}

\begin{Rmk}\label{R:Gamma(2e)_is_max} Let $e^{\prime}$ be the even element in the set $\{e,e+1\}$. A more detailed analysis of the Hilbert symbol should show that $\s_{\Gamma(e^{\prime})}$ is a splitting of $\Gamma(e^{\prime})$. However, one can readily see, by comparing the relations \eqref{R3} and \eqref{R4} with \eqref{MR3} and \eqref{MR4}, that  the image of any splitting of $\Gamma(2e-1)$ cannot be normal in $\Kt$, and thus $\Gamma(2e)$ is the maximal principal congruence subgroup of $K$ which splits with normal image in $\Kt$.
\end{Rmk}

Now, since the group $\s(\Gamma(2e))$ is normal in $\Kt$, we can let
\begin{equation}\label{E:def_finite_K_tilde}
\widetilde{{\Kf}}:=\Kt/\s(\Gamma(2e)).
\end{equation}
Then this finite group is a nontrivial $\mu_{2}$-central extension of $K/\Gamma(2e)$; namely we have
\begin{equation}\label{Mod4CExt}
	1\rightarrow \mu_{2}\rightarrow \widetilde{{\Kf}}\rightarrow K/\Gamma(2e) \rightarrow 1.
\end{equation}
We know the extension is nontrivial because $\Tt_{0}/\s(T_{2e})$ is nonabelian.

\begin{Rmk} Suppose that $\mathbb{G}$ is an $\Ocal$-model for $G$ such that $K=\mathbb{G}(\Ocal)$. Then the group $K/\Gamma(2e)$ is isomorphic to $\mathbb{G}(\Ocal/4\Ocal)$. Chevalley groups over local rings were studied by Abe \cite{A69}. Central extensions of such groups, which include $\widetilde{{\Kf}}$, have been studied by Stein \cite{S73}.
\end{Rmk}



%%%%%%%%%%%%%%%%%%%%%%%%%%%%%%%%%%%%%%%%%%%%%%%%%%%%%%%%%%%%%%%%%%%%%%%

\subsection{Parahoric induction}\label{SS:ParaInd}

%%%%%%%%%%%%%%%%%%%%%%%%%%%%%%%%%%%%%%%%%%%%%%%%%%%%%%%%%%%%%%%%%%%%%%

In this subsection we review the parahoric induction of \cite{CMO19}, which generalizes the parabolic induction of \cite{D09}, and prove some basic results.

For this subsection we let $G$ denote an arbitrary profinite group instead of our Chevalley group, and let $\mathcal{M}(G)$ be the category of smooth representations of $G$.

Let us recall a couple of definitions from \cite{D09,CMO19}.

\begin{Def}
	The group $G$ is said to have a {\it virtual Iwahori decomposition} $(U,T,V)$ if $U,T,V\subseteq G$ are closed subgroups such that 
	\begin{itemize}
		\item $T$ normalizes $U$ and $V$,
		\item the multiplication map $U\times T\times V\rightarrow G$ is an open embedding,
		\item $G$ contains arbitrarily small open normal subgroups $K$ such that the multiplication map $(U\cap K)\times (T\cap K)\times (V\cap K)\rightarrow K$ is a homeomorphism.
	\end{itemize}
	A virtual Iwahori decomposition is call an {\it Iwahori decomposition} if the multiplication map $U\times T\times V\rightarrow G$ is a homeomorphism.
\end{Def}

Let $\Hcal(G)$ be the space of locally constant $\mathbb{C}$-valued functions on $G$. The convolution product gives $\Hcal(G)$ the structure of an idempotented $\mathbb{C}$-algebra, and each smooth representation is viewed as an $\Hcal(G)$-module, so that the category $\mathcal{M}(G)$ is identified with the category of smooth nondegenerate $\Hcal(G)$-modules. We write $e_G\in\Hcal(G)$ for the idempotent of the trivial representation of $G$. We then define parahoric induction and restriction from \cite{D09,CMO19} as follows.

\begin{Def}
	Let $(U,T,V)$ be a virtual Iwahori decomposition of a profinite group $G$. Let $\tau$ be a smooth $T$-module and let $\sigma$ be a smooth $G$-module. Let $e_{U}\in\Hcal(U)$ and $e_{V}\in\Hcal(V)$ be the idempotents of the trivial representation of $U$ and $V$ respectively. 
	
	Define functors $i_{U,V}:\mathcal{M}(T)\rightarrow \mathcal{M}(G)$ and $r_{U,V}:\mathcal{M}(G)\rightarrow \mathcal{M}(T)$ by
	\begin{align*}
		i_{U,V}:=&\Hcal(G)e_{U}e_{V}\otimes_{\Hcal(T)}(-);\\	
		r_{U,V}:=&e_{U}e_{V}\Hcal(G)\otimes_{\Hcal(G)}(-).
	\end{align*}
    We may write $i_{U,V}^{G}=i_{U,V}$ and $r_{U,V}^{G}=r_{U,V}$ if we want to emphasize $G$.
\end{Def}

Let us list a couple of important properties of these functors proved in \cite{CMO19,D09}.
\begin{Prop}\label{P:two_functions_are_adjoint}
Let $(U, T, V)$ be a virtual Iwahori decomposition of a profinite group $G$. 
\begin{enumerate}[(1)]
\item There are natural isomorphisms 
\[
i_{U, V}\cong i_{V, U}\qand r_{U, V}\cong r_{V, U}.
\]
\item $i_{U, V}$ is both left and right adjoint to $r_{U, V}$; namely we have
\[
\Hom_{\Hcal(G)}(\sigma, i_{U, V}\tau)\cong \Hom_{\Hcal(T)}(r_{U, V}\sigma, \tau)
\qand
\Hom_{\Hcal(G)}(i_{U, V}\tau, \sigma)\cong\Hom_{\Hcal(T)}(\tau, r_{U,V}\sigma)
\]
for all $\sigma\in\mathcal{M}(G)$ and $\tau\in\mathcal{M}(T)$.
\end{enumerate}
\end{Prop}
\begin{proof}
    This is \cite[Theorem 2.18 (1), (2) and (3)]{CMO19}.
\end{proof}



If $(U, T, V)$ is an Iwahori decomposition (instead of merely virtual), the above two functors have particularly nice properties as follows.
\begin{Prop}\label{P:two_functors_for_Iwahori_decomp}
Suppose that $(U, T, V)$ is an Iwahori decomposition of $G$. Then we have the following:
\begin{enumerate}[(1)]
\item $i_{U, V}$ sends $\Irr(T)$ to $\Irr(G)$, and $r_{U, V}$ sends $\Irr(G)$ to $\Irr(T)\cup\{0\}$.
\item $r_{U, V}\circ i_{U, V}\cong id_{\mathcal{M}(T)}$.\label{P:two_functors_for_Iwahori_decomp2}
\item If $\sigma\in\Irr(G)$ is such that $r_{U, V}\sigma\neq 0$, then $i_{U, V}r_{U, V}\sigma\cong\sigma$.
\item If $\sigma\in\Irr(G)$, then
\[
\dim_{\C}\Hom_T(\sigma^U, \sigma^V)\leq 1.
\]
Further, if this dimension is zero then $r_{U, V}\sigma=0$, and if it is 1, then the $\Hom$ space is generated by an isomorphism $e_Ve_U$ and $r_{U, V}\sigma\cong\sigma^U\cong\sigma^V$.
\end{enumerate}
\end{Prop}
\begin{proof}
This is \cite[Theorem 2.23]{CMO19}.
\end{proof}


The following definition is motivated by the well known result that we will recall in Lemma \ref{WCompat}.
\begin{Def}
	Suppose that $G$ has two virtual Iwahori decompositions $(U,T,V)$ and $(U^{\prime},T,V^{\prime})$. We call the two virtual decompositions {\it compatible} if the following four identities hold. 
	\begin{align*}
		U=&(U\cap U^{\prime})(U\cap V^{\prime}),\\
		V=&(V\cap \,U^{\prime})(V\cap V^{\prime}),\\
		U^{\prime}=&(U^{\prime}\cap U)(U^{\prime}\cap V),\\
		V^{\prime}=&(V^{\prime}\cap U)(V^{\prime}\cap V).
	\end{align*}
\end{Def}


The following is an important property of compatible decompositions.
\begin{Prop}\label{P:compatible_Iwahori_decomp}
Assume $(U, T, V)$ and $(U', T, V')$ are compatible virtual Iwahori decompositions of $G$. Then we have
\[
i_{U, V}\cong i_{U', V'}\qand r_{U, V}\cong r_{U', V'}.
\]
\end{Prop}
\begin{proof}
    This is \cite[Theorem 2.18 (4)]{CMO19}.
\end{proof}




\begin{Lem}\label{IsoTwist}
	Let $H_1$ and $H_2$ be profinite groups and assume each $H_i$ has a virtual Iwahori decomposition $(U_{i},T_{i},V_{i})$. Let $\phi:H_{1}\rightarrow H_{2}$ be a group isomorphism such that $\phi(U_{1})=U_{2}$, $\phi(T_{1})=T_{2}$, and $\phi(V_{1})=V_{2}$. Let $\tau\in\Irr(T_{1})$. Then as $H_{2}$-modules
	\begin{equation*}
		\,^{\phi}i_{U_1,V_1}\tau\cong i_{\phi(U_1),\phi(V_1)}\,^{\phi}\tau. 
	\end{equation*}
\end{Lem}
\begin{proof}
	One can directly check that there is a well-defined linear map $\,^{\phi}i_{U_1,V_1}\tau\rightarrow i_{\phi(U_1),\phi(V_1)}\,^{\phi}\tau$ defined by $f(g)e_{U_{1}}e_{V_{1}}\otimes v\mapsto f(\phi^{-1}(g))e_{U_{2}}e_{V_{2}}\otimes v$, where $g\in G$ and $v\in \tau$, and this map is an isomorphism of $H_{2}$-modules. 
\end{proof}


\begin{Prop}\label{TwistAut}
	Let $G$ be a profinite group with an Iwahori decomposition $(U,T,V)$. Let $\phi\in\Aut(G)$ be such that $\phi(T)=T$. Let $\tau\in\Irr(T)$.
	Suppose that 
	\begin{itemize}
		\item $\,^{\phi}\tau\cong \tau$, 
		\item the two Iwahori factorizations $(U,T,V)$ and $(\phi(U),T,\phi(V))$ are compatible.
	\end{itemize} 
	Then
	\begin{equation*}
		\,^{\phi}i_{U,V}\tau\cong i_{U,V}\tau.
	\end{equation*}
\end{Prop}

\begin{proof}
	By Lemma \ref{IsoTwist} we have $\,^{\phi}i_{U,V}\tau\cong i_{\phi(U),\phi(V)}\,^{\phi}\tau$. Since $\,^{\phi}\tau\cong \tau$ we have $i_{\phi(U),\phi(V)}\,^{\phi}\tau\cong i_{\phi(U),\phi(V)}\tau$. But since $(U, T, V)$ and  $(\phi(U),T,\phi(V))$ are compatible by our assumption, the proposition follows from Proposition \ref{P:compatible_Iwahori_decomp}.
\end{proof}

\begin{Lem}\label{L:IndRes2x}
    Let $G$ be a profinite group with an Iwahori decomposition $(U,T,V)$. Let $H\subseteq G$ be an open subgroup with an Iwahori decomposition $(U^{\prime}=H\cap U,T^{\prime}=H\cap T,V)$. Let $\tau$ be an irreducible representation of $T$ that remains irreducible when restricted to $T^{\prime}$.
    
    The $\Tt^{\prime}$-module homomorphism $\phi:r_{U^{\prime},V}^{H}i_{U,V}^{G}\tau\rightarrow r_{U,V}^{G}i_{U,V}^{G}\tau$ defined by
    \begin{equation*}
        e_{U^{\prime}}e_{V}f\otimes_{\mathcal{H}(H)} v \mapsto e_{U}e_{V}f\otimes_{\mathcal{H}(G)} v
    \end{equation*}
    is a well defined isomorphism of $T^{\prime}$-modules.
\end{Lem}  

\begin{proof}

    To begin we show that the map $\phi$ is well defined by constructing it directly. We do this in several steps.

    Because $H\subseteq G$ is of finite index, extension by $0$ defines a $(H,\mathcal{H}(H))$-bimodule inclusion $\mathcal{H}(H)\hookrightarrow \mathcal{H}(G)$. (Here we are fixing a Haar measure on $G$ and restricting it to a Haar measure on $H$.) Thus the map $A:e_{U^{\prime}}e_{V}\mathcal{H}(H)\rightarrow e_{U}e_{V}\mathcal{H}(G)$ defined by $e_{U^{\prime}}e_{V} f\mapsto e_{U}e_{V}f$ (i.e. left convolution by $e_{U}$) is a $(T^{\prime},\mathcal{H}(H))$-bimodule homomorphism. We tensor by $\otimes_{H}i_{U,V}^{G}\tau$ and get a map $A\otimes id_{i_{U,V}^{G}\tau}:e_{U^{\prime}}e_{V}\mathcal{H}(H)\otimes_{\mathcal{H}(H)}i_{U,V}^{G}\tau\rightarrow e_{U}e_{V}\mathcal{H}(G)\otimes_{\mathcal{H}(H)}i_{U,V}^{G}\tau$ by functoriality. Finally since the right $\mathcal{H}(G)$-module action on $e_{U}e_{V}\mathcal{H}(G)$ and the left $\mathcal{H}(G)$ action on $i_{U,V}^{G}\tau$ extends the respective $\mathcal{H}(H)$ action, we have a surjective map $B:e_{U}e_{V}\mathcal{H}(G)\otimes_{\mathcal{H}(H)}i_{U,V}^{G}\tau\twoheadrightarrow e_{U}e_{V}\mathcal{H}(G)\otimes_{\mathcal{H}(G)}i_{U,V}^{G}\tau$ defined by $e_{U}e_{V}f\otimes_{\mathcal{H}(H)} v\mapsto e_{U}e_{V}f\otimes_{\mathcal{H}(G)} v$. Then $\phi=B \circ(A\otimes_{\mathcal{H}(H)} id_{i_{U,V}^{G}\tau})$.

    Let $\Gamma$ be the kernel of $i_{U,V}^{G}\tau$, so that $\Gamma$ is a normal subgroup which acts trivially on $i_{U,V}^{G}\tau$. Note that in $\mathcal{H}(G)$ we have $e_{\Gamma\cap H}*e_{\Gamma}=e_{\Gamma}$. Thus under the map $\phi$ we have
    \begin{equation*}
        e_{U^{\prime}}e_{V}e_{\Gamma\cap H}\otimes_{\mathcal{H}(H)} v\mapsto e_{U}e_{V}e_{\Gamma}\otimes_{\mathcal{H}(G)} v.
    \end{equation*}

    Since $r_{U_{0},V}^{\It_{2}}i_{U,V}^{G}\tau$ is spanned by vectors of the form $e_{U}e_{V}e_{\Gamma}\otimes v$, where $v\in i_{U,V}^{G}\tau$, it follows that the map $\phi$ is surjective. By Proposition \ref{P:two_functors_for_Iwahori_decomp} \eqref{P:two_functors_for_Iwahori_decomp2} $r_{U,V}^{G}i_{U,V}^{G}\tau\cong \tau$, so the map is nonzero. By Lemma \ref{L:ResDimBound} and Proposition \ref{P:two_functors_for_Iwahori_decomp}, $\dim r_{U^{\prime},V}^{H}i_{U,V}^{G}\tau\leq \dim\tau$. Thus the map is an isomorphism of $T^{\prime}$-modules. 

    
\end{proof}

The next proposition provides a sufficient condition for when parahoric induction is preserved under restriction. 

\begin{Prop}\label{ResGamma1}
	Let $G$ be a profinite group with an Iwahori decomposition $(U,T,V)$. Let $H\subseteq G$ be an open subgroup with an Iwahori decomposition $(U^{\prime}=H\cap U,T^{\prime}=H\cap T,V)$. Suppose that 
	\begin{itemize}
		\item $\tau$ is a representation of $T$ such that $\tau$ restricted to $T^{\prime}$ remains irreducible;
		\item $(i_{U,V}\tau)^{U^{\prime}}=(i_{U,V}\tau)^{U}$.
	\end{itemize}
	
	Then as $H$-modules
    \[
    (i_{U,V}\tau)|_H\cong  i_{U^{\prime},V}\tau.
    \]

\end{Prop}

\begin{proof}
	To begin we show that $\Hom_{H}(i_{U^{\prime},V}\tau,i_{U,V}\tau)\cong \mathbb{C}$. By Proposition \ref{P:two_functions_are_adjoint},
	\begin{equation*}
		\Hom_{H}(i_{U^{\prime},V}\tau,\, i_{U,V}\tau)\cong \Hom_{T^{\prime}}(\tau,\, r_{U^{\prime},V}i_{U,V}\tau).
	\end{equation*}
	
	By Lemma \ref{L:IndRes2x}, as $T^{\prime}$-modules $r_{U^{\prime},V}i_{U,V}\tau\cong r_{U,V}i_{U,V}\tau$. By Proposition \ref{P:two_functors_for_Iwahori_decomp} \eqref{P:two_functors_for_Iwahori_decomp2}, we know $r_{U,V}i_{U,V}\tau\cong \tau$ as $T$-modules. This is also true as $T^{\prime}$-modules since $\tau$ restricted to $T^{\prime}$ remains irreducible. Thus $r_{U^{\prime},V}i_{U,V}\tau\cong \tau$ as $T^{\prime}$-modules. Therefore $\Hom_{H}(i_{U^{\prime},V}\tau,\, i_{U,V}\tau)\cong \mathbb{C}$.
	
	Now let $\phi\in \Hom_{H}(i_{U^{\prime},V}\tau,\, i_{U,V}\tau)$ be nonzero. We want to show that $\phi$ is an $H$-module isomorphism. The map $\phi$ is injective because $i_{U^{\prime},V}\tau$ is an irreducible $H$-module and $\phi\neq 0$. Thus it remains to show that $\phi$ is surjective.
	
	By the Iwahori factorization, the space $i_{U,V}\tau$ is generated as a $V$-module by elements of the form $e_{U}e_{V}\otimes w$, where $w\in \tau$. So, in particular $i_{U,V}\tau$ is generated as an $H$-module by the subspace $(i_{U,V}\tau)^{U}$. By assumption, $(i_{U,V}\tau)^{U}=(_{U,V}\tau)^{U^{\prime}}$, so $i_{U,V}\tau$ is $H$-generated by $(i_{U,V}\tau)^{U^{\prime}}$.
	
	We apply the $U^{\prime}$-invariants functor to the injection
	\begin{equation*}
		i_{U^{\prime},V}\tau \xhookrightarrow{\;\phi\;} i_{U,V}\tau
	\end{equation*}
	to get 
	\begin{equation*}
		\left(i_{U^{\prime},V}\tau\right)^{U^{\prime}}\xhookrightarrow{\;\phi^{U^{\prime}}\;} \left(i_{U,V}\tau\right)^{U^{\prime}}=\left(i_{U,V}\tau\right)^{U}.
	\end{equation*}
	By Proposition \ref{P:two_functors_for_Iwahori_decomp} (2) and (4), this injection must be an isomorphism of $T^{\prime}$-modules, as the domain and codomain are isomorphic to $\tau$. Thus the image of $\phi$ contains $\left(i_{U,V}\tau\right)^{U}$. Therefore as an $H$-module the image of $\phi$ generates $i_{U,V}\tau$. But $\phi$ is an $H$-modules homomorphism, so $\phi$ is surjective.
\end{proof}

\begin{Lem}\label{L:ResDimBound}
    Let $G$ be a profinite group with an Iwahori decomposition $(U,T,V)$. Let $\pi$ be a finite-dimensional representation of $G$. Then
    \begin{equation*}
        \dim r_{U,V}\pi\leq \dim \pi^{V}.
    \end{equation*}
\end{Lem}

\begin{proof}
    Since $\pi$ is finite dimensional there exists a normal subgroup $N\subseteq G$ such that $N$ acts trivially on $\pi$. Note that for any $f\in\mathcal{H}$ we have $e_{N}*f=f*e_{N}$.
    
    By definition $r_{U,V}\pi$ is spanned by vectors of the form $e_{U}e_{V}f\otimes v$, where $f\in \mathcal{H}(G)$ and $v\in\pi$. But $e_{U}e_{V}f\otimes v=e_{U}e_{V}e_{N}\otimes e_{V}f\cdot v$. Now the result follows because $e_{V}f\cdot v\in \pi^{V}$.
\end{proof}

If $G$ is finite, then it is convenient to make the identification
\[
\Hcal(G)\cong\C[G], \quad \delta_g\mapsto g,
\]
of algebras, where $\C[G]$ is the group algebra of $G$ and $\delta_g$ is the delta function at $g\in G$. Under this identification, we have
\begin{equation}\label{E:e_H_idempotent}
e_H=\frac{1}{|H|}\sum_{h\in H}h
\end{equation}
for each subgroup $H\subseteq G$. Note that for each $g\in G$ such that $gHg^{-1}=H$ we have $ge_H=e_Hg$. Also if $h\in H$, then $he_H=e_Hh=e_H$; namely $e_H$ ``absorbs" $h$. Note also that the third condition for virtual Iwahori decomposition is trivial because one can just choose $K$ to be the trivial group.

Let us prove a few facts for finite groups.
\begin{Lem}\label{L:spanning_vectors}
Assume $G$ is a finite group with an Iwahori decomposition $(U, T, V)$ and $\tau\in\mathcal{M}(T)$. Then $i_{U, V}\tau$ (as a $\C$-vector space) is spanned by the vectors of the form
\[
ye_Ue_V\otimes v\qquad (y\in V,\; v\in\tau).
\]
\end{Lem}
\begin{proof}
By definition, the space $i_{U, V}\tau$ is generated by the vectors of the form $ge_Ue_V\otimes v$ for $g\in G$ and $v\in\tau$. Now, by the Iwahori decomposition, each $g^{-1}$ is written as $g^{-1}=x t y$ for some $x\in U, t\in T, y\in V$, so that $g=y^{-1}t^{-1} x^{-1}$. But $x^{-1}e_U=e_U$ by absorbing $x^{-1}$, and 
\[
t^{-1}e_Ue_V\otimes v=e_Ue_Vt^{-1}\otimes v=e_Ue_V\otimes \tau(t^{-1})v.
\]
Hence we can write $ge_Ue_V\otimes v=y^{-1}e_Ue_V\otimes \tau(t^{-1})v$. The lemma follows.
\end{proof}

The final lemma in this subsection relates parahoric induction with induction in stages where one step involves standard induction.

\begin{Lem}\label{LHoms}
	Let $G$ be a finite group with an Iwahori decomposition $(U,T,V)$. Let $H\subseteq G$ be a subgroup with an Iwahori decomposition $(U^{\prime}=H\cap U,T,V)$. Let $\tau$ be a representation of $T$. Then
	\begin{equation*}
		\Hom_{G}(\Ind_{H}^{G}i_{V,U^{\prime}}\tau,\, i_{V,U}\tau)\neq 0.
	\end{equation*}
\end{Lem}

\begin{proof} The surjective $(G,T)$-bimodule map $\Hcal(G)\otimes_{H}\Hcal(H)e_{V}e_{U^{\prime}} \rightarrow \Hcal(G)e_{V}e_{U}$ defined by $g\otimes_{H} he_{V}e_{U^{\prime}}\mapsto gh e_{V}e_{U}$ induces a surjective $G$-module map 
	\begin{equation*}
		\Hcal(G)\otimes_{H}\Hcal(H)e_{V}e_{U^{\prime}}\otimes_{T}\tau \longrightarrow \Hcal(G)e_{V}e_{U}\otimes_{T}\tau=i_{V,U}\tau.
	\end{equation*}
	Since $i_{V,U}\tau\neq 0$ by \cite[Theorem 2.18 (6)]{CMO19}, the result follows.
\end{proof}

%\textbf{Might want to strengthen the previous lemma. It is enough for the image of $\phi$ to contain some nonzero vector of the form $e_{U}e_{V}\otimes w_{0}$, because once such a vector is in the image we can use the $\Tt_{0}$ action to get all elements of the form $e_{U}e_{V}\otimes w$.}





%%%%%%%%%%%%%%%%%%%%%%%%%%%%%%%%%%%%%%%%%%%%%%%%%%%%%%%%%%%%%%%%%%%%%%%%%%%%%%%%%%%%%%%%%%

\subsection{Construction of Iwahori types}\label{IwahoriType}

%%%%%%%%%%%%%%%%%%%%%%%%%%%%%%%%%%%%%%%%%%%%%%%%%%%%%%%%%%%%%%%%%%%%%%%%%%%%%%%%%%%%%%%%%%

Now, we apply the above described theory of parahoric induction to our pseudo-spherical representations $\tau$. The basic idea is that we will construct a representation of the Iwahori subgroup $\It$ by ``parahorically inducing" $\tau$ to $\It$, and prove that this representation behaves well with respect to the Weyl group by using the machinery of parahoric induction. The main results are stated in Proposition \ref{IRep}. (We will later prove in Theorem \ref{IType} that this representation is an Iwahori type of $\Gt$ in the sense of Bushnell-Kutzko \cite[(4.1) Definition]{BK98}.)

Recall from \S \ref{SSecSplittings} that we have a splitting $\s:\Gamma(2e)\to\Kt$ whose image is normal in $\Kt$. The kernel of the representation we will construct will contain $\s(\Gamma(2e))$. Hence we almost always work with ``representations modulo ${\Gamma(2e)}$", so that we view these representations as representations of subgroups of the finite quotient $\widetilde{\Kf}=\Kt\slash \s(\Gamma(2e))$. 

For this purpose, let us introduce the following notational convention. For each subgroup $H\subseteq K$, we denote its image in $K\slash\Gamma(2e)$ by the corresponding sans-serif font $\Hf$. To be precise
\[
\Hf:=H\slash (\Gamma(2e)\cap H)\subseteq K\slash \Gamma(2e)=\Kf.
\]
Further, denote the image of $\Ht$ in $\Kft$ by $\widetilde{\Hf}$, so that we have
\[
\begin{tikzcd}
    \Kt\arrow[r, phantom, "\supseteq"]\ar[d]&\Ht\arrow[r, two heads]\ar[d]&\widetilde{\Hf}\arrow[r, phantom, "\subseteq"]\ar[d]&\makebox[0pt][l]{$\widetilde{\Kf}=\Kt\slash\s(\Gamma(2e))$}\ar[d, shift left=.5em]\\
    K\arrow[r, phantom, "\supseteq"]&H\arrow[r, two heads]&\Hf\arrow[r, phantom, "\subseteq"]&\makebox[0pt][l]{$\Kf=K\slash\Gamma(2e)$\rlap{\; .}}
\end{tikzcd}
\]
In particular, we often use the groups
\[
\Uf_{j}^{\pm}\qand  \Uf_{\alpha, j}\qquad (j\in\Z_{\geq 0}),
\]
where the corresponding $U_j^{\pm}$ and $U_{\alpha, j}$ are recalled from \eqref{E:notation_U_alpha_j}. We often set
\[
\Uf_j=\Uf_j^{+}.
\]
 We are particularly interested in the case when $j=0, 1$.
These groups are usually viewed as subgroups of $\widetilde{\Kf}$ via the splittings $\s^{\pm}:U_0^{\pm}\to\Kt$, so that we have
\[
\s^{\pm}:\Uf_0^{\pm}\xhookrightarrow{\qquad}\widetilde{\Kf}.
\]
The images of the elements $\xt_{\alpha}(u)$ and $\htt_{\alpha}(t)$ in $\widetilde{\Kf}$ are to be denoted by the same symbol.

We should mention that the three groups 
\[
\widetilde{\Gammaf}(1),\quad \widetilde{\If}=\widetilde{\Gammaf}_0(1),\quad \widetilde{\If}_2=\widetilde{\Gammaf}_0(2)
\]
have Iwahori decompositions
\[
(\Uf_1, \widetilde{\Tf}_1, \Uf_1^-),\quad (\Uf_0, \widetilde{\Tf}_0, \Uf_1^-),\quad (\Uf_0, \widetilde{\Tf}_0, \Uf_2^-),
\]
respectively. Also the group $\widetilde{\Kf}$ has a virtual Iwahori decomposition $(\Uf_0,\widetilde{\Tf}_0, \Uf_0^-)$.

\begin{comment}
Recall the group ${\Kf}=K/\Gamma(2e)$ from subsection \ref{SSecSplittings}. We will be working with the following subgroups. Let $j\in\mathbb{Z}_{\geq 0}$. We define 
\begin{align*}
	{\Tf}_{j}:=&T_{j}\Gamma(2e)/\Gamma(2e)\\
	{\Tf}:=&{\Tf}_{0}\\	
	{\Uf}^{\pm}_{j}:=&U^{\pm}_{j}\Gamma(2e)/\Gamma(2e)\\
	{\Uf}^{\pm}:=&{\Uf}^{\pm}_{0}\\
	{\Bf}^{\pm}:=&{\Tf}{\Uf}^{\pm}.\\
\end{align*}
We may write ${\Bf}={\Bf}^{+}$. And for $\alpha\in \Phi$ we write ${\Uf}_{\alpha,j}:=U_{\alpha,j}\Gamma(2e)/\Gamma(2e)$

Let ${\If}$ be the subgroup of ${\Kf}$ generated by ${\Bf}$ and ${\Uf}^{-}_{1}$. Let ${\Kf}(1)$ be the subgroup generated by the three subgroups ${\Uf}^{\pm}_{1}$ and ${\Tf}_{1}$.

Note that by the results of \S \ref{SSecSplittings} and in particular sequence \eqref{Mod4CExt} we know that $\widetilde{{\Kf}}$ is a nontrivial $\mu_{2}$-central extension of ${\Kf}$. For any subgroup ${\Hf}\subseteq {\Kf}$ we write $\widetilde{{\Hf}}$ for its preimage in $\widetilde{{\Kf}}$. Note that the canonical splitting $\s^{\pm}:U^{\pm}_{0}\rightarrow \Kt$ induces a splitting of ${\Uf}^{\pm}$ into $\widetilde{{\Kf}}$. We identify ${\Uf}^{\pm}$ with its image under this splitting in $\widetilde{{\Kf}}$.

\end{comment}

\quad

Now, let $\tau$ be a pseudo-spherical representation of $\Tt_0$. Since $T_{2e}\subseteq T_{R}$ and $T_R$ is in the kernel of $\tau$, we view $\tau$ as a representation of 
\[
\widetilde{\Tf}_0=\Tt_0\slash T_{2e}=\Tt_0\slash \left(\Tt_0\cap \s(\Gamma(2e))\right).
\]
We continue to write $\tau$ for this representation of $\widetilde{{\Tf}}_0$. 

We then parahorically induce $\tau$ to a representation of $\widetilde{\If}$; namely we set
\[
\sigma:=i_{{\Uf},{\Uf}^{-}_{1}}^{\Ift}\tau,
\]
which is an irreducible representation of $\widetilde{{\If}}$ by Proposition \ref{P:two_functors_for_Iwahori_decomp}. We also consider 
\[
\sigma_{1}:=i_{{\Uf}_{1},{\Uf}^{-}_{1}}^{\Gammaft(1)}\tau,
\]
which is an irreducible representation of $\widetilde{\Gammaf}(1)$ by Proposition \ref{P:two_functors_for_Iwahori_decomp}. Note that $\widetilde{\Gammaf}(1)\subseteq\widetilde{\If}$. One of the main results in this section is to show that
\[
\sigma|_{\widetilde{\Gammaf}(1)}\cong\sigma_1,
\]
which will be proven in Proposition \ref{IRep} among other things. 


We begin with a few technical lemmas.


\begin{Lem}\label{WCompat}
	Let $w\in W$. Then $({\Uf}_{0},{\Tf}_{0},{\Uf}^{-}_{0})$ and $(\,^{w}{\Uf}_{0},{\Tf}_{0},\,^{w}{\Uf}^{-}_{0})$ are compatible virtual Iwahori decompositions of ${\Kf}$. Specifically, 
	\begin{align*}
		{\Uf}_0=&({\Uf}_0\cap \,^{w}{\Uf}_0)({\Uf}_0\cap \,^{w}{\Uf}_0^{-}),\\
		{\Uf}_0^{-}=&({\Uf}_0^{-}\cap \,^{w}{\Uf}_0)({\Uf}_0^{-}\cap \,^{w}{\Uf}_0^{-}),\\
		\,^{w}{\Uf}_0=&(\,^{w}{\Uf}_0\cap {\Uf}_0)(\,^{w}{\Uf}_0\cap {\Uf}_0^{-}),\\
		\,^{w}{\Uf}_0^{-}=&(\,^{w}{\Uf}_0^{-}\cap {\Uf}_0)(\,^{w}{\Uf}_0^{-}\cap {\Uf}_0^{-}).\\
	\end{align*}
	The same also holds for ${\Gammaf}(1)$ with the Iwahori factorization $({\Uf}_{1},{\Tf}_{1},{\Uf}^{-}_{1})$.
\end{Lem}

\begin{proof}
	This follows because ${\Uf}_0\cap \,^{w}{\Uf}_0=\prod_{\alpha\in \Phi^{+}\cap w\Phi^{+}}{\Uf}_{\alpha,0}$, along with similar identities for the other intersections.
\end{proof}


\begin{comment}
Let $\alpha\in \Delta$. For $j\in\mathbb{Z}_{\geq0}$ let ${\Uf}_{\Phi^{-}_{\alpha},j}\subseteq {\Kf}$ be the subgroup generated by the subgroups ${\Uf}_{\beta,j}$, where $\beta\in \Phi^{-}\smallsetminus\{-\alpha\}$. Note that the multiplication map defines a set bijection ${\Uf}_{\Phi^{-}_{\alpha},j}\times {\Uf}_{-\alpha,j}\cong {\Uf}_{j}^{-}$, and the order of the factors can be reversed. Thus one can readily see from \eqref{E:e_H_idempotent} that the idempotent $e_{{\Uf}_{j}^{-}}$ can be factored as
\begin{equation}\label{IdemFactor}
	e_{{\Uf}_{j}^{-}}=e_{{\Uf}_{\Phi^{-}_{\alpha},j}}e_{{\Uf}_{-\alpha,j}}=e_{{\Uf}_{-\alpha,j}}e_{{\Uf}_{\Phi^{-}_{\alpha},j}}.
\end{equation}
\end{comment}


\begin{Lem}\label{LemVarIndRootMove1}
	Let $\alpha\in\Delta$ and $u\in\mathcal{O}\smallsetminus\{0\}$. Then in the group algebra $\mathbb{C}[\;\Ift\;]$
	\begin{equation*}
		e_{{\Uf}_{-\alpha,1}}\xt_{\alpha}(u)e_{{\Uf}_{1}^{-}}e_{{\Uf}_{0}}=e_{{\Uf}_{1}^{-}}\frac{1}{\left|\varpi\Ocal\slash 4\Ocal\right|}\sum_{v\in \varpi\mathcal{O}/4\mathcal{O}}(1+uv,-u)\htt_{\alpha}(1+uv) e_{{\Uf}_{0}}.
	\end{equation*}
If $u\in \Ocal^{\times}$, note that $\frac{1}{\left|\varpi\Ocal\slash 4\Ocal\right|}\sum_{v\in \varpi\mathcal{O}/4\mathcal{O}}(1+uv,-u)\htt_{\alpha}(1+uv)=X_{\alpha,-u,1}$. (Recall \S \ref{SSSqRoot}.)
\end{Lem}

\begin{proof}
Let ${\Uf}_{\Phi^{-}_{\alpha},1}\subseteq {\Kf}$ be the subgroup generated by the subgroups ${\Uf}_{\beta,1}$, where $\beta\in \Phi^{-}\smallsetminus\{-\alpha\}$. Note that the multiplication map defines a set bijection ${\Uf}_{\Phi^{-}_{\alpha},1}\times {\Uf}_{-\alpha,1}\cong {\Uf}_{1}^{-}$, and the order of the factors can be reversed. Thus one can readily see from \eqref{E:e_H_idempotent} that the idempotent $e_{{\Uf}_{1}^{-}}$ can be factored as
\[
	e_{{\Uf}_{1}^{-}}=e_{{\Uf}_{\Phi^{-}_{\alpha},1}}e_{{\Uf}_{-\alpha,1}}=e_{{\Uf}_{-\alpha,1}}e_{{\Uf}_{\Phi^{-}_{\alpha},1}}.
\]
Hence we have
	\begin{equation*}
		e_{{\Uf}_{-\alpha,1}}\xt_{\alpha}(u)e_{{\Uf}_{\Phi^{-}_{\alpha},1}}=e_{{\Uf}_{-\alpha,1}}e_{{\Uf}_{\Phi^{-}_{\alpha},1}}\xt_{\alpha}(u)=e_{{\Uf}_{1}^{-}}\xt_{\alpha}(u),
	\end{equation*}
    where we used that $\xt_{\alpha}(u)$ normalizes ${\Uf}_{\Phi^{-}_{\alpha},1}$.
	
	Using Proposition \ref{PMRootSwap} we have
	\begin{equation*}
		\xt_{\alpha}(u)e_{{\Uf}_{-\alpha,1}}=\frac{1}{\left|\varpi\Ocal\slash 4\Ocal\right|}\sum_{v\in \varpi\mathcal{O}/4\mathcal{O}}(1+uv,-u)\xt_{-\alpha}(\frac{v}{1+uv})\htt_{\alpha}(1+uv) \xt_{\alpha}(\frac{u}{1+uv}).
	\end{equation*}
	Now the result follows because 
	\begin{align*}
		&e_{{\Uf}_{-\alpha,1}}\sum_{v\in \varpi\mathcal{O}/4\mathcal{O}}(1+uv,-u)\xt_{-\alpha}(\frac{v}{1+uv})\htt_{\alpha}(1+uv) \xt_{\alpha}(\frac{u}{1+uv})e_{{\Uf}_{0}}\\
		&=e_{{\Uf}_{-\alpha,1}}\sum_{v\in \varpi\mathcal{O}/4\mathcal{O}}(1+uv,-u)\htt_{\alpha}(1+uv)e_{{\Uf}_{0}}
	\end{align*}
    by letting $e_{{\Uf}_{-\alpha,1}}$ absorb $\xt_{-\alpha}(\frac{v}{1+uv})$.
\end{proof}

\begin{Lem}\label{LemVarIndRootMove2}
	Let $\alpha\in\Delta$, $u\in\mathcal{O}$ and $w\in \tau$. Then in the $\widetilde{{\If}}$-module $\sigma=i_{{\Uf}_{1}^{-},{\Uf}_{0}}\tau$
	\begin{equation*}
		e_{{\Uf}_{1}}\xt_{\alpha}(u)e_{{\Uf}_{1}^{-}}e_{{\Uf}_{0}}\otimes w = e_{{\Uf}_{1}}e_{{\Uf}_{1}^{-}}e_{{\Uf}_{0}}\otimes w,
	\end{equation*}
    where we recall that the tensor product is over $\C[\Tft_0]$.
\end{Lem}
\begin{proof}
        This is trivial if $u\in\varpi\mathcal{O}$. So assume $u\in\mathcal{O}^{\times}$. Let $X=X_{-\alpha,u,1}=X_{\alpha, -u, 1}\in\mathbb{C}[\Tft_{0}]$ be as in \eqref{E:X_alpha_u_j_def}. By Corollary \ref{XInvert}, the operator $\tau(X)$ is invertible. Thus
	\begin{align*}
		&e_{{\Uf}_{1}}\xt_{\alpha}(u)e_{{\Uf}_{1}^{-}}e_{{\Uf}_{0}}\otimes w\\
        = & e_{{\Uf}_{1}}\xt_{\alpha}(u)e_{{\Uf}_{1}^{-}}e_{{\Uf}_{0}}X\otimes \tau(X)^{-1}w\\
        = & e_{{\Uf}_{1}}\xt_{\alpha}(u)Xe_{{\Uf}_{1}^{-}}e_{{\Uf}_{0}}\otimes \tau(X)^{-1}w\\
		= & e_{{\Uf}_{1}}\frac{1}{\left|\varpi\Ocal\slash 4\Ocal\right|}\sum_{v\in \varpi\Ocal/4\Ocal}(1+uv,u)\xt_{\alpha}(u)\htt_{-\alpha}(1+uv)e_{{\Uf}_{1}^{-}}e_{{\Uf}_{0}}\otimes \tau(X)^{-1}w,
	\end{align*}
    where for the second equality we used $e_{{\Uf}_{1}^{-}}e_{{\Uf}_{0}}X=Xe_{{\Uf}_{1}^{-}}e_{{\Uf}_{0}}$, noting that $\htt_{-\alpha}(1+uv)$ normalizes both $\Uf_1^-$ and $\Uf_0$. 
    
    Note that
    \[
     e_{\Uf_1}\xt_{\alpha}(u)=e_{\Uf_1}\xt_{\alpha}(\frac{u}{1+uv}),
    \]
    because for $v\in\varpi\Ocal$ we have $u\equiv \frac{u}{1+uv}$ modulo $\varpi$. Also
    \[
    e_{{\Uf}_{1}^{-}}=\xt_{-\alpha}(\frac{v}{1+uv})e_{{\Uf}_{1}^{-}}
    \]
    by absorbing $\xt_{-\alpha}(\frac{v}{1+uv})$. Hence we can write
	\begin{align*}
		&e_{{\Uf}_{1}}\sum_{v\in \varpi\Ocal/4\Ocal}(1+uv,u)\xt_{\alpha}(u)\htt_{-\alpha}(1+uv)e_{{\Uf}_{1}^{-}}e_{{\Uf}_{0}}\otimes \tau(X)^{-1}w \\
		= & e_{{\Uf}_{1}}\sum_{v}(1+uv,u)\xt_{\alpha}(\frac{u}{1+uv})\htt_{-\alpha}(1+uv)\xt_{-\alpha}(\frac{v}{1+uv})e_{{\Uf}_{1}^{-}}e_{{\Uf}_{0}}\otimes \tau(X)^{-1}w\\
        = & e_{{\Uf}_{1}}\left(\sum_{v}\xt_{-\alpha}(v)\right)\xt_{\alpha}(u)e_{{\Uf}_{1}^{-}}e_{{\Uf}_{0}}\otimes \tau(X)^{-1}w\quad (\text{Proposition \ref{PMRootSwap}})\\
		= & \left|\varpi\Ocal\slash 4\Ocal\right|e_{{\Uf}_{1}}e_{{\Uf}_{-\alpha,1}}\xt_{\alpha}(u)e_{{\Uf}_{1}^{-}}e_{{\Uf}_{0}}\otimes \tau(X)^{-1}w.
	\end{align*}
	
	Now we apply Lemmas \ref{LemVarIndRootMove1} and \ref{Xpmalpha} to get 
		\begin{equation*}
		e_{{\Uf}_{1}}e_{{\Uf}_{-\alpha,1}}\xt_{\alpha}(u)e_{{\Uf}_{1}^{-}}e_{{\Uf}_{0}}\otimes \tau(X)^{-1}w = e_{{\Uf}_{1}}e_{{\Uf}_{1}^{-}}e_{{\Uf}_{0}}\otimes w.
	\end{equation*}
    The lemma is proven.
\end{proof}

This lemma implies
\begin{Prop}\label{UInvar}
	Let $\alpha\in \Delta$. Then we have 
	\begin{equation*}
		\sigma^{{\Uf_0}}=\sigma^{{\Uf_0}\cap \,^{w_{\alpha}}\widetilde{{\If}}}=\sigma^{{\Uf}_{1}}.
	\end{equation*}
\end{Prop}
\begin{proof}
Since $\Uf_0\supseteq (\Uf_0\cap \,^{w_{\alpha}}\widetilde{{\If}}) \supseteq \Uf_1$, the inclusions $\sigma^{{\Uf_{0}}}\subseteq\sigma^{{\Uf_{0}}\cap \,^{w_{\alpha}}\widetilde{{\If}}}\subseteq \sigma^{{\Uf}_{1}}$ are immediate. It remains to show $\sigma^{{\Uf}_{1}}\subseteq \sigma^{{\Uf_0}}$. Note that $\sigma=i_{\Uf_0, \Uf_1^-}\tau\cong i_{{\Uf}^{-}_{1},{\Uf_0}}\tau$ by Proposition \ref{P:two_functions_are_adjoint} (1). By the Iwahori factorization and Lemma \ref{L:spanning_vectors}, $(i_{{\Uf}^{-}_{1},{\Uf_0}}\tau)^{{\Uf}_{1}}$  is spanned by elements of the form $e_{{\Uf}_{1}}u e_{{\Uf}_{1}^{-}}e_{{\Uf}_{0}}\otimes w$, where $u\in {\Uf}_{0}$ and $w\in \tau$. We also know that ${\Uf}_{0}$ is generated by the elements of the form $x_{\alpha}(u)$, where $\alpha\in \Delta$ and $u\in \mathcal{O}$. The result follows by Lemma \ref{LemVarIndRootMove2}, induction and the identity
	\begin{equation*}
		e_{{\Uf}_{1}}xx_{\alpha}(u)e_{{\Uf}_{1}^{-}}e_{{\Uf}_{0}}\otimes w =xe_{{\Uf}_{1}}x_{\alpha}(u)e_{{\Uf}_{1}^{-}}e_{{\Uf}_{0}}\otimes w
	\end{equation*}
	for $x\in {\Uf}_{0}$.
\end{proof}

Now we can state and prove the main theorem of this subsection. Let 
\[
\sigma_{\alpha}:=i^{\widetilde{{\If}}\cap\,^{w_{\alpha}}\widetilde{{\If}}}_{{\Uf}_{0}\cap \,^{w_{\alpha}}\widetilde{{\If}},{\Uf}^{-}_{1}}\tau,
\]
and recall that 
\[
\sigma=i_{{\Uf_{0}},{\Uf}^{-}_{1}}^{\Ift}\tau\qand \sigma_{1}=i_{{\Uf}_{1},{\Uf}^{-}_{1}}^{\Gammaft(1)}\tau,
\]
all of which are irreducible representations of $\widetilde{{\If}}\cap\,^{w_{\alpha}}\widetilde{{\If}}$, $\widetilde{\If}$ and $\widetilde{\Gammaf}(1)$, respectively, by Proposition \ref{P:two_functors_for_Iwahori_decomp}.

\begin{Prop}\label{IRep} Let $\alpha\in\Delta$.
	\begin{enumerate}[(1)]
		\item The $\widetilde{{\Gammaf}}(1)$-module $\sigma_{1}$ is $W$-invariant, \ie
        $^w\sigma_1\cong\sigma_1$ for all $w\in W$.\label{sigWInvar}
		\item As $\widetilde{{\Gammaf}}(1)$-modules, $\sigma_{1}\cong\sigma|_{\widetilde{{\Gammaf}}(1)}$.\label{ResSig}
		\item As $\widetilde{{\If}}\cap\,^{w_{\alpha}}\widetilde{{\If}}$-modules, $\sigma_{\alpha}\cong\sigma|_{\widetilde{{\If}}\cap\,^{w_{\alpha}}\widetilde{{\If}}}$.\label{ResSigAlpha}
		\item As $\widetilde{{\If}}\cap\,^{w_{\alpha}}\widetilde{{\If}}$-modules , $\,^{w_{\alpha}}\sigma_{\alpha}\cong\sigma_{\alpha}$.\label{SigAWInv}
	\end{enumerate}
\end{Prop}

\begin{proof}	
	\eqref{sigWInvar} and \eqref{SigAWInv} follow from Proposition \ref{TwistAut} by setting $\phi$ to be the conjugation by the Weyl group element. Indeed, the two conditions of Proposition \ref{TwistAut} are satisfied by the $W$-invariance of $\tau$ and by Lemma \ref{WCompat}.
	
	\eqref{ResSig} and \eqref{ResSigAlpha} follow from Proposition \ref{ResGamma1}. Indeed, the two conditions of Proposition \ref{ResGamma1} are satisfied by Lemma \ref{RedtoU1} and Proposition \ref{UInvar}.
\end{proof}


%%%%%%%%%%%%%%%%%%%%%%%%%%%%%%%%%%%%%%%%%%%%%%%%%%%%%%%%%%%%%%%%%%%%%%%%%%%%%%%%%%%%%%%%%%

\subsection{Extending $\sigma$ to $\widetilde{\Kf}$}\label{S:KType}

%%%%%%%%%%%%%%%%%%%%%%%%%%%%%%%%%%%%%%%%%%%%%%%%%%%%%%%%%%%%%%%%%%%%%%%%%%%%%%%%%%%%%%%%%%
Next we show that the $W$-invariance of the $\widetilde{{\Gammaf}}(1)$-module $\sigma_{1}$ allows us to prove
\begin{Thm}\label{KRep}
	The $\widetilde{\Gammaf}(1)$-module $\sigma_{1}$ extends uniquely to a $\widetilde{{\Kf}}$-module.
\end{Thm}

Since we know that $\sigma_1$ is $W$-invariant, this is actually almost immediate by using the following known results on the simply-connected simply-laced Chevalley groups over finite fields of characteristic $2$, except for a few low rank cases.
\begin{Prop}\label{SchurMult}
	Let $\kappa$ be a finite field of characteristic $2$. Let $G_\kappa$ be the $\kappa$-points of a simply-connected simply-laced Chevalley group defined over $\kappa$ with an irreducible root system. Then $G_\kappa$ is centrally closed except the following five cases:
\begin{equation}\label{E:exceptional_low_rank_cases}
\SL_{2}(\mathbb{F}_{4}),\quad \SL_{3}(\mathbb{F}_{2}),\quad \SL_{3}(\mathbb{F}_{4}),\quad \SL_{4}(\mathbb{F}_{2}),\quad D_{4}(\mathbb{F}_{2}). 
\end{equation}
Hence except these five cases, $G_\kappa$ has no nontrivial double cover. 
\end{Prop}
\begin{proof}
This follows from \cite[Corollary 2 and Remark (b), page 58]{S16}. 
\end{proof}

\begin{Rmk}
The Schur multiplier of the above five groups can be found in \cite{S81}. Note that \cite[Theorem 1.1]{S81} contains a typo. The group $A_{3}(2)\cong \SL_{4}(\mathbb{F}_{2})\cong A_{8}$ should be listed under (a) in that theorem. Also note other isomorphisms $\SL_{2}(\mathbb{F}_{4})\cong A_{5}$ and $\SL_{3}(\mathbb{F}_{2})\cong \mathrm{PSL}_{3}(\mathbb{F}_{7})$.
\end{Rmk}

Let $\kappa=\Ocal/\varpi\Ocal$ and let $G_{\kappa}$ be the group of $\kappa$-points of the simply-connected Chevalley group over $\kappa$ of the same type as $G$. Note that $G_{\kappa}\cong K/\Gamma(1)\cong{\Kf}/{\Gammaf}(1)$, where the isomorphisms are induced by the natural maps on Chevalley generators. We also write $T_{\kappa}\subseteq B_{\kappa}\subseteq G_{\kappa}$ for the $\kappa$-points of a maximal torus and Borel subgroup of the Chevalley group. We choose these groups so that they are compatible with the Chevalley pinning of $G$.

With this said, we can give the proof of the above theorem
\begin{proof}[Proof of Theorem \ref{KRep}]
 Since we are assuming that $\Phi\neq A_{1}$ or $D_2$ and it is simply-laced, we can see that $\Kft$ is perfect from the relation \eqref{MR2}.  Hence the uniqueness follows.
    
	From Proposition \ref{IRep} \eqref{sigWInvar} and \eqref{ResSig}, we know that the irreducible $\widetilde{{\Gammaf}}(1)$-module $\sigma_{1}$ is invariant under conjugation by $\widetilde{{\Kf}}$, because $\widetilde{{\Kf}}$ is generated by $\widetilde{{\If}}$ and $W$. Hence, if $G_\kappa$ is not among the five exceptional cases of \eqref{E:exceptional_low_rank_cases}, one can just invoke \cite[(11.45) Proposition]{CR81}, which immediately proves the theorem.
    
    It remains to show that the representation extends for the five exceptional cases of \eqref{E:exceptional_low_rank_cases}. Our argument for these exceptional cases is inspired by Savin \cite[Proposition 4.1]{S12}. The basic idea is to embed $G$ to the group $G'$ of the same type with one rank larger. Then since we already know that the result holds for $G'$, we will obtain the result for $G$ by restriction.

    To be precise, by Proposition \ref{SchurMult}, we may assume $G$ is the simply-laced simply-connected Chevalley group of type $A_r$ or $D_r$. (Though we have only to consider when $G=\SL_3, \SL_4$ or $D_4$, our argument works in general.) Define $G'$ to be the simply-laced simply-connected Chevalley group of the same type as that of $G$ with rank $r+1$, so that we have the natural embedding $G\subseteq G'$. This embedding is obtained as follows. Let $\Phi'$ be the root system of $G'$, where the set of simple roots $\Delta'$ is so that $\Delta\subseteq\Delta'$ and the unique root in $\Delta'\smallsetminus\Delta$ is the root at the left endpoint of the Dynkin diagram. Then $\Phi\subseteq\Phi'$ is a subroot system, and the group $G$ viewed as a subgroup of $G$ is the group generated by $\xt_{\alpha}(t)$ for $\alpha\in\Phi$ and $t\in F$. (If $G=\SL_r$ then $G$ is embedded in $G'=\SL_{r+1}$ in the lower right corner.) We denote by $T'$, $K'$, $\Gamma(k)'$, etc, the corresponding groups for $G'$.

    Let $\Nf\subseteq\Kft'$ be the group generated by the subgroups $\Uf_{\alpha, 1}$ for $\alpha\in{\Phi'}^+\smallsetminus\Phi^+$. (If $G=\SL_r$ then $\Nf$ is the unipotent group of the ``first rows" in $\SL_{r+1}$.)  Similarly we let $\Nf^-\subseteq\Kft'$ be the group generated by the subgroups $\Uf_{\alpha, 1}$ for $\alpha\in{\Phi'}^-\smallsetminus\Phi^-$. Let $\sigma$ be our pseudo-spherical representation of $\Gammaft(1)$ constructed from a pseudo-spherical representation $\tau$ of $\Tt_1$. We know from Proposition \ref{PseudoSphericalBranch} that there is a pseudo-spherical representation $\tau'$ of $\Tt_1'$ such that $\tau'|_{\Tt_1}$ contains $\tau$ with multiplicity one. Let $\sigma'$ be the pseudo-spherical representation of $\Gammaf(1)'$ constructed from $\tau'$. We claim that
    \begin{equation}\label{ExtThmIso1}
        \Hom_{\Gammaft(1)}(\sigma,\,{\sigma'}^{\Nf})\cong \mathbb{C}.
    \end{equation}
    
    
    
    %We fix a pinning of $G$ with respect to the Borel subgroup $B$ with maximal torus $T$. Let $\Phi$ be the root system of $(G,T)$. Using the pinning, define $G^{\prime}\subseteq G$ to be the Levi subgroup associated to a subdiagram of the Dynkin diagram of type $A_{r-1}$ or $D_{r-1}$, respectively. For any $H\subseteq G$, let $H^{\prime}=H\cap G^{\prime}$. (\textbf{Probably need a bit more detail here.}) Let $\Phi^{\prime}$ be the root system of $(G^{\prime},T^{\prime})$.

    %Let $\Nf^{\pm}\subseteq \Kft$ be the subgroup generated by the subgroup $\Uf_{\alpha,1}$, where $\alpha\in\Phi^{\pm}-\Phi^{\prime}$. We will write $\Nf=\Nf^{+}$. Let $\sigma^{\prime}$ be a pseudospherical representation of $\Gammaft(1)^{\prime}$ constructed from a pseudospherical $\Tt_{1}^{\prime}$-representation $\tau^{\prime}$. Let $\tau$ be a pseudospherical $\Tt_{1}$-representation containing $\tau^{\prime}$, which exists by Proposition \ref{PseudoSphericalBranch}. Let $\sigma$ be the pseudospherical representation of $\Gammaft(1)$ constructed from $\tau$.

    %We claim that
    %\begin{equation}\label{ExtThmIso1}
    %    \Hom_{\Gammaft(1)^{\prime}}(\sigma^{\prime},\sigma^{\Nf})\cong \mathbb{C}.
    %\end{equation}


    Note that $\Gammaft(1)'$ has an Iwahori factorization relative to the triple $(\Nf^{-},\Tft_{1}'\Gammaft(1),\Nf)$. Thus as $\Tft_{1}'\Gammaft(1)$-modules we have $\sigma^{\Nf}\cong r_{\Nf^{-},\Nf}^{\Tft_{1}'\Gammaft(1)}(\sigma')$ by Proposition \ref{P:two_functors_for_Iwahori_decomp}. On the other hand, as $\Gammaft(1)'$-modules, we have $\sigma'\cong i_{\Nf^{-},\Nf}^{\Gammaft(1)'}\left(i_{\Uf^{-}_{1},\Uf_{1}}^{\Tft_{1}'\Gammaft(1)}\tau'\right)$ by \cite[Theorem 2.17 (7)]{CMO19}. If we combine the last two isomorphisms and apply Proposition \ref{P:two_functors_for_Iwahori_decomp} \eqref{P:two_functors_for_Iwahori_decomp2}, we know ${\sigma'}^{\Nf}\cong i_{\Uf^-_{1},\Uf_{1}}^{\Tft_{1}'\Gammaft(1)}\tau'$ as $\Tft_{1}'\Gammaft(1)$-modules. Thus by Proposition \ref{P:two_functions_are_adjoint},  
    \begin{equation*}
        \Hom_{\Gammaft(1)}(\sigma,{\sigma'}^{\Nf})\cong \Hom_{\Tft_{1}}(\tau,r_{\Uf_{1}^{-},\,\Uf_{1}}^{\Tft_{1}}\left(i_{\Uf^-_{1},\Uf_{1}}^{\Tft_{1}'\Gammaft(1)}\tau'\right)).
    \end{equation*}
    Now as $\Tft_{1}$-modules we have $r_{\Uf_{1},\Uf_{1}}^{\Tft_{1}}\left(i_{\Uf^-_{1},\Uf_{1}}^{\Tft_{1}'\Gammaft(1)}\tau'\right)\cong \tau'$ by Lemma \ref{L:IndRes2x} and Proposition \ref{P:two_functors_for_Iwahori_decomp} \eqref{P:two_functors_for_Iwahori_decomp2}. The isomorphism \eqref{ExtThmIso1} follows because $\tau$ appears in $\tau'|_{\Tft_1}$ with multiplicity one.

    By isomorphism \eqref{ExtThmIso1}, we can identify $\sigma$ with the $\sigma$-isotypic subspace of ${\sigma'}^{\Nf}$, up to scaling. Then for any $k\in\Kft$ we can consider the $\Gammaft(1)$-submodule $k\sigma\subseteq {\sigma'}^{\Nf}$. As $\Gammaft(1)$-modules we have $k\sigma\cong {}^{k}\!\sigma$. But we know that ${}^k\!\sigma\cong \sigma$ as $\Gammaft(1)$-modules by Proposition \ref{IRep}. Since $\sigma$ is identified with the $\sigma$-isotypic subspace of ${\sigma'}^{\Nf}$, it follows that $k\sigma=\sigma$ as subspaces of ${\sigma'}^{\Nf}$. Thus $\sigma\subseteq{\sigma'}^{\Nf}$ is a $\Kft$-submodule. This proves the theorem.
\end{proof}

Recall we have the inclusions
\[
\widetilde{\Gammaf}(1)\subseteq \widetilde{\If}\subseteq \widetilde{\Kf},
\]
and the $\widetilde{\If}$-module $\sigma$ is an extension of $\sigma_1$. Now we prove that the extension of $\sigma_{1}$ to $\widetilde{{\Kf}}$ is also an extension of $\sigma$, namely
\begin{Prop}\label{ExtICompat}
   Let $\sigma_{\widetilde{{\Kf}}}$ be an arbitrary extension of $\sigma_{1}$ to $\widetilde{\Kf}$. Then 
   \[
   \sigma_{\widetilde{{\Kf}}}|_{\widetilde{{\If}}}\cong \sigma.
   \]
\end{Prop}

\begin{proof}
First note that the set of extensions of $\sigma_{1}$ to $\widetilde{{\If}}$ is a torsor under the group $\Hom(\widetilde{{\If}}/\widetilde{{\Gammaf}}(1),\mathbb{C}^{\times})$. Since $\Phi$ is simply-laced, any such character must factor through $T_{\kappa}$. Thus every extension of $\sigma_{1}$ to $\widetilde{{\If}}$ is of the form $\chi\otimes \sigma$ for some $\chi\in \Hom(T_{\kappa},\mathbb{C}^{\times})$. 

Hence, the restriction of $\sigma_{\widetilde{{\Kf}}}$ to $\widetilde{{\If}}$ is of the form $\chi\otimes \sigma$ for some $\chi\in \Hom(T_{\kappa},\mathbb{C}^{\times})$. We will show $\chi=1$. Let $\alpha\in\Delta$. We restrict further to $\widetilde{{\If}}\cap\,^{w_{\alpha}}\widetilde{{\If}}$ and apply Proposition \ref{IRep} \eqref{ResSigAlpha} to see that the restriction of $\sigma_{\widetilde{{\Kf}}}$ to $\widetilde{{\If}}\cap\,^{w_{\alpha}}\widetilde{{\If}}$ is isomorphic to $\chi\otimes \sigma_{\alpha}$. Since $\sigma_{\widetilde{{\Kf}}}$ and $\sigma_{\alpha}$ are $w_{\alpha}$ invariant by Theorem \ref{IRep} \eqref{sigWInvar} and \eqref{SigAWInv}, it follows that $\chi$ must be $w_{\alpha}$ invariant. Since $\Phi$ is simply-laced with rank $>1$ the character $\chi$ must be trivial on elements of the form $h_{\alpha}(t)$ for all $t\in \Ocal^{\times}/(1+\varpi^{2e}\Ocal)$, which implies that $\chi$ is trivial since $G$ is simply-connected.
\end{proof}


%%%%%%%%%%%%%%%%%%%%%%%%%%%%%%%%%%%%%%%%%%%%%%%%%%%%%%%%%%%%%%%%%%%%%%%

\subsection{Additional properties of $\sigma$}\label{SigProps}

%%%%%%%%%%%%%%%%%%%%%%%%%%%%%%%%%%%%%%%%%%%%%%%%%%%%%%%%%%%%%%%%%%%%%%%

The $\widetilde{{\If}}$-module $\sigma=i_{{\Uf}_{0},{\Uf}^{-}_{1}}\tau$ exhibits some additional properties that will be useful in our investigation of a Hecke algebra on $\Gt$. We record these in this subsection. Recall from \eqref{E:I_and_I_2_def} that ${\If}_{2}\subseteq {\If}$ is the subgroup generated by the subgroups ${\Tf}_{0}$, ${\Uf}_{0}$ and ${\Uf}^{-}_{2}$. Let
\[
\sigma_{2}:=i_{{\Uf}_{0},{\Uf}^{-}_{2}}^{\Ift_2}\tau,
\]
which is an irreducible $\widetilde{\If}_2$-module by Proposition \ref{P:two_functors_for_Iwahori_decomp}.

\begin{Lem} \label{+-IdemSwap}

	Let $\alpha\in\Phi^{+}$ and $t\in \mathcal{O}$. Then in the group algebra of $\widetilde{{\Kf}}$, we have
	\begin{equation*}
		e_{{\Uf}_{0}}e_{{\Uf}^{-}_{2}}x_{\alpha}(\varpi^{2e-1}t)=e_{{\Uf}_{0}}e_{{\Uf}^{-}_{2}}.
	\end{equation*}
\end{Lem}

\begin{proof}
	 For all $y\in {\Uf}^{-}_{2}$ the commutator $[y,x_{\alpha}(\varpi^{2e-1}t)]$ is trivial mod $\Gamma(2e)$, which gives
     \[
     e_{{\Uf}^{-}_{2}}x_{\alpha}(\varpi^{2e-1}t)=x_{\alpha}(\varpi^{2e-1}t)e_{{\Uf}^{-}_{2}}.
     \]
     The result follows because $x_{\alpha}(\varpi^{2e-1}t)\in {\Uf}_{0}$.
\end{proof}

\begin{comment}
\begin{Lem} \label{+-IdemSwap}

	Let $\alpha\in\Phi^{+}$ and $t\in \mathcal{O}$. Let $\Sigma^{\pm}\subseteq\Phi$ a positive (negative) system of roots (not necessarily $\Phi^{\pm}$). 
	
	Suppose that $({\Uf}_{\Sigma^{+}},T_{0},{\Uf}_{\Sigma^{-}})$ is the Iwahori factorization of $\widetilde{{\If}}_{2}$ associated to the positive system of roots defined by $\Sigma^{+}$. 
	
	Then
	\begin{equation*}
		e_{{\Uf}_{\Sigma^{+}}}e_{{\Uf}_{\Sigma^{-}}}x_{\alpha}(\varpi^{2e-1}t)=e_{{\Uf}_{\Sigma^{+}}}e_{{\Uf}_{\Sigma^{-}}}.
	\end{equation*}
\end{Lem}

\begin{proof}
	Note that we have the following (set) factorizations into subgroups.
	\begin{align*}
		{\Uf}_{\Sigma^{-}} \cong& {\Uf}_{\Phi^{+}\cap \Sigma^{-},0}\times {\Uf}_{\Phi^{-}\cap \Sigma^{-},2}\\
		{\Uf}^{+}_{0} \cong& {\Uf}_{\Phi^{+}\cap \Sigma^{+},0}\times {\Uf}_{\Phi^{+}\cap \Sigma^{-},0}\\
		{\Uf}_{\Sigma^{+}}\cong& {\Uf}_{\Phi^{-}\cap \Sigma^{+},2}\times {\Uf}_{\Phi^{+}\cap \Sigma^{+},0}
	\end{align*}
	Thus we have the following factorizations of idempotents.
	\begin{align*}
		e_{{\Uf}_{\Sigma^{-}}} \cong& e_{{\Uf}_{\Phi^{+}\cap \Sigma^{-},0}}e_{{\Uf}_{\Phi^{-}\cap \Sigma^{-},2}}\\
		e_{{\Uf}^{+}_{0}} \cong& e_{{\Uf}_{\Phi^{+}\cap \Sigma^{+},0}} e_{{\Uf}_{\Phi^{+}\cap \Sigma^{-},0}}\\
		e_{{\Uf}_{\Sigma^{+}}}\cong& e_{{\Uf}_{\Phi^{-}\cap \Sigma^{+},2}} e_{{\Uf}_{\Phi^{+}\cap \Sigma^{+},0}}
	\end{align*}
	
	Thus we have
	\begin{equation*}
		e_{{\Uf}_{\Sigma^{+}}}e_{{\Uf}_{\Sigma^{-}}}= e_{{\Uf}_{\Phi^{-}\cap \Sigma^{+},2}} e_{{\Uf}^{+}_{0}} e_{{\Uf}_{\Phi^{-}\cap \Sigma^{-},2}}
	\end{equation*}
	
	Now the result follows because for all $y\in U^{-}(\varpi^{2}\mathcal{O})$ the commutator $[y,x_{\alpha}(\varpi^{2e-1}t)]$ is trivial mod $4$, and $x_{\alpha}(\varpi^{2e-1}t)\in U^{+}_{0}$.
\end{proof}
\end{comment}

\begin{Prop}\label{Gamma02IrrInd}
	The $\widetilde{{\If}}$-module $\Ind_{\widetilde{{\If}}_{2}}^{\widetilde{{\If}}}\sigma_{2}$ is irreducible.
\end{Prop}

\begin{proof}
It suffices to show $\End_{\widetilde{{\If}}}(\Ind_{\widetilde{{\If}}_{2}}^{\widetilde{{\If}}}\sigma_{2})\cong \mathbb{C}$. But by the Frobenius reciprocity and Proposition \ref{P:two_functions_are_adjoint} (2), we have
	\begin{align*}
		\End_{\widetilde{{\If}}}(\Ind_{\widetilde{{\If}}_{2}}^{\widetilde{{\If}}}\sigma_{2})\cong& \Hom_{\widetilde{{\If}}_{2}}(\sigma_{2},\Ind_{\widetilde{{\If}}_{2}}^{\widetilde{{\If}}}\sigma_{2})\\
		\cong& \Hom_{\widetilde{{\Tf}}_{0}}(\tau, r_{{\Uf}^{+}_{0},{\Uf}^{-}_{2}}^{\widetilde{{\If}}_{2}}\Ind_{\widetilde{{\If}}_{2}}^{\widetilde{{\If}}}\sigma_{2}).
	\end{align*}
	
	By definition
	\begin{equation*}
		r_{{\Uf}^{+}_{0},{\Uf}^{-}_{2}}^{\widetilde{{\If}}_{2}}\Ind_{\widetilde{{\If}}_{2}}^{\widetilde{{\If}}}\sigma_{2}=e_{{\Uf}^{+}_{0}}e_{{\Uf}^{-}_{2}}\mathcal{H}(\widetilde{{\If}}_{2})\otimes_{\widetilde{{\If}}_{2}} \mathcal{H}(\widetilde{{\If}})\otimes_{\widetilde{{\If}}_{2}}\mathcal{H}(\widetilde{{\If}}_{2})e_{{\Uf}^{+}_{0}}e_{{\Uf}^{-}_{2}}\otimes_{\widetilde{{\Tf}}_{0}}\tau.
	\end{equation*}
Thus by the Iwahori factorization and Lemma \ref{L:spanning_vectors}, $r_{{\Uf}^{+}_{0},{\Uf}^{-}_{2}}^{\widetilde{{\If}}_{2}}\Ind_{\widetilde{{\If}}_{2}}^{\widetilde{{\If}}}\sigma_{2}$ is spanned by elements of the form 
\begin{equation*}
	e_{{\Uf}^{+}_{0}}e_{{\Uf}^{-}_{2}}\otimes_{\widetilde{{\If}}_{2}} y\otimes_{\widetilde{{\If}}_{2}}e_{{\Uf}^{+}_{0}}e_{{\Uf}^{-}_{2}}\otimes_{\widetilde{{\Tf}}_{0}} w,
\end{equation*}	
where $y\in {\Uf}^{-}_{1}$ and $w\in\tau$.

We claim that if $y\in {\Uf}^{-}_{1}\smallsetminus{\Uf}^{-}_{2}$, then $e_{{\Uf}^{+}_{0}}e_{{\Uf}^{-}_{2}}\otimes _{\widetilde{{\If}}_{2}}y\otimes_{\widetilde{{\If}}_{2}}e_{{\Uf}^{+}_{0}}\otimes_{\widetilde{{\Tf}}_{0}} w=0$. To prove the claim, let us write $y=y^{\prime}x_{-\alpha}(u)$, where $u\in\varpi\mathcal{O}^{\times}$ and $y^{\prime}\in \prod_{\beta}{\Uf}_{-\beta,1}$, where the product is over $\beta\in \Phi^{+}\smallsetminus\{\alpha\}$. Since $v\varpi\mathcal{O}\subseteq 4\mathcal{O}$, there exists $v\in \varpi^{2e-1}\mathcal{O}^{\times}$ such that $(1+ uv,\varpi)=-1$ by Corollary \ref{C:Hilbert_symbol_unique_element_integral_radical}. Because $v\in \varpi^{2e-1}\mathcal{O}$ and $y^{\prime}\in \prod_{\beta}{\Uf}_{-\beta,1}$ where $\beta\in \Phi^{+}\smallsetminus\{\alpha\}$, we have $y^{\prime}\xt_{\alpha}(v)=\xt_{\alpha}(v)y^{\prime}$.

We then have
\begin{align*}
	&e_{{\Uf}^{+}_{0}}e_{{\Uf}^{-}_{2}}\otimes_{\widetilde{{\If}}_{2}}y\otimes_{\widetilde{{\If}}_{2}}e_{{\Uf}^{+}_{0}}e_{{\Uf}^{-}_{2}}\otimes_{\widetilde{{\Tf}}_{0}} w\\
    =&e_{{\Uf}^{+}_{0}}e_{{\Uf}^{-}_{2}}\otimes_{\widetilde{{\If}}_{2}}y\otimes_{\widetilde{{\If}}_{2}}\xt_{\alpha}(v)e_{{\Uf}^{+}_{0}}e_{{\Uf}^{-}_{2}}\otimes_{\widetilde{{\Tf}}_{0}} w\quad(\text{because $\xt_{\alpha}(v)e_{{\Uf}^{+}_{0}}=e_{{\Uf}^{+}_{0}}$})\\
	=&e_{{\Uf}^{+}_{0}}e_{{\Uf}^{-}_{2}}\otimes_{\widetilde{{\If}}_{2}}y^{\prime}\xt_{-\alpha}( u)\xt_{\alpha}(v)\otimes_{\widetilde{{\If}}_{2}}e_{{\Uf}^{+}_{0}}e_{{\Uf}^{-}_{2}}\otimes_{\widetilde{{\Tf}}_{0}} w\\
	=&-e_{{\Uf}^{+}_{0}}e_{{\Uf}^{-}_{2}}\otimes_{\widetilde{{\If}}_{2}}y^{\prime}\xt_{\alpha}(v)\htt_{\alpha}(1+uv)\xt_{-\alpha}(u)\otimes_{\widetilde{{\If}}_{2}}e_{{\Uf}^{+}_{0}}e_{{\Uf}^{-}_{2}}\otimes_{\widetilde{{\Tf}}_{0}} w\quad (\text{Lemma \ref{+-IdemSwap}})\\
	=&-e_{{\Uf}^{+}_{0}}e_{{\Uf}^{-}_{2}}\otimes_{\widetilde{{\If}}_{2}}\xt_{\alpha}(v)y^{\prime}\xt_{-\alpha}(u)\otimes_{\widetilde{{\If}}_{2}}e_{{\Uf}^{+}_{0}}e_{{\Uf}^{-}_{2}}\otimes_{\widetilde{{\Tf}}_{0}} w\quad (\text{by $y^{\prime}\xt_{\alpha}(v)=\xt_{\alpha}(v)y^{\prime}$})\\
	=&-e_{{\Uf}^{+}_{0}}e_{{\Uf}^{-}_{2}}\otimes_{\widetilde{{\If}}_{2}}y\otimes_{\widetilde{{\If}}_{2}}e_{{\Uf}^{+}_{0}}e_{{\Uf}^{-}_{2}}\otimes_{\widetilde{{\Tf}}_{0}} w,
\end{align*}
where for the last equality we used
\[
e_{{\Uf}^{+}_{0}}e_{{\Uf}^{-}_{2}}\xt_{\alpha}(v)=e_{{\Uf}^{+}_{0}}e_{{\Uf}^{-}_{2}},
\]
which holds because $\xt_{\alpha}(v)$ normalizes $\Uf_2^-$ and $\xt_{\alpha}(v)\in\Uf_0^+$. Thus $e_{{\Uf}^{+}_{0}}e_{{\Uf}^{-}_{2}}\otimes_{\widetilde{{\If}}_{2}}ye_{{\Uf}^{+}_{0}}e_{{\Uf}^{-}_{2}}\otimes_{\widetilde{{\Tf}}_{0}} w=0$, and so we have $r_{{\Uf}^{+}_{0},{\Uf}^{-}_{2}}^{\widetilde{{\If}}_{2}}\Ind_{\widetilde{{\If}}_{2}}^{\widetilde{{\If}}}\sigma_{2}\cong \tau$ as $\widetilde{{\Tf}}_{0}$-modules.

Therefore $\End_{\widetilde{{\If}}}(\Ind_{\widetilde{{\If}}_{2}}^{\widetilde{{\If}}}\sigma_{2})\cong \mathbb{C}$ and the result follows.
\end{proof}

\begin{Prop}\label{IndSig2}
	As $\widetilde{\If}$-modules
	\begin{equation*}
		\sigma \cong \Ind_{\widetilde{\If}_2}^{\widetilde{\If}}\sigma_{2}.
	\end{equation*}
\end{Prop}
\begin{proof}
	By Lemma \ref{LHoms} we know that $\Hom_{\widetilde{\If}}(\sigma,\Ind_{\widetilde{\If}_2}^{\widetilde{\If}}\sigma_{2})\neq 0$. By Lemma \ref{Gamma02IrrInd} we know that $\Ind_{\widetilde{\If}_{2}}^{\widetilde{\If}}\sigma_{2}$ is irreducible. Since $\sigma$ is also irreducible the result follows by Schur's Lemma.
\end{proof}


\begin{Lem}\label{pmUniComm}
	Let $\ell\in \mathbb{Z}$ be such that $0\leq \ell\leq e$. Then the groups ${\Uf}_{2e-2\ell}$ and ${\Uf}^{-}_{2\ell}$ commute.
\end{Lem}
\begin{proof}
Let $\alpha,\beta\in\Phi^{+}$ be such that $\beta\neq \alpha$. Then for each $u,v\in\mathcal{O}$, one can readily derive the following two identities.
	\begin{align*}
		[x_{\alpha}(\varpi^{2e-2\ell}u),x_{-\beta}(\varpi^{2\ell}v)]=&1 \quad (\text{by \eqref{E:MR2'} viewed mod $\Gamma(2e)$});\\
		x_{\alpha}(\varpi^{2e-2\ell}u)x_{-\alpha}(\varpi^{2\ell}v)=&x_{-\alpha}(\varpi^{2\ell}v)x_{\alpha}(\varpi^{2e-2\ell}u)\quad (\text{by Proposition \ref{+-IdemSwap}}).
	\end{align*}
	We note that in the second identity there is no Hilbert symbol $(1+\varpi^{2e}uv,\varpi^{2\ell}v)$ because $1+\varpi^{2e}uv\in R$ and if $v\in \varpi\Ocal$ then $1+\varpi^{2e}uv\in 1+\varpi^{2e+1}\Ocal$. 
\end{proof}

%\textbf{NOTE: \cite[Example 2.15]{CMO19}. Can we use this?}


\begin{Prop}\label{Sig2UInvar} In the $\widetilde{{\If}}_{2}$-module $\sigma_{2}$ we have
	\begin{equation*}
		\sigma_{2}^{{\Uf}_{2e-2}}=\sigma_{2}.
	\end{equation*}
\end{Prop}

\begin{proof}
	By the Iwahori factorization and Lemma \ref{L:spanning_vectors}, $\sigma_{2}$ is spanned by elements of the form $ye_{{\Uf}_{0}}e_{{\Uf}^{-}_{2}}\otimes_{\Tt_{0}} w$, where $y\in {\Uf}^{-}_{2}$ and $w\in \tau$. Let $x\in \Uf_{2e-2}$, so that $xy=yx$ by Lemma \ref{pmUniComm}. Hence 
    \[
    xye_{{\Uf}_{0}}e_{{\Uf}^{-}_{2}}\otimes_{\Tt_{0}} w=yxe_{{\Uf}_{0}}e_{{\Uf}^{-}_{2}}\otimes_{\Tt_{0}} w=ye_{{\Uf}_{0}}e_{{\Uf}^{-}_{2}}\otimes_{\Tt_{0}} w,
    \]
    because $e_{\Uf_0}$ absorbs $x$. The result follows.
\end{proof}

\begin{comment}

\begin{Prop}
	The subspace $\sigma_{2}^{{\Uf}_{-\alpha,2e-2}}$ is a $ I_2^{\alpha}:= I_2\cap \,^{w_{\alpha}} I_2$-module and it is isomorphic to $\sigma_{2}^{\alpha}:=i_{U_{0}^{\alpha},U^{-}_{2}}\tau$.
\end{Prop}

\textbf{Might be able to generalize in a few directions. Should definitely try for the affine simple reflections; more general weyl group elements; other valuation values.}

\begin{proof}
	\textbf{To Do.}
\end{proof}

\end{comment}

%%%%%%%%%%%%%%%%%%%%%%%%%%%%%%%%%%%%%%%%%%%%%%%%%%%%%%%%%%%%%%%%%%%%%%%%%

\subsection{Brief wrap-up}

%%%%%%%%%%%%%%%%%%%%%%%%%%%%%%%%%%%%%%%%%%%%%%%%%%%%%%%%%%%%%%%%%%%%%%%%%

Let us briefly summarize what we have achieved so far. We have first constructed a pseudo-spherical representation $\tau$ of the compact group $\Tt_0$, which is a finite dimensional irreducible representation of $\Tt_0$ that is trivial on $T_R$. We showed that the pseudo-spherical representations are Weyl invariant. 

Next using parahoric induction, we obtained an irreducible representation $\sigma:=i_{\Uf, \Uf_1^-}^{\Ift}\tau$ of the finite group $\Ift\cong \It/\s(\Gamma(2e))$, which we can inflate to the Iwahori subgroup $\It$. This representation $\sigma$ has two important properties: first, it remains irreducible when restricted to $\Gammat(1)$, and in fact can be viewed as being parahorically induced; second, $\sigma$ extends to $\Kft=\Kt/\s(\Gamma(2e))$, which we also call a pseudo-spherical representation and denote it by the same symbol $\sigma$. We can inflate $\sigma$ to a representation of the compact group $\Kt$. 

As in Remark \ref{R:Gamma(2e)_is_max}, the group $\Gamma(2e)$ is the maximal congruence subgroup of the hyperspecial maximal compact subgroup $K$ that splits to $\Kt$ with normal image. Hence the representation $\sigma$ of $\Kt$ can be considered as the covering analog of the trivial representation. Hence it makes sense to call $\sigma$ a ``pseudo-spherical" representation of $\Kt$. Indeed, $\sigma$ does play a role of the trivial representation in the representation theory of covering groups.


%%%%%%%%%%%%%%%%%%%%%%%%%%%%%%%%%%%%%%%%%%%%%%%%%%%%%%%%%%%%%%%%%%%%%%%%%

\subsection{Finite Shimura correspondence}\label{SS:finite_Shimura}

%%%%%%%%%%%%%%%%%%%%%%%%%%%%%%%%%%%%%%%%%%%%%%%%%%%%%%%%%%%%%%%%%%%%%%%%%

A first instance where one can see how our pseudo-spherical representation plays a role of the trivial representation is a ``finite Shimura correspondence" originally established by Savin in \cite[Theorem 4.1]{S12}, when $F/\mathbb{Q}_{2}$ is unramified. In this section, we generalize it to an arbitrary $F$. Indeed, this correspondence establishes an equivalence of the categories of representations of the finite group $K\slash\Gamma(1)$ with a certain category of representations of $\Kft$, and this equivalence of categories is precisely given by tensoring the pseudo-spherical representation as $\pi\mapsto \pi\otimes\sigma$ for a representation $\pi$ of $K\slash\Gamma(1)$.


To be precise, let $\mathcal{M}(G_\kappa)$ be the category of representations of $G_\kappa$, where we recall $\kappa$ is the residue field of our $2$-adic field $F$. Let $\mathcal{M}_{\sigma}(\widetilde{{\Kf}})$ be the full subcategory of smooth $\widetilde{{\Kf}}$-modules $\widetilde{\pi}$ such that $\widetilde{\pi}|_{\widetilde{\Gammaf}(1)}$ is a nontrivial $\sigma$-isotypic $\widetilde{\Gammaf}(1)$-module, or $\pit=0$. Let us then define two functors
\begin{align*}
\mathcal{F}_{\sigma}&:\mathcal{M}(G_{\kappa})\rightarrow \mathcal{M}_{\sigma}(\widetilde{{\Kf}}),\qquad \pi\mapsto \sigma\otimes \pi;\\
\mathcal{G}_{\sigma}&:\mathcal{M}_{\sigma}(\widetilde{{\Kf}})\rightarrow \mathcal{M}(G_{\kappa}),\qquad \widetilde{\pi}\mapsto \Hom_{\widetilde{\Gammaf}(1)}(\sigma,\widetilde{\pi}).
\end{align*}
Here, $\widetilde{{\Kf}}$ acts diagonally on $\sigma\otimes \pi$. Also $G_\kappa=\widetilde{\Kf}\slash\widetilde{\Gammaf}(1)$ acts on $\Hom_{\widetilde{\Gammaf}(1)}(\sigma,\widetilde{\pi})$ by ``conjugation", namely for each $g\in \widetilde{\Kf}\slash\widetilde{\Gammaf}(1)$ and $\phi\in \Hom_{\widetilde{\Gammaf}(1)}(\sigma,\widetilde{\pi})$
\[
(g\cdot \phi)(v)=g\phi(g^{-1}v)
\] 
for all $v\in \widetilde{\pi}$, which is indeed well-defined.

The main theorem of this subsection is the following ``finite Shimura correspondence".

\begin{Thm}\label{FiniteShimCor}
	The functors $\mathcal{F}_{\sigma}$ and $\mathcal{G}_{\sigma}$ are mutual quasi-inverses and hence define an equivalence of categories between $\mathcal{M}(G_{\kappa})$ and $\mathcal{M}_{\sigma}(\widetilde{{\Kf}})$.
\end{Thm}

We first need a couple of elementary lemmas.

\begin{Lem}\label{FinRkLem}
	Let $V_{1},V_{2},W$ be vector spaces. If $V_{1},V_{2}$ are finite dimensional, then the map 
    \[
    \Hom(V_{1},V_{2})\otimes W\rightarrow \Hom(V_{1},V_{2}\otimes W),\quad f\otimes w\mapsto (v\mapsto f(v)\otimes w),
    \]
    is an isomorphism of vector spaces. 
	
	Furthermore, suppose that $G$ is a group and $V_{1},V_{2},W$ are $G$-modules. Let $G$ act diagonally on tensor products and on Hom spaces via``conjugation''. Then this isomorphism is a $G$-module isomorphism. 
\end{Lem}	
\begin{proof}
Elementary exercise.
\end{proof}

\begin{comment}
\begin{proof}
    
	The map $\Hom(V_{1},V_{2})\otimes W\rightarrow \Hom(V_{1},V_{2}\otimes W)$ defined by $f\otimes w\mapsto (v\mapsto f(v)\otimes w)$ is always an injection, as can be seen by choosing a basis for $W$. Note that if $W$ is finite dimensional $F$ must be an isomorphism since the domain and codomain have the same dimension. 
	
	Now since $V_{1}$ is finite dimensional the image of any $f\in \Hom(V_{1},V_{2}\otimes W)$ must be finite dimensional and thus there exists a finite dimensional $W_{f}\subseteq W$ such that $f\in \Hom(V_{1},V_{2}\otimes W_{f})$. Now since $F$ restricted to $\Hom(V_{1},V_{2})\otimes W_{f}$ defines an isomorphism with $\Hom(V_{1},V_{2}\otimes W_{f})$ we see that $f$ is in the image of $F$ and so $F$ is surjective. 
	
	The second claim follows by a direct computation.
    
\end{proof}
\end{comment}



\begin{Lem}\label{HOMotimesSig}
    Let $M \subseteq L$ be finite groups. Let $\pi_1, \pi_2$ be (not necessarily finite dimensional) $L$-modules which are trivial on $M$, and let $\sigma_1, \sigma_2$ be irreducible (hence necessarily finite dimensional) $L$-modules which remain irreducible when restricted to $M$. If $\phi:\sigma_1\to\sigma_2$ is an $L$-module isomorphism, then the map
    \[
    \Hom_L(\pi_1,\pi_2)\longrightarrow\Hom_L(\pi_1\otimes\sigma_1, \pi_2\otimes\sigma_2),\quad f\mapsto f\otimes \phi,
    \]
    is an isomorphism of vector spaces.
\end{Lem}
\begin{proof}
    By the tenor-Hom adjunction, we have
    \[
     \Hom(\pi_1\otimes\sigma_1, \pi_2\otimes\sigma_2)\cong \Hom(\pi_1, \Hom(\sigma_1,\pi_2\otimes\sigma_2)),
    \]
    which gives
    \[
     \Hom_L(\pi_1\otimes\sigma_1, \pi_2\otimes\sigma_2)\cong \Hom_L(\pi_1, \Hom(\sigma_1,\pi_2\otimes\sigma_2)).
    \]
    Since $M$ acts trivially on $\pi_1$, the image of each element in $\Hom_L(\pi_1, \Hom(\sigma_1,\pi_2\otimes\sigma_2))$ lies in the $M$-invariant subspace of $\Hom(\sigma_1,\pi_2\otimes\sigma_2)$, implying
    \[
     \Hom_L(\pi_1\otimes\sigma_1, \pi_2\otimes\sigma_2)\cong \Hom_L(\pi_1, \Hom_M(\sigma_1,\pi_2\otimes\sigma_2)).
    \]
    
    Noting that both $\sigma_1$ and $\sigma_2$ are finite dimensional, Lemma \ref{FinRkLem} gives
    \[
    \Hom_M(\sigma_1,\pi_2\otimes\sigma_2)\cong \Hom_M(\sigma_1, \sigma_2)\otimes \pi_2\cong\C\otimes\pi_2\cong\pi_2,
    \]
    where we used $\Hom_M(\sigma_1, \sigma_2)\cong\C$ by Schur's lemma. Therefore, we have
    \[
    \Hom_L(\pi_1\otimes\sigma_1, \pi_2\otimes\sigma_2)\cong \Hom_L(\pi_1, \pi_2).
    \]

    All of the preceding isomorphisms are explicit once $\phi$ is used to normalize the isomorphism $\Hom_M(\sigma_1, \sigma_2)\cong\C$ coming from Schur's lemma, and it can be checked directly that $f\in \Hom_L(\pi_{1},\pi_{2})$ corresponds to $f\otimes \phi\in \Hom_L(\pi_{1}\otimes \sigma_{1},\pi_{2}\otimes \sigma_{2})$.
\end{proof}


\begin{comment}

\begin{Lem}\label{HOMotimesSig}
	Let $H\subseteq \widetilde{{\Kf}}$ be a subgroup such that $\widetilde{\Gamma}(1)\subseteq H$. Let $\pi_{1},\pi_{2}$ be $H$-modules trivial on $\widetilde{\Gamma}(1)$. Let $\sigma_{1},\sigma_{2}$ be two irreducible $\widetilde{{\Kf}}$-modules which remain irreducible when restricted to $\widetilde{\Gamma}(1)$. 
	
	If $\phi:\sigma_{1}\rightarrow \sigma_{2}$ is a $\widetilde{{\Kf}}$-module isomorphism, then the map $\Hom_{H}(\pi_{1},\pi_{2})\rightarrow \Hom_{H}(\pi_{1}\otimes \sigma_{1},\pi_{2}\otimes \sigma_{2})$ defined by $f\mapsto f\otimes \phi$ is an isomorphism of vector spaces.
\end{Lem}

\begin{proof}
	By the tensor-Hom adjunction we have
	\begin{equation*}
		\Hom(\pi_{1}\otimes \sigma_{1},\pi_{2}\otimes \sigma_{2})\cong \Hom(\pi_{1},\Hom(\sigma_{1},\pi_{2}\otimes \sigma_{2})).
	\end{equation*}
	Taking $\widetilde{\Gamma}(1)$-invariants and using that $\widetilde{\Gamma}(1)$ acts trivially on $\pi_{j}$ we get 
	\begin{equation*}
		\Hom_{\widetilde{\Gamma}(1)}(\pi_{1}\otimes \sigma_{1},\pi_{2}\otimes \sigma_{2})\cong \Hom(\pi_{1},\Hom_{\widetilde{\Gamma}(1)}(\sigma_{1},\pi_{2}\otimes \sigma_{2})).
	\end{equation*}
	Now taking $H$-invariants we get
	\begin{equation*}
		\Hom_{H}(\pi_{1}\otimes \sigma_{1},\pi_{2}\otimes \sigma_{2})\cong \Hom_{H}(\pi_{1},\Hom_{\widetilde{\Gamma}(1)}(\sigma_{1},\pi_{2}\otimes \sigma_{2})).
	\end{equation*}
	Now since $\sigma_{1},\sigma_{2}$ are finite dimensional we apply Lemma \ref{FinRkLem} to get
	\begin{equation*}
		\Hom_{H}(\pi_{1},\Hom_{\widetilde{\Gamma}(1)}(\sigma_{1},\pi_{2}\otimes \sigma_{2}))\cong \Hom_{H}(\pi_{1},\Hom_{\widetilde{\Gamma}(1)}(\sigma_{1},\sigma_{2})\otimes \pi_{2}).
	\end{equation*}
	Since $\sigma_{1},\sigma_{2}$ are an irreducible $\widetilde{\Gamma}(1)$-modules, Schur's Lemma implies 
	\begin{equation*}
		\Hom(\pi_{1},\Hom_{\widetilde{\Gamma}(1)}(\sigma_{1},\sigma_{2})\otimes \pi_{2})\cong\Hom(\pi_{1},\pi_{2}).
	\end{equation*}
	All of the preceding isomorphisms are explicit once $\phi$ is used to normalize the isomorphism coming from Schur's Lemma, and it can be checked directly that $f\in \Hom(\pi_{1},\pi_{2})$ corresponds to $f\otimes \phi\in \Hom_{\widetilde{\Gamma}(1)}(\pi_{1}\otimes \sigma_{1},\pi_{2}\otimes \sigma_{2})$.
\end{proof}

\end{comment}

Now, we are ready to prove the theorem.

\begin{proof}[Proof of Theorem \ref{FiniteShimCor}]
Noting that $G_\kappa=\widetilde{\Kf}\slash \widetilde{\Gammaf}(1)$, we can identify the category $\mathcal{M}(G_\kappa)$ with the category of $\widetilde{\Kf}$-modules which are trivial on $\widetilde{\Gammaf}(1)$. Then the above lemma (with $L=\widetilde{\Kf}$, $M=\widetilde{\Gammaf}(1)$ and $\sigma_1=\sigma_2=\sigma$) means that the functor $\mathcal{F}_\sigma$ is fully faithful. (Note that we necessarily have $\Hom_L(\pi_1,\pi_2)=\Hom_{L\slash M}(\pi_1,\pi_2)$ in the above lemma.)

Hence to show $\mathcal{F}_\sigma$ is an equivalence of the categories, it suffices to show that the functor $\mathcal{F}_\sigma$ is essentially surjective on the objects. But for each $\pit\in\mathcal{M}_\sigma(\widetilde{\Kf})$, we have
\[
\mathcal{F}_\sigma\circ\mathcal{G}_\sigma(\pit)=\sigma\otimes \Hom_{\widetilde{\Gammaf}(1)}(\sigma, \pit).
\]
Now, we have the natural $\widetilde{\Kf}$-isomorphism
\[
\sigma\otimes \Hom_{\widetilde{\Gammaf}(1)}(\sigma, \pit)\longrightarrow \pit,\quad  v\otimes f\mapsto f(v),
\]
because $\pit|_{\widetilde{\Gammaf}(1)}$ is $\sigma$-isotypic. Hence the functor $\mathcal{F}_\sigma$ is essentially surjective on the objects.

Now that we have shown that $\mathcal{F}_\sigma\circ\mathcal{G}_\sigma\cong \mathrm{id}_{\mathcal{M}_\sigma(\widetilde{\Kf})}$ and $\mathcal{F}_\sigma$ is an equivalence of the categories, we know that $\mathcal{G}_\sigma$ is a quasi-inverse of $\mathcal{F}_\sigma$.


\begin{comment}
	First we show that the composition $\mathcal{G}_{\sigma}\circ\mathcal{F}_{\sigma}$ is naturally isomorphic to the identity functor $\mathrm{id}_{\mathcal{M}(G_{\kappa})}$. We will just check this on objects and omit the compatibility on morphisms, which can also be checked directly.
	
	Let $\pi$ be a $G_{\kappa}$-module. By definition $\mathcal{G}_{\sigma}\circ\mathcal{F}_{\sigma}(\pi)=\Hom_{\widetilde{\Gamma}(1)}(\sigma,\sigma\otimes \pi)$. By Lemma \ref{FinRkLem} we have a $\widetilde{{\Kf}}$-module isomorphism $\Hom(\sigma,\sigma)\otimes \pi \rightarrow \Hom(\sigma,\sigma\otimes \pi)$ defined by $f \otimes v \mapsto (w \mapsto f(w)\otimes v)$. Note that the action of $\widetilde{{\Kf}}$ factors through ${\Kf}$.
	
	Taking $\widetilde{\Gamma}(1)$-invariant gives a $G_{\kappa}$-module isomorphism $\Hom_{\widetilde{\Gamma}(1)}(\sigma,\sigma)\otimes \pi \cong \Hom_{\widetilde{\Gamma}(1)}(\sigma,\sigma\otimes \pi)$. Now since $\sigma$ is an irreducible $\widetilde{\Gamma}(1)$-module that extends to $\widetilde{{\Kf}}$, Schur's Lemma implies that $\Hom_{\widetilde{\Gamma}(1)}(\sigma,\sigma)\otimes \pi\cong \pi$ as $G_{\kappa}$-modules.
	
	Second we show that the composition $\mathcal{F}_{\sigma}\circ\mathcal{G}_{\sigma}$ is naturally isomorphic to the identity functor $\mathrm{id}_{\mathcal{M}_{\sigma}(\widetilde{{\Kf}})}$. Again we omit the compatibility on morphisms.
	
	Let $\widetilde{\pi}$ be an object in $\mathcal{M}_{\sigma}(\widetilde{{\Kf}})$. By definition $\mathcal{F}_{\sigma}\circ\mathcal{G}_{\sigma}(\widetilde{\pi})=\sigma \otimes \Hom_{\widetilde{\Gamma}(1)_{4}}(\sigma,\widetilde{\pi})$. The natural map $\sigma \otimes \Hom_{\widetilde{\Gamma}(1)_{4}}(\sigma,\widetilde{\pi})\rightarrow \widetilde{\pi}$ defined by $v\otimes f\mapsto f(v)$ is an $\widetilde{{\Kf}}$-isomorphism, since $\widetilde{\pi}$ restricted to $\widetilde{\Gamma}(1)$ is $\sigma$-isotypic.
\end{comment}
    
\end{proof}

The next proposition shows that the equivalence of categories of Theorem \ref{FiniteShimCor} commutes with induction in certain cases.

\begin{Prop}\label{SCInd}
Assume $\Hft$ is such that $\Gammaft(1)\subseteq\Hft\subseteq\Kft$. Let $\pi$ be a representation of $\Hft$ that is trivial on $\widetilde{{\Gammaf}}(1)$, so that $\pi$ is really a representation of $\Hft/\widetilde{{\Gammaf}}(1)\hookrightarrow G_{\kappa}=\Kft\slash\Gammaft(1)$.
	\begin{enumerate}
		\item The $\widetilde{{\Kf}}$-module $\Ind_{\Hft}^{\widetilde{{\Kf}}}(\pi\otimes\sigma)$ is an object in $\mathcal{M}_{\sigma}(\widetilde{{\Kf}})$.\label{IndInCat}
		%\item The map $\phi:\mathbb{C}[\widetilde{{\Kf}}]\otimes_{H}(\pi\otimes\sigma)\rightarrow \mathcal{F}_{\sigma}(\mathbb{C}[\widetilde{{\Kf}}]\otimes_{H} \pi)=(\mathbb{C}[\widetilde{{\Kf}}]\otimes_{H} \pi)\otimes \sigma$ defined by $k\otimes_{H}(v\otimes w)\mapsto (k\otimes_{H}v)\otimes kw$ is an isomorphism of $\widetilde{{\Kf}}$-modules.\label{cIndSigIso} %\textbf{This map does not look linear in $k$, but this should be okay as the group elements make up a basis.}
		\item The map 
        \[
        \phi:\left(\Ind_{\Hft}^{\widetilde{{\Kf}}}\pi\right)\otimes \sigma\longrightarrow \Ind_{\Hft}^{\widetilde{{\Kf}}}(\pi\otimes \sigma),\quad f\otimes w\mapsto (k\mapsto f(k)\otimes kw),
        \]
        is a surjective map of $\widetilde{{\Kf}}$-modules. Furthermore, if $\pi$ is finite dimensional, then both sides have the same dimension, so that $\pi$ is an isomorphism.\label{IndSigIso}
		\item The map 
        \[\End_{G_{\kappa}}(\Ind_{\Hft}^{\widetilde{{\Kf}}}\pi)\rightarrow \End_{\widetilde{{\Kf}}}(\Ind_{\Hft}^{\widetilde{{\Kf}}}(\pi\otimes \sigma)),\quad F\mapsto \phi\circ(F\otimes id_{\sigma})\circ \phi^{-1},
        \]
        is an isomorphism of $\mathbb{C}$-algebras. \label{EndIso}
	\end{enumerate}
\end{Prop}

\begin{proof}
	\eqref{IndInCat} follows from the Mackey theory and $\widetilde{{\Gammaf}}(1)\cap \,^{g}\Hft=\widetilde{{\Gammaf}}(1)$ for any $g\in \widetilde{{\Kf}}$.
	
	We prove \eqref{IndSigIso}. A direct computation shows that the map is well defined and a $\widetilde{{\Kf}}$-module homomorphism. We will show this is surjective. Let $W$ be the subspace of $\Ind_{\Hft}^{\widetilde{{\Kf}}}\pi$ consisting of functions with support contained in $\Hft$. Let $W^{\prime}$ be the subspace of $\Ind_{\Hft}^{\widetilde{{\Kf}}}(\pi\otimes \sigma)$ consisting of functions with support contained in $\Hft$. Recall that as $\Hft$-modules $W\cong \pi$ and $W^{\prime}\cong \pi\otimes \sigma$. In both cases the isomorphism is given by evaluation at the identity. Under these identifications, one can see that $\phi:W\otimes \sigma\rightarrow W^{\prime}$ becomes the identity map. Thus $W^{\prime}$ is in the image of $\phi$. Since the index $[\widetilde{{\Kf}}:\Hft]$ is finite and $\Ind_{\Hft}^{\widetilde{{\Kf}}}(\pi\otimes \sigma)$ is generated by $W^{\prime}$ as a $\widetilde{{\Kf}}$-module, it follows that $\phi$ is surjective. 
	
	%The map is well defined because for any $h\in H$ we have $(kh\otimes_{H}v)\otimes khw=(k\otimes_{H}hv)\otimes khw$. We can also see from the definition that the map is surjective. As the dimension of the domain and codomain is $[\widetilde{{\Kf}}:H](\dim\pi)(\dim\sigma)$ the map is a linear isomorphism. Finally a direct computation shows that the map is a $\widetilde{{\Kf}}$-module homomorphism. \textbf{Don't want to assume that $\pi$ is finite dimensional. Should just be able to write down the inverse map.}
	
	\eqref{EndIso} follows from \eqref{IndSigIso} and Theorem \ref{FiniteShimCor}. 
\end{proof}

The isomorphism of Proposition \ref{SCInd} \eqref{EndIso} has a very important property in that it preserves support, when viewed as an isomorphism of Hecke algebras. We return to this point in \S \ref{FinShimHecke} after we introduce the notation for Hecke algebras.  

%%%%%%%%%%%%%%%%%%%%%%%%%%%%%%%%%%%%%%%%%%%%%%%%%%%%%%%%%%%%%%%%%%%%%%%

\subsection{Whittaker Invariants}\label{SS:WhitInvar}

%%%%%%%%%%%%%%%%%%%%%%%%%%%%%%%%%%%%%%%%%%%%%%%%%%%%%%%%%%%%%%%%%%%%%%%


In this subsection we record a few simple results on twisted ${\Uf}_{0}$-invariants for representations in $\mathcal{M}_{\sigma}(\widetilde{{\Kf}})$. (Though we will not use it for our main theorems, it is certainly of some interest.) Recall that $\sigma^{{\Uf}_{0}}\cong \tau$ as $\Tft_{0}$-modules.


\begin{Cor}
	Let $\psi\in\Hom({\Uf}_{0},\mathbb{C}^{\times})$ be a nontrivial character that is trivial on ${\Uf}_{1}$. Then
	\begin{equation*}
		\sigma^{({\Uf}_{0},\psi)}=0. 
	\end{equation*}
\end{Cor}

\begin{proof}
	This follows directly from Lemma \ref{UInvar} because $\sigma^{({\Uf}_{0},\psi)}\subseteq \sigma^{\Uf_1}$.
\end{proof}

Thus we see that for $\sigma$ to support a $\psi$-Whittaker functional the conductor of $\psi$ must be strictly larger than $1$.


\begin{Cor}\label{WhitMult}
	Let $\pi$ be a $G_{\kappa}$-module and let $\psi\in\Hom({\Uf}_{0},\mathbb{C}^{\times})$. Then 
	\begin{equation*}
		(\pi\otimes \sigma)^{({\Uf}_{0},\psi)}\supseteq \bigoplus_{\psi=\psi_{1}\psi_{2}}\pi^{({\Uf}_{0},\psi_{1})}\otimes \sigma^{({\Uf}_{0},\psi_{2})},
	\end{equation*}
    where the direct sum is over all $\psi_{1},\psi_{2}\in\Hom({\Uf}_{0},\mathbb{C}^{\times})$ such that $\psi=\psi_{1}\psi_{2}$ and $\psi_{1}$ is trivial on ${\Uf}_{1}$. 
    In particular, if $\psi_{1}$ is nondegenerate, $\psi_{2}$ is trivial, and $\pi$ is $\psi_{1}$-generic, then
	\begin{equation*}
		\dim(\pi\otimes \sigma)^{({\Uf}_{0},\psi)}\geq \dim\tau=\left|Y\slash \Yt\right|^{\frac{ef}{2}}.
	\end{equation*}
\end{Cor}

From Corollary \ref{WhitMult} we see that multiplicity one for Whittaker functionals fails for the finite group $\Kft$. Thus it would be interesting to study the associated Gelfand-Graev representation in this finite group setting. Such an investigation will likely been needed to generalize the results of \cite{GGK24,GGK25} to the wild double covers studied in this paper.

%%%%%%%%%%%%%%%%%%%%%%%%%%%%%%%%%%%%%%%%%%%%%%%%%%%%%%%%%%%%%%%%%%%%%%%

\section{Hecke algebras}\label{S:Hecke}

%%%%%%%%%%%%%%%%%%%%%%%%%%%%%%%%%%%%%%%%%%%%%%%%%%%%%%%%%%%%%%%%%%%%%%%

Once we have constructed the Iwahori type $(\It, \sigma)$, which will be shown to be indeed a type in the sense of Bushnell-Kutzko, we can form the Iwahori Hecke algebra
\[
\Hcal(\Gt, \It; \sigma).
\]
In this section, we will prove that the support of this Hecke algebra is $\Wt_{\ea}$ (see \S \ref{SS:affine_Weyl_covering_group} for $\Wt_{\ea}$) and establish its IM-presentation. 

\subsection{Basic theory of Hecke algebras}\label{SS:HeckeReview}

In this subsection, we catalog basics of representations of compact groups and Hecke algebras to the extent necessary for our purposes, and fix our notation and conventions.

Let $G$ be a locally compact group and $K\subseteq G$ an open compact subgroup. Let $(\pi, V)$ be a finite dimensional (not necessarily irreducible) representation of $K$. We define the $\pi$-spherical Hecke algebra of $G$ by
\[
\Hcal(G, K; \pi):=\{f\in C_c^\infty(G, \End_{\C}(\pi^\vee))\st f(k_1gk_2)=\pi^\vee(k_1)f(g)\pi^\vee(k_2)\text{ for all $k_1, k_2\in K$}\}.
\]
This is an algebra under convolution. Specifically, for $f_{1},f_{2}\in \Hcal(G, K; \pi)$ we have
\begin{equation*}
	f_{1}*f_{2}:=\int_{G}f_{1}(h)f_{2}(h^{-1}g)dg=\int_{G}f_{1}(gh^{-1})f_{2}(h)dg,
\end{equation*}
where the Haar measure on $G$ is normalized so that the measure of $K$ is $1$.

Let $(\rho,W)$ be a $G$-module. The Hecke algebra $\Hcal(G, K; \pi^\vee)$ acts on $\Hom_{K}(V,W)$ on the right in the following way. Let $f\in \Hcal(G, K; \pi^\vee)$ and $A\in\Hom_{K}(V,W)$. Then
\begin{equation*}
	(A*f)(v):=\int_{G}\rho(g^{-1})A(f(g)v)dg.
\end{equation*}


The Hecke algebra $\Hcal(G, K; \pi^\vee)$ also has a a canonical left action on $\ind_{K}^{G}\pi$ defined by
    \begin{equation*}
        f*F(g):=\int_{G}f(x)F(x^{-1}g)dx,
    \end{equation*}
    where $f\in \Hcal(G, K; \pi^\vee)$ and $F\in \ind_{K}^{G}\pi$. The following lemma is well-known. (See, for example, \cite[(2.6)]{BK98}.)

\begin{Lem}\label{HeckeAsEnd}
    The canonical left action of $\Hcal(G, K; \pi^\vee)$ on $\ind_{K}^{G}\pi$ induces an isomorphism of $\mathbb{C}$-algebras $\Hcal(G, K; \pi)\cong \End_{G}(\ind_{K}^{G}\pi)$. 
\end{Lem}

For each $g\in G$, we let
\[
\Hcal(G, K; \pi)_g
\]
be the vector space of all $f\in \Hcal(G, K; \pi)$ that vanish outside the double coset $KgK$.

A very important fact we repeatedly use in this paper is 
\begin{Lem}\label{L:Bushnell_Kutzko_equivalence}
For each $g\in G$, we have canonical isomorphisms
\[
\Hom_{K\cap{^gK}}(\;{^g\pi},\pi)\cong\Hcal(G, K; \pi^\vee)_g\cong\Hcal(G, K; \pi)_{g^{-1}}
\]
as vector spaces.
\end{Lem}
\begin{proof}
See \cite[(4.1.1) Proposition, p.144]{Bushnell-Kutzko-Book}.
\end{proof}


\begin{Lem}\label{L:Bushnell_Kutzko_induced}
Let $K\subseteq G$ be as above and let $H$ be an open compact group such that $H\subseteq K\subseteq G$. Let $\rho$ be a finite dimensional representation of $H$. We have a natural isomorphism
\[
\Hcal(G, K; \Ind_H^K\rho)\cong \Hcal(G, H; \rho)
\]
such that for each $g\in G$
\[
\Hcal(G, K; \Ind_H^K\rho)_g\cong\bigoplus_{\substack{h\in H\backslash G \slash H\\KhK=KgK}} \Hcal(G, H; \rho)_h,
\]
where the sum is over all the double cosets $HhH$ such that $KhK=KgK$.
\end{Lem}
\begin{proof}
See \cite[(4.1.5) Corollary, p.146]{Bushnell-Kutzko-Book}.
\end{proof}

The next lemma will be useful for relating some of the Hom-spaces introduced above.

\begin{Lem}\label{HomIsoRel}
	Let $G$  and $H$ be groups and let $\phi:G\rightarrow H$ be a group isomorphism. Let $\sigma,\tau$ be $G$-modules. Then the Hom-spaces $\Hom_{G}(\sigma,\tau)$ and $\Hom_{H}(\,^{\phi}\sigma,\,^{\phi}\tau)$ are equal as subspaces of $\Hom_{\mathbb{C}}(\sigma,\tau)$.
\end{Lem}

Next we describe some results on twisting Hecke algebras by automorphisms.

\begin{Lem}\label{L:HeckeAutTwist}
	Let $\phi\in\Aut(G)$. Then
	\begin{enumerate}
		\item[(1)] The map $\Hcal(G,K,\pi)\rightarrow \Hcal(G,\phi(K),\,^{\phi}\pi)$ defined by $f\mapsto f\circ\phi^{-1}$ is an isomorphism of $\mathbb{C}$-algebras.\label{HeckeAutTwist}
    \end{enumerate}
	Suppose further that $\phi(K)=K$ and $\,^{\phi}\pi\cong \pi$, and let $A_{\phi}\in \GL(\pi^\vee)$ be such that $A_{\phi}\pi^\vee(k)=\pi^\vee(\phi^{-1}(k))A_{\phi}$. Then
    \begin{enumerate}
		\item[(2)]  $\Hcal(G,K,\,^{\phi}\pi)=A_{\phi}\Hcal(G,K,\pi)A_{\phi}^{-1}$ (we emphasize that this is an equality of spaces of functions and not just an isomorphism of algebras);\label{HeckeEqual}
		\item[(3)] $\Hcal(G,K,\,\pi)_{\phi(g)}\neq 0$ if and only if $\Hcal(G,K,\pi)_{g}\neq 0$. \label{HeckeSuppRel}
	\end{enumerate}	
\end{Lem}



%%%%%%%%%%%%%%%%%%%%%%%%%%%%%%%%%%%%%%%%%%%%%%%%%%%%%%%%%%%%%%%%%%%%%%%

\subsection{Finite Hecke algebra Shimura Correspondence}\label{FinShimHecke}

%%%%%%%%%%%%%%%%%%%%%%%%%%%%%%%%%%%%%%%%%%%%%%%%%%%%%%%%%%%%%%%%%%%%%%%

In this subsection we prove that the finite Shimura correspondence of Theorem \ref{FiniteShimCor} preserves the support of Hecke algebras. Let us set up the notation first. Let $H\subseteq G$ be such that $\Gamma(1)\subseteq H\subseteq K$, so that $\Gammaft(1)\subseteq \Hft \subseteq \Kft$. Denote by $H_\kappa$ the image of $H$ in $G_{\kappa}=K\slash\Gamma(1)$, so that we have
    \[
    \begin{tikzcd}
        \Kt\slash \s(\Gamma(2e))\ar[r,phantom, "{=}"]&\Kft\ar[r, two heads]&G_\kappa \ar[r, phantom, "{=}"]&K\slash\Gamma(1)\\
        \Ht\slash \s(\Gamma(2e))\ar[r,phantom, "{=}"]&\Hft\ar[r, two heads]\ar[u,phantom, "{\subseteq}" sloped, description] &H_\kappa \ar[r, phantom, "{=}"]\ar[u,phantom, "{\subseteq}" sloped, description]&H\slash\Gamma(1) \rlap{\,.}
    \end{tikzcd}
    \]
Fix a representation $\pi$ of $H_\kappa$, which we also view as a representation of $\Hft$ by inflating via the above surjection $\Hft\twoheadrightarrow H_{\kappa}$.

For each $x\in G_{\kappa}$, we let $\xt\in\Kft$ be a preimage under the surjection $\Kft\twoheadrightarrow G_{\kappa}$. The choice of $\xt$ is not unique, but if $\xt'$ is another choice, then
\begin{equation}\label{E:double_coset_equal_finite_Hecke}
    \Hft\,\xt\,\Hft=\Hft\,\xt'\,\Hft.
\end{equation}
because $\xt'=\xt \gamma$ for some $\gamma\in\Gammaft(1)$.

We then have
\begin{Thm}\label{T:finite_Hecke_algebra_isom}
    With the above notation, we have a $\C$-algebra isomorphism 
    \[
    \Hcal(G_{\kappa}, H_{\kappa}; \pi)\cong \Hcal(\Kft, \Hft; \pi\otimes\sigma)
    \]
    which restricts to an isomorphism
    \[
    \Hcal(G_{\kappa}, H_{\kappa}; \pi)_x\cong \Hcal(\Kft, \Hft; \pi\otimes\sigma)_{\xt}
    \]
    of $\C$-vector spaces for all $x\in G_{\kappa}$; namely the isomorphism is a support-preserving. (Note that the space $\Hcal(\Kft, \Hft; \pi\otimes\sigma)_{\xt}$ is independent of the choice of $\xt$ thanks to \eqref{E:double_coset_equal_finite_Hecke}.)
\end{Thm}
\begin{proof}
We already know the isomorphisms
\begin{equation}\label{E:composition_of_three_isomorphism_fintie_Hecke}
\Hcal(G_{\kappa},H_{\kappa};\pi)\xrightarrow{\;\sim\;} \End_{G_{\kappa}}(\Ind_{H_{\kappa}}^{G_{\kappa}}\pi)
\xrightarrow{\;\sim\;}\End_{\Kft}(\Ind_{\Hft}^{\Kft}(\pi\otimes \sigma))\xrightarrow{\;\sim\;}\Hcal(\Kft,\Hft;\pi\otimes \sigma),
\end{equation}
where the first and the third isomorphisms are as in Lemma \ref{HeckeAsEnd}, and the second isomorphism is by Theorem \ref{FiniteShimCor} and Proposition \ref{SCInd} \eqref{EndIso}. Hence we have the isomorphism of the two Hecke algebras.

By Lemma \ref{L:Bushnell_Kutzko_equivalence}, we know that
\[
\Hcal(G_{\kappa}, H_{\kappa}; \pi)_{x^{-1}}\cong\Hom_{H_\kappa\cap{\,^{x}H_\kappa}}(\,^{x}\pi, \pi)\qand
\Hcal(\Kft, \Hft; \pi\otimes\sigma)_{\xt^{-1}}\cong\Hom_{\Hft\cap{\,^{\xt}\Hft}}(\,^{\xt}\pi\otimes\sigma, \pi\otimes\sigma).
\]
Since $\Hom_{H_\kappa\cap{\,^{x}H_\kappa}}(\,^{x}\pi, \pi)\cong \Hom_{\Hft_\kappa\cap{\,^{x}\Hft_\kappa}}(\,^{x}\pi, \pi)$, Lemma \ref{HOMotimesSig} (with $M=\Gammaft(1)$ and $L=\Hft_\kappa\cap{\,^{x}\Hft_\kappa}$) gives
\[
\Hom_{H_\kappa\cap{\,^{x}H_\kappa}}(\,^{x}\pi, \pi)\cong \Hom_{\Hft\cap{\,^{\xt}\Hft}}(\,^{\xt}\pi\otimes\sigma, \pi\otimes\sigma).
\]
Here, it should be mentioned that $^{\xt}\Gammaft(1)=\Gammaft(1)$, so $\Gammaft(1)\subseteq \Hft_\kappa\cap{\,^{x}\Hft_\kappa}$. Hence we have
\[
\Hcal(G_{\kappa}, H_{\kappa}; \pi)_{x^{-1}}\cong \Hcal(\Kft, \Hft; \pi\otimes\sigma)_{\xt^{-1}}.
\]

It remains to show that this isomorphism is induced by the composite of the three isomorphisms of \eqref{E:composition_of_three_isomorphism_fintie_Hecke}.  But this is a matter of book keeping since these three isomorphisms are explicit, and is left to the reader.
\end{proof}


\begin{comment}

By Lemma \ref{HeckeAsEnd} we have isomorphisms
\begin{equation*}
	A:\Hcal(G_{\kappa},H,\pi))\rightarrow \End_{G_{\kappa}}(\Ind_{H}^{G_{\kappa}}\pi),\quad T\mapsto (f\mapsto T*f)
\end{equation*} 
and
\begin{equation*}
	C^{-1}:\Hcal(\Kt,\Ht,\pi\otimes \sigma)\rightarrow \End_{\Kt}(\Ind_{\Ht}^{\Kt}(\pi\otimes \sigma)),
    \quad T\mapsto (f\mapsto T*f).
\end{equation*} 
From Theorem \ref{FiniteShimCor} and Proposition \ref{SCInd} \eqref{EndIso} we have an isomorphism
\begin{equation*}
	B:\End_{G_{\kappa}}(\Ind_{H}^{G_{\kappa}}\pi)\rightarrow\End_{\Kt}(\Ind_{\Ht}^{\Kt}(\pi\otimes \sigma)),
\end{equation*}
which is best described in two steps. First, by Theorem \ref{FiniteShimCor} there is an isomorphism $B^{\prime}:\End_{G_{\kappa}}(\Ind_{H}^{G_{\kappa}}\pi)\rightarrow \End_{\Kt}((\Ind_{H}^{G_{\kappa}}\pi)\otimes \sigma)$ defined by $f\mapsto \mathcal{F}_{\sigma}(f)=f\otimes id_{\sigma}$. Second, recall the isomorphism $b^{\prime\prime}:(\Ind_{H}^{G_{\kappa}}\pi)\otimes \sigma\rightarrow \Ind_{\Ht}^{\Kt}(\pi\otimes \sigma)$ from Proposition \ref{SCInd} \eqref{IndSigIso} which induces an isomorphism $B^{\prime\prime}:\End_{\Kt}((\Ind_{H}^{G_{\kappa}}\pi)\otimes \sigma)\rightarrow\End_{\Kt}(\Ind_{\Ht}^{\Kt}(\pi\otimes \sigma))$ defined by $\phi\mapsto b^{\prime\prime}\circ \phi\circ (b^{\prime\prime})^{-1}$. Now $B=B^{\prime\prime}\circ B^{\prime}$.

Let $f\otimes v\in (\Ind_{H}^{G_{\kappa}}\pi)\otimes \sigma$. Define $F_{f,v}:=b^{\prime\prime}f\in \Ind_{\Ht}^{\Kt}(\pi\otimes \sigma)$. The vectors of this form span $\Ind_{\Ht}^{\Kt}(\pi\otimes \sigma)$. Thus the elements of $\End_{\Kt}(\Ind_{\Ht}^{\Kt}(\pi\otimes \sigma))$ are determined by their action on these vectors.

Now we can describe the map $B$. Let $f\in \Ind_{H}^{G_{\kappa}}\pi$, $v\in\sigma$, and $\phi\in \End_{G_{\kappa}}(\Ind_{H}^{G_{\kappa}}\pi)$. Then
\begin{equation*}
	B(\phi)(F_{f,v})=b^{\prime\prime}\circ \phi\otimes id_{\sigma}\circ (b^{\prime\prime})^{-1}(b^{\prime\prime}(f\otimes v))=F_{\phi(f),v}.
\end{equation*}

\begin{Thm}\label{FiniteHeckeSuppThm}
	Let $\Ht\subseteq \Kt$ be a subgroup such that $\widetilde{\Gamma}(1)\subseteq \Ht$. Let $\pi$ be a representation of $\Ht$ that is trivial on $\widetilde{\Gamma}(1)$, so that $\pi$ is really a representation of $H:=\Ht/\widetilde{\Gamma}(1)\hookrightarrow G_{\kappa}$. Let $\kt\in\Kt$ and let $k$ be its image in $G_{\kappa}$. 
	
	The composition of $\mathbb{C}$-algebra isomorphisms 
	\begin{equation}\label{FiniteHeckeIso}
		\Hcal(G_{\kappa},H,\pi)\stackrel{A}{\rightarrow}\End_{G_{\kappa}}(\Ind_{H}^{G_{\kappa}}\pi)\stackrel{B}{\rightarrow}\End_{\Kt}(\Ind_{\Ht}^{\Kt}(\pi\otimes \sigma))\stackrel{C}{\rightarrow}\Hcal(\Kt,\Ht,\pi\otimes \sigma)
	\end{equation} 
	restricts to give an isomorphism of vector spaces
	\begin{equation*}
		\Hcal(G_{\kappa},H,\pi)_{k}\cong \Hcal(\Kt,\Ht,\pi\otimes \sigma)_{\kt}.
	\end{equation*}
\end{Thm}

\begin{proof}
	We will prove this as follows. By Lemmas \ref{L:Bushnell_Kutzko_equivalence} and \ref{HOMotimesSig} we know that $\Hcal(G_{\kappa},H,\pi)_{k}\cong \Hcal(\Kt,\Ht,\pi\otimes \sigma)_{\kt}$ as vector spaces so the content of the theorem is that this isomorphism is induced by those of \eqref{FiniteHeckeIso}. To verify this we will take $\Phi_{k}\in \Hcal(G_{\kappa},H,\pi)_{k}$ and a corresponding (by Lemmas \ref{L:Bushnell_Kutzko_equivalence} and \ref{HOMotimesSig}) $\widetilde{\Phi}_{\kt}\in \Hcal(\Kt,\Ht,\pi\otimes \sigma)_{\kt}$, and check that $B\circ A(\Phi_{k})=C^{-1}\widetilde{\Phi}_{\kt}$ as elements of $\End_{\Kt}(\Ind_{\Ht}^{\Kt}(\pi\otimes \sigma))$.

	By Lemma \ref{L:Bushnell_Kutzko_equivalence} we have a natural isomorphism of vector spaces $\Hom_{H\cap \,^{k}H}(\pi,\,^{k}\pi)\cong\Hcal(G_{\kappa},H,\pi)_{k}$ defined by $\phi_{k}\mapsto \Phi_{k}$, where $\supp\Phi_{k}=HkH$ and $\Phi_{k}(h_{1}kh_{2})=\pi(h_{1})\phi_{k}\pi(h_{2})$. Similarly we have a natural isomorphism of vector spaces $\Hom_{\Ht\cap \,^{\kt}\Ht}(\pi\otimes\sigma,\,^{\kt}(\pi\otimes \sigma))\cong\Hcal(\Kt,\Ht,\pi\otimes \sigma)_{\kt}$ defined by $\widetilde{\phi}_{\kt}\mapsto \widetilde{\Phi}_{\kt}$, where $\supp\widetilde{\Phi}_{\kt}=\Ht\kt\Ht$ and $\widetilde{\Phi}_{\kt}(\htt_{1}\kt\htt_{2})=(\pi\otimes \sigma)(\htt_{1})\widetilde{\phi}_{k}(\pi\otimes)(\htt_{2})$.
	
	Now given $f\in \Ind_{H}^{G_{\kappa}}\pi$ the action of $\Phi_{k}$ on $f$ via $A$ is
	\begin{equation*}
		[\Phi_{k}*f](g)=\sum_{h\in H/H\cap\,^{k}H} \pi(h)\phi_{k}f(k^{-1}h^{-1}g).
	\end{equation*}
	Thus for $f\in \Ind_{H}^{G_{\kappa}}\pi$, $v\in \sigma$ we have
	\begin{equation*}
		[B\circ A(\Phi_{k})](F_{f,v})=F_{\Phi_{k}*f,v}.
	\end{equation*}
	
	Let $\widetilde{\phi}_{\kt}=\phi\otimes \sigma(\kt)\in \Hom_{\Ht\cap \,^{\kt}\Ht}(\pi\otimes\sigma,\,^{\kt}(\pi\otimes \sigma))$ be the element corresponding to $\phi_{k}$ under the isomorphism of Lemma \ref{HOMotimesSig}. We must compute $[C^{-1}(\widetilde{\Phi}_{\kt})](F_{f,v})=\widetilde{\Phi}_{\kt}*F_{f,v}$.
	
	\begin{align*}
		(\widetilde{\Phi}_{\kt}*F_{f,v})(\gt)=&\sum_{\htt\in \Ht/\Ht\cap\,^{\kt}\Ht}\widetilde{\Phi}_{\kt}(\htt\kt)F_{f,v}(\kt^{-1}\htt^{-1}\gt)\\
		=&\sum_{\htt\in \Ht/\Ht\cap\,^{\kt}\Ht}[\pi(\htt)\otimes\sigma(\htt)]\circ[\phi_{k}\otimes \sigma(\kt)](f(\kt^{-1}\htt^{-1}\gt)\otimes \kt^{-1}\htt^{-1}\gt v)\\
		=&\sum_{\htt\in \Ht/\Ht\cap\,^{\kt}\Ht}[\pi(\htt)\circ\phi_{k}f(\kt^{-1}\htt^{-1}\gt)]\otimes \gt v\\
		=&F_{\Phi_{k}f,v}.
	\end{align*}
	Thus it follows that $[B\circ A(\Phi_{k})]=C^{-1}\widetilde{\Phi}_{\kt}$.
\end{proof}

\end{comment}

This support preserving property has significant consequences when we specialize to the case of parabolic induction. Let
\[
P_\kappa=M_\kappa N_\kappa\subseteq G_\kappa
\]
be a parabolic subgroup of $G_\kappa$, where $M_\kappa$ and $N_\kappa$ are the Levi part and the unipotent subgroup, respectively. Then there exists a subgroup $P\subseteq K$ such that $P\supseteq\Gamma (1)$ and the corresponding $\Pft\subseteq\Kft$ is the preimage of $P_\kappa$ under the surjection $\Kft\twoheadrightarrow G_\kappa$.  We then have

\begin{Cor}\label{ParaHeckeIso}
Let $\pi$ be a representation of $M_\kappa$ which we inflate to a representation of $P_\kappa$. We then have a support-preserving $\C$-algebra algebra isomorphism 
\[
\Hcal(G_{\kappa},P_{\kappa};\pi)\cong\Hcal(\Kft,\Pft;\pi\otimes \sigma).
\]
In particular, if $B_\kappa$ is the Borel subgroup, we have
\[
\Hcal(G_{\kappa},B_{\kappa};\pi)\cong\Hcal(\Kft,\Ift;\pi\otimes \sigma).
\]
\end{Cor}
\begin{proof}
    Just let $H=P$ in the above theorem.
\end{proof}

Hecke algebras of the form $\Hcal(G_{\kappa},P_{\kappa};\pi)$, where $\pi$ is cuspidal, were thoroughly studied in \cite{HL80}. We highlight a couple of their results. (Here we will be translating the statements of \cite{HL80}, which primarily use the language of endomorphism algebras, into the language of Hecke algebras.) First, \cite[Proposition (3.9)]{HL80} (for which they refer to Springer \cite[pp. 635-636]{S71}) states that the Hecke algebra $\Hcal(G_{\kappa},P_{\kappa};\pi)$ has a basis indexed by a subset of $P_{\kappa}$-double cosets such that the support of a basis function $T_{w}$ is the corresponding $P_{\kappa}$-double coset $P_{\kappa}\,w\,P_{\kappa}$.  Second, in terms of this basis, \cite[Theorem (4.14)]{HL80} describes a presentation for $\Hcal(G_{\kappa},P_{\kappa}; \pi)$ (roughly braid relations and quadratic relations). The isomorphism of Corollary \ref{ParaHeckeIso} then implies that $\Hcal(\Kft,\Pft; \pi\otimes \sigma)$ inherits the presentation of $\Hcal(G_{\kappa},P_{\kappa}; \pi)$, but this isomorphism also preserves support. Thus if $\Tt_{\wt}\in \Hcal(\Kft, \Pft; \pi\otimes \sigma)$ corresponds to $T_{w}\in \Hcal(G_{\kappa}, P_{\kappa}; \pi)$, then we know that $\supp(\Tt_{\wt})=\Pft\,\wt\, \Pft$ and under the surjection $\Kft\twoheadrightarrow G_{\kappa}$ we have $\wt\mapsto w$.


\begin{comment}
Now, let us specialize to the situation where $P_\kappa=B_\kappa$, the Borel subgroup, and $\pi=\one$, the trivial representation; namely consider the Hecke algebra $\Hcal(G_{\kappa}, B_{\kappa}; \one)$, which is just equal to the space of $\mathbb{C}$-valued functions on $B_{\kappa}\backslash G_{\kappa}/B_{\kappa}$. By the Bruhat decomposition the $B_{\kappa}$-double cosets are in bijection with $W$. Thus for each $w\in W$ define $T_{w}\in \Hcal(G_{\kappa},B_{\kappa};\one)$ to be the unique function supported on $B_{\kappa}\,w\, B_{\kappa}$ such that $T_{w}(\dot{w})=1$, where $\dot{w}\in N_{G_{\kappa}}(T_{\kappa})$ represents $w$. Note that this definition does not depend on the choice of $\dot{w}$. The set $\{T_{w}\st w\in W\}$ is a basis for $\Hcal(G_{\kappa},B_{\kappa};\one)$.

As a $\mathbb{C}$-algebra, $\Hcal(G_{\kappa},B_{\kappa},1)$ is generated by the elements $T_{\alpha}:=T_{w_{\alpha}}$, where $\alpha\in\Delta$ subject to the following relations. Let $\ell_{\Delta}$ be the length function of $W$ relative to the simple roots $\Delta$. For $\alpha\in \Delta$
\begin{equation}\label{WQuadRel}
	T_{\alpha}^{2}=q+(q-1)T_{\alpha}.
\end{equation}
For $w_{1},w_{2}\in W$ such that $\ell_{\Delta}(w_{1}w_{2})=\ell_{\Delta}(w_{1})+\ell_{\Delta}(w_{2})$
\begin{equation}\label{WBraidRel}
	T_{w_{1}}T_{w_{2}}=T_{w_{1}w_{2}}.
\end{equation}

Now if we apply Corollary \ref{ParaHeckeIso} we get a basis along with algebra generators and relations for the Hecke algebra $\Hcal(\Kft,\Bft;\sigma)$. Let us also write $T_{w}$ for the element of $\Hcal(\Kft,\Bft; \sigma)$ corresponding to $T_{w}\in \Hcal(G_{\kappa},B_{\kappa}; \one)$ under the isomorphism of Corollary \ref{ParaHeckeIso}. We note that by the support preserving property we have $\supp(T_{w})=\Bft\,w\,\Bft$. Note that the above relations on $T_w$'s imply that $T_{w}$ is invertible for all $w\in W$.

%, and by Lemma \ref{} we know that $T_{w}(\wt)=\sigma(\wt)$, where $\wt\in\Wcal$ represents $w$.

\end{comment}

%%%%%%%%%%%%%%%%%%%%%%%%%%%%%%%%%%%%%%%%%%%%%%%%%%%%%%%%%%%%%%%%%%%%%%%

\subsection{Two isomorphic Hecke algebras $\Hcal$ and $\Hcal_2$}\label{HandH2}

%%%%%%%%%%%%%%%%%%%%%%%%%%%%%%%%%%%%%%%%%%%%%%%%%%%%%%%%%%%%%%%%%%%%%%%

\begin{comment}
In \S \ref{SSecSplittings}, Proposition \ref{Gamma12eSplitting} we proved that the map $\s_{\Gamma_{1}(2e)}$ is a splitting of the congruence subgroup $\Gamma_{1}(2e)$. Furthermore, the image of this splitting is normal in $\widetilde{\Gamma}_{0}(2e)$, as can be seen from the Steinberg relations. By the Iwahori factorization and a direct calculation it follows that the inclusion $\Tt_{0}\hookrightarrow \widetilde{\Gamma}_{0}(2e)$ induces an isomorphism $\Tt_{0}/T_{2e}\cong \widetilde{\Gamma}_{0}(2e)/\s_{\Gamma_{1}(2e)}(\Gamma_{1}(2e))$.

Since $\tau$ is a representation of $\Tt_{0}/T_{R}$ and $T_{2e}\subseteq T_{R}$, the representation $\tau$ inflates to a representation of $\Tt_{0}/T_{2e}\cong \widetilde{{\Tf}}_{0}$ and thus also to $\widetilde{\Gamma}_{0}(2e)$. We abuse notation and refer to these representations as $\tau$.
\end{comment}


Now, let us introduce the Iwahori Hecke algebra we will study. Recall that we have first constructed a pseudo-spherical representation $\tau$ of the compact group $\Tt_0$. The kernel of $\tau$ contains $T_R$, which allows us to view $\tau$ as a representation of the finite group $\Tft=\Tt_0\slash T_R$, which we also denote by $\tau$. By parahoric induction, we have constructed a representation  $\sigma=i_{{\Uf}_{0},{\Uf}_{1}^{-}}^{\Ift}\tau$ of the finite group $\Ift=\It\slash \s(\Gamma(2e))$. By inflating $\sigma$ by the surjection $\It\twoheadrightarrow \widetilde{{\If}}$, we view $\sigma$ as a representation of $\It$, which we also call $\sigma$. We then define 
\begin{equation*}
	\mathcal{H}:=\mathcal{H}(\Gt,\widetilde{I};\sigma).
\end{equation*}

For our analysis, however, it will be convenient to also consider another Hecke algebra constructed using the congruence subgroup 
\[
I_{2}:=\Gamma_{0}(2).
\]
Note the inclusions $I_2\subseteq I\subseteq K$. By using the Iwahori decomposition $(U_0, T_0, U_2^-)$ of $I_2$, we consider another parahoric induction
\[
\sigma_{2}:=i_{{\Uf}_{0},{\Uf}_{2}^{-}}^{\Ift_2}\tau,
\]
which we often view as a representation of $\It_2$ by inflating via $\It_2\twoheadrightarrow \Ift_2$. We then define 
\begin{equation*}
	\Hcal_{2}:=\mathcal{H}(\Gt,\It_2;\sigma_{2}).
\end{equation*}



These two Hecke algebras are actually isomorphic:
\begin{Prop}\label{HIsoH2}
	As $\mathbb{C}$-algebras 
	\begin{equation*}
		\Hcal\cong \Hcal_{2}.
	\end{equation*}
	%Furthermore, for all $\alpha\in\Delta$ we have $\Hcal_{w_{\alpha}}\cong (\Hcal_{2})_{w_{\alpha}}\cong \mathbb{C}$.
\end{Prop}

\begin{proof}
	The claim follows directly from Proposition \ref{IndSig2} and Lemma \ref{L:Bushnell_Kutzko_induced}. %The second claim follows from Proposition \ref{CompatSupp}.
\end{proof}

\begin{Rmk}
We work with $\Hcal_{2}$ primarily because the group $\It_2$ behaves better with the affine reflection $w_{\af}$ than $\It$. For example, the group generated by $\It_2$ and $w_{\af}$ is compact, whereas the one generated by $\It$ and $w_{\aff}$ is not. Accordingly, when we work with $w_{\af}$, we often work with $\Hcal_2$.
\end{Rmk}

%%%%%%%%%%%%%%%%%%%%%%%%%%%%%%%%%%%%%%%%%%%%%%%%%%%%%%%%%%%%%%%%%%%%%%%

\subsection{Support of $\Hcal$: multiplicity one}\label{SS:MultOne}

%%%%%%%%%%%%%%%%%%%%%%%%%%%%%%%%%%%%%%%%%%%%%%%%%%%%%%%%%%%%%%%%%%%%%%%

In this subsection we begin our study of the support of $\Hcal$. Let us start with

\begin{Prop}\label{HMult1}
	For each $g\in G$, we have
	\begin{equation*}
		\dim\Hcal_{g}\leq 1.
	\end{equation*}
\end{Prop}

\begin{proof}
	First note that by Lemma \ref{L:Bushnell_Kutzko_equivalence} it suffices to compute the dimension of the Hom-space $\Hom_{\It\cap\,^{g}\It}(\,^{g}\sigma,\sigma)$. Second, note that if $\It g\It=\It g'\It$ then $\dim\Hcal_{g^{-1}}=\dim\Hcal_{{g'}^{-1}}$, and hence by the Iwahori-Bruhat decomposition, we may assume $g\in N_{\Gt}(\Tt)$.
	
	Let $\Gamma:=\Gamma(1)$. Since $ \Gamma\subseteq I$ we have $\Hom_{\It\cap\,^{g}\It}(\,^{g}\sigma,\sigma)\subseteq \Hom_{\widetilde{\Gamma}\cap\,^{g}\widetilde{\Gamma}}(\,^{g}\sigma,\sigma)$. Thus it suffices to prove 
    \[
    \dim\Hom_{\widetilde{\Gamma}\cap\,^{g}\widetilde{\Gamma}}(\,^{g}\sigma,\sigma)\leq 1,
    \]
    where we recall from Proposition \ref{IRep} \eqref{ResSig} that $\sigma|_{\Gammat}$ is still irreducible. (The advantage of $\Gamma$ over $I$ is that $\Gamma$ is normalized by $\Wcal$.)
	
	As $g\in N_{\Gt}(\Tt)$ there exists $w\in \Wcal$ and $t\in \Tt$ such that $g=tw$. Then by Proposition \ref{IRep} we know $^w\sigma\cong\sigma$ as $\Gammat$-modules, which implies that $\Hom_{\widetilde{\Gamma}\cap\,^{g}\widetilde{\Gamma}}(\,^{g}\sigma,\sigma)\cong \Hom_{\widetilde{\Gamma}\cap\,^{t}\widetilde{\Gamma}}(\,^{t}\sigma,\sigma)$. Hence we may assume $g=t\in\Tt$.
	
	Next, note that there exists $w_{0}\in\Wcal$ such that 
    \[
    (w_0tw_0)U_j(w_0tw_0)^{-1}\subseteq U_j
    \]
    for any $j\in\Z$, where we recall the notation $U_j=U_j^+$ from \eqref{E:notation_U_alpha_j}. By Lemma \ref{HomIsoRel} and Proposition \ref{IRep} \eqref{sigWInvar} and \eqref{ResSig}, we have
	\begin{equation*}
		\Hom_{\widetilde{\Gamma}\cap\,^{t}\widetilde{\Gamma}}(\,^{t}\sigma,\sigma)\cong \Hom_{\widetilde{\Gamma}\cap\,^{w_{0}tw_{0}^{-1}}\widetilde{\Gamma}}(\,^{w_{0}tw_{0}^{-1}}\sigma,\sigma).
	\end{equation*}
	Thus we may assume $g=t\in \Tt$ such that $tU_{j}t^{-1}\subseteq U_{j}$ for any $j\in \mathbb{Z}$. 
	
	In this case, we see that $\widetilde{\Gamma}\cap\,^{t}\widetilde{\Gamma}$ has an Iwahori decomposition of the form
	\begin{equation*}
		\widetilde{\Gamma}\cap\,^{t}\widetilde{\Gamma}\cong U^{-}_{1}\times \Tt_{1} \times \,^{t}U_{1}.
	\end{equation*}
	Specifically, we have $U^{-}_{1}\subseteq \widetilde{\Gamma}\cap\,^{t}\widetilde{\Gamma}$, and $\,^{t}U_{1}\subseteq U_{1}$.
	
	Let $A\subseteq \,^{t}\sigma$ be the span of the elements of the form $e_{U_{1}}e_{U^{-}_{1}}\otimes_{\Tt_{1}} w$, where $w\in\tau$. Note that $A$ is a $\Tt_{1}$-module such that $A\subseteq \,^{t}\sigma^{\,^{t}U_{1}}=\sigma^{U_{1}}$ because the idempotent $e_{U_1}$ absorbs the elements in $U_1$. Also note that, as a $\widetilde{\Gamma}\cap\,^{t}\widetilde{\Gamma}$-module, $\,^{t}\sigma$ is generated by $A$, because of the above Iwahori decomposition of $\widetilde{\Gamma}\cap\,^{t}\widetilde{\Gamma}$ and the containment $\widetilde{\Gamma}\cap\,^{t}\widetilde{\Gamma}\supset U^{-}_{1}$. Thus by restriction from $\,^{t}\sigma$ to $A$ we get an injective map
	\begin{equation*}
		\Hom_{\widetilde{\Gamma}\cap\,^{t}\widetilde{\Gamma}}(\,^{t}\sigma,\sigma)\hookrightarrow\Hom_{\Tt_{1}}(A,\sigma^{U_{1}}).
	\end{equation*}
        Now, as $\Tt_{1}$-modules, $\sigma^{U_{1}}\cong \tau$ by Proposition \ref{P:two_functors_for_Iwahori_decomp} (4). Therefore $\dim\Hom_{\Tt_{1}}(A,\sigma^{U_{1}})\leq 1$, since $\dim A\leq \dim\tau$, proving the proposition.
\end{proof}

The following gives a sufficient condition for $\mathcal{H}_{g}$ to be $0$.

\begin{Prop}\label{HVanish}
	Let $f\in \Hcal$, $w\in \Wcal$ and $y\in Y\smallsetminus\Yt$. Let $\yt(\varpi)\in \Tt$ be any preimage of $y(\varpi)\in T$. Then $f(\yt(\varpi) w)=0$. In particular, 
	\begin{equation*}
		\Hcal_{\yt(\varpi) w}=0.
	\end{equation*}
\end{Prop}
\begin{proof}
	Since $y\notin\Yt$, there exists a simple root $\alpha\in\Delta$ such that $\la \alpha, y\ra\notin 2\Z$. Now, let $t\in (1+\varpi^{2e}\Ocal)\smallsetminus \Ocal^{\times 2}$. Then since $1+\varpi^{2e}\Ocal\subseteq R$, we know $(t, \varpi)=-1$ by Corollary \ref{C:Hilbert_symbol_unique_element_integral_radical}. Thus
	\[
	\htt_{\alpha}(t) \yt(\varpi)=(t, \varpi^{\la\alpha, y\ra})\yt(\varpi)\htt_{\alpha}(t)=-\yt(\varpi)\htt_{\alpha}(t).
	\]
	Recalling that $\Kt$ normalizes $\Gamma(2e)$, we have
	\[
	w^{-1}\htt_{\alpha}(t)w\in T_{2e}.
	\]
	Since $T_{2e}$ acts trivially on $\sigma$, we have
	\[
	f(\yt(\varpi) w)=f(\htt_{\alpha}(t)\yt(\varpi) w)=-f(\yt(\varpi)ww^{-1}\htt_{\alpha}(t)w)=-f(\yt(\varpi)w).
	\]
	Hence $f(\yt(\varpi) w)=0$.
\end{proof}


\begin{comment}
Our next task will be to construct a basis for $\Hcal$ and prove that these elements satisfy certain relations. Unfortunately we must address these questions simultaneously.

By Proposition \ref{HMult1} we see that a basis for $\Hcal$ is indexed by a subset of $I\backslash G/I$. By the Iwahori-Bruhat decomposition $I\backslash G/I$ can be identified with $W_{\ea}:=N_{G}(T)/T_{0}$ the extended affine Weyl group. (Notation for affine Weyl groups can be found in \S \eqref{RootSys}.)

The group $W_{\ea}$ acts faithfully on $Y\otimes \mathbb{R}$ by affine transformations and it contains the affine Weyl group $W_{\af}$. We use the convention such that the elements $h_{\alpha}(\varpi^{j})$ act on $Y\otimes \mathbb{R}$ by $(y\otimes c\mapsto (y-j\alpha^{\vee})\otimes c)$. This identifies the subgroup of translations $T/T_{0}\subseteq W_{\ea}$ with $Y$. Since $G$ is simply-connected we have $W_{\ea}=W_{\af}$. 

The stabilizer of $0\in Y\otimes \mathbb{R}$ in $W_{\ea}$ is isomorphic to the Weyl group $W$ and this subgroup normalizes the subgroup of translations $Y$, and gives a decomposition $W_{\ea}=W_{\af}\cong W\ltimes Y$.

Proposition \ref{HVanish} shows that a basis for $\Hcal$ is indexed by a subset of $\Wt_{\ea}:=W\ltimes \Yt\subseteq W_{\af}=W_{\ea}$. We will show later that $\Wt_{\ea}$ exactly indexes a basis for $\Hcal$ (Proposition \ref{HLineDecompAndInvert}). The group $\Wt_{\ea}$ contains as a subgroup of finite index the affine Weyl group $\Wt_{\af}$ associated with the affine root system attached to $\frac{1}{2}\Phi$. The group $\Wt_{\ea}$ also acts on the set of alcoves defined by this affine root system. Let $\Omega\subseteq \Wt_{\ea}$ be the subgroup stabilizing (set-wise not point-wise) the alcove $\widetilde{\mathfrak{A}}$. Then $\Wt_{\ea}$ has a decomposition $\Wt_{\ea}\cong \Omega\ltimes \Wt_{\af}$.

With this decomposition we can extend the length function $\ell$ on $\Wt_{\af}$ to $\Wt_{\ea}$ such that for $(\omega,w)\in \Wt_{\ea}$ we have $\ell((\omega,w)):=\ell(w)$. Thus the subset $\Omega$ can now also be identified as the subset of length $0$ elements in $\Wt_{\ea}$. We note that $\Omega\cong \Yt/2Y$.


\end{comment}

%%%%%%%%%%%%%%%%%%%%%%%%%%%%%%%%%%%%%%%%%%%%%%%%%%%%%%%%%%%%%%%%%%%%%%%

\subsection{Linear Hecke algebra for $I_2$}\label{SS:LinearHecke}

%%%%%%%%%%%%%%%%%%%%%%%%%%%%%%%%%%%%%%%%%%%%%%%%%%%%%%%%%%%%%%%%%%%%%%%

Here we take an interlude to study a linear Hecke algebra which will help us investigate several aspects of $\Hcal_{2}$. (A similar idea is used by Savin \cite{S04} in the tame case.) In particular, we record generalizations of the results of \cite[Section 5]{Karasiewicz} to the present setting. All the results generalize, but we only record those that we directly use later. The proofs carry over without change once the correct generalization is used, so we will be brief and refer the reader to \cite[Section 5]{Karasiewicz} for more details. The main point is that the mod $4$ congruence conditions in the $F=\mathbb{Q}_{2}$ setting of \cite{Karasiewicz} should become mod $\varpi^{2}$ congruence conditions in the present setting of an arbitrary $2$-adic field (as opposed to mod $\varpi^{2e}$). Note also that the congruence subgroup $I_2:=\Gamma_{0}(2)$ in the current paper corresponds to $\Gamma_{0}(4)$ in \cite{Karasiewicz}.

Let $\underline{\Hcal}:=C_{c}^{\infty}(I_2\backslash G/I_2)$, the algebra of $I_2$-biinvariant compactly supported functions on $G$. The space $\underline{\Hcal}$ has multiplication given by convolution where we normalize the Haar measure so that $I_2$ has measure $1$. 

We consider a length function $\ell_{2}:G\rightarrow\mathbb{R}_{\geq 0}$ defined by
\begin{equation*}
	q^{\ell_{2}(g)}:=[I_2gI_2:I_2].
\end{equation*}
By composing with any set theoretic section $W_{\af}\rightarrow N_{G}(T)$, we can consider the function
\[
\ell_{2}:W_{\af}\rightarrow N_G(T)\rightarrow\mathbb{R}_{\geq 0},
\]
which we also denote by $\ell_2$, because one can see that this composite does not depend on the choice of the section because $T_{0}\subseteq  I_2$.

Now we state relevant generalizations as follows, for which we recall the notation $\Omegat, \Wt_{\aff}, \Wt_{\ea}$ and $\tilde{\ell}$ from \S \ref{SS:affine_Weyl_covering_group}.

\begin{Lem}\label{Len0elts}
    Let $w_{1}\in \Omegat$ and $w_{2}\in \widetilde{W}_{\af}$. Then $\ell_{2}(w_{1}w_{2})=\ell_{2}(w_{2})$ and $\s(w_{1})\in N_{G}( I_2)$.
\end{Lem}

\begin{proof}
	This generalizes \cite[Lemma 5.4]{Karasiewicz} and can be proved in the same way.
\end{proof}

\begin{Prop}\label{CosetProp}
	Let $g_{1},g_{2},g,\delta\in G$ be such that $\ell_{2}(g_{1}g_{2})=\ell_{2}(g_{1})+\ell_{2}(g_{2})$. Suppose that the coset $\delta I_2$ satisfies the two conditions
	\begin{itemize}
		\item $\delta I_2\subseteq  I_2g_{1} I_2$;
		\item $(\delta I_2)^{-1}g\subseteq  I_2g_{2} I_2$.
	\end{itemize} 
	Then
	\begin{itemize}
		\item $ I_2g I_2=  I_2g_{1}g_{2} I_2$;
		\item $\delta I_2= g_{1} I_2$.
	\end{itemize} 
\end{Prop}

\begin{proof}
	This follows by adapting the proof of \cite[Proposition 5.2]{Karasiewicz} and the subsequent remark. We note that this proof uses $\underline{\Hcal}$.
\end{proof}

\begin{Thm}\label{ProportionalLengths}
	Let $w\in \widetilde{W}_{\ea}$. Then
	\begin{equation*}
		\ell_{2}(w)=2\tilde{\ell}(w).
	\end{equation*}
\end{Thm}

\begin{proof}
	This is proved as in \cite[Theorem 5.5]{Karasiewicz} by studying the structure of $\underline{\Hcal}$.
\end{proof}





%%%%%%%%%%%%%%%%%%%%%%%%%%%%%%%%%%%%%%%%%%%%%%%%%%%%%%%%%%%%%%%%%%%%%%%

\subsection{Support: affine simple reflections}\label{SuppAfSimpRel}

%%%%%%%%%%%%%%%%%%%%%%%%%%%%%%%%%%%%%%%%%%%%%%%%%%%%%%%%%%%%%%%%%%%%%%%

Recall that the group $\Wt_{\af}$ is generated by the simple affine reflections $w_{\af}, w_{\alpha_1}, \dots, w_{\alpha_n}$, where $w_{\af}$ is defined as in \eqref{E:w_aff_definition}. In this section, we will show that our Hecke algebra $\Hcal$ is supported on these elements, namely

\begin{Prop} \label{affSupp_H}
	We have
	\begin{equation*}
		\Hcal_{w_{\alpha}}\cong \Hom_{\It\cap{^{w_{\alpha}}\It}}(\,^{w_{\alpha}}\sigma,\;\sigma)\cong \mathbb{C}
	\end{equation*}
    for all $\alpha\in\Delta\cup\{\af\}$.
\end{Prop}

Since we already know $\dim\Hcal_{w_{\alpha}}\leq 1$ (Proposition \ref{HMult1}) and $\Hcal\cong\Hcal_2$ (Proposition \ref{HIsoH2}), and since $\It w_{\alpha}^{-1}\It=\It w_{\alpha}\It$ and so $\Hcal_{w_{\alpha}}=\Hcal_{w_{\alpha}^{-1}}$, we know from Lemma \ref{L:Bushnell_Kutzko_induced} that Proposition \ref{affSupp_H} above is implied by the following proposition.
\begin{Prop} \label{affSupp}
	We have
	\begin{equation*}
		(\Hcal_2)_{w_{\alpha}}\cong \Hom_{\It_2\cap{^{w_{\alpha}}\It_2}}(\,^{w_{\alpha}}\sigma_{2},\;\sigma_{2})\cong \mathbb{C}
	\end{equation*}
    for all $\alpha\in\Delta\cup\{\af\}$.
\end{Prop}


Before we begin the proof we will set up some preliminary results. We focus on the case of $\alpha=\af$. The case of $\alpha\in\Delta$ can be proved similarly. We omit these details because Proposition \ref{affSupp_H} also follows directly from the finite Shimura correspondence. Alternatively, since we know that the $\It$-module $\sigma$ extends to a $\Kt$-module, we may apply the method of \cite{Takeda_Wood} to prove Proposition \ref{affSupp_H} for $\alpha\in\Delta$.


It will be convenient to have a slightly different model for the representation $\sigma_{2}$. Let 
\begin{equation}\label{E:profinite_model_sigma_2}
\sigma_{2}^{p}:=i_{U_{0},U^{-}_{2}}^{\It_2}\tau
\end{equation}
be the $\It_2$-module constructed by parahoric induction, where $\tau$ is the representation $\Tt_{0}$. Here, the superscript $p$ is for ``profinite" as opposed to ``finite". (Recall that $\sigma_{2}$ was first defined as a representation of the {\it finite} group $\Ift_{2}$ via the finite group version of parahoric induction and then inflated to the {\it profinite} group $\It_2$. Here, $\sigma_2^p$ is defined by the profinite version of parahoric induction.) But this is another model of $\sigma_2$. Namely, we have

\begin{Lem}
As $\It_2$-modules 
\[
\sigma_{2}\cong\sigma_{2}^{p}.
\]
\end{Lem}
\begin{proof}
This is a consequence of the following identity in $\sigma_{2}^{p}$
    \begin{equation*}
        f(g)e_{U_{0}}e_{U^{-}_{2}}\otimes_{\Tt_{0}}w=f(g)e_{U_{0}}e_{\textbf{s}(\Gamma_{2e})}e_{U^{-}_{2}}\otimes_{\Tt_{0}}w=f(g)e_{\textbf{s}(\Gamma_{2e})}e_{U_{0}}e_{U^{-}_{2}}\otimes_{\Tt_{0}}w,
    \end{equation*}
where $w\in\tau$. The main input needed for this identity is Proposition \ref{P:SplitNormIm} \eqref{Gam2eSplitNorm}.
\end{proof}

Hence by abuse of notation, we often denote this profinite model simply by $\sigma_2$

\begin{Rmk}
The reason for switching to this profinite model is that conjugation by $w_{\af}$ does not stabilize $\Gamma(2e)$ and thus does not interact well with the quotients by $\s(\Gamma(2e))$ that we used previously.
\end{Rmk}

\begin{comment}
Let 
\[
I_{2}^{\af}:= I_2\cap \,^{w_{\af}} I_2\qand\sigma_{2}^{\af}:=i_{U^{-}_{2}\cap I_2^{\af},U_{0}}^{ I_2^{\af}}\tau.
\]

\begin{Lem}\label{L:GammaAfRep} $ $
    \begin{enumerate}
        \item $\,^{w_{\af}}\sigma_{2}^{\af}\cong \sigma_{2}^{\af}$.\label{WafTwistIso}
        
        Now suppose that $e>1$.
        \item The subspace $\sigma_{2}^{U_{-\alpha_{0},2e-4}}\subseteq \sigma_{2}$ is a $\It_2^{\af}$-module.\label{UInvarSub}
        \item $\Hom_{\It_2^{\af}}(\sigma_{2}^{\af},\sigma_{2})\cong \Hom_{\It_2^{\af}}(\sigma_{2}^{\af},\sigma_{2}^{U_{-\alpha_{0},2e-4}})\cong \mathbb{C}$. \label{GammaAfIsotypic}
    \end{enumerate}
\end{Lem}

\begin{proof}
    Item \eqref{WafTwistIso} follows from Proposition \eqref{TwistAut} because $\tau$ is $W$-invariant and $Z(\Tt)$-invariant.
    Item \eqref{UInvarSub} can be checked using the Steinberg relations.
    Item \eqref{GammaAfIsotypic} follows from Proposition \ref{P:two_functions_are_adjoint} and Item \eqref{WafTwistIso}, because $U_{\alpha_{0},2e}$ acts trivially on $\sigma_{2}^{\af}$.
\end{proof}

%\textcolor{red}{\textbf{At some point I thought $\sigma_{2}^{\af}\cong\sigma_{2}^{U_{-\alpha_{0},2e-4}}$, but I cannot seem to prove this at the moment.}

\end{comment}


\begin{proof}[Proof of Theorem \ref{affSupp}] 
As mentioned above, we only prove the proposition for $\alpha=\af$. Let
\[
I_{2}^{\af}:= I_2\cap \,^{w_{\af}} I_2\qand\sigma_{2}^{\af}:=i_{U^{-}_{2}\cap I_2^{\af},U_{0}}^{ I_2^{\af}}\tau.
\]
Since $\tau$ is $W$-invariant and $Z(\Tt)$-invariant, Proposition \ref{TwistAut} implies
\[
\,^{w_{\af}}\sigma_{2}^{\af}\cong \sigma_{2}^{\af}.
\]
We also know by Proposition \ref{P:two_functions_are_adjoint} and Lemma \ref{L:IndRes2x} that
\[
\Hom_{\It_2^{\af}}(\sigma_{2}^{\af},\sigma_{2})\cong \mathbb{C}.
\]

Hence, we can identify $\sigma_{2}^{\af}$ with the $\sigma_{2}^{\af}$-isotypic subspace of $\sigma_{2}$, and this identification is unique up to nonzero scaling. Let $A:\sigma_{2}^{\af}\hookrightarrow \sigma_{2}$ be the canonical $\It_2^{\af}$-inclusion of this isotypic subspace. Let $C:\sigma_{2}\twoheadrightarrow \sigma_{2}^{\af}$ be the canonical $\It_2^{\af}$-projection onto this isotypic subspace. Let $\,^{w_{\af}}C:\,^{w_{\af}}\sigma_{2}\twoheadrightarrow \,^{w_{\af}}\sigma_{2}^{\af}$ be the $\It_2^{\af}$-homomorphism obtained by twisting $C$ by $w_{\af}$. Note that $\,^{w_{\af}}C=C$ as elements of $\End(\sigma_{2})$. Let $B:\,^{w_{\af}}\sigma_{2}^{\af}\rightarrow \sigma_{2}^{\af}$ be an $\It_2^{\af}$-isomorphism, which is also unique up to nonzero scaling.

Now by construction $A\circ B\circ C\in\Hom_{\It_2^{\af}}(\,^{w_{\alpha}}\sigma_{2},\sigma_{2})$ and this map is nonzero. (We note though that this map may not be an isomorphism.)
\end{proof}

\begin{Rmk}
    Suppose that $e\geq 2$. Since $U_{\alpha_{0},2e}$ acts trivially on $\sigma_{2}^{\af}$ and $\,^{w_{\af}}\sigma_{2}^{\af}\cong \sigma_{2}^{\af}$ the canonical image of $\sigma_{2}^{\af}$ in $\sigma_{2}$ actually must be contained in $\sigma_{2}^{U_{-\alpha_{0},2e-4}}$. (When $e=1$ the group $U_{-\alpha_{0},2e-4}$ is not contained in $I_{2}$, but $U_{-\alpha_{0},0}$ already acts trivially on $\sigma_{2}$.)
    
\end{Rmk}


\begin{comment}
\textbf{Remark:} The isomorphism $\sigma_{2}^{\alpha}\cong \,^{w_{\alpha}}\sigma_{2}^{\alpha}$ in the proof of Lemma \ref{affSupp} can be made explicit. With this explicit generator in $\Hom_{ I_2^{\alpha}}(\sigma_{2},\,^{w_{\alpha}}\sigma_{2})$ it seems like it may be possible to directly compute the quadratic relation for a nonzero element of $(\Hcal_{2})_{w_{\alpha}}$. However, at this point we have not been able to carry this out.
\end{comment}

%%%%%%%%%%%%%%%%%%%%%%%%%%%%%%%%%%%%%%%%%%%%%%%%%%%%%%%%%%%%%%%%%%%%%%%

\subsection{Support: length $0$ elements}\label{SuppLen0}

%%%%%%%%%%%%%%%%%%%%%%%%%%%%%%%%%%%%%%%%%%%%%%%%%%%%%%%%%%%%%%%%%%%%%%%

In this subsection we show that $\Hcal_{2}$ is supported on the length $0$ elements $\Omegat\subseteq \Wt_{\ea}$. Let us first mention
\begin{Lem}\label{L:lenght_0_element_representative}
    The elements of $\Omegat$ are represented by elements in $N_{\Gt}(\Tt)\cap N_{\Gt}(\It_2)$.
\end{Lem}
\begin{proof}
This readily follows from  Lemma \ref{Len0elts} and Theorem \ref{ProportionalLengths}.
\end{proof}

We then have

\begin{Prop}\label{Len0Supp}
	For each $n\in N_{\Gt}(\Tt)\cap N_{\Gt}(\It_2)$, we have
	\begin{equation*}
		(\Hcal_{2})_{n}\cong\Hom_{\It_2\cap{^{n}\It_{2}}}(\,^{n}\sigma_{2},\sigma_{2})\cong \mathbb{C}\qand \Hcal_n\cong\Hom_{\It\cap{^n\It}}(\,^n\sigma, \sigma)\cong\C.
	\end{equation*}
    Hence both $\Hcal_2$ and $\Hcal$ are supported on the length 0 elements.
\end{Prop}

\begin{proof}
For the proof, we need to use the profinite model \eqref{E:profinite_model_sigma_2} for $\sigma_2$. Then the argument for $\Hcal_2$ is similar to Proposition \ref{IRep} \eqref{SigAWInv} along with the observation that modulo $\Tt_{0}$ the element $n$ represents an element of $\Omegat\subseteq \widetilde{W}_{\ea}\cong W\ltimes \Yt$. Again we apply Proposition \ref{TwistAut}. The first bullet is satisfied because $\tau$ is $W$-invariant and $\widetilde{Y}$-invariant. The second bullet is satisfied by a minor variant of Lemma \ref{WCompat}. The assertion for $\Hcal$ follows from Proposition \ref{HMult1}, Proposition \ref{HIsoH2} and Lemma \ref{L:Bushnell_Kutzko_induced}.
\end{proof}

%%%%%%%%%%%%%%%%%%%%%%%%%%%%%%%%%%%%%%%%%%%%%%%%%%%%%%%%%%%%%%%%%%%%%%%

\subsection{Quadratic relations}\label{SS:QuadRel}

%%%%%%%%%%%%%%%%%%%%%%%%%%%%%%%%%%%%%%%%%%%%%%%%%%%%%%%%%%%%%%%%%%%%%%%

In this subsection we prove that the elements of $\Hcal$ constructed in \S \ref{SuppAfSimpRel} and \ref{SuppLen0} satisfy quadratic relations. Given our current tool set, we must consider several cases. We begin with the quadratic relations for the elements corresponding to simple linear reflections. For these elements we work with the Iwahori model $\Hcal$.

Let $\alpha\in \Delta$. By Lemma \ref{affSupp} we know that $\Hcal_{w_{\alpha}}\cong \mathbb{C}$. We need to identify a suitably normalized generator of this one dimensional vector space. To do this we will use the finite Shimura correspondence.

Let $\Hcal_{\Kt}\subseteq\Hcal$ be the subalgebra of functions supported on $\Kt$. Pulling back along the surjective group homomorphism $\Kt\twoheadrightarrow \widetilde{{\Kf}}$ defines a $\mathbb{C}$-algebra isomorphism $\Hcal(\widetilde{{\Kf}},\widetilde{{\If}},\sigma)\stackrel{\sim}{\rightarrow} \Hcal_{\Kt}$ such that it is support-preserving with respect to the surjection $\Kt\twoheadrightarrow \widetilde{{\Kf}}$. Now we can apply the finite Shimura correspondence of Hecke algebras to describe the structure of $\Hcal(\widetilde{{\Kf}},\widetilde{{\If}},\sigma)$. We make this precise in the next proposition.

\begin{Prop}\label{QuadBraidSimpLinReflect}
	For every $\alpha\in\Delta$ there exists $T_{\alpha}\in \Hcal_{w_{\alpha}}$ such that the follow identities hold. For all $\alpha,\beta\in \Delta$ we have 
	\begin{enumerate}
		\item $T_{\alpha}^{2}=(q-1)T_{\alpha}+q$;\label{QuadSimpLin}
		\item if $\langle\alpha,\beta\rangle=0$, then $T_{\alpha}T_{\beta}=T_{\beta}T_{\alpha}$;\label{CommSimpLin}
		\item if $\langle\alpha,\beta\rangle=-1$, then $T_{\alpha}T_{\beta}T_{\alpha}=T_{\beta}T_{\alpha}T_{\beta}$.\label{BraidSimpLin}
	\end{enumerate}
In particular, for each $\alpha\in\Delta$, the element $T_{\alpha}$ is invertible.
\end{Prop}

\begin{proof}
	Corollary \ref{ParaHeckeIso} describes a support preserving isomorphism 
	\begin{equation*}
		\Hcal(\widetilde{{\Kf}},\widetilde{{\If}},\sigma)\cong \Hcal(G_{\kappa},B_{\kappa},1).
	\end{equation*}
	Iwahori \cite{I64} proved that for each $\alpha\in\Delta$ there exists a unique element $t_{\alpha}\in \Hcal(G_{\kappa},B_{\kappa},1)_{w_{\alpha}}$ satisfying the identities \eqref{QuadSimpLin}, \eqref{CommSimpLin}, \eqref{BraidSimpLin}. For each $\alpha\in\Delta$, define $T_{\alpha}\in \Hcal_{\Kt}$ to be the image of $t_{\alpha}$ under the chain of algebra isomorphisms $\Hcal_{\Kt}\cong\Hcal(\widetilde{{\Kf}},\widetilde{{\If}},\sigma)\cong \Hcal(G_{\kappa},B_{\kappa},1)$ prescribed above. As the two isomorphisms preserve support, it follows that $T_{\alpha}\in \Hcal_{w_{\alpha}}$.
\end{proof}

It remains to consider the simple affine reflection $w_{\af}$. Here it will be convenient to work with the $I_2$-model $\Hcal_{2}$. Let $T_{\af}\in (\Hcal_{2})_{w_{\af}}$ be any nonzero element. First we will prove that $T_{\af}$ satisfies a weak quadratic relation, which will be valid for all Cartan types. Then we will refine this in types $A_r,D_r,E_{6}, E_{7}$, by normalizing $T_{\af}$ so that it satisfies the same quadratic relation that appears in Proposition \ref{QuadBraidSimpLinReflect}. The same quadratic relation is expected to hold in type $E_{8}$, but we are not able to prove it in this paper. We will make some comments about this at the end of this subsection.

\begin{Prop}\label{WeakQuadSimpAff}
	Let $T_{\af}\in (\Hcal_{2})_{w_{\af}}$ be nonzero. Then there exists $a,b\in\mathbb{C}$ with $b\neq 0$ such that
	\begin{equation*}
		T_{\af}^{2}=aT_{\af}+b.
	\end{equation*}
        In particular, $T_{\af}$ is invertible.
\end{Prop}
 
\begin{proof}
In this proof we will occasionally abuse notation and write $w_{\af}$ for the element $\wt_{\alpha_0}(\varpi^2)\in N_{\Gt}(\Tt)$.
 
By the definition of convolution and the Steinberg relations, 
	\begin{align*}
		\supp \,T_{\af}^{2}\subseteq&  \It_2\wt_{\alpha_{0}}(\varpi^{2}) \It_2\wt_{\alpha_{0}}(\varpi^{2}) \It_2\\
		=& \It_2\wt_{\alpha_{0}}(\varpi^{2})U_{\alpha_{0},2}\wt_{\alpha_{0}}(-\varpi^{2}) \It_2\\
		=& \It_2U_{-\alpha_{0},-2} \It_2\\
		=& \It_2\cup \It_2\wt_{\alpha_{0}}(\varpi) \It_2\cup \It_2\wt_{\alpha_{0}}(\varpi^{2}) \It_2.
	\end{align*}
    By Proposition \ref{HVanish} and the support preserving isomorphism $\Hcal\cong \Hcal_{2}$, we know that $T_{\af}^{2}(g)=0$ for any $g\in \It_2\wt_{\alpha_{0}}(\varpi) \It_2$. Therefore there exists $a,b\in\mathbb{C}$ such that
	\begin{equation*}
		T_{\af}^{2}=aT_{\af}+b.
	\end{equation*}
	
	We claim that $T_{\af}^{2}(1)=b\neq 0$. First, for arbitrary $g\in G$, we have
	
	\begin{align*}
		T_{\af}^{2}(g)=&\int_{\It_2w_{\af}\It_2}T_{\af}(h)T_{\af}(h^{-1}g)dh\\
		=&\sum_{h\in\It_2w_{\af}\It_2/\It_2}T_{\af}(h)T_{\af}(h^{-1}g)\\
		=&\sum_{h\in U_{\alpha_{0},2}/U_{\alpha_{0},4}}T_{\af}(hw_{\af})T_{\af}(w_{\af}^{-1}h^{-1}g).
	\end{align*}
	Now we take $g=1$. Then
	\begin{align*}
		T_{\af}^{2}(1)=&\sum_{h\in U_{\alpha_{0},2}/U_{\alpha_{0},4}}T_{\af}(hw_{\af})T_{\af}(w_{\af}^{-1}h^{-1})\\
		=&\sum_{h\in U_{\alpha_{0},2}/U_{\alpha_{0},4}}\sigma_{2}(h)T_{\af}(w_{\af})T_{\af}(w_{\af}^{-1})\sigma_{2}(h^{-1})\\
        =& \sum_{h\in U_{\alpha_{0},2}/U_{\alpha_{0},4}}\sigma_{2}(h)T_{\af}(w_{\af})T_{\af}(w_{\af}\htt_{\alpha_{0}}(-1))\sigma_{2}(h^{-1})\\
        =&\sum_{h\in U_{\alpha_{0},2}/U_{\alpha_{0},4}}\sigma_{2}(h)T_{\af}(w_{\af})^{2}\sigma_{2}(\htt_{\alpha_{0}}(-1)h^{-1}),
	\end{align*}
    where we used $w_{\af}^{-2}=\htt_{\alpha_{0}}(-1)$ for the second equality.
    
        Next we apply Proposition \ref{affSupp} to compute $T_{\af}(w_{\af})^{2}\in\End(\sigma_{2})$. From the proof of Proposition \ref{affSupp} we know that the map $T_{\af}(w_{\af})$ factors into three pieces, $T_{\af}(w_{\af})=A\circ B\circ C$, where $A,B,C$ are as in the proof of Proposition \ref{affSupp}. Thus we have $T_{\af}(w_{\af})^{2}=A\circ B\circ C\circ A\circ B\circ C$. Now by definition $C\circ A = id_{\sigma_{2}^{\af}}$. Thus $A\circ B\circ C\circ A\circ B\circ C=A\circ B^{2}\circ C$. Now $B\in \Hom_{ \It_2^{\af}}(\,^{w_{\af}}\sigma_{2}^{\af},\sigma_{2}^{\af})$, where we have put $I_{2}^{\af}= I_2\cap \,^{w_{\af}} I_2$. Thus 
        \begin{align*}
            B^{2}\in\Hom_{ \It_2^{\af}}(\,^{w_{\af}}\sigma_{2}^{\af},\,^{w_{\af}^{-1}}\sigma_{2}^{\af})=&\Hom_{ \It_2^{\af}}(\sigma_{2}^{\af},\,^{w_{\af}^{-2}}\sigma_{2}^{\af})\\
            =&\Hom_{ \It_2^{\af}}(\sigma_{2}^{\af},\,^{\htt_{\alpha_{0}}(-1)}\sigma_{2}^{\af})\\
            =&\mathbb{C}\sigma_{2}^{\af}(\htt_{\alpha_{0}}(-1)).
        \end{align*}
        Therefore $T_{\af}(w_{\af})^{2} = c A\circ \sigma_{2}^{\af}(\htt_{\alpha_{0}}(-1))\circ C$ for some $c\in\mathbb{C}^{\times}$. Since $\sigma_{2}^{\af}(h_{\alpha_{0}}(-1))\circ C=C\circ \sigma_{2}(\htt_{\alpha_{0}}(-1))$ we have
        \begin{align*}
            T_{\af}(w_{\af})^{2}=&\sum_{h\in U_{\alpha_{0},2}/U_{\alpha_{0},4}}\sigma_{2}(h)T_{\af}(w_{\af})^{2}\sigma_{2}(\htt_{\alpha_{0}}(-1)h^{-1})\\
            =& c\sum_{h\in U_{\alpha_{0},2}/U_{\alpha_{0},4}}\sigma_{2}(h)\circ A\circ \sigma_{2}^{\af}(\htt_{\alpha_{0}}(-1))\circ C\circ\sigma_{2}(\htt_{\alpha_{0}}(-1)h^{-1})\\
            =&(-1,-1)c\sum_{h\in U_{\alpha_{0},2}/U_{\alpha_{0},4}} \sigma_{2}(h)\circ A\circ C\circ\sigma_{2}(h^{-1}).
        \end{align*}

        By taking the trace and noting that $A\circ C$ is just a projection operator we see that 
        \begin{equation*}
            \Tr(T_{\af}^{2}(1))=(-1,-1,)c\left|\Ocal\slash\varpi\Ocal\right|^{2}\Tr(A\circ C)\neq 0.
        \end{equation*}
        Thus $b=T_{\af}^{2}(1)\neq 0$. This implies $T_{\af}$ is invertible.
    \end{proof}

Now we specialize to the Cartan types $A_r,D_r,E_{6},E_{7}$ and prove an exact quadratic relation. The feature of these types that we exploit is the existence of an inner automorphism of $\Gt$ with good properties, essentially coming from a length $0$ element in $\Wt_{\ea}$. 

Let us make this inner automorphism explicit. When $\Phi$ is simply-laced and not of type $E_{8}$, then there exists a simple root $\alpha_{j_{0}}$ such that in the expression $-\alpha_{0}=\sum c_{j}\alpha_{j}$, we have $c_{j_{0}}=1$. The set of possible $\alpha_{j_{0}}$ is exactly the orbit of $\alpha_{0}$ under the automorphisms of the affine Dynkin diagram induced by the length $0$ elements of the extended affine group $\Wt_{\ea}$, excluding $\alpha_{0}$ itself. Alternatively, this set is in bijection with the nontrivial cosets in $\Yt/2Y$. In particular, the simple root $\alpha_{j_{0}}$ corresponds to the coset $2\omega_{j_{0}}^{\vee}+2Y\in\Yt/2Y$, where $\omega_{j_0}^\vee$ is the fundamental coweight associated with $\alpha_{j_0}$. (See \cite[Ch.\ 6, \S 2.3, Prop.\ 6]{BourLGLA46} for details.)




\begin{Lem}\label{StabGam2}
	Suppose that $\Phi$ is of type $A_r,D_r$, $E_{6}$, or $E_{7}$. Let $\varepsilon\in\Omegat$ be nontrivial. Let $\alpha_{j_{0}}:=\varepsilon(\alpha_{0})\in \Delta$. Then there exists $n\in N_{\Gt}(\Tt)\cap N_{\Gt}( I_2)$ representing $\varepsilon\in\Omegat$ such that
	\begin{enumerate}
		\item ${}^n\Tt=\Tt$,\label{TInv}
		\item ${}^n\It_2=\It_2$,\label{Gamma2Inv}
		\item ${}^n\It_2\wt_{\alpha_{j_{0}}}\It_2=\It_2w_{\af}\It_2$.\label{Gam2AfDoubCoset}
	\end{enumerate}
\end{Lem}

\begin{proof}
	All three properties follow directly by construction.
\end{proof}

With this lemma, we can prove

\begin{Prop}\label{QuadSimpAff} 
Let $\Phi$ be of type $A_r$,$D_r$, $E_{6}$ or $E_{7}$. Then there exists $T_{\af}\in(\Hcal_{2})_{w_{\af}}$ such that 
\begin{equation*}
		T_{\af}^{2}=(q-1)T_{\af}+q.
	\end{equation*}
\end{Prop}

\begin{proof}
Let $n\in N_{\Gt}(\Tt)\cap N_{\Gt}( I_2)$ represent a nontrivial element in $\Omegat$. Such $n$ exists by Lemma \ref{L:lenght_0_element_representative}. Since $\,^{n}\sigma_{2}\cong \sigma_{2}$ by Proposition \ref{Len0Supp}, there exists $A_{n}\in \GL(\sigma_2)$ such that $A_{n}\,^{n}\sigma_{2}(g)=\sigma_{2}(g)A_{n}$ for all $g\in  I_2$.

Define
\[
T_{\af}:=A_{n}(T_{\alpha_{j_{0}}}\circ\Int(n^{-1}))A_{n}^{-1}\in\Hcal_{2},
\]
by which we mean $T_{\af}$ is such that $T_{\af}(g)=A_{n}(T_{\alpha_{j_{0}}}\circ\Int(n^{-1})(g))A_{n}^{-1}$. One can then see from Lemmas \ref{L:HeckeAutTwist} and \ref{StabGam2} that
\[
T_0\in(\Hcal_2)_{w_{\af}}.
\]
The quadratic relation follows from the quadratic relation for $T_{\alpha_{j_{0}}}$ (Proposition \ref{QuadBraidSimpLinReflect}).
\end{proof}

Finally we prove that the length $0$ elements satisfy a quadratic relation.

\begin{Prop}\label{QuadLen0}
	Let $\varepsilon\in \Omegat$. There exists $T_{\varepsilon}\in (\Hcal_{2})_{\varepsilon}$ (unique up to sign) such that 
	\begin{equation*}
		T_{\varepsilon}^{2}=1.
	\end{equation*}
	In particular, $T_{\varepsilon}$ is invertible.
\end{Prop}

\begin{proof}
	Let $t_{\varepsilon}\in (\Hcal_{2})_{\varepsilon}$ be nonzero. Note that $\Omegat\cong \widetilde{Y}/2Y$ is an elementary abelian $2$-group. With this and the definition of convolution we have $\supp(t_{\varepsilon}^{2})\subseteq  \It_2 \varepsilon^{2}= \It_2$. Thus we see that $t_{\varepsilon}^{2}=C$ for some $C\in \mathbb{C}$.
	
	Note that the subgroup $\langle \It_2, \varepsilon\rangle\subseteq\Gt$ generated by $\It_2$ and $\varepsilon$ has $\It_2$ as a subgroup of index $2$. Thus by Proposition \ref{Len0Supp} and Clifford theory we know that $\Ind_{\It_2}^{\langle \It_2, \varepsilon\rangle}(\sigma_{2})$ decomposes into two distinct irreducible representations $\pi_{1},\pi_{2}$ extending $\sigma_{2}$. Then because $t_{\varepsilon}$ can be viewed as an element of the semisimple algebra $\End_{\langle \It_2, \varepsilon\rangle}(\Ind_{\It_2}^{\langle \It_2, \varepsilon\rangle}(\sigma_{2}))$ it follows that $C\neq 0$. Thus if we  define $T_{\varepsilon}:=(C)^{-\frac{1}{2}}t_{\varepsilon}$ (unique up to sign), then $T_{\varepsilon}^{2}=1$.
\end{proof}


%%%%%%%%%%%%%%%%%%%%%%%%%%%%%%%%%%%%%%%%%%%%%%%%%%%%%%%%%%%%%%%%%%%%%%%

\subsection{Braid relations}\label{SS:BraidRel}

%%%%%%%%%%%%%%%%%%%%%%%%%%%%%%%%%%%%%%%%%%%%%%%%%%%%%%%%%%%%%%%%%%%%%%%

In this subsection we investigate the braid relations. We begin by proving the braid relations up to scaling.

For $w\in \Wt_{\ea}$, we let $t_{w}\in (\Hcal_{2})_{w}$. By Proposition \ref{HMult1} the element $t_{w}$ is unique up to scaling. Later we will be concerned with normalizing the elements $t_{w}$, but for now this is not important.

\begin{Prop}\label{Prop:GenBraid}
	Let $w_{1},w_{2}\in \Wt_{\ea}$. If $\ell(w_{1}w_{2})=\ell(w_{1})+\ell(w_{2})$, then in $\Hcal_{2}$ we have 
    \[
    t_{w_{1}}*t_{w_{2}}=C(w_{1},w_{2})t_{w_{1}w_{2}},
    \]
    for some $C(w_{1},w_{2})\in\mathbb{C}$. Furthermore, if $t_{w_{1}}$ is invertible and $t_{w_{2}}\neq 0$, then $C(w_{1},w_{2})\in \mathbb{C}^{\times}$.
	
	%\textbf{$C(w_{1},w_{2})$ should be a $2$-cocycle. Add this.}
\end{Prop}

\begin{proof}
	Our proof follows \cite[Proposition 6.8]{Karasiewicz}, which adapted the argument of \cite[Proposition 6.2]{S04}.
	
	For any $g\in \Gt$ we need to compute $t_{w_{1}}*t_{w_{2}}(g)$. By the definition of convolution we need to identify the cosets $\delta\It_2\subseteq \Gt$ such that
	\begin{align*}
		\delta\It_2\subseteq& \It_2w_{1}\It_2\qand\\
		(\delta\It_2)^{-1}g\subseteq& \It_2w_{2}\It_2.
	\end{align*}
	These containments only depend on $G$. Thus by Theorem \ref{ProportionalLengths} and Proposition \ref{CosetProp}, which uses the length hypothesis, it follows that $\supp\,t_{w_{1}}*t_{w_{2}}\subseteq \It_2w_{1}w_{2}\It_2$.
	
	By Proposition \ref{HMult1} there exists $C\in\mathbb{C}$ such that $t_{w_{1}}*t_{w_{2}}=Ct_{w_{1}w_{2}}$. If we further assume that $t_{w_{1}}$ is invertible and $t_{w_{2}}\neq 0$, then $C$ cannot be $0$.
	\end{proof}
	
	Using Proposition \ref{Prop:GenBraid} we can define a basis of $\Hcal_{2}$ such that each element is invertible. 
	
	\begin{Lem}\label{SuppComplete}
		Let $w\in \widetilde{W}_{\ea}$. Let $\varepsilon\in\Omegat$ and let $\beta_{1},\ldots,\beta_{k}\in \Delta\cup\{\alpha_{0}\}$ be a finite sequence such that
		\begin{equation*}
			w=\varepsilon w_{\beta_{1}}\cdots w_{\beta_{k}}
		\end{equation*}
		is a reduced expression for $w$. Then the element $T_{\varepsilon}T_{\beta_{1}}\cdots T_{\beta_{k}}\in (\Hcal_{2})_{w}$ is nonzero and invertible.
	\end{Lem}
	
	\begin{proof}
		This follows by induction on the length of $w$ and by Propositions \ref{QuadBraidSimpLinReflect}, \ref{QuadSimpAff}, \ref{QuadLen0}, and \ref{Prop:GenBraid}. 
	\end{proof}
	
	\begin{Prop}\label{HLineDecompAndInvert}
		The set of subspaces $(\Hcal_{2})_{w}$, where $w\in \widetilde{W}_{\ea}$, is a line decomposition for $\Hcal_{2}$, namely
        \[
        \Hcal_2=\bigoplus_{w\in\Wt_{\ea}}(\Hcal_{2})_{w}.
        \]
        Furthermore, any nonzero element of $(\Hcal_{2})_{w}$ is invertible.
	\end{Prop}
	
	\begin{proof}
		This follows from Theorem \ref{HMult1}, Proposition \ref{HVanish} and Lemma \ref{SuppComplete}, along with the fact that the distinct $(\Hcal_{2})_{w}$'s have disjoint support.
	\end{proof}
	
	For the case of type $A_r,D_r$, $E_{6}$, $E_{7}$, we can refine Proposition \ref{HLineDecompAndInvert} to produce a normalized basis for $\Hcal_{2}$ for which the usual braid relations are satisfied.
	
	From \S \ref{SS:QuadRel} we have the following normalized elements. For $\alpha\in\Delta\cup\{\af\}$ there is a unique $T_{\alpha}\in (\Hcal_{2})_{w_{\alpha}}$ such that $T_{\alpha}^{2}=(q-1)T_{\alpha}+q$. For $\varepsilon\in\Omegat$ there is a unique up to sign $T_{\varepsilon}\in (\Hcal_{2})_{\varepsilon}$ such that $T_{\varepsilon}^{2}=1$. To begin we show that the elements indexed by $\Delta\cup\{\af\}$ satisfy the braid relations. 
    
    Let us first prove the following elementary fact.	
	\begin{Lem}\label{BraidLem}
		Let $a,b,c\in \mathbb{C}^{\times}$. Let $\mathcal{A}$ be an associative $\mathbb{C}$-algebra with unit. Suppose that $X,Y\in \mathcal{A}$ are such that
		\begin{itemize}
			\item $1,X,Y\in \mathcal{A}$ are linearly independent;
			\item $X^{2}=aX+b$;
			\item $Y^{2}=aY+b$;
			\item $XYX=cYXY$, or $XY=cYX$.
		\end{itemize}
		Then $c=1$.
	\end{Lem}
	
	\begin{proof}
		This follows by a direct calculation. We will prove the case where $XYX=cYXY$. The other case is easier.
		
		Since $b\neq 0$, the elements $X$ and $Y$ are invertible. Thus we know that $X=cYXYX^{-1}Y^{-1}$. Squaring and applying the quadratic relations gives
		\begin{equation*}
			aX+b = c^{2}YX(aY+b)X^{-1}Y^{-1}.
		\end{equation*}
		The right hand side simplifies to 
		\begin{equation*}
			c^{2}YX(aY+b)X^{-1}Y^{-1}=c^{2}(aYXYX^{-1}Y^{-1}+b)=acX+c^{2}b.
		\end{equation*}
		Now since $X$ and $1$ are linearly independent we have $a=ac$ and $b=bc^{2}$. Since $a\neq 0$ it follows that $c=1$.
 	\end{proof}
%The hypotheses in Lemma \ref{BraidLem} are not minimal. This presentation is just to avoid stating two totally separate cases.
	
	\begin{Prop}\label{CompleteBraidRel}
		Let $\Phi$ be of type $A_r,D_r$, $E_{6}$ or $E_{7}$. Let $\alpha,\beta\in \Delta\cup\{\af\}$ be distinct.
		\begin{enumerate}
			\item If $\langle\alpha,\beta\rangle=0$, then $T_{\alpha}T_{\beta}=T_{\beta}T_{\alpha}$.
			\item If $\langle\alpha,\beta\rangle=-1$, then $T_{\alpha}T_{\beta}T_{\alpha}=T_{\beta}T_{\alpha}T_{\beta}$.
		\end{enumerate}
	\end{Prop}
	
	\begin{proof}
		This follows from Proposition \ref{Prop:GenBraid} and Lemma \ref{BraidLem}.
	\end{proof}
	
	\begin{Prop}\label{P:TwDef}
		Let $\Phi$ be of type $A_r,D_r$, $E_{6}$ or $E_{7}$. Let $w\in\widetilde{W}_{\af}$ and let $w=w_{1}\cdots w_{k}$ be a reduced expression for $w$ in terms of the affine simple reflections $w_{\alpha}$, where $\alpha\in\Delta\cup\{\af\}$. Define $T_{w}:=T_{w_{1}}\cdots T_{w_{k}}\in \Hcal_{2}$. Then the element $T_{w}$ is independent of the reduced expression for $w$.
	\end{Prop}
	
	\begin{proof}
		This follows from \cite[Th\'{e}or\`{e}me 3]{T69} and Proposition \ref{CompleteBraidRel}.
        %\cite[Theorem 7.1, Section 7.3 exercise 2]{H90} That might not be enough. 
	\end{proof}
	
	\begin{Lem}\label{OmegaAut}
		Let $\Phi$ be of type $A_r,D_r$, $E_{6}$ or $E_{7}$. Let $\varepsilon\in\Omegat$ and $\alpha\in\Delta\cup\{\af\}$. Then $T_{\varepsilon}T_{\alpha}T_{\varepsilon}^{-1}=T_{\varepsilon(\alpha)}$.
	\end{Lem}
	
	\begin{proof}
		By Proposition \ref{Prop:GenBraid} we know that $T_{\varepsilon}T_{\alpha}T_{\varepsilon}^{-1}=cT_{\varepsilon(\alpha)}$ for some $c\in\mathbb{C}^{\times}$. Because $T_{\alpha}$ and $T_{\varepsilon(\alpha)}$ satisfy the same quadratic relation with nonzero linear coefficient it follows that $c=1$ by Lemma \ref{BraidLem}.
	\end{proof}


    \subsection{IM-presentation}\label{SS:IMPres}
    
    Finally we can prove that when $\Phi$ has Cartan type $A_r,D_{2r+1},E_{6},E_{7}$ the Hecke algebra $\Hcal$ has an Iwahori-Matsumoto presentation (IM-presentation for short). The same result is expected to hold for $D_{2r}$ and $E_{8}$, though we are not able to prove it in this paper.

        Let us first define an abstract algebra derived from the extended Weyl group $\Wt_{\ea}$ defined via generators and relations as in \cite[Section 3.1]{IM}. Let 
        \[
        H_{\af}=\la 1, t_{\alpha}\st \alpha\in\Delta\cup\{\af\}\ra
        \]
        be the unital $\mathbb{C}$-algebra generated by the symbols $t_{\alpha}$, where $\alpha\in\Delta\cup\{\af\}$, subject to the following relations:
        \begin{align*}
            &t_{\alpha}^{2}=(q-1)t_{\alpha}+q;\\
            &t_{\alpha}t_{\beta}=t_{\beta}t_{\alpha},\quad \text{if } \langle\alpha,\beta\rangle=0;\\
            &t_{\alpha}t_{\beta}t_{\alpha}=t_{\beta}t_{\alpha}t_{\beta},\quad \text{if } \langle\alpha,\beta\rangle=-1.
        \end{align*}
        Note that for $w\in \Wt_{\af}$ with a reduced expression $w=w_{\alpha_{i_{1}}}\cdots w_{\alpha_{i_{k}}}$ we can define $t_{w}=t_{\alpha_{i_{1}}}\cdots t_{\alpha_{i_{k}}}$. This definition is independent of the choice of reduced word by the same argument as in Proposition \ref{P:TwDef}.
        
        The group $\Omegat\subseteq \Wt_{\ea}$ acts on $H_{\af}$ by algebra automorphisms as follows. Recall that the elements of $\Omegat$ act on $\Delta\cup\{\af\}$ by permutation. Then for $\varepsilon\in\Omegat$ define $\varepsilon(t_{\alpha})=t_{\varepsilon(\alpha)}$. This induces an algebra automorphism of $H_{\af}$. Moreover, this is compatible with the action of $\Omegat$ on $\Wt_{\af}$ in the sense that $\varepsilon(t_{w})=t_{\varepsilon w\varepsilon^{-1}}$. Let 
        \[
        H_{\ea}:=\mathbb{C}[\Omegat]\otimes H_{\af}.
        \]
        This $\mathbb{C}$ vector space can be given the structure of a $\mathbb{C}$-algebra as follows. For $\varepsilon_{j}\in \Omegat$ and $w_{j}\in \Wt_{\af}$ define
        \begin{equation*}
            (\varepsilon_{1}\otimes t_{w_{1}})\cdot (\varepsilon_{2}\otimes t_{\alpha_{2}}):=\varepsilon_{1}\varepsilon_{2}\otimes t_{\varepsilon_{2}^{-1}(\alpha_{1})}t_{\alpha_{2}}
        \end{equation*}

        
        
        \begin{Thm}\label{IMPres} Assume that $G$ has Cartan type $A_r, D_{2r+1}, E_{6},E_{7}$. Then the linear map $H_{\ea}\rightarrow \Hcal$ defined by $\varepsilon\otimes t_{w}\mapsto T_{\varepsilon}T_{w}$, where $\varepsilon\in\Omegat$ and $w\in \Wt_{\af}$, is an isomorphism of $\mathbb{C}$-algebras.
        \end{Thm}

        \begin{proof}
            The map is a linear isomorphism because the set $\{\varepsilon\otimes t_{w}\st\varepsilon\in\Omegat,\,w\in\Wt_{\af}\}$ is a basis for $H_{\ea}$ and the set $\{T_{\varepsilon}T_{w}\st\varepsilon\in \Omegat,\,w\in\Wt_{\af}\}$ is a basis for $\Hcal$. The map is a $\mathbb{C}$-algebra homomorphism because of Propositions \ref{QuadBraidSimpLinReflect}, \ref{QuadSimpAff}, \ref{QuadLen0}, \ref{CompleteBraidRel}, and Lemma \ref{OmegaAut}.
        \end{proof}

        Theorem \ref{IMPres} and the work of \cite{IM} imply that $\Hcal$ is isomorphic to the Iwahori-Hecke algebra of a linear group. Specifically, let 
        \[
        G^{\prime}:=G/Z_{2},
        \]
        where $Z_{2}$ is the subgroup of $2$-torsion in the center of $G$. (The significance of $Z_{2}$ is that the fundamental group of $G^{\prime}$ is isomorphic to $\Yt/2Y$. Thus the dual root datum of $G^{\prime}$ agrees with the dual root datum of $\Gt$.) Let $I^{\prime}$ be an Iwahori subgroup of $G^{\prime}$ and let $\Hcal(G^{\prime},I^{\prime})$ be the Iwahori-Hecke algebra of locally constant $I^{\prime}$-biinvariant functions on $G^{\prime}$.

        \begin{Cor}\label{HeckeShimCorr}
            As $\mathbb{C}$-algebras, we have
            \[
            \Hcal(G^{\prime},I^{\prime})\cong\Hcal\cong H_{\ea}.
            \]
        \end{Cor}

        \begin{proof}
            This follows directly from Theorem \ref{IMPres} and \cite[Section 3]{IM}. 
        \end{proof}
	
We end this subsection with a few comments on the cases of $E_{8}$ and $D_{2r}$. We expect the same type of Iwahori-Matsumoto presentation to hold in the case of $D_{2r}$ and $E_{8}$ as well. When $F$ is unramified over $\mathbb{Q}_{2}$ the approach of \cite{Karasiewicz} can be generalized by incorporating some ideas from this current paper. We will not go into this because we believe that these cases for any $F$ will follow as we refine and expand the techniques used in this paper. 

For example, for $E_{8}$ the only obstruction is proving that the the simple affine reflection satisfies the desired quadratic relation. This should follow from an analysis along the lines of \S \ref{SSecSplittings} for other vertices in the building of $G$ and an analogous finite Shimura correspondence at these vertices. For $D_{2r}$ the difficulty arises because in this case $\Omegat$ is not a cyclic group. The resolution of this issue is related to the question of extending $\sigma_{2}$ from $\It_2$ to its normalizer in $\Gt$. However, at the moment, certain basic properties of $\sigma_{2}$ remain unknown, such as its dimension. As we further investigate the basic properties of $\sigma_{2}$, this will naturally put us in a better position to address the extension question.

%%%%%%%%%%%%%%%%%%%%%%%%%%%%%%%%%%%%%%%%%%%%

\subsection{Satake Isomorphism}\label{SS:Satake}

%%%%%%%%%%%%%%%%%%%%%%%%%%%%%%%%%%%%%%%%%%%%

In this subsection we state the Satake isomoprhism for our double cover $\Gt$. Everything in this subsection essentially follows from a Bernstein presentation of $\Hcal$, which can be deduced from an IM presentation. So, it applies when $G$ has Cartan type $A_{r},D_{2r+1},E_{6}$, $E_{7}$; it will apply to the other types when the remaining gaps in the IM presentation are closed. 

Let $Z(\Hcal)$ be the center of the Hecke algebra $\Hcal$. Recall that since $\Hcal$ satisfies an IM presentation deforming the extended affine Weyl group $\Wt_{ea}\cong W\ltimes \Yt$ it is known that $Z(\Hcal)\cong \mathbb{C}[\Yt]^{W}$ (e.g. \cite{L89} or \cite[Section 2]{HKP10}).

Let 
\[
e_{\sigma}:=\frac{1}{\sum_{w\in W}q^{\ell(w)}}\sum_{w\in W} T_{w}.
\]
Note that since the $\It$-representation $\sigma$ extends to $\Kt$ we have $\Hcal(\Gt,\Kt,\sigma)=e_{\sigma}*\Hcal*e_{\sigma}$.

Now using the known structure of affine Hecke algebras one can prove the following corollary.

\begin{Cor}\label{C:Satake}
    Suppose that $\Hcal$ admits an IM-presentation. Then the map $Z(\Hcal)\rightarrow \Hcal(\Gt,\Kt,\sigma)$ defined by $z \mapsto z*e_{\sigma}$ defines an isomorphism of $\mathbb{C}$-algebras. In particular, $\Hcal(\Gt,\Kt,\sigma)$ is commutative.
\end{Cor}

%%%%%%%%%%%%%%%%%%%%%%%%%%%%%%%%%%%%%%%%%%%%

\section{Iwahori types}\label{S:IwahoriType}

%%%%%%%%%%%%%%%%%%%%%%%%%%%%%%%%%%%%%%%%%%%%

In this section we show that $(\It,\sigma)$ is a type in the sense of \cite[(4.1) Definition]{BK98}. We emphasize that this result is valid for any Cartan type. As a corollary, we derive a local Shimura correspondence when $G$ is of Cartan type $A_r, D_{2r+1}, E_6$ and $E_7$.

The following arguments essentially follow \cite{BK98}, which build on the Borel-Casselman theorem. However, there is one notable difference. Bushnell-Kutzko introduce the notion $G$-cover \cite[(8.1) Definition]{BK98} as a way to axiomatize the construction of types for a reductive group $G$ from types of its Levi subgroups. Their main theorem for constructing types using $G$-covers is \cite[(8.3) Theorem]{BK98}. The definition of $G$-cover refers back to \cite[(6.1) Definition]{BK98} which contains two conditions (i) and (ii). Our proposed type $(\It,\sigma)$ satisfies condition (ii) by Proposition \ref{HLineDecompAndInvert}, but it does not satisfy condition (i). Specifically, $\sigma$ is neither trivial on $U_{0}$ nor $U^{-}_{1}$, when $e>1$. Fortunately, we are able to modify the argument of Bushnell-Kutzko using special properties of $\sigma$. For example, Lemma \ref{VeryPosSupp} plays a particularly important role. So, we include the details to make clear the necessary changes, in the hope that they might point toward a valuable extension of the notion of $G$-cover.

We begin with a few lemmas.

\begin{Lem}\label{L3to1}
	Let $G$ be a finite group with an Iwahori factorization $(V,T,U)$. Let $\pi\in \Irr(G)$ be such that $r_{U,V}\pi\neq 0$. The element $e_{U}e_{V}e_{U}\in \End({\pi})$ induces a linear isomorphism of the space $\pi^{U}$ and thus is a generator of the one-dimensional space $\Hom_{T}(\pi^{U},\pi^{U})=\mathbb{C}e_{U}$.
\end{Lem}

\begin{proof}
	By Proposition \ref{P:two_functors_for_Iwahori_decomp} the element $e_{U}e_{V}\in \Hom_{T}(\pi^{V},\pi^{U})$ is an isomorphism, and $r_{U,V}\pi\cong \pi^{U}\cong \pi^{V}$ is an irreducible $T$-module. Similarly, $e_{V}e_{U}\in \Hom_{T}(\pi^{U},\pi^{V})$ is an isomorphism. Hence $e_{U}e_{V}e_{U}=(e_Ue_V)(e_Ve_U)\in \Hom_{T}(\pi^{U},\pi^{U})=\mathbb{C}e_{U}$ is nonzero and the result follows.
\end{proof}

In the next lemma we prove that the Hecke algebra $\Hcal$ contains elements supported on certain positive elements of $Z(\Tt)$ with a natural normalization. We note that this lemma does not rely on any of the previous support computations in $\Hcal$. 

\begin{Lem}\label{VeryPosSupp}
	Let $z\in Z(\Tt)$ be such that $zU_{0}z^{-1}\subseteq U_{2e}$. Then the element $e_{\Uf_{0}}e_{\Uf^{-}_{1}}\in\Hcal(\,\Ift\,)$ defines a nonzero element in $\Hom_{\It\cap\,^{z}\It}(\sigma,\,^{z}\sigma)$.
\end{Lem}

\begin{proof}
	Note that since $zU_{0}z^{-1}\subseteq U_{2e}$, it also follows that $z^{-1}U^{-}_{1}z\subseteq U^{-}_{2e+1}$. This implies that $\It\cap\,^{z}\It$ has an Iwahori decomposition
	\begin{equation*}
		\It\cap\,^{z}\It\cong \,^{z}U_{0}\times\Tt_{0}\times U^{-}_{1}.
	\end{equation*} 
	 We have to show $e_{\Uf_{0}}e_{\Uf^{-}_{1}}\in \Hom_{\It\cap\,^{z}\It}(\sigma,\,^{z}\sigma)$. For this it suffices to show that it is $\Tt_{0}$-, $U_0$- and $U_1^-$-intertwining. But one can readily see that it is $\Tt_0$-intertwining because $z\in Z(\Tt)$. Next, since $\,^{z}U_{0}\subseteq U_{2e}\subseteq \ker\sigma$, we see that $\,^{z}U_{0}$ acts trivially on the domain. Now the image of $e_{\Uf_{0}}e_{\Uf^{-}_{1}}$ is contained in $\sigma^{U_{0}}=(\,^{z}\sigma)^{\,^{z}U_{0}}$. Thus $\,^{z}U_{0}$ also acts trivially on the codomain. Thus $e_{\Uf_{0}}e_{\Uf^{-}_{1}}$ intertwines the $\,^{z}U_{0}$ action. The case of $U^{-}_{1}$ is similar.
		
	The map $e_{\Uf_{0}}e_{\Uf^{-}_{1}}$ is nonzero by Proposition \ref{P:two_functors_for_Iwahori_decomp} (4) because $\Hom_{\Tt}(\sigma^{\Uf_0},\sigma^{\Uf_1^-})\neq 0$.
\end{proof}

\begin{Rmk}
    Lemma \ref{VeryPosSupp} and Proposition \ref{HMult1} together imply that $\Hom_{\It\cap\,^{z}\It}(\sigma,\,^{z}\sigma)=\mathbb{C}e_{\Uf_{0}}e_{\Uf^{-}_{1}}$.
\end{Rmk} 

The next two lemmas are based on \cite[(7.9) Theorem]{BK98}. We cannot invoke this result directly because in our case the pair $(\It,\sigma)$ does not satisfy hypothesis (a) of this theorem, because $(\It,\sigma)$ does not satisfy condition (i) in \cite[(6.1) Definition]{BK98} as mentioned at the beginning of this section.

\begin{Lem}\label{LemSurj}
Let $q:\pi\to\pi_U$ be the canonical surjection. The map $\Hom(\sigma,\pi)\rightarrow \Hom(\sigma^{U_{0}},\pi_{U})$ defined by $f\mapsto q\circ f|_{\sigma^{U_{0}}}$ is a homomorphism of $(\It\cap \widetilde{B},\widetilde{B})$-bimodules. The map induced by restriction $\Psi:\Hom_{\It}(\sigma,\pi)\rightarrow \Hom_{\Tt_{0}}(\sigma^{U_{0}},\pi_{U})$ is a surjective linear map.
\end{Lem}

\begin{proof}
	Let $f\in \Hom_{\Tt_{0}}(\sigma^{U_{0}},\pi_{U})$. Note that by definition $f\circ e_{{\Uf}_{0}}=f$. When convenient we use this to identify $f$ with $f\circ e_{{\Uf}_{0}}\in \Hom_{\Tt_{0}}(\sigma,\pi_{U})$.
	
	Now since the category of $\Tt_{0}$-modules is semisimple we can fix a $\Tt_{0}$-module splitting $s:\pi_{U}\rightarrow \pi$ of the $\Tt_{0}$-module surjection $q:\pi\rightarrow \pi_{U}$. Define 
    \[
    F:=s\circ f\circ e_{{\Uf}_{0}}\in \Hom_{\Tt_{0}}(\sigma,\pi),
    \]
    so that $q\circ F=f\circ e_{{\Uf}_{0}}$.
	
	Since $\sigma$ is finite dimensional, there exists $j\in \mathbb{Z}_{\geq 1}$ such that $\Im(F)\subseteq \pi^{U^{-}_{j}}$. Let $z\in Z(\Tt)$ be such that $z^{-1}U^{-}_{1}z\subseteq U^{-}_{j}\cap U^{-}_{2e+1}$.

    Define
    \[
    F_{1}:=\int_{\It}\pi(x)\pi(z)F\circ \sigma(x^{-1})dx\in\Hom_{\It}(\sigma,\pi).
    \]
    Let us compare $F_1$ to $f$. Let $w\in \sigma$. Then by the Iwahori factorization, we have
	\begin{align*}
		q\circ F_1(e_{{\Uf}_{0}}w)=&q\circ \int_{\It} \pi(x)\pi(z)F\circ \sigma(x^{-1}) e_{{\Uf}_{0}}w\,dx\\
		=&q\circ \int_{U_{0}\Tt_{0}U^{-}_{1}} \pi(utv)\pi(z)F\circ \sigma((utv)^{-1}) e_{{\Uf}_{0}}w \,du\,dt\,dv\\
		=&q\circ \int_{U_{0}\Tt_{0}U^{-}_{1}} \pi(utvz)F\circ \sigma((utv)^{-1}) e_{{\Uf}_{0}}w \,du\,dt\,dv.
	\end{align*}
	Using the idempotent $e_{{\Uf}_{0}}$ and the fact that $F$ is a $\Tt_{0}$-module homomorphism we have 
	\begin{align*}
		&q\circ \int_{U_{0}\Tt_{0}U^{-}_{1}} \pi(utvz)F\circ \sigma((utv)^{-1}) e_{U_{0}}w \,du\, dt\, dv\\
        \equiv&q\circ \int_{\Tt_{0}U^{-}_{1}} \pi(tvz)F\circ \sigma((tv)^{-1}) e_{{\Uf}_{0}}w \,dt\,dv\\
		\equiv&q\circ \int_{U^{-}_{1}} \pi(vz)F\circ \sigma(v^{-1}) e_{{\Uf}_{0}}w \,dv,
	\end{align*}
    where by $\equiv$ we mean equality up to a positive scalar. (The discrepancy comes from the choice of measures. We will use the same convention in what follow.)
	Now because $\Im(F)\subseteq \pi^{U^{-}_{j}}$ and  $z^{-1}U^{-}_{1}z\subseteq U^{-}_{j}$, we have
	\begin{align*}
		q\circ \int_{U^{-}_{1}} \pi(vz)F\circ \sigma(v^{-1}) e_{{\Uf}_{0}}w \,dv
        \equiv&q\circ \int_{U^{-}_{1}} \pi(z)\pi(z^{-1}vz)F\circ \sigma(v^{-1}) e_{{\Uf}_{0}}w \,dv\\
		\equiv&q\circ \pi(z)F\circ e_{{\Uf}^{-}_{1}}e_{{\Uf}_{0}}w\\
		\equiv&\pi_{U}(z)f \circ e_{{\Uf}_{0}}e_{{\Uf}^{-}_{1}}e_{{\Uf}_{0}}w,
	\end{align*}
    where the last equality is by the definition of $F$. Finally we can apply Lemma \ref{L3to1} to get 
	\begin{equation*}
		\pi_{U}(z)f \circ e_{{\Uf}_{0}}e_{{\Uf}^{-}_{1}}e_{{\Uf}_{0}}w\equiv\pi_{U}(z)f \circ e_{{\Uf}_{0}}w.
	\end{equation*}
	Thus we have shown that
	\begin{equation}\label{SurjEq1}
		q\circ F_1(e_{{\Uf}_{0}}w)\equiv\pi_{U}(z)f \circ e_{{\Uf}_{0}}w.
	\end{equation}
	
	Recall that for $T\in \Hcal$ and $A\in \Hom_{\It}(\sigma,\pi)$, we define $A*T\in\Hom_{\It}(\sigma,\pi)$ by
    \[
    (A*T)(w):=\int_{\Gt}\pi(h^{-1})A(T(h)w)\,dh.
    \]
    Now suppose that $T\in \Hcal_{z^{-1}}$ is such that $T(z^{-1})=e_{{\Uf}_{0}}e_{{\Uf}^{-}_{1}}$. Such an element exists by Lemma \ref{VeryPosSupp} because $z^{-1}U^{-}_{1}z\subseteq U^{-}_{2e+1}$. Then by definition 
	\begin{align*}
		q[A*T(w)] =&q[ \int_{\Gt}\pi(h^{-1})A(T(h)w)\,dh]\\
		=& q[\int_{\It z^{-1}\It}\pi(h^{-1})A(T(h)w)\,dh]\\
		=& q[\sum_{\It\backslash \It z^{-1}\It}\pi(h^{-1})A(T(h)w)].
	\end{align*}
	Now because $z^{-1}U^{-}_{1}z\subseteq U^{-}_{j}\cap U^{-}_{2e+1}$,
	\begin{align*}
		q[\sum_{\It\backslash \It z^{-1}\It}\pi(h^{-1})A(T(h)w)] =& q[\sum_{zU_{0}z^{-1}\backslash U_{0}}\pi(u^{-1}z)A(T(z^{-1})\sigma(u)w)]\\
		=& q[\sum_{zU_{0}z^{-1}\backslash U_{0}}\pi(z)A(e_{{\Uf}_{0}}e_{{\Uf}^{-}_{1}}\sigma(u)w)]\\
		\equiv&\pi_{U}(z)q[A(e_{{\Uf}_{0}}e_{{\Uf}^{-}_{1}}e_{{\Uf}_{0}}w)].
	\end{align*}
	By Lemma \ref{L3to1},
	\begin{equation*}
	\pi_{U}(z)q[A(e_{{\Uf}_{0}}e_{{\Uf}^{-}_{1}}e_{{\Uf}_{0}}w)]\equiv\pi(z)q[A(e_{{\Uf}_{0}}w)].
	\end{equation*}
	Thus we have shown
	\begin{equation}\label{SurjEq2}
		q[A*T(w)] \equiv\pi_{U}(z)q[A(e_{{\Uf}_{0}}w)].
	\end{equation}
	
	Now we set $A=F_1*T^{-1}$ and combine equations \eqref{SurjEq1} and \eqref{SurjEq2} to get 
    \begin{equation*}
		\pi_{U}(z)f \circ e_{{\Uf}_{0}}w\stackrel{\eqref{SurjEq1}}{\equiv}q\circ F_1(e_{{\Uf}_{0}}w)\equiv q\circ F_1*T^{-1}*T(e_{{\Uf}_{0}}w)\stackrel{\eqref{SurjEq2}}{\equiv}\pi_{U}(z)q[[F*T^{-1}](e_{{\Uf}_{0}}w)].
	\end{equation*}
	
	Thus 
	\begin{equation*}
		q\circ [F*T^{-1}]\circ e_{{\Uf}_{0}}\equiv f\circ e_{{\Uf}_{0}}
	\end{equation*}
	and the result follows.
\end{proof}

\begin{Lem}\label{LemInj}
	The map $\Hom(\sigma,\pi)\rightarrow \Hom(\sigma^{U_{0}},\pi_{U})$ defined by $f\mapsto q\circ f|_{\sigma^{U_{0}}}$ is a homomorphism of $(\It\cap \widetilde{B},\widetilde{B})$-bimodules. The map induced by restriction $\Psi:\Hom_{\It}(\sigma,\pi)\rightarrow \Hom_{\Tt_{0}}(\sigma^{U_{0}},\pi_{U})$ is an injective linear map.
\end{Lem}

\begin{proof}
	Suppose that $F\in \Hom_{\It}(\sigma,\pi)$ is such that $q\circ F\circ e_{{\Uf}_{0}}=0$. Since $F(\sigma^{U_{0}})$ is finite dimensional there exists $j\in\mathbb{Z}$ such that for all $w\in \sigma^{U_{0}}$ we have $\int_{U_{j^{\prime}}}\pi(u) F(w)du=0$ for any $j^{\prime}\leq j$.
	
	Let $z\in Z(\Tt)$ be such that $zU_{j}z^{-1}\subseteq U_{0}\cap U_{j+2e}$. Let $T\in\Hcal_{z^{-1}}$ be such that $T(z^{-1})=e_{{\Uf}_{0}}e_{{\Uf}^{-}_{1}}$, which exists by Lemma \ref{VeryPosSupp}. Then since $zU_{j}z^{-1}\subseteq U_{0}$ for $w\in \sigma^{U_{0}}$, we have 
	\begin{align*}
		F*T(w)=& \int_{\It z^{-1}\It}\pi(h^{-1})F(T(h)w)\,dh\\
		=&\int_{U_{0}}\pi(u^{-1}z)F(T(z^{-1})\sigma(u)w)\,du\\
		=&\int_{U_{0}}\pi(u^{-1}z)F(e_{{\Uf}_{0}}e_{{\Uf}^{-}_{1}}\sigma(u)w)\,du\\
		=&\pi(z)\int_{U_{0}}\pi(z^{-1}u^{-1}z)F(e_{{\Uf}_{0}}e_{{\Uf}^{-}_{1}}\sigma(u)w)\,du\\
		=&0,
	\end{align*}
        where the last equality follows because $zU_{j}z^{-1}\subseteq U_{0}$.
    
	Thus $F*T$ is zero on $\sigma^{U_{0}}$. Since $F*T$ is an $\It$-module homomorphism and $\sigma^{U_{0}}$ generates $\sigma$ as an $\It$-module, it follows that $F*T=0$. Since $T$ is invertible we see that $F=0$.
\end{proof}

Now, let us recall from \S \ref{TtReps} that we defined the representation
\[
\tau_\eta:=\eta\otimes\tau
\]
of $Z(\Tt)\Tt_0$, where $\eta$ is a genuine $W$-invariant character and $\tau$ is a pseudo-spherical representation of $\Tt_0$, and then for each unramified character $\chi:T\to\C^\times$ we defined the representation 
\[
i_{\eta,\tau}(\chi):=\chi\otimes\Ind_{Z(\Tt)\Tt_0}^{\Tt}\tau_{\eta}
\]
of $\Tt$, which is irreducible. We then have

\begin{Prop}\label{PSeriesConstits}
    Let $\pi$ be a smooth irreducible representation of $\Gt$. Then the $\sigma$-isotypic subspace $\pi^{\sigma}$ is nonzero if and only if $\Hom_{\Gt}(\pi,\Ind_{\Bt}^{\Gt}i_{\eta,\tau}(\chi))\neq 0$ for some unramified character $\chi:T\rightarrow \mathbb{C}^{\times}$. 
\end{Prop}

\begin{proof}
    We begin with a few preliminary remarks. The $\Tt$-representation $\pi_{U}$ is finitely generated, because $\pi$ is irreducible, and thus $\pi_{U}$ has an irreducible quotient. By Frobenius reciprocity it is equivalent to proving the following. The $\sigma$-isotypic subspace $\pi^{\sigma}$ is nonzero if and only if $\Hom_{\Tt}(\pi_{U},i_{\eta,\tau}(\chi))\neq 0$ for some unramified character $\chi:T\rightarrow \mathbb{C}^{\times}$. Now we start the proof.
    
    Suppose that $\pi^{\sigma}\neq 0$. This implies that $\Hom_{\It}(\sigma,\pi)\neq 0$. Thus by Lemma \ref{LemInj} it follows that $\Hom_{\Tt_{0}}(\sigma^{U_{0}},\pi_{U})\neq 0$. Since $\sigma^{U_{0}}\cong\tau$ as $\Tt_{0}$-representations we see that $\pi^{\tau}\neq0$. 

    Since $\tau$ is a $\Tt$-type and $\pi_{U}^{\tau}\neq 0$ it follows that $\pi_{U}$ has an irreducible quotient with a nonzero $\tau$-isotypic subspace. Thus by Proposition \ref{TTypes}, $\Hom_{\Tt}(\pi_{U},i_{\eta,\tau}(\chi))\neq 0$ for some unramified character $\chi:T\rightarrow \mathbb{C}^{\times}$.
    
    Now we prove the converse. Suppose that $\Hom_{\Tt}(\pi_{U},i_{\eta,\tau}(\chi))\neq 0$ for some unramified character $\chi:T\rightarrow \mathbb{C}^{\times}$. Since $i_{\eta,\tau}(\chi)$ is irreducible, there is a surjective $\Tt$-homomorphism $\pi_{U}\twoheadrightarrow i_{\eta,\tau}(\chi)$. Thus $\pi_{U}^{\tau}\neq 0$, because $i_{\eta,\tau}(\chi)^{\tau}\neq 0$. Finally we apply Lemma \ref{LemSurj} to see that $\pi^{\sigma}\neq0$.
\end{proof}

Now we can prove our main theorem. 

\begin{Thm}\label{TypesThm}\label{IType}
Let $\mathcal{M}(\Gt,\sigma)$ be the full subcategory of $\mathcal{M}(\Gt)$ consisting of objects generated by their $\sigma$-isotypic subspace.
	\begin{enumerate}
		\item The category $\mathcal{M}(\Gt,\sigma)$ is closed relative to subquotients in $\mathcal{M}(\Gt)$.\label{SubQuot}
		\item For every irreducible object $(\pi,V)$ in $\mathcal{M}(\Gt,\sigma)$, there is an unramified character $\chi:T\rightarrow \mathbb{C}^{\times}$ such that $(\pi,V)$ is isomorphic to an irreducible subquotient of $\Ind_{\widetilde{B}}^{\Gt}(i(\chi))$.\label{BernComp}
	\end{enumerate}
In particular, $(\It, \sigma)$ is a type.
\end{Thm}

\begin{proof} 
\eqref{SubQuot} can be proved as in \cite[Theorem 7.1]{Karasiewicz}, with \cite[Proposition 7.4]{Karasiewicz} being replaced by Proposition \ref{PSeriesConstits}. This argument essentially follows \cite[(8.3) Theorem]{BK98}.

\eqref{BernComp} is already part of Proposition \ref{PSeriesConstits}.
\end{proof}

As a corollary, we have the following local Shimura correspondence.

\begin{Cor}\label{C:local_Shimura_Correspondence}
    Assume the Cartan type of $G$ is $A_r, D_{2r+1}, E_6$ or $E_7$. Then the category $\mathcal{M}(\Gt,\sigma)$ is equivalent to the Bernstein block of the Iwahori spherical representations of $G'=G\slash Z_2$.
\end{Cor}
\begin{proof}
    This is immediate from the above theorem and Corollary \ref{HeckeShimCorr}. Also see the proof of \cite[Theorem 8.1]{Karasiewicz}.
\end{proof}




\appendix

%%%%%%%%%%%%%%%%%%%%%%%%%%%%%%%%%%%%%%%%%%%%

\section{Hilbert symbol over $2$-adic fields}\label{S:Appendix_Hilbert_symbol}

%%%%%%%%%%%%%%%%%%%%%%%%%%%%%%%%%%%%%%%%%%%%

In this appendix, we establish various properties of the quadratic Hilbert symbol for our $2$-adic field $F$. Everything in this appendix is likely known, but we include proofs because we have not been able to locate the proofs in the literature.

Since $F$ is $2$-adic, we first have
\begin{Lem}
The quadratic Hilbert symbol is not trivial on $\Ocal^\times\times\Ocal^\times$.
\end{Lem}
\begin{proof}
If $a\in\Ocal^\times$ is such that $(a, b)=1$ for all $a\in\Ocal^\times$ then $F(\sqrt{a})$ is an unramified extension of $F$ by \cite[Exercise 5, p.209]{Serre}. But certainly there is a unit $a\in\Ocal^\times$ such that $F(\sqrt{a})$ is ramified over $F$. 
\end{proof}


We need to analyze the Hilbert symbol restricted to $\Ocal^\times\times\Ocal^\times$. For this purpose, define
\[
R=\{a\in\Ocal^\times\st \text{$(a, -)$ is trivial on $\Ocal^\times$}\},
\]
which we call the {\it integral radical} of the Hilbert symbol. By the above lemma, $R\subsetneq\Ocal^\times$. Clearly we have $\Ocal^{\times2}\subseteq R$. 


Let us mention the following ``one-dimensionality" of the integral radical.
\begin{Lem}\label{DimRad}
Let us view $F^\times\slash F^{\times 2}$ as an $\F_2$-vector space. Then
\[
\dim_{\mathbb{F}_{2}}(RF^{\times 2}/F^{\times 2})=1,
\]
where we view $RF^{\times 2}\slash F^{\times 2}$ as a subspace of $F^\times\slash F^{\times 2}$.
\end{Lem}
\begin{proof}
The $\mathbb{F}_{2}$ linear map $F^{\times}/F^{\times 2}\rightarrow \Hom_{\mathbb{F}_{2}}(F^{\times}/F^{\times 2},\mathbb{F}_{2})$ defined by $t\mapsto (u\mapsto (t,u))$ is an isomorphism, because the Hilbert symbol is nondegenerate. Moreover, by definition, this map identifies $RF^{\times 2}/F^{\times 2}$ with the annihilator $\mathrm{ann}(\mathcal{O}^{\times}F^{\times 2}/F^{\times 2})$. Since $\mathcal{O}^{\times}F^{\times 2}/F^{\times 2}\subseteq F^{\times }/F^{\times 2}$ is a hyperplane, the lemma follows.
\end{proof}



%Let us also mention the following.
%\begin{Lem}
%Assume $p\neq 2$. Then
%\[
%(\varpi, \varpi)=\begin{cases}1&\text{if $\frac{q-1}{2}$ is even};\\
%-1&\text{if $\frac{q-1}{2}$ is odd}.
%\end{cases}
%\]
%If $p=2$ then there exists a uniformizer $\varpi\in \Ocal$ such that
%\[
%(\varpi, \varpi)_F=1.
%\]
%\end{Lem}
%\begin{proof}
%The case for $p\neq 2$ is well-known. (See, for example, \cite[p.336]{Neukirch}.) Assume $p=2$. Let $\varpi$ be any uniformizer. If $(\varpi, \varpi)=1$ then we are done. If $(\varpi, \varpi)=-1$ then pick a unit $u\in\Ocal^\times$ such that $(u, u)_F=-1$. (Such $u$ exists by the above lemma.) Then
%\begin{align*}
%(u\varpi, u\varpi)&=(u, u)\cdot(u, \varpi)\cdot(\varpi, u)\cdot(\varpi, \varpi)\\
%&=(-1)\cdot (u, \varpi)^2\cdot (-1)\\
%&=1.
%\end{align*}
%The lemma is proven.
%\end{proof}


For each $k\in\Z_{\geq 1}$, we set
\[
U_k=1+\varpi^k\Ocal.
\]
We often use the well-known
\begin{Lem}\label{L:square_root_exits_in_redidue_field}
Every element in $\Ocal\slash\varpi\Ocal$ has a square root, and hence $\Ocal^\times=\Ocal^{\times 2}U_1$.
\end{Lem}
\begin{proof}
The Frobenius map (the squaring map) on $\Ocal\slash\varpi\Ocal$ is an automorphism, and hence every element in $\Ocal\slash\varpi\Ocal$ has a square root. Then every element in $\Ocal^\times$ is written as $a^2+b\varpi$ for some nonzero $a\in\Ocal\slash\varpi\Ocal$ and some $b\in \Ocal\slash\varpi\Ocal$. We can then write $a^2+b\varpi=a^2(1+a^{-2}b\varpi)$.
\end{proof}

We then have
\begin{Lem}\label{L:2-adic_arithmetic}
\quad
\begin{enumerate}[(a)]
\item $U_{2e+1}\subseteq \Ocal^{\times 2}$.
\item\label{Val2eSqrs} For each $u\in\Ocal$, we have $1+4u\in\Ocal^{\times 2}$ if and only if there exists $v\in \Ocal$ such that $v+v^2=u$ in $\Ocal\slash\varpi\Ocal$.
\item $U_{2e}\nsubseteq\Ocal^{\times 2}$.
\item $U_{2e}\slash (U_{2e}\cap\Ocal^{\times 2})\cong\Z\slash 2\Z$.
\end{enumerate}
\end{Lem}
\begin{proof}
\quad
\begin{enumerate}[(a)]
\item This follows from \cite[(5.5) Proposition, page 137]{Neukirch} because the exponential map $\varpi^{e+1}\Ocal\to U_{e+1}$ is an isomorphism.

\item Let $u\in\Ocal$. Assume $1+4u\in\Ocal^{\times 2}$, so $1+4u=w^2$ for some $w\in\Ocal^\times$. One can then write $w=1+\varpi^{\ell}v$ for some $\ell\in\Z_{\geq 1}$ and $v\in \Ocal^\times$. So, we have $1+4u=(1+\varpi^{\ell} v)^2$. A short calculation shows that $\ell=e$, so $\varpi^{\ell}v\in 2\Ocal^\times$. Hence by re-choosing $v$ (if necessary), we have $1+4u=(1+2v)^2$, which implies $v+v^2=u$. Conversely, assume there exists $v\in \Ocal$ such that $v+v^2=u$ in $\Ocal\slash\varpi\Ocal$, namely $u=v+v^2+\varpi v'$ for some $v'\in\Ocal$. We then have $1+4u=(1+2v)^2+4\varpi v'$. But
\[
(1+2v)^2+4\varpi v'=(1+2v)^2(1+\frac{4\varpi v'}{(1+2v)^2})\qand 1+\frac{4\varpi v'}{(1+2v)^2}\in U_{2e+1}.
\]
By part (a), we know $U_{2e+1}\subseteq \Ocal^{\times 2}$, which implies $1+4u\in\Ocal^{\times 2}$.

\item Consider the group homomorphism $\sigma:\F_{q}\to\F_{q}$ given by $x\mapsto x+x^2$. The kernel is exactly $\mathbb{F}_{2}$. Hence there exists $u\in\Ocal$ such that $u=x+x^2$ has no solution. For such $u$, we know $1+4u\notin\Ocal^{\times 2}$ by part (b).

\item Recall that there is a group isomorphism $\Ocal/\varpi\Ocal \rightarrow U_{2e}/U_{2e+1}$ defined by $u+\varpi\Ocal \mapsto (1+4u)U_{2e+1}$. Under this map we see that the squares in $U_{2e}$ correspond to the elements of $\Ocal/\varpi\Ocal\cong \mathbb{F}_{q}$ in the image of the map $\sigma$ of part (c). Since the image of $\sigma$ has index $2$ in $\mathbb{F}_{q}$ the result follows.
\end{enumerate}
\end{proof}


In what follows, we will show $R=U_{2e}\Ocal^{\times 2}$. Let us start with one inclusion.
\begin{Lem}
We have $U_{2e}\subseteq R$.
\end{Lem}
\begin{proof}
Let $1-4u\in U_{2e}$ for $u\in\Ocal$. If $u\notin\Ocal^\times$, then $1-4u\in U_{2e+1}$ and we already know $U_{2e+1}\subseteq\Ocal^{\times 2}\subseteq R$. Hence we assume $u\in\Ocal^\times$. We have to show $(1-4u, a)=1$ for all $a\in\Ocal^\times$. Let $v\in\Ocal$ be such that $1-4v\in U_{2e}\cap\Ocal^{\times 2}$. Then for all $a\in F^\times$ we have
\begin{align*}
(1-4u, a)&=((1-4u)(1-4v), a)=(1-4(u+v)+16uv,a)\\
&=((1-4(u+v))(1+\frac{16uv}{1-4(u+v)}),a)=(1-4(u+v),a),
\end{align*}
where we used $1+\frac{16uv}{1-4(u+v)}\in U_{2e+1}\subseteq \Ocal^{\times 2}$ by Lemma \ref{L:2-adic_arithmetic} (a). Also $(1-4(u+v), u+v)=(1-4(u+v), 4(u+v))=1$, where we used $(1-x, x)=1$ for all $x\in F^\times$. Hence by setting $a=u+v$, we obtain
\begin{equation}\label{E:(1-4u, u+v)=1}
(1-4u, u+v)=(1-4(u+v),u+v)=1.
\end{equation}

Now, for $a\in\Ocal^\times$, one can write $au^{-1}=r^2(1+\varpi s)$ for some $r\in\Ocal^\times$ and $s\in\Ocal$ by Lemma \ref{L:square_root_exits_in_redidue_field}. Then
\[
(1-4u, a)=(1-4u, ur^2(1+\varpi s))=(1-4u, u+\varpi(us)).
\]
But $1-4\varpi(us)\in U_{2e+1}\subseteq\Ocal^{\times 2}$, and hence by taking $v=\varpi(us)$ in \eqref{E:(1-4u, u+v)=1}, we obtain $(1-4u, a)=1$.
\end{proof}



%\begin{Lem}
%\quad
%\begin{enumerate}[(a)]
%\item $U_{2e}\subseteq R$.
%\item $U_{2e-1}\nsubseteq R$.
%\end{enumerate}
%\end{Lem}
%\begin{proof}
%\quad
%\begin{enumerate}[(a)]
%\item Let $1-4u\in U_{2e}$. We have to show $(1-4u, a)=1$ for all $a\in\Ocal^\times$. Let $v\in\Ocal$ be such that $1-4v\in U_{2e}\cap\Ocal^{\times 2}$. Then for all $a\in F^\times$ we have
%\begin{align*}
%(1-4u, a)&=((1-4u)(1-4v), a)=(1-4(u+v)+16uv,a)\\
%&=((1-4(u+v))(1+\frac{16uv}{1-4(u+v)}),a)=(1-4(u+v),a),
%\end{align*}
%where we used $\frac{16uv}{1-4(u+v)}\in U_{2e+1}\subseteq \Ocal^{\times 2}$ by Lemma \ref{L:2-adic_arithmetic} (a). Also $(1-4(u+v), u+v)=(1-4(u+v), 4(u+v))=1$, where we used $(1-x, x)=1$ for all $x\in F^\times$. Hence by setting $a=u+v$.
%\begin{equation}\label{E:(1-4u, u+v)=1}
%(1-4u, u+v)=(1-4(u+v),u+v)=1.
%\end{equation}
%
%Now, for $a\in\Ocal^\times$, one can write $au^{-1}=r^2(1+\varpi s)$ for some $r\in\Ocal^\times$ and $s\in\Ocal$ because every element in $\Ocal\slash\varpi\Ocal$ has a square root. Then
%\[
%(1-4u, a)=(1-4u, ur^2(1+\varpi s))=(1-4u, u+\varpi(us)).
%\]
%But $1-4\varpi(us)\in U_{2e+1}\subseteq\Ocal^{\times 2}$ and hence by taking $v=\varpi(us)$ in \eqref{E:(1-4u, u+v)=1}, we obtain $(1-4u, a)=1$.
%
%\item Let $1-\varphi^{2e-1}v\in U_{2e-1}\cap R$, where we assume $v\in\Ocal^\times$. Since $(x, 1-x)=1$ for all $x\in F^\times$, we have
%\[
%1=(1-\varphi^{2e-1}v, \varpi^{2e-1}v)=(1-\varphi^{2e-1}v, \varpi),
%\]
%which implies that $(1-\varphi^{2e-1}v, -)$ is trivial on $F^\times$. Hence we must have $1-\varphi^{2e-1}v\in\Ocal^{\times 2}$. But we already know $U_{2e-1}\nsubseteq\Ocal^{\times 2}$. Hence $U_{2e-1}\cap R\subsetneq U_{2e-1}$.
%\end{enumerate}
%\end{proof}

\begin{comment}
Next we have
\begin{Lem}
Let $k\leq 2e$. Then 
\[
U_{k}\cap\Ocal^{\times 2}\subsetneq U_{k-1}\cap\Ocal^{\times 2}.
\]
\end{Lem}
\begin{proof}
First note that since $k\geq 2e$ we know
\[
U_{k}\slash (U_{k}\cap\Ocal^{\times 2})\cong U_{k}\Ocal^{\times 2}\slash\Ocal^{\times 2}\neq 0
\]
by Lemma \ref{L:2-adic_arithmetic} (d). Also we have the inclusion
\[
0\neq U_{k}\slash (U_{k}\cap\Ocal^{\times 2})\cong U_{k}\Ocal^{\times 2}\slash\Ocal^{\times 2}\xhookrightarrow{\quad}
U_{k-1}\Ocal^{\times 2}\slash\Ocal^{\times 2}\cong U_{k-1}\slash (U_{k-1}\cap\Ocal^{\times 2}).
\]
Since $[U_{k}: U_{k-1}]=2^f$, we must have
\[
[U_{k-1}\slash (U_{k-1}\cap\Ocal^{\times 2}):U_{k}\slash (U_{k}\cap\Ocal^{\times 2})]< 2^{f},
\]
where the inequality is strict because $U_{k}\slash (U_{k}\cap\Ocal^{\times 2})\neq 0$.


Now, assume $U_{k}\cap\Ocal^{\times 2}=U_{k-1}\cap\Ocal^{\times 2}$. Then
\begin{align*}
U_{k-1}\slash (U_{k-1}\cap\Ocal^{\times 2})\Big\slash U_{k}\slash (U_{k}\cap\Ocal^{\times 2})
&\cong U_{k-1}\slash (U_{k}\cap\Ocal^{\times 2})\Big\slash U_{k}\slash (U_{k}\cap\Ocal^{\times 2})\\
&=U_{k-1}\Big\slash U_{k},
\end{align*}
so
\[
[U_{k-1}\slash (U_{k-1}\cap\Ocal^{\times 2}):U_{k}\slash (U_{k}\cap\Ocal^{\times 2})]= 2^f,
\]
which is a contradiction.
\end{proof}
\end{comment}

We can now prove
\begin{Prop}
We have $U_{2e}\Ocal^{\times 2}= R$, and hence $R\slash\Ocal^{\times 2}=\Z\slash 2\Z$.
\end{Prop}
\begin{proof}
We already know $\Ocal^{\times 2}\subseteq U_{2e}\Ocal^{\times 2}\subseteq R$ and $[U_{2e}\Ocal^{\times 2}:\Ocal^{\times 2}]=2$. By Lemma \ref{DimRad} we know that $[R:\Ocal^{\times 2}]=2$. Thus $R=U_{2e}\Ocal^{\times 2}$.
\end{proof}


\begin{comment}
To show the other inclusion, let $u\in R$. Then there exists $k$ such that $u\in U_k\Ocal^{\times 2}$. If $k\geq 2e$, then clearly $u\in U_{2e}\Ocal^{\times 2}$. So assume $k<2e$. 

Assume $1\leq k <2e$. Then we may assume $k$ is odd. Indeed, if $k$ is even, then by the above lemma there exists $v\in (U_{k+1}\Ocal^\times\smallsetminus U_{k})\cap \Ocal^{\times 2}$. Then $uv\in U_{k+1}\Ocal^{\times 2}$, and still $uv\in R$ because $v\in \Ocal^{\times 2}\subseteq R$. Let $k$ be the minimum such odd integer, so that we can write $u=(1-\varpi^{k}x)w$ for some $x\in\Ocal^\times$ and $w\in\Ocal^{\times 2}$. Note then that $1-\varpi^{k}x\in R$. Recall that $(1-a, a)=1$ for all $a\in F^\times$. Hence we have
\[
1=(\varpi^kx, 1-\varpi^kx)=(\varpi, 1-\varpi^kx)\cdot (x, 1-\varpi^kx)=(\varpi, 1-\varpi^kx),
\]
where the second equality follows because $k$ is odd and the third one follows because $1-\varpi^{k}x\in R$. But this implies $(a, 1-\varpi^kx)=1$ for all $a\in F^\times$, so $1-\varpi^kx\in F^{\times 2}$, which implies $u\in\Ocal^{\times 2}\subseteq U_{2e}\Ocal^{\times 2}$. Assume $u\in\Ocal^\times\smallsetminus U_1$. Then we already know $u$ is a square by Lemma \ref{L:square_root_exits_in_redidue_field}. 

Hence we have shown $R=U_{2e}\Ocal^{\times 2}$, which gives $R\slash\Ocal^{\times 2}=\Z\slash 2\Z$ by Lemma \ref{L:2-adic_arithmetic} (d).
\end{comment}


The above discussion can be summarized as follows:
\[
\begin{tikzcd}
&\Ocal^{\times 2}\ar[r, phantom, "\subseteq"]&R=U_{2e}\Ocal^{\times 2}&\\
U_{2e+1}\ar[r, phantom, "\subseteq"]&U_{2e}\cap\Ocal^{\times 2}\ar[r, phantom, "\subseteq"]\ar[u, phantom, "{\rotatebox{90}{$\subseteq$}}"]&U_{2e}\ar[r, phantom, "\subseteq"]\ar[u, phantom, "{\rotatebox{90}{$\subseteq$}}"]&U_{2e-1}
\rlap{\, ,}
\end{tikzcd}
\]
where
\[
U_{2e}\slash (U_{2e}\cap\Ocal^{\times 2})=\Z\slash 2\Z.
\]

%\[
%U_{2e+1}\Ocal^{\times 2}=\Ocal^{\times 2}\subsetneq U_{2e}\Ocal^{\times 2}=R\subsetneq U_{2e-1}\Ocal^{\times 2},
%\]
%where
%\[
%U_{2e}\Ocal^{\times 2}\slash \Ocal^{\times 2}= U_{2e}\slash (U_{2e}\cap\Ocal^{\times 2})=\Z\slash 2\Z.
%\]


Let us also mention
\begin{Cor}\label{C:dimension_integral_radical}
\[
[\Ocal^{\times}:R]=2^{ef}.
\]
\end{Cor}
\begin{proof}
It is well-known that
\[
[\Ocal^\times: \Ocal^{\times 2}]=2^{ef+1}.
\]
(See, for example, \cite[(5.8) Proposition, p.142]{Neukirch}.) Since $R\slash\Ocal^{\times 2}=\Z\slash 2\Z$, the corollary follows.
\end{proof}


\begin{comment}
%%%%%%%%%%%%%%%%%%%%%%%%%%%%%%%%%%%%%%%%%%%%

\section{More on the support}

%%%%%%%%%%%%%%%%%%%%%%%%%%%%%%%%%%%%%%%%%%%%



This Lemma along with \cite[Proposition 2.16]{CMO19} should imply that the $2$-cocycle in $\Hcal_{2}$ is trivial for sufficiently large strictly dominant coweights.

We can also prove a slight variant on the last lemma. There are fewer assumptions on $y$, but the conclusion is also weaker. Can this be improved?  

\begin{Lem}
	Let $y\in \Yt$ such that $y$ is positive. Then for any $f\in \Hom_{ I_2\cap\,^{y(\varpi)} I_2}(\sigma_{2},\,^{y(\varpi)}\sigma_{2})$ we have
	\begin{equation*}
		f=fe_{{\Uf}^{-}_{2}\cap\ker\,^{y(\varpi)}\sigma_{2}}=e_{{\Uf}_{0}\cap\,^{y(\varpi)^{-1}}\ker\sigma_{2}}f.
	\end{equation*}
\end{Lem}



%%%%%%%%%%%%%%%%%%%%%%%%%%%%%%%%%%%%%%%%%%%%

\section{To Do}

%%%%%%%%%%%%%%%%%%%%%%%%%%%%%%%%%%%%%%%%%%%%

\begin{itemize}
	\item Add signposting
	\item Case of $E_{8}$?
	\item Can we compute the Whittaker dimension for the basic $\widetilde{{\Kf}}$-type?
	\item Check Lemma \ref{LemVarIndRootMove2} very carefully.
	\item Check Prop \ref{HMult1} very carefully.
	\item $\dim \sigma_{2}=?$
	\item Improve Prop \ref{PseudoSphericalBranch} and Theorem \ref{KRep}.
	\item Satake Iso
	\item The Shimura Corr Hecke alg iso depends on the choice of signs for elts in $\Omega\cong \widetilde{Y}/2Y$. This seems to be a bit different from what I wrote at the end of my ANT paper. Need to think about what is the correct set of choices to track. 
	\item In finite Shim Corr we can use different $\tau$ to construct $\sigma$. Do these yield disjoint pieces of the category?
        \item Need to see where Marty's work on covers of tori fits in.
\end{itemize}
\end{comment}

%\nocite{}
\bibliographystyle{amsalpha}
\bibliography{Hecke_algebra2_bio}


\end{document}





